% Poly "Probabilités" -- Mines Paris, UE11
% Compilation : pdflatex poly.tex (x2)   ou   latexmk -pdf poly.tex
% Contenu : conversion des sources markdown (md/ProbaI.md ... md/ProbaV.md) en chapitres LaTeX (chapitres/ProbaI.tex ... ProbaV.tex),
% corrigés regroupés en fin de poly (chapitres/ProbaX_corrections.tex). Même mise en forme que le poly d'Equations différentielles.
\documentclass[french,a4paper,11pt,open=any,headings=small]{scrbook}
\usepackage{../cdis} % mise en forme commune aux polys (packages, environnements, macros communes, corrections)
\hypersetup{
  pdftitle={Probabilités},
  pdfauthor={UE11, Mines Paris - PSL}}

% ---- Macros de Probabilités ----
\providecommand{\P}{}\renewcommand{\P}{\mathbb{P}}
\newcommand{\A}{\mathcal{A}}
\newcommand{\E}{\mathcal{E}}
\newcommand{\B}{\mathcal{B}}
\providecommand{\L}{}\renewcommand{\L}{\mathcal{L}}
\providecommand{\No}{}\renewcommand{\No}{\mathcal{N}}
\newcommand{\Esp}{\mathbb{E}}
\newcommand{\V}{\mathbb{V}}
\newcommand{\cov}{\text{Cov}}
% niveaux de difficulté des exercices
\newcommand{\zero}{$\mathord{\boldsymbol{\circ}}$}
\newcommand{\one}{$\mathord{\bullet}$}
\newcommand{\two}{$\mathord{\bullet}\mathord{\bullet}$}
\newcommand{\three}{$\mathord{\bullet}\mathord{\bullet}\mathord{\bullet}$}
\newcommand{\four}{$\mathord{\bullet}\mathord{\bullet}\mathord{\bullet}\mathord{\bullet}$}

% Les solutions de ce poly sont repérées par un \label (et non un \hypertarget)
\renewcommand{\cdissollink}[2]{\hyperref[#1]{#2}}

\title{Probabilités}
\author{UE11, Maths 2 \\
	Mines Paris - PSL\\~\\
	Thomas Romary et Mike Pereira\footnote{Ce document est un des produits
		du projet \href{https://github.com/}{\(\mbox{\faGithub}\)
			\texttt{paulinebernard/CDIS}} issu de la collaboration de
		\href{mailto:pauline.bernard@minesparis.psl.eu}{Pauline Bernard}
		(CAS), \href{mailto:thomas.romary@minesparis.psl.eu}{Thomas
			Romary} et \href{mailto:mike.pereira@minesparis.psl.eu}{Mike Pereira} (STIM). 
		%Il dérive du projet
		%  \href{https://github.com/}{\(\mbox{\faGithub}\)
			%  \texttt{boisgera/CDIS}}, initié par la collaboration de
		%  \href{mailto:sebastien.boisgerault@mines-paristech.fr}{(S)ébastien
			%  Boisgérault} (CAOR),
		%  \href{mailto:thomas.romary@mines-paristech.fr}{(T)homas Romary} et
		%  \href{mailto:emilie.chautru@mines-paristech.fr}{(E)milie Chautru}
		%  (GEOSCIENCES),
		%  \href{mailto:pauline.bernard@mines-paristech.fr}{(P)auline Bernard}
		%  (CAS), avec la contribution de
		%  \href{mailto:gabriel-stolz@mines-paristech.fr}{Gabriel Stoltz} (Ecole
		%  des Ponts ParisTech, CERMICS). 
		Il est mis à disposition selon les
		termes de \href{http://creativecommons.org/licenses/by-nc-sa/}{la
			licence Creative Commons ``attribution -- pas d'utilisation
			commerciale -- partage dans les mêmes conditions'' 4.0
			internationale}.}}
\date{\today}


\begin{document}
\maketitle

\setcounter{tocdepth}{2}
\tableofcontents

\part{Probabilités et variables aléatoires}
\chapter{Introduction}\label{I-introduction}

Le but de ce cours est de consolider et compléter les connaissances en
théorie des probabilités acquises en CPGE mais surtout de permettre
d'acquérir le raisonnement probabiliste. En effet, les probabilités
peuvent être vues comme un outil de modélisation de phénomènes qui ont
la caractéristique d'être aléatoires. L'aléatoire peut intervenir de
différentes manières dans ces phénomènes.

\begin{itemize}
\tightlist
\item
  Dans les cas d'école que sont les jeux de pile ou face ou de lancés de
  dés, la différence entre les résultats, si l'on réitère l'expérience,
  peut être liée à l'impulsion initiale communiquée au dé et à d'autres
  facteurs environnementaux comme le vent, la rugosité de la table, etc.
  Le hasard intervient du fait de la méconnaissance des conditions
  initiales, car la pièce ou le dé ont des trajectoires parfaitement
  définies par la mécanique classique.
\item
  Dans beaucoup de cas de figure, on fait intervenir l'aléatoire dans la
  modélisation du fait d'une connaissance incomplète des phénomènes. On
  parle alors de modélisation de l'incertitude. C'est le cas par exemple
  en sciences du climat, dont nous discutons ci-dessous.
\item
  Dans certains domaines, tels la physique quantique, l'aléatoire fait
  intrinsèquement partie de la théorie.
\end{itemize}

Les probabilités sont aussi un préalable indispensable pour aborder
l'analyse statistique des données et les méthodes d'apprentissage
automatique.

En CPGE, les probabilités ont été vues dans le cadre de phénomènes
aléatoires qui admettent un nombre au plus dénombrable de résultats
possibles. Ce cadre restreint est supposé connu. On pourra se reporter
au chapitre 3 du cours de Garnier
(\protect\hyperlink{refI-polygarnier}{2018}) ou aux deux premiers
chapitres du cours de Jourdain (\protect\hyperlink{refI-polyponts}{2016})
pour une éventuelle mise à niveau.

\section{Exemple de la modélisation du
climat}\label{I-exemple-de-la-moduxe9lisation-du-climat}

Les prédictions de la météo et les projections du climat proviennent
généralement de modèle numériques qui simulent les différents processus
physiques à l'oeuvre. Les incertitudes dans la construction et
l'application de ces modèles sont variées et peuvent être réparties en
quatre groupes~: les conditions initiales (on ne connaît jamais
parfaitement l'ensemble des variables climatiques en tout point du
globe), les conditions aux limites (par exemple lorsqu'on travaille à
l'échelle d'un continent ou d'un pays), les valeurs des paramètres
intervenant dans les modèles (constantes issues d'observations
diverses), les incertitudes structurelles enfin qui relèvent des choix
de modélisation. Pour tenir compte de ces incertitudes, les
climatologues effectuent des ensembles de simulations, où les
différentes quantités incertaines sont échantillonnées selon des modèles
probabilistes (voir en particulier Tebaldi and Knutti
(\protect\hyperlink{refI-tebaldi}{2007})).

L'incertitude sur les conditions initiales est particulièrement
influente aux faibles échelles de temps. La météo est un système
chaotique, les prévisions sont extrêmement sensibles aux variations des
conditions initiales utilisées pour initialiser les modèles. Ces
dernières sont en revanche connues de manière imparfaite. Il est ainsi
nécessaire de tenir compte de cette incertitude. Ceci est rendu possible
par la modélisation probabiliste. Les conditions aux limites font
intervenir des \textbf{variables continues} (température, pression,
vitesse du vent, etc.), ce qui représente la principale nouveauté par
rapport au programme de CPGE, et il est nécessaire de caractériser leurs
relations de \textbf{dépendance}. Il convient par ailleurs de tenir
compte des observations (satellites, station de mesure, \ldots) en
incluant cette information dans la modélisation probabiliste. On parle
de \textbf{conditionnement} aux données. Il s'agit enfin d'en générer
les valeurs via des algorithmes de \textbf{simulation stochastique}. La
validité de l'approche est assurée par les \textbf{théorèmes limites},
qui garantissent la représentativité des ensembles générés.

\section{Plan du cours}\label{I-plan-du-cours}

Le cours est organisé en 5 amphis et abordera consécutivement les
notions évoquées ci-dessus, à savoir les probabilités définies sur
\(\mathbb{R}\), les variables et vecteurs aléatoires réels,
l'indépendance et le conditionnement de variables aléatoires, l'étude
des suites de variables aléatoires et, enfin, les méthodes de simulation
stochastique.

\section{Historique}\label{I-historique}

Avant que l'étude des probabilités soit considérée comme une science,
l'observation du hasard dans les événements naturels a amené les
philosophes et les scientifiques à réfléchir sur la notion de liens
entre événements, causes et conséquences, et lois de la nature. Les jeux
de hasard, les situations météorologiques ou les trajectoires des astres
ont fait partie des domaines étudiés. Les explications données sont
alors liées au destin, à une colère céleste ou à une présence divine.

Il est communément admis que le début de la science des probabilités se
situe au XVIe siècle avec l'analyse de jeux de hasard par Jérôme Cardan
et au XVIIe siècle avec les discussions entre Pierre de Fermat et Blaise
Pascal au sujet de paradoxes issus de ces jeux, notamment posés par
Antoine Gombaud, chevalier de Méré. Cette nouvelle théorie est nommée
géométrie aléatoire par le chevalier de Méré en 1654, elle est appelée
par la suite calcul conjectural, arithmétique politique et plus
communément aujourd'hui théorie des probabilités. Cette théorie, dite
des probabilités modernes, est alors étudiée par de nombreux penseurs
jusqu'au XIXe siècle~: Kepler, Galilée, Leibniz, Huygens, Halley,
Buffon, les frères Bernoulli, Moivre, Euler, D'Alembert, Condorcet,
Laplace, Fourier. Elle est principalement basée sur les événements
discrets et la combinatoire.

Des considérations analytiques ont forcé l'introduction de variables
aléatoires continues dans la théorie. Cette idée prend tout son essor
dans la théorie moderne des probabilités, dont les fondations ont été
posées par Andreï Nikolaevich Kolmogorov. Kolmogorov combina la notion
d'univers, introduite par Richard von Mises et la théorie de la mesure
pour présenter son système d'axiomes pour la théorie des probabilités en
1933. Très vite, son approche devint la base incontestée des
probabilités modernes.

Le XXe siècle voit également le développement de l'application de la
théorie des probabilités dans plusieurs sciences.

Avec la mécanique newtonienne, la théorie du champ électromagnétique ou
la thermodynamique, la physique classique est la théorie utilisée
jusqu'à la fin du XIXe siècle. En 1925, Erwin Schrödinger étudie
l'équation qui détermine l'évolution d'une onde au cours du temps~:
l'équation de Schrödinger. Max Born utilise cette équation pour décrire
une collision entre des particules telles que des électrons ou des
atomes. Les observations de ces expériences l'amènent à supposer que le
module de la fonction d'onde est la probabilité que la particule soit
détectée en un point de l'espace. C'est le début d'une nouvelle approche
de la physique quantique.

En 1900, Louis Bachelier fut un des premiers mathématiciens à modéliser
les variations de prix boursiers grâce à des variables aléatoires. « le
marché n'obéit qu'à une seule loi~: la loi du hasard ». Bachelier
utilise alors le calcul stochastique pour étudier les variations
boursières au cours du temps. En 1970, Fischer Black et Myron Scholes
reprennent les idées de Bachelier pour modéliser les rendements d'une
action.

L'utilisation des probabilités en biologie a pris un essor dans les
années 1970, notamment dans l'étude de l'évolution des espèces. La
reproduction des individus est modélisée par un choix aléatoire des
gènes transmis ainsi que des mutations apparaissant de manière aléatoire
sur les individus. L'extinction des espèces ou des gènes est alors
étudiée en fonction des effets stochastiques.

De nos jours, l'Ecole française de Probabilités est très active. La
première Médaille Fields décernée à un probabiliste a été attribuée à
Wendelin Werner en 2006. Les probabilités se développent de plus en
plus, alimentées en particulier de manière essentielle par la physique,
le développement des réseaux de télécommunications, la finance, la
biologie, la médecine\ldots{} Elles permettent de construire des modèles
mathématiques, qui peuvent être validés par les données suivant la
théorie statistique, et fournissent également des possibilités
d'expérimentations fictives dans de multiples domaines d'applications.

\chapter{Objectifs d'apprentissage}\label{I-objectifs-dapprentissage}

\runinhead{Probabilités des
événements}\phantomsection\label{I-probabilituxe9s-des-uxe9vuxe9nements}

\begin{itemize}
\item
  connaître les définitions de base des probabilités~:

  \begin{itemize}
  \tightlist
  \item
    $\mathord{\boldsymbol{\circ}}$ espace fondamental \(\Omega\)
  \item
    $\mathord{\boldsymbol{\circ}}$ événement \(A \subset \Omega\)
  \item
    $\mathord{\boldsymbol{\circ}}$ tribu (de parties)
  \item
    $\mathord{\bullet}$ probabilité
  \item
    $\mathord{\bullet}$ propriété presque sûre
  \end{itemize}
\item
  connaître, savoir démontrer et exploiter les résultats suivants~:

  \begin{itemize}
  \tightlist
  \item
    $\mathord{\bullet}$ les propriétés élémentaires d'une probabilité
  \item
    $\mathord{\bullet}$ le théorème de continuité monotone
  \end{itemize}
\item
  indépendance et conditionnement

  \begin{itemize}
  \tightlist
  \item
    $\mathord{\bullet}$ connaître la définition de la probabilité
    conditionnelle
  \item
    $\mathord{\bullet}$ connaître et savoir exploiter la formule des
    probabilités totales
  \item
    $\mathord{\bullet}$ connaître et savoir exploiter la formule de
    Bayes
  \item
    $\mathord{\bullet}$ connaître la définition de l'indépendance
    d'événements
  \end{itemize}
\end{itemize}

\runinhead{Variables
aléatoires}\phantomsection\label{I-variables-aluxe9atoires}

\begin{itemize}
\item
  variables aléatoires réelles

  \begin{itemize}
  \tightlist
  \item
    $\mathord{\bullet}\mathord{\bullet}$ connaître la définition de la
    tribu borélienne
  \item
    $\mathord{\bullet}\mathord{\bullet}$ savoir manipuler la notion de
    tribu engendrée
  \item
    $\mathord{\bullet}\mathord{\bullet}$ savoir que la tribu borélienne
    est engendrée par les intervalles de la forme
    \(]-\infty, x],~ x\in \mathbb{R}\)
  \item
    $\mathord{\bullet}$ connaître la définition d'une variable aléatoire
    réelle et de sa loi de probabilité
  \end{itemize}
\item
  loi des variables aléatoires réelles

  \begin{itemize}
  \tightlist
  \item
    $\mathord{\bullet}$ savoir que la fonction de répartition
    caractérise la loi d'une variable aléatoire réelle
  \item
    $\mathord{\bullet}\mathord{\bullet}$ connaître les propriétes de
    base de la fonction de répartition
  \item
    $\mathord{\bullet}\mathord{\bullet}$ savoir calculer la fonction de
    répartition d'une variable aléatoire
  \end{itemize}
\item
  densité de probabilité

  \begin{itemize}
  \tightlist
  \item
    $\mathord{\bullet}\mathord{\bullet}$ connaître la définition d'une
    densité de probabilité
  \item
    $\mathord{\bullet}\mathord{\bullet}$ connaître la définition d'une
    variable aléatoire réelle à densité
  \item
    $\mathord{\bullet}\mathord{\bullet}$ savoir identifier si une
    variable aléatoire admet une densité
  \end{itemize}
\end{itemize}

\chapter{Probabilités des
événements}\label{I-probabilituxe9s-des-uxe9vuxe9nements-1}

\section{Phénomènes aléatoires et
événements}\label{I-phuxe9nomuxe8nes-aluxe9atoires-et-uxe9vuxe9nements}

L'objet de la théorie des probabilités est l'analyse mathématique de
phénomènes dans lesquels le hasard intervient. Les phénomènes aléatoires
résultent d'expériences dont le résultat ne peut être prédit à l'avance
et qui \textbf{peut} varier si on répète l'expérience dans des
conditions identiques.

Il est aisé de trouver des exemples de tels phénomènes.

\begin{exemple}\phantomsection\label{I-exemples}
\begin{enumerate}
\def\labelenumi{\arabic{enumi}.}
\tightlist
\item
  Jeu de Pile ou Face
\item
  Lancé de dés
\item
  Durée de vie d'une ampoule électrique
\item
  Température demain à 12h au sommet de la tour Eiffel
\item
  Évolution de la vitesse d'une molécule dans un gaz raréfié sur un
  intervalle de temps \([t_1,t_2]\)
\end{enumerate}

La théorie des probabilités vise à fournir un modèle mathématique pour
décrire ces phénomènes. Elle repose sur trois ingrédients essentiels
dont on donne ici les définitions.
\end{exemple}

\addcontentsline{toc}{subsection}{L'espace fondamental}
\begin{definition}[{L'espace
fondamental}]\phantomsection\label{I-lespace-fondamental}

Noté habituellement \(\Omega\), l'\emph{espace fondamental} (ou encore
l'\emph{espace d'état} ou \emph{univers}) contient l'ensemble de tous
les résultats possibles d'un phénomène aléatoire. Un résultat possible
d'une expérience sera noté \(\omega \in \Omega\).

Si on reprend les exemples précédents, on peut facilement définir les
univers associés.
\end{definition}

\begin{exemple}\phantomsection\label{I-exemples-1}
\begin{enumerate}
\def\labelenumi{\arabic{enumi}.}
\tightlist
\item
  Jeu de Pile ou Face, \(\Omega = \{\text{pile}, \text{face}\}\)
\item
  Lancé de dés, \(\Omega = \{1,2,3,4,5,6\}\)
\item
  Durée de vie d'une ampoule électrique, \(\Omega = [0,+\infty [\)
\item
  Température demain à 12h au sommet de la tour Eiffel (en degrés
  Kelvin), \(\Omega = [0,+\infty [\)
\item
  Évolution de la vitesse d'une molécule dans un gaz raréfié sur un
  intervalle de temps \([t_1,t_2]\), \(\Omega\)~: ensemble des
  application continues sur \([t_1,t_2]\) à valeurs dans
  \(\mathbb{R}^3\)
\end{enumerate}

Cette liste d'exemples montre que l'espace \(\Omega\) peut varier
énormément dans sa structure, d'une expérience à l'autre. Cela permet de
réaliser la richesse de la théorie qu'il faut mettre en place, pour
créer un modèle qui englobe tous ces cas. Nous verrons également
ultérieurement que le modèle abstrait que nous allons construire
permettra de s'affranchir du fait que \(\Omega\) décrit précisément tous
les résultats possibles de l'expérience.
\end{exemple}

\addcontentsline{toc}{subsection}{Evénement}
\begin{definition}[{Evénement}]\phantomsection\label{I-evuxe9nement}

Un \emph{événement} est une propriété qui est vérifiée ou non une fois
l'expérience réalisée. On identifie un événement \(A\) à un
sous-ensemble ou \emph{partie} de \(\Omega\),
i.e.~\(A = \{\omega \in \Omega : A \text{ est vérifiée pour } \omega \}\).
\end{definition}

\begin{exemple}\phantomsection\label{I-exemples-2}
\begin{enumerate}
\def\labelenumi{\arabic{enumi}.}
\tightlist
\item
  Jeu de Pile ou Face~: \(A = \{\text{pile}\}\).
\item
  Lancé de dés~: \(A = \{1,3,5\}\).
\item
  Durée de vie d'une ampoule électrique~:
  \(A = [t_1,t_2] \subset \mathbb{R}_+\).
\item
  Température demain à 12h au sommet de la tour Eiffel (en degrés
  Kelvin)~: \(A = [T_1,T_2]\cup[T_3,T_4] \subset \mathbb{R}_+\).
\item
  Évolution de la vitesse d'une molécule dans un gaz raréfié sur un
  intervalle de temps \([t_1,t_2] \subset \mathbb{R}_+\)~:
  \(A = \{ f \in C([t_1,t_2], \mathbb{R}^3) : \|f-g\|_{\infty} \leq a \}\),
  où \(g \in C([t_1,t_2],\mathbb{R}^3)\) et \(a \in \mathbb{R}_+\).
\end{enumerate}

Les événements étant des ensembles, les opérations ensemblistes
classiques admettent une interprétation probabiliste.
\end{exemple}

\subsection*{Correspondance entre opérations logiques et
ensemblistes}\label{I-correspondance-entre-opuxe9rations-logiques-et-ensemblistes}

\begin{longtable}[]{@{}ccc@{}}
\toprule
Terminologie probabiliste & Terminologie ensembliste &
Notation\tabularnewline
\midrule
\endhead
événement certain & ensemble entier & \(\Omega\)\tabularnewline
événement impossible & ensemble vide & \(\varnothing\)\tabularnewline
événement contraire & complémentaire & \(A^c\)\tabularnewline
événement atomique & singleton & \(\{\omega\}\)\tabularnewline
implication & inclusion & \(\subset\)\tabularnewline
et & intersection & \(\cap\)\tabularnewline
ou & réunion & \(\cup\)\tabularnewline
événements incompatibles & ensembles disjoints &
\(A_1\cap A_2 = \varnothing\)\tabularnewline
\bottomrule
\end{longtable}

On doit maintenant répondre à la question de savoir quels sont les
événements dont on va vouloir évaluer la probabilité d'occurence. On va
ainsi regrouper les événements en un ensemble \(\mathcal{A}\) qui
constitue une collection de sous-ensembles de \(\Omega\). On va
souhaiter en particulier pouvoir combiner des événements au sein de
\(\mathcal{A}\) par les opérations ensemblistes courantes.

Ceci conduit à la notion de \emph{tribu de parties} de \(\Omega\).

\addcontentsline{toc}{subsection}{Tribu}
\begin{definition}[{Tribu}]\phantomsection\label{I-deftribu}

Une \emph{tribu} (ou \emph{\(\sigma\)-algèbre}) \(\mathcal{A}\) est une
collection de sous-ensembles de \(\Omega\) tels que~:

\begin{enumerate}
\def\labelenumi{\arabic{enumi}.}
\tightlist
\item
  \(\Omega \in \mathcal{A}\),
\item
  \(A \in \mathcal{A}\Rightarrow A^c \in \mathcal{A}\),
\item
  \(\forall\, n \in \mathbb{N}, A_n \in \mathcal{A}\Rightarrow \bigcup_n A_n \in \mathcal{A}\).
\end{enumerate}
\end{definition}

Le couple \((\Omega, \mathcal{A})\) est appelé \emph{espace
probabilisable}. On notera qu'il s'agit d'un espace mesurable dans la
terminologie du cours d'intégration.

\begin{exemple}\phantomsection\label{I-exemples-3}
\begin{enumerate}
\def\labelenumi{\arabic{enumi}.}
\tightlist
\item
  \(\mathcal{A}= \{\varnothing,\Omega\}\) est la tribu \emph{grossière}
  ou \emph{triviale}~: c'est la plus petite tribu de \(\Omega\).
\item
  Dans le cas où \(\Omega\) est au plus dénombrable, on choisit
  systématiquement l'ensemble \(\mathcal{P}(\Omega)\) des parties de
  \(\Omega\) dont on vérifie aisément qu'il s'agit d'une tribu\,; c'est
  le cadre étudié en CPGE. On verra plus loin que cette tribu est trop
  grande dans le cas où \(\Omega\) est infini non dénombrable.
\item
  Si \(\Omega = \mathbb{R}\), on peut le munir de la tribu formée des
  ensembles mesurables de \(\mathbb{R}\), dite \emph{tribu de Lebesgue}.
\item
  On verra plus loin qu'il est aisé de définir une tribu sur tout espace
  topologique.
\end{enumerate}
\end{exemple}

\section{Probabilité}\label{I-probabilituxe9}

\addcontentsline{toc}{subsection}{Probabilité}
\begin{definition}[{Probabilité}]\phantomsection\label{I-defproba}

Une \emph{probabilité} sur l'espace \((\Omega, \mathcal{A})\) est une
application \(\mathbb{P}: \mathcal{A}\rightarrow [0,1]\), telle que~:

\begin{enumerate}
\def\labelenumi{\arabic{enumi}.}
\tightlist
\item
  \(\mathbb{P}(\Omega) = 1\),
\item
  pour toute suite (dénombrable) \((A_n)_{n\in\mathbb{N}^\star}\)
  d'éléments de \(\mathcal{A}\) \textbf{deux à deux disjoints}, on a
  \begin{equation*}
   \mathbb{P}\left(\bigcup_{n\in\mathbb{N}^\star} A_n\right) = \sum_{n\in\mathbb{N}^\star} \mathbb{P}(A_n).
   \end{equation*}
\end{enumerate}
\end{definition}

Le triplet \((\Omega, \mathcal{A}, \mathbb{P})\) est appelé \emph{espace
probabilisé}. La modélisation probabiliste consiste ainsi à décrire une
expérience aléatoire par la donnée d'un espace probabilisé.

On notera que dans la terminologie du cours d'intégration,
\((\Omega, \mathcal{A}, \mathbb{P})\) est un espace mesuré. On vérifie
aisément que \(\mathbb{P}\) est une mesure qui a la particularité d'être
finie, c'est-à-dire telle que \(\mathbb{P}(\Omega) =1\). On parle aussi
de \emph{mesure de probabilité}.

La définition suivante est fondamentale en théorie des probabilités.
Elle introduit une notion de ``vrai ou faux'' qui dépend de la
probabilité choisie sur l'espace fondamental.

\addcontentsline{toc}{subsection}{Propriété presque-sûre}
\begin{definition}[{Propriété
presque-sûre}]\phantomsection\label{I-propriuxe9tuxe9-presque-suxfbre}

Soit \((\Omega, \mathcal{A}, \mathbb{P})\) un espace probabilisé. On dit
qu'un événement \(A\in\mathcal{A}\) se réalise
\emph{\(\mathbb{P}\)-presque sûrement} (en abrégé \(\mathbb{P}\)-p.s.)
si \(\mathbb{P}(A) = 1\).
\end{definition}

\addcontentsline{toc}{subsection}{Propriétés élémentaires}
\begin{proposition}[{Propriétés
élémentaires}]\phantomsection\label{I-elemprop}

\begin{enumerate}
\def\labelenumi{\arabic{enumi}.}
\tightlist
\item
  \(\forall\, A \in \mathcal{A}\), \(\mathbb{P}(A) \in [0,1]\) et
  \(\mathbb{P}(A^c)= 1-\mathbb{P}(A)\).
\item
  \(\forall\, A,B \in \mathcal{A}\),
  \(A \subset B \Rightarrow \mathbb{P}(A) \leq \mathbb{P}(B)\).
\item
  \(\forall\, A,B \in \mathcal{A}\),
  \(\mathbb{P}(A \cup B ) = \mathbb{P}(A) + \mathbb{P}(B) - \mathbb{P}(A \cap B)\).
\item
  Inégalité de Boole~: \(\forall\, n \in \mathbb{N}^\ast\),
  \(\forall\, (A_i)_{1 \leq i \leq n} \in \mathcal{A}\),
  \[\mathbb{P}\left(\bigcup_{i=1}^n A_i\right) \leq \sum_{i=1}^n \mathbb{P}(A_i).\]
\item
  Formule de Poincaré~: \(\forall\, n \in \mathbb{N}^\ast\),
  \(\forall\, (A_i)_{1 \leq i \leq n} \in \mathcal{A}\)
  \[ \mathbb{P}\left(\bigcup_{i=1}^n A_i\right) = \sum_{i=1}^n \mathbb{P}(A_i) - \sum_{1 \leq i < j \leq n} \mathbb{P}(A_i \cap A_j) + \ldots + (-1)^n \mathbb{P}\left(\bigcap_{i=1}^n A_i\right).\]
\end{enumerate}
\end{proposition}

\runinhead{Exercice -- Démonstration (\(\mathord{\bullet}\))}\phantomsection\addcontentsline{toc}{subsubsection}{Démonstration}\label{I-propelem}
\solref{I-answer-propelem}

\addcontentsline{toc}{subsection}{Théorème de la continuité monotone}
\begin{theoreme}[{Théorème de la continuité
monotone}]\phantomsection\label{I-continuitemonotone}

Dans le cas d'une suite \((A_n)_{n\in\mathbb{N}^\ast}\) d'éléments de
\(\mathcal{A}\) croissante, on a
\[ \mathbb{P}\left(\bigcup_{n\in\mathbb{N}^\ast} A_n\right) = \lim_{n \rightarrow \infty} \mathbb{P}(A_n).\]
\end{theoreme}

\runinhead{Exercice -- Démonstration (\(\mathord{\bullet}\))}\phantomsection\addcontentsline{toc}{subsubsection}{Démonstration}\label{I-contmon}
\solref{I-answer-contmon}

\runinhead{Exercice -- Une définition alternative de la
probabilité (\(\mathord{\bullet}\))}\phantomsection\addcontentsline{toc}{subsubsection}{Une
définition alternative de la probabilité}\label{I-altdef}
Soit \((\Omega, \mathcal{A})\) un espace probabilisable. Supposons que
\(\mathbb{P}: \mathcal{A}\to [0,1]\) vérifie~:

\begin{enumerate}
\def\labelenumi{\arabic{enumi}.}
\tightlist
\item
  \(\mathbb{P}(\Omega) = 1\),
\item
  Pour \(A, B \in \mathcal{A}\), tels que \(A\cap B = \varnothing\),
  \(\mathbb{P}(A\cup B)= \mathbb{P}(A) + \mathbb{P}(B)\) (additivité),
\item
  Pour toute suite \((A_n)_{n\in\mathbb{N}^\ast}\) d'éléments de
  \(\mathcal{A}\) croissante
  \[\mathbb{P}\left(\bigcup_{n\in\mathbb{N}^\ast} A_n\right) = \lim_{n \rightarrow \infty} \mathbb{P}(A_n).\]
\end{enumerate}

Montrer que \(\mathbb{P}\) vérifie la propriété de
\protect\hyperref[I-defproba]{\(\sigma\)-additivité (p.~\pageref*{I-defproba})}.\solref{I-answer-altdef}

\addcontentsline{toc}{subsection}{Suite décroissante}
\begin{remarque}[{Suite
décroissante}]\phantomsection\label{I-suite-duxe9croissante}

Dans le cas d'une suite décroissante, on a
\[ \mathbb{P}\left(\bigcap_{n \in \mathbb{N}^\ast} A_n\right) = \lim_{n \rightarrow \infty} \mathbb{P}(A_n).\]

Le second point de la \protect\hyperref[I-defproba]{définition de la
probabilité (p.~\pageref*{I-defproba})} donne la probabilité de la
réunion \(\cup_n A_n\) en fonction des \(\mathbb{P}(A_n)\) lorsque les
événements sont deux à deux disjoints. Si ce n'est pas le cas, on a tout
de même la majoration suivante~:
\end{remarque}

\addcontentsline{toc}{subsection}{Une majoration bien utile}
\begin{proposition}[{Une majoration bien
utile}]\phantomsection\label{I-une-majoration-bien-utile}

Soit \(\mathbb{P}\) une probabilité et soit
\((A_n)_{n\in \mathbb{N}^\ast}\) une famille dénombrable d'événements.
On a alors
\[ \mathbb{P}(\cup_n A_n) \leq \sum_{n\in\mathbb{N}^\ast} \mathbb{P}(A_n) \]
\end{proposition}

\runinhead{Exercice -- Démonstration (\(\mathord{\bullet}\))}\phantomsection\addcontentsline{toc}{subsubsection}{Démonstration}\label{I-boolinf}
\solref{I-answer-boolinf}

\section{Probabilité
conditionnelle}\label{I-probabilituxe9-conditionnelle}

La construction d'un modèle probabiliste repose sur l'information connue
\textbf{a priori} sur l'expérience aléatoire. Ce modèle permet de
quantifier les probabilités de réalisation de certains résultats de
l'expérience. Il est fondamental de remarquer que si l'information
change, les probabilités de réalisation changent.

\addcontentsline{toc}{subsection}{Information a priori}
\begin{exemple}[{Information a
priori}]\phantomsection\label{I-information-a-priori}

On cherche pour un lancer de deux dés, la probabilité de l'événement
``la somme est supérieure ou égale à 10''. Elle vaut 1/6 sans
information supplémentaire, 1/2 si l'on sait que le résultat d'un des
dés est 6, 0 si l'on sait a priori que le résultat d'un des dés est 2.
Pour obtenir ces résultats, on a calculé dans chaque cas le rapport du
nombre de résultats favorables sur le nombre de cas possibles. Il est
ainsi indispensable de bien définir l'espace de probabilité lié à
l'expérience munie de l'information a priori. On remarque également que
l'information a priori a changé la valeur de la probabilité de
l'événement.

L'outil qui va nous permettre d'introduire de l'information est la
probabilité conditionnelle dont nous donnons ici la définition.
\end{exemple}

\addcontentsline{toc}{subsection}{Probabilité conditionnelle}
\begin{definition}[{Probabilité
conditionnelle}]\phantomsection\label{I-defprobacond}

Soient \((\Omega, \mathcal{A}, \mathbb{P})\) un espace probabilisé,
\(A, B \in \mathcal{A}\) tels que \(\mathbb{P}(B)>0\). La
\emph{probabilité conditionnelle} de \(A\) sachant \(B\) est le nombre
\begin{equation*}
\mathbb{P}(A|B) = \frac{\mathbb{P}(A\cap B)}{\mathbb{P}(B)}.
\end{equation*}
\end{definition}

Cela définit une probabilité comme le montre la proposition suivante.

\addcontentsline{toc}{subsection}{Conséquences}
\begin{proposition}[{Conséquences}]\phantomsection\label{I-consuxe9quences}

\begin{enumerate}
\def\labelenumi{\arabic{enumi}.}
\tightlist
\item
  Soient \((\Omega, \mathcal{A}, \mathbb{P})\) un espace probabilisé et
  \(B \in \mathcal{A}\) tel que \(\mathbb{P}(B)>0\). Alors l'application
  de \(\mathcal{A}\) dans \([0, 1]\) qui à \(A\) associe
  \(\mathbb{P}(A|B)\) définit une nouvelle probabilité sur \(\Omega\),
  appelée probabilité conditionnelle sachant \(B\).
\item
  Si \(\mathbb{P}(A) > 0\) et \(\mathbb{P}(B) > 0\) , nous avons
  \[\mathbb{P}(A|B)\, \mathbb{P}(B) = \mathbb{P}(A \cap B) = \mathbb{P}(B|A)\, \mathbb{P}(A).\]
\end{enumerate}
\end{proposition}

\begin{preuve}
Il est clair que \(0 \leq \mathbb{P}(A|B) \leq 1\). Par ailleurs, les
deux propriétés de la \protect\hyperref[I-defproba]{définition de la
probabilité (p.~\pageref*{I-defproba})} pour \(\mathbb{P}(\cdot|B)\)
proviennent des mêmes propriétés pour \(\mathbb{P}\) et des remarques
suivantes~: \(\Omega \cap B = B\), et
\((\bigcup_{n\in\mathbb{N}^\ast} A_n ) \cap B = \bigcup_{n\in\mathbb{N}^\ast} (A_n \cap B)\).
De plus, si \(A\) et \(C\) sont disjoints, il en est de même de
\(A \cap B\) et \(C \cap B\). L'assertion 2 est évidente, d'après la
définition de la \protect\hyperref[I-defprobacond]{Probabilité
conditionnelle (p.~\pageref*{I-defprobacond})}.\hfill$\blacksquare$
\end{preuve}

\addcontentsline{toc}{subsection}{Formule des probabilités totales}
\begin{proposition}[{Formule des probabilités
totales}]\phantomsection\label{I-formprobatot}

Soit \((B_n)_{n\in\mathbb{N}^\ast}\) une partition finie ou dénombrable
d'événements de \(\Omega\) (i.e.~telle que
\(\bigcup_{n\in\mathbb{N}^\ast} B_n = \Omega\) et les \(B_n\) sont
deux-à-deux disjoints), telle que \(\mathbb{P}(B_n ) > 0\) pour tout
\(n\in\mathbb{N}^\ast\). Pour tout \(A \in \mathcal{A}\), on a alors
\begin{equation*}
P(A) = \sum_{n\in\mathbb{N}^\ast} P(A \cap B_n ) = \sum_{n\in\mathbb{N}^\ast} P(A|B_n)\, P(B_n).
\end{equation*}
\end{proposition}

\begin{preuve}
Nous avons \(A = \bigcup_{n\in\mathbb{N}^\ast} (A\cap B_n)\). Par
hypothèse, les ensembles \((A\cap B_n)\) sont deux-à-deux disjoints et
de plus \(\mathbb{P}(A\cap B_n) = \mathbb{P}(A|B_n)\,\mathbb{P}(B_n)\).
Le résultat découle du deuxième point de la
\protect\hyperref[I-defproba]{définition de la
probabilité (p.~\pageref*{I-defproba})}.\hfill$\blacksquare$
\end{preuve}

\addcontentsline{toc}{subsection}{Formule de Bayes}
\begin{proposition}[{Formule de Bayes}]\phantomsection\label{I-bayes}

Selon les mêmes hypothèses que ci-dessus et si \(\mathbb{P}(A) > 0\), on
a

\begin{equation*}
\forall\, i \in \{1,\dots,n\},\  \mathbb{P}(B_i | A) = \dfrac{\mathbb{P}(A | B_i)\, \mathbb{P}(B_i)}{\sum_{n\in\mathbb{N}^\ast} \mathbb{P}(A | B_n)\, \mathbb{P}(B_n)}.
\end{equation*}
\end{proposition}

\begin{preuve}
Le dénominateur vaut \(\mathbb{P}(A)\) d'après la
\protect\hyperref[I-formprobatot]{Formule des probabilités
totales (p.~\pageref*{I-formprobatot})}. La définition de la
\protect\hyperref[I-defprobacond]{probabilité
conditionnelle (p.~\pageref*{I-defprobacond})} implique~:
\begin{equation*}
\mathbb{P}(B_i | A) = \dfrac{\mathbb{P}(A \cap B_i)}{\mathbb{P}(A)} = \dfrac{\mathbb{P}(A | B_i) \mathbb{P}(B_i)}{\mathbb{P}(A)}.
\end{equation*}\hfill$\blacksquare$
\end{preuve}

\begin{remarque}\phantomsection\label{I-remarque}
La formule de Bayes, simple conséquence des axiomes et de la définition
de la probabilité conditionnelle, tient une place à part dans le calcul
des probabilités en raison de son importance pratique considérable et
des controverses auxquelles son application pratique a donné lieu~: elle
est à la base de toute une branche de la statistique appelée statistique
bayésienne.
\end{remarque}

\runinhead{Exercice -- Test (\(\mathord{\bullet}\))}\phantomsection\addcontentsline{toc}{subsubsection}{Test}\label{I-test}
Un individu est tiré au hasard dans une population où l'on trouve une
proportion \(10^{-4}\) de séropositifs. On lui fait passer un test de
détection de la séropositivité. Par ailleurs, des expérimentations
antérieures ont permis de savoir que les probabilités d'avoir un
résultat positif lors de l'application du test si l'individu est
séropositif, ou s'il ne l'est pas, sont respectivement égales à 0,99
(c'est la sensibilité du test) et à 0,001 (0,999 = 1 - 0,001 est la
spécificité du test). Sachant que le test donne un résultat positif,
quelle est la probabilité pour que l'individu soit effectivement
séropositif\,?\solref{I-answer-test}

\section{Indépendance des
événements}\label{I-induxe9pendance-des-uxe9vuxe9nements}

La notion d'indépendance est absolument fondamentale en probabilités et
nous verrons par la suite toutes ses implications dans la modélisation
de l'aléatoire.

Intuitivement, deux événements \(A\) et \(B\) sont indépendants si le
fait de savoir que \(A\) est réalisé ne donne aucune information sur la
réalisation de \(B\) et réciproquement.

Si \(B\) est un événement de probabilité strictement positive, \(A\)
sera dit indépendant de \(B\) si
\[\mathbb{P}(A | B) = \frac{\mathbb{P}(A\cap B)}{\mathbb{P}(B)} = \mathbb{P}(A)\]
On remarque que cette formule se symétrise et la notion d'indépendance
se définit finalement comme suit.

\addcontentsline{toc}{subsection}{Indépendance de deux événements}
\begin{definition}[{Indépendance de deux
événements}]\phantomsection\label{I-induxe9pendance-de-deux-uxe9vuxe9nements}

Deux événements \(A\) et \(B\) sont \emph{indépendants} si et seulement
si \begin{equation*}
\mathbb{P}(A\cap B) = \mathbb{P}(A)\, \mathbb{P}(B).
\end{equation*}
\end{definition}

\begin{remarque}\phantomsection\label{I-remarques}
\begin{itemize}
\tightlist
\item
  La probabilité de voir \(A\) réalisé ne dépend pas de la réalisation
  de \(B\), et réciproquement.
\item
  Cette notion est liée au choix de la probabilité \(\mathbb{P}\) et
  n'est pas une notion ensembliste. Cela n'a en particulier rien à voir
  avec le fait que \(A\) et \(B\) soient disjoints ou non.
\item
  Si \(\mathbb{P}(A)>0\) et \(\mathbb{P}(B)>0\), alors \begin{equation*}
   \mathbb{P}(A\cap B) = \mathbb{P}(A)\, \mathbb{P}(B) \Leftrightarrow \mathbb{P}(A) = \mathbb{P}(A|B) \Leftrightarrow \mathbb{P}(B) = \mathbb{P}(B|A).
   \end{equation*}
\end{itemize}
\end{remarque}

\begin{proposition}\phantomsection\label{I-proposition}
Si les événements \(A\) et \(B\) sont indépendants, alors il en est de
même des couples \((A^c,B)\), \((A,B^c)\) et \((A^c,B^c)\).
\end{proposition}

\runinhead{Exercice -- Démonstration (\(\mathord{\bullet}\))}\phantomsection\addcontentsline{toc}{subsubsection}{Démonstration}\label{I-indep}
\solref{I-answer-indep}

\runinhead{Exercice -- Auto-indépendant\,? (\(\mathord{\bullet}\))}\phantomsection\addcontentsline{toc}{subsubsection}{Auto-indépendant\,?}\label{I-autoindep}
À quelle condition un événement \(A\) est-il indépendant de
lui-même\,?\solref{I-answer-autoindep}

\begin{exemple}\phantomsection\label{I-exemples-4}
\begin{enumerate}
\def\labelenumi{\arabic{enumi}.}
\tightlist
\item
  On lance 3 fois un dé. Si \(A_i\) est un événement qui ne dépend que
  du \(i^\text{ème}\) lancer, alors \(A_1\) , \(A_2\) , \(A_3\) sont
  indépendants.
\item
  On tire une carte au hasard dans un jeu de 52 cartes. Soit \(A\) =
  \{la carte est une dame\} et \(B\) = \{la carte est un coeur\}. Il est
  facile de voir que \(\mathbb{P}(A) = 4/52\) et
  \(\mathbb{P}(B) = 13/52\) et
  \(\mathbb{P}(A \cap B) = \mathbb{P}( \{\text{la carte est la dame de coeur}\} ) = 1/52 = \mathbb{P}(A) \mathbb{P}(B)\).
  Ainsi, les événements \(A\) et \(B\) sont indépendants pour la
  probabilité uniforme \(\mathbb{P}\).
\item
  On suppose maintenant que le jeu de cartes soit trafiqué. Soit
  \(\widetilde{\mathbb{P}}\) la nouvelle probabilité correspondant au
  tirage de cartes. On suppose également que
  \[ \widetilde{\mathbb{P}}(\{\text{As de trèfle}\} ) = \frac{1}{2}, \,\,\,\,\, \widetilde{\mathbb{P}}(\{\text{autre carte}\}) = \frac{1}{2} \frac{1}{51} = \frac{1}{102},\]
  alors
  \[\widetilde{\mathbb{P}}(A \cap B) = \frac{1}{102} \neq \widetilde{\mathbb{P}}(A)\widetilde{\mathbb{P}}(B) = \frac{2}{51}\frac{13}{102}.\]
  Les événements \(A\) et \(B\) ne sont pas indépendants sous la
  probabilité \(\widetilde{\mathbb{P}}\).
\end{enumerate}
\end{exemple}

\addcontentsline{toc}{subsection}{Indépendance de \(n\) événements}
\begin{definition}[{Indépendance de \(n\)
événements}]\phantomsection\label{I-induxe9pendance-de-n-uxe9vuxe9nements}

\(n\) événements \(A_1, \dots, A_n\) sont \emph{mutuellement
indépendants} ou simplement \emph{indépendants} si et seulement si pour
toute partie \(I \subset \lbrace 1,\dots, n\rbrace\) \begin{equation*}
\mathbb{P}(\cap_{i\in I} A_i) = \prod_{i\in I} \mathbb{P}(A_i).
\end{equation*}
\end{definition}

\begin{remarque}\phantomsection\label{I-remarque-1}
Attention, l'indépendance des événements deux à deux n'implique pas
l'indépendance mutuelle des événements.
\end{remarque}

\section{Réflexion sur le concept de
probabilité}\label{I-remarque-ruxe9flexion-sur-le-concept-de-probabilituxe9}

La théorie mathématiques des probabilités ne dit pas quelle loi de
probabilité choisir sur un espace \((\Omega,\mathcal{A})\) parmi
l'infinité de lois possibles. Ce problème qui concerne ceux qui veulent
appliquer le calcul des probabilités, renvoie à la nature ``physique''
du concept de probabilité qui formalise et quantifie le sentiment
d'incertitude vis-à-vis d'un événement. Ce problème d'ordre conceptuel
oppose deux écoles de pensée, la conception objectiviste et la
conception subjectiviste.

Pour les tenants du premier point de vue, la probabilité d'un événement
peut être déterminée de manière unique. Dans la vision dite classique,
héritée des jeux de hasard, \(\Omega\) est fini et on donne à chaque
événement élémentaire la même probabilité. Le calcul des probabilités se
résume alors à un problème de dénombrement et la probabilité d'un
événement est le rapport du nombre de cas favorables sur le nombre de
cas possibles. Dans le cas infini, la vision fréquentiste repose sur la
loi des grands nombres~: si on répète un grand nombre de fois
l'expérience, la proportion de fois où un événement sera réalisé va
converger vers la probabilité de cet événement. Dans ce cadre, il est
impossible de donner une valeur et même un sens à un événement non
répétable comme ``pleuvra-t-il demain\,?''. En outre, la répétition à
l'infini d'une même expérience étant physiquement irréalisable, la loi
des grands nombres étant un résultat qui suppose défini le concept de
probabilité, la vision fréquentiste est logiquement intenable.

Dans la conception subjectiviste, la probabilité objective d'un
événement n'existe pas et n'est donc pas une grandeur mesurable analogue
à la masse d'un corps, par exemple. C'est simplement une mesure
d'incertitude qui reflète un degré de croyance pouvant varier avec les
circonstances et l'observateur, donc subjective, la seule exigence étant
qu'elle satisfasse aux axiomes du calcul des probabilités. Des méthodes
ont alors été proposées pour passer d'un simple pré-ordre sur les
événements, à une probabilité. Puisque la répétition n'est plus
nécessaire, on peut probabiliser des événements non répétables et
étendre ainsi le domaine d'application du calcul des probabilités,
notamment pour orienter des prises de décisions. On notera que la
\protect\hyperref[I-bayes]{formule de Bayes (p.~\pageref*{I-bayes})}
permet d'intégrer facilement de l'information a priori, dans la mesure
où celle-ci est probabilisée.

On arrête ici ces quelques remarques sans prendre parti dans une
querelle qui dure encore. L'un ou l'autre point de vue sera adopté selon
les ouvrages rencontrés. Dans tous les cas, les outils mathématiques
développés dans ce cours seront adaptés. On rappelle tout de même que la
modélisation probabiliste a prouvé son efficacité dans de nombreuses
applications mais que, comme tout modèle, ce n'est qu'une représentation
simplificatrice de la réalité et que ses hypothèses doivent être mises à
l'épreuve des faits. À ce titre, on citera Georges Matheron qui dans son
essai sur la pratique des probabilités Estimer et Choisir (Matheron
(\protect\hyperlink{refI-matheron}{1978})) écrit fort justement~: ``Il
n'y a pas de probabilités en soi. Il n'y a que des modèles
probabilistes''.

\chapter{Variables aléatoires}\label{I-variables-aluxe9atoires-1}

En théorie moderne des probabilités, on préfère prendre un point de vue
fonctionnel plutôt qu'ensembliste, et utiliser les variables aléatoires
plutôt que les événements. Une variable aléatoire est une grandeur qui
dépend du résultat de l'expérience. Par exemple,

\begin{itemize}
\tightlist
\item
  le nombre de 6 obtenus dans un lancer de 3 dés,
\item
  le nombre d'appels dans un central téléphonique pendant une heure,
\item
  la distance du point d'atteinte d'une flèche au centre de la cible,
\item
  la valeur maximale d'un prix d'actif sur un intervalle de temps donné,
\end{itemize}

sont des variables aléatoires.

\addcontentsline{toc}{subsection}{Variable aléatoire}
\begin{definition}[{Variable
aléatoire}]\phantomsection\label{I-variable-aluxe9atoire}

Soient \((\Omega,\mathcal{A})\) et \((E,\mathcal{E})\) deux espaces
probabilisables (mesurables). Une \emph{variable aléatoire} \(X\) à
valeurs dans \(E\) est une application mesurable de
\((\Omega,\mathcal{A})\) dans \((E,\mathcal{E})\), c'est-à-dire telle
que \(\forall B \in \mathcal{E}\), \(X^{-1}(B) \in \mathcal{A}\).
\end{definition}

En pratique, l'ensemble \(E\) pourra être un ensemble fini ou
dénombrable ou \(\mathbb{R}\) ou \(\mathbb{R}^d\) ou encore un espace
plus sophistiqué tel que l'ensemble \(C(\mathbb{R}_+ , \mathbb{R}^d)\)
des fonctions continues de \(\mathbb{R}_+\) dans \(\mathbb{R}^d\).

\addcontentsline{toc}{subsection}{Vocabulaire}
\begin{remarque}[{Vocabulaire}]\phantomsection\label{I-vocabulaire}

La terminologie, consacrée par l'usage, peut être trompeuse. Une
variable aléatoire n'est pas une variable (au sens de l'analyse) mais
une fonction. Cette terminologie est apparentée à la notion de variable
en physique ou en sciences humaines où on désigne volontiers par
``variable'' la valeur prise par une fonction de l'état du système
étudié.

L'intérêt principal de travailler avec des variables aléatoires est de
pouvoir substituer à l'espace abstrait \(\Omega\) des résultats de
l'expérience l'espace \(E\), mieux connu dans la pratique. Ainsi, grâce
à une variable aléatoire \(X\), nous pouvons transporter la structure
abstraite du modèle probabiliste \((\Omega, \mathcal{A}, \mathbb{P})\)
sur l'espace d'arrivée \(E\), en posant pour \(B \in \mathcal{E}\)
\[\mathbb{P}_X (B) = \mathbb{P}(X^{-1}(B)) = \mathbb{P}(\{\omega, X(\omega)\in B\})\]

Cette formule définit une nouvelle probabilité, notée \(\mathbb{P}_X\)
et définie sur \((E,\mathcal{E})\), qui s'appelle la \emph{loi de la
variable} \(X\). On dit aussi que \(\mathbb{P}_X\) est la \emph{mesure
(de probabilité) image} de \(\mathbb{P}\) par \(X\).

Pour se convaincre que \(\mathbb{P}_X\) définit bien une probabilité sur
\((E,\mathcal{E})\), il suffit de remarquer que \begin{align*}
& X^{-1}(\varnothing) = \varnothing, X^{-1}(E) = \Omega, X^{-1}(B^c) = X^{-1}(B)^c \\
& X^{-1}(\cap_i A_i) = \cap_i X^{-1}(A_i), X^{-1}(\cup_i A_i) = \cup_i X^{-1}(A_i)
\end{align*}
\end{remarque}

\addcontentsline{toc}{subsection}{Notation}
\begin{remarque}[{Notation}]\phantomsection\label{I-notation}

Il est usuel de noter l'ensemble
\(X^{-1}(B) = \{\omega \in \Omega, X(\omega) \in B\}\) par
\(\{X \in B\}\), ce qui allège les écritures. On se rappelera néanmoins
que cette notation désigne un sous-ensemble de \(\Omega\).
\end{remarque}

\(\mathbb{P}_X\) sera plus facile à caractériser que \(\mathbb{P}\)
puisque \(E\) est un ensemble connu (on pourra en particulier utiliser
ses propriétés topologiques) alors que \(\Omega\) est un espace
abstrait. Les variables que nous rencontrerons dans ce cours seront soit
à valeurs dans un ensemble dénombrable, soit à valeurs dans
\(\mathbb{R}\) ou dans \(\mathbb{R}^d\). Nous les appellerons
respectivement des variables aléatoires discrètes, réelles ou des
vecteurs aléatoires. Leurs lois seront alors des probabilités
respectivement sur un ensemble dénombrable, sur \(\mathbb{R}\) ou sur
\(\mathbb{R}^d\). Le cas discret est considéré connu.

\section{Variables aléatoires
réelles}\label{I-variables-aluxe9atoires-ruxe9elles}

Dans le cas où \(E = \mathbb{R}\), on travaille généralement avec la
\emph{tribu des boréliens}, que l'on notera \(\mathcal{B}(\mathbb{R})\),
qui est très utile comme on le verra pour identifier les lois de
probabilité. On rappelle ici quelques propriétés sur les tribus, déjà
vues dans le cours d'intégration.

\addcontentsline{toc}{subsection}{Tribu engendrée}
\begin{definition}[{Tribu
engendrée}]\phantomsection\label{I-tribu-engendruxe9e}

Si \(C \subset \mathcal{P}(\Omega)\), on appelle tribu engendrée par
\(C\) la plus petite tribu contenant \(C\). Elle existe toujours, car
d'une part \(\mathcal{P}(\Omega)\) est une tribu contenant \(C\), et
d'autre part l'intersection d'une famille quelconque de tribus est une
tribu. Ainsi, la tribu engendrée par \(C\) est l'intersection de toutes
les tribus contenant \(C\).
\end{definition}

\addcontentsline{toc}{subsection}{Tribus engendrées}
\begin{exemple}[{Tribus
engendrées}]\phantomsection\label{I-tribus-engendruxe9es}

\begin{itemize}
\tightlist
\item
  La tribu engendrée par un ensemble \(A \subset \Omega\) est
  \({\varnothing, A, A^c , \Omega}\).
\item
  si \((A_i )_{i \in I}\) est une partition finie ou dénombrable de
  \(\Omega\) (i.e.~les \(A_i\) sont deux-à-deux disjoints et leur
  réunion est \(\Omega\)), la tribu engendrée par \(\{A_i , i \in I\}\)
  est l'ensemble des réunions \(B_J = \cup_{i \in J} A_i\), où \(J\)
  décrit la classe de toutes les parties de I.
\end{itemize}
\end{exemple}

\addcontentsline{toc}{subsection}{Tribu borélienne sur \(\mathbb{R}\)}
\begin{definition}[{Tribu borélienne sur
\(\mathbb{R}\)}]\phantomsection\label{I-bor}

Si \(\Omega = \mathbb{R}\), on appelle tribu borélienne, que l'on note
\(\mathcal{B}(\mathbb{R})\), la tribu engendrée par la classe des
ouverts de \(\mathbb{R}\).
\end{definition}

À titre d'exercice de maniement des tribus, on donne en détail la
démonstration du résultat suivant~:

\addcontentsline{toc}{subsection}{Définition alternative}
\begin{proposition}[{Définition
alternative}]\phantomsection\label{I-altbor}

La tribu borélienne de \(\mathbb{R}\) est la tribu engendrée par les
intervalles de la forme \(] - \infty, a]\) pour \(a \in \mathbb{Q}\).
\end{proposition}

\begin{preuve}
On rappelle que toute tribu est stable par passage au complémentaire,
par réunion ou intersection dénombrable.

Soit \(C_1\) la classe des intervalles ouverts de \(\mathbb{R}\),
\(C_2\) la classe des intervalles \(] - \infty, a]\) pour
\(a \in \mathbb{Q}\) et \(\mathcal{B}(\mathbb{R})\) la tribu borélienne.

Tout ouvert \(A\) est réunion dénombrable d'intervalles ouverts, en
effet, on a \(A = \cup_{(q,n)\in B} ]q -\frac{1}{n}, q -\frac{1}{n}[\),
où \(B\) est l'ensemble (dénombrable) des couples \((q,n)\) avec
\(q \in \mathbb{Q}\), \(n \in \mathbb{N}\) et
\(]q -\frac{1}{n}, q -\frac{1}{n}[ \in A\). La tribu engendrée par les
intervalles ouverts est donc identique à la tribu borélienne par la
stabilité par réunion dénombrable et le fait que
\(C_1 \subset \mathcal{B}(\mathbb{R})\).

Réciproquement, soit \(]x, y[\) un intervalle ouvert de \(\mathbb{R}\).
Soit \((x_n)_n\) une suite de rationnels décroissant vers \(x\) et
\((y_n)_n\) une suite de rationnels croissant strictement vers y. On a~:
\[ ]x, y[= \bigcup_n (] - \infty, y_n ]\bigcap ] - \infty, x_n ]^c ).\]
Donc \(C_1\) et par conséquent \(\mathcal{B}(\mathbb{R})\) sont inclus
dans la tribu engendrée par \(C_2\). Finalement, comme tout fermé est
dans \(\mathcal{B}(\mathbb{R})\) par complémentarité, on a aussi
\(C_2 \subset \mathcal{B}(\mathbb{R})\), d'où le
résultat.\hfill$\blacksquare$
\end{preuve}

On peut maintenant définir une variable aléatoire réelle~:

\addcontentsline{toc}{subsection}{Variable aléatoire réelle}
\begin{definition}[{Variable aléatoire
réelle}]\phantomsection\label{I-defvar}

Soit l'espace d'état \(\Omega\) muni de la tribu \(\mathcal{A}\) des
événements. Une application \(X\) de \(\Omega\) dans \(\mathbb{R}\) est
une \emph{variable aléatoire réelle} si \(X^{-1}(B) \in \mathcal{A}\)
pour tout \(B \in \mathcal{B}(\mathbb{R})\), autrement dit si \(X\) est
une application mesurable de \((\Omega, \mathcal{A})\) dans
\((\mathbb{R},\mathcal{B}(\mathbb{R}))\).
\end{definition}

\addcontentsline{toc}{subsection}{Loi d'une variable aléatoire réelle}
\begin{definition}[{Loi d'une variable aléatoire
réelle}]\phantomsection\label{I-defloivar}

La probabilité \(\mathbb{P}_X\), définie sur
\((\mathbb{R},\mathcal{B}(\mathbb{R}))\) par
\(\mathbb{P}_X (B) = \mathbb{P}(X^{-1}(B))\) pour
\(B \in \mathcal{B}(\mathbb{R})\) est appelée \emph{loi de la variable
\(X\)}, ou \emph{distribution} de \(X\).
\end{definition}

On peut voir \(\mathbb{P}_X\) comme une transposition de \(\mathbb{P}\)
sur \(\mathbb{R}\). On a alors le résultat très utile suivant~:

\addcontentsline{toc}{subsection}{Composition}
\begin{proposition}[{Composition}]\phantomsection\label{I-composition}

Si \(X_1, \ldots, X_n\) sont des variables aléatoires réelles et si
\(g\) est une fonction mesurable de \(\mathbb{R}^n\) dans \(\mathbb{R}\)
(munis de leur tribu borélienne), alors \(Y = g(X_1,\ldots,X_n)\) est
une variable aléatoire réelle.
\end{proposition}

\begin{preuve}
On note \(X = (X_1, \ldots, X_n)\). \(X\) est mesurable de
\((\Omega,\mathcal{A})\) dans
\((\mathbb{R}^n,\mathcal{B}(\mathbb{R}^n))\) puisque ce vecteur est
constitué de composantes mesurables de \((\Omega,\mathcal{A})\) dans
\((\mathbb{R},\mathcal{B}(\mathbb{R}))\) et que
\(\mathcal{B}(\mathbb{R}^n) = \mathcal{B}(\mathbb{R})\otimes \ldots \otimes \mathcal{B}(\mathbb{R})\).
Soit \(B \in \mathcal{B}(\mathbb{R})\), alors
\(Y^{-1}(B) = (g \circ X)^{-1}(B) = X^{-1}(g^{-1}(B))\). Puisque \(g\)
et \(X\) sont mesurables, \(g^{-1}(B) \in \mathcal{B}(\mathbb{R}^n)\) et
\(X^{-1}(g^{-1}(B))\in \mathcal{A}\).\hfill$\blacksquare$
\end{preuve}

Comme application de ce résultat, on a les propriétés suivantes~:

\addcontentsline{toc}{subsection}{Conséquences}
\begin{proposition}[{Conséquences}]\phantomsection\label{I-consuxe9quences-1}

Soient \(X\), \(Y\) et \((X_n)_{n \in \mathbb{N}^\star}\) des variables
aléatoires réelles. On a

\begin{enumerate}
\def\labelenumi{\arabic{enumi}.}
\item
  \(X + Y\), \(XY\), \(\frac{X}{Y}\) si \(Y \neq 0\), sont des variables
  aléatoires.
\item
  \(\sup_{1\leq p \leq n} X_p\), \(\inf_{1\leq p \leq n} X_p\), sont des
  variables aléatoires.
\item
  \(\sup_{n\geq 1} X_n\), \(\inf_{n\geq 1} X_n\), sont des variables
  aléatoires.
\item
  Si \(X_n(\omega) \xrightarrow[n \to \infty]{} Z(\omega)\),
  \(\forall \omega\), alors la limite \(Z\) est une variable aléatoire.
\item
  \(Z = 1_A\) est une variable aléatoire \(\Leftrightarrow\)
  \(A \in \mathcal{A}\).
\end{enumerate}
\end{proposition}

\section{Loi des variables aléatoires
réelles}\label{I-loi-des-variables-aluxe9atoires-ruxe9elles}

Nous avons vu précédemment la définition générale d'une probabilité
\(\mathbb{P}\) sur un espace quelconque \(\Omega\) muni d'une tribu
\(\mathcal{A}\). Un problème fondamental est de construire et de
caractériser ces probabilités. La résolution de ce problème lorsque
\(\Omega\) est fini ou dénombrable est supposée connue. Puisque les
variables aléatoires réelles sont définies sur \(\mathbb{R}\), nous
allons étudier leur loi de probabilité \(\mathbb{P}_X\) définie sur
\(\mathbb{R}\) muni de la tribu \(\mathcal{B}(\mathbb{R})\).

\addcontentsline{toc}{subsection}{Fonction de répartition}
\begin{definition}[{Fonction de
répartition}]\phantomsection\label{I-deffdr}

Soit \(X\) une variable aléatoire réelle et \(\mathbb{P}_X\) sa loi. La
\emph{fonction de répartition} de \(X\) est la fonction
\begin{equation*}
F_X(x) = \mathbb{P}_X(\left]-\infty, x\right]) = \mathbb{P}(X \leq x),\ x \in \mathbb{R}.
\end{equation*}
\end{definition}

\addcontentsline{toc}{subsection}{La fonction de répartition caractérise
la probabilité}
\begin{theoreme}[{La fonction de répartition caractérise la
probabilité}]\phantomsection\label{I-carac}

La fonction de répartition \(F\) caractérise la probabilité
\(\mathbb{P}_X\) sur (\(\mathbb{R},\mathcal{B}(\mathbb{R})\)).
\end{theoreme}

\begin{preuve}
On veut montrer que la connaissance de \(F\) telle que définie
\protect\hyperref[I-deffdr]{ci-dessus (p.~\pageref*{I-deffdr})}
détermine \(\mathbb{P}_X\) de manière unique, c'est-à-dire que si il
existe une autre probabilité \(\mathbb{P}_Y\) telle que
\(G(x) = \mathbb{P}_Y(]-\infty,x])\) pour \(x \in \mathbb{R}\) et si
\(F = G\), alors \(\mathbb{P}_X = \mathbb{P}_Y\).

Soit \(\mathcal{B}_0\) l'ensemble de toutes les réunions finies
d'intervalles disjoints de la forme \(]x,y]\) ou
\(-\infty \leq x < y \leq +\infty\). Cet ensemble contient
\(\mathbb{R}\), \(\varnothing\) et est stable par passage au
complémentaire et par réunion finie, c'est donc une algèbre. De manière
évidente, il est aussi stable par intersection finie et est ainsi un
\(\pi\)-système. Par ailleurs, on a déjà vu que
\(\sigma(\mathcal{B}_0)\) (la tribu engendrée par \(\mathcal{B}_0\)) est
la tribu borélienne \(\mathcal{B}(\mathbb{R})\).

D'après \protect\hyperref[I-deffdr]{la définition de la fonction de
répartition (p.~\pageref*{I-deffdr})}, on a
\(\mathbb{P}_X(]x,y]) = F(y)-F(x)\) pour tous \(x< y\). Par conséquent,
si \(B = \cup_{i=1}^n ]x_i,y_i]\), avec \(x_i < y_i < x_{i+1}\), on a
\[\mathbb{P}_X(B) = \sum_{i=1}^n \mathbb{P}_X(]x_i,y_i]) = \sum_{i=1}^n F(y_i)-F(x_i),\]
car les intervalles sont disjoints.

Si \(\mathbb{P}_Y\) est une autre probabilité telle que
\(F(x) = \mathbb{P}_Y(]-\infty,x])\), alors ce qui précède montre que
\(\mathbb{P}_X = \mathbb{P}_Y\) sur \(\mathcal{B}_0\) qui engendre
\(\mathcal{B}(\mathbb{R})\). \(\mathbb{P}_X\) et \(\mathbb{P}_Y\) étant
des mesures finies, la proposition C.3.12 du cours de calcul intégral
nous permet d'assurer que \(\mathbb{P}_X = \mathbb{P}_Y\) pour tous les
éléments de \(\mathcal{B}(\mathbb{R})\), autrement dit que ces deux
mesures de probabilités sont identiques sur
\((\mathbb{R},\mathcal{B}(\mathbb{R}))\).\hfill$\blacksquare$
\end{preuve}

\addcontentsline{toc}{subsection}{Caractérisation de la fonction de
répartition}
\begin{theoreme}[{Caractérisation de la fonction de
répartition}]\phantomsection\label{I-theofdr}

Une fonction \(F\) est la fonction de répartition d'une unique
probabilité \(\mathbb{P}_X\) sur
\((\mathbb{R},\mathcal{B}(\mathbb{R}))\) si et seulement si elle vérifie
les trois conditions suivantes~:

\begin{enumerate}
\def\labelenumi{\arabic{enumi}.}
\tightlist
\item
  elle est croissante,
\item
  elle est continue à droite,
\item
  \(\lim\limits_{x \to -\infty} F(x) = 0, \lim\limits_{x \to +\infty} F(x) = 1\).
\end{enumerate}
\end{theoreme}

\begin{preuve}
La première assertion est immédiate d'après la
\protect\hyperref[I-deffdr]{définition (p.~\pageref*{I-deffdr})}. Pour
la seconde, on remarque que si \(x_n\) décroît vers \(x\), alors
\(]-\infty,x_n]\) décroît vers \(]-\infty,x]\) et donc \(F(x_n)\)
décroît vers \(F(x)\) par le
\protect\hyperref[I-continuitemonotone]{théorème de la continuité
monotone (p.~\pageref*{I-continuitemonotone})}. La troisième assertion
se montre de manière analogue en remarquant que \(]-\infty,x]\) décroît
vers \(\varnothing\) (resp. croît vers \(\mathbb{R}\)) lorsque \(x\)
décroît vers \(-\infty\) (resp. croît vers \(+\infty\)).

Pour la réciproque, on se reportera au chapitre 7 de Jacod and Protter
(\protect\hyperlink{refI-Jacod}{2003}).\hfill$\blacksquare$
\end{preuve}

\addcontentsline{toc}{subsection}{Quelques propriétés}
\begin{remarque}[{Quelques
propriétés}]\phantomsection\label{I-quelques-propriuxe9tuxe9s}

Comme \(F\) est croissante, elle admet une limite à gauche en chaque
point notée \(F(x^-)\). En remarquant que
\(]-\infty,y[\, = \lim\limits_{n \to +\infty}]-\infty,y_n[\) si \(y_n\)
tend vers \(y\) par valeurs décroissantes, on obtient pour \(x < y\)~:

\begin{itemize}
\tightlist
\item
  \(\mathbb{P}_X(]x,y]) = \mathbb{P}( x < X \leq y) = F(y) - F(x)\)
\item
  \(\mathbb{P}_X(]x,y[) = \mathbb{P}( x < X < y) = F(y-) - F(x)\)
\item
  \(\mathbb{P}_X([x,y]) = \mathbb{P}( x \leq X \leq y) = F(y) - F(x^-)\)
\item
  \(\mathbb{P}_X([x,y[) = \mathbb{P}( x \leq X < y) = F(y-) - F(x^-)\)
\end{itemize}

En particulier, \(\mathbb{P}_X(\{x\}) = F(x) - F(x^-)\) est le
\textbf{saut} de la fonction \(F\) au point \(x\). On a donc
\(\mathbb{P}_X(\{x\}) = 0\) pour tout \(x\) si et seulement si \(F\) est
continue en tout point.
\end{remarque}

\addcontentsline{toc}{subsection}{Nécessité des tribus}
\begin{remarque}[{Nécessité des
tribus}]\phantomsection\label{I-nuxe9cessituxe9-des-tribus}

Le \protect\hyperref[I-carac]{théorème
ci-dessus (p.~\pageref*{I-carac})} explique pourquoi, d'un point de vue
strictement mathématique, il est nécessaire d'introduire les tribus en
probabilités, malgré la complexité que cela engendre.

Plus concrètement, considérons l'exemple suivant~: soit
\(\Omega = [0,1]\) et \(\mathbb{P}\) telle que
\(\mathbb{P}(]a,b]) = b-a\) pour \(0\leq a\leq b\leq 1\) (il s'agit de
la loi uniforme sur {[}0,1{]}). C'est une probabilité naturelle qui
assigne à tout intervalle sa longueur comme probabilité. Supposons
maintenant que l'on souhaite étendre de manière unique \(\mathbb{P}\)
aux \(2^{[0,1]}\) éléments de \(\mathcal{P}([0,1])\) de manière à ce que
\(\mathbb{P}(\Omega) =1\) et
\(\mathbb{P}\left(\cup_{n\in\mathbb{N}^\star} A_n\right) = \sum_{n\in\mathbb{N}^\star} \mathbb{P}(A_n)\)
pour toute suite \((A_n)_{n\in\mathbb{N}^\star}\) tels que
\(\mathcal{A}_n\cap A_m = \varnothing\) pour \(n \neq m\). On peut
prouver qu'un tel \(\mathbb{P}\) n'existe pas. \(\mathcal{P}([0,1])\)
est trop ``grand'' pour définir un tel \(\mathbb{P}\). Il contient en
particulier des ensembles non mesurables.

Si l'on voulait travailler avec la tribu
\(\mathcal{A}= \mathcal{P}(\mathbb{R})\), il n'existerait que très peu
de probabilités sur \(\mathbb{R}\), à savoir les probabilités discrètes
que l'on décrit rapidement ci-dessous.
\end{remarque}

%\addcontentsline{toc}{subsection}{Probabilités discrètes}
\begin{exemple}[{Probabilités
discrètes}]\phantomsection\label{I-ex.discret}

\begin{enumerate}
\def\labelenumi{\arabic{enumi}.}
\item
  Les masses de Dirac (ou \textbf{mesures} de Dirac).

  \(X\) est identiquement égale à \(a \in \mathbb{R}\). Alors sa loi est
  la mesure de Dirac en \(a\), la probabilité \(\mathbb{P}_X\) sur
  \(\mathbb{R}\) qui vérifie pour \(A \in \mathcal{P}(\mathbb{R})\)
  \begin{equation*}
      \mathbb{P}_X(A) = \left\{ \begin{array}{ll}
      1  &\text{si } a \in A, \\
      0 &\text{sinon.}
      \end{array}
      \right.
  \end{equation*} Sa fonction de répartition est
  \(F(x) = 1_{[a,+\infty[}(x)\).
\item
  Les probabilités portées par \(\mathbb{N}\).

  \(X\) est à valeurs dans \(\mathbb{N}\). Comme \(\mathbb{N}\) est une
  partie de \(\mathbb{R}\), toute probabilité sur \(\mathbb{N}\) peut
  être considérée comme une probabilité sur \(\mathbb{R}\) qui ne
  ``charge'' que \(\mathbb{N}\). Plus précisément, si \(Q\) est la loi
  de probabilité de \(X\) sur \(\mathbb{N}\), on définit son
  ``extension'' \(\mathbb{P}_X\) à \(\mathbb{R}\) en posant
  \(\mathbb{P}_X(A) = Q(A\cap \mathbb{N})\). Si \(q_n = Q(\{n\})\) pour
  \(n \in \mathbb{N}\), la fonction de répartition \(F\) de
  \(\mathbb{P}_X\) est \begin{equation*}
  F(x) = \left\{ \begin{array}{ll}
  0  &\text{si } x <0, \\
  \sum_{i=0}^{\lfloor x \rfloor} q_i &\text{sinon,}
  \end{array}
  \right.
  \end{equation*} où \(\lfloor \cdot \rfloor\) désigne la partie
  entière. A titre d'exemple, on représente ci-dessous les fonctions de
  répartitions de la loi binomiale et de la loi de Poisson.
  \begin{center}
    \includegraphics{CdfBinom.tex.pdf}
  \end{center}
  \begin{center}
    \includegraphics{CdfPoisson.tex.pdf}
  \end{center}  
    
\item
  Les probabilités discrètes.

  Plus généralement, si \(X\) prend ses valeurs dans une partie finie ou
  dénombrable \(E\) de \(\mathbb{R}\), la loi \(\mathbb{P}_X\) de \(X\)
  est caractérisée pour tout \(i \in E\) par
  \(q_i = \mathbb{P}_X({i}) = \mathbb{P}(X=i)\). La fonction de
  répartition \(F\) de \(X\) est alors
  \[F(x) = \sum_{\substack{i \in E \\ i \leq x}} q_i,\] avec la
  convention qu'une somme ``vide'' vaut 0. On retrouve bien l'exemple 2
  si \(E = \mathbb{N}\). On voit que \(F\) est \textbf{purement
  discontinue} au sens où elle est complètement caractérisée par ses
  sauts \(\triangle F(x) = F(x) - F(x^-)\)~:
  \[F(x) = \sum_{y\leq x} \triangle F(y).\]
\end{enumerate}

Il existe bien d'autres probabilités, non discrètes, sur \(\mathbb{R}\).
Le paragraphe suivant est consacré à un exemple très important, celui
des variables aléatoires dont la loi admet une densité.
\end{exemple}

\begin{figure}
\centering
\includegraphics{CdfBinom.tex.pdf}
\caption{fonction de répartition de la loi
binomiale}\label{I-fig-CdfBinom.tex}
\end{figure}

\begin{figure}
\centering
\includegraphics{CdfPoisson.tex.pdf}
\caption{fonction de répartition de la loi de
Poisson}\label{I-fig-CdfPoisson.tex}
\end{figure}

\section{Variables aléatoires réelles à
densité}\label{I-variables-aluxe9atoires-ruxe9elles-uxe0-densituxe9}

\addcontentsline{toc}{subsection}{Densité de probabilité}
\begin{definition}[{Densité de
probabilité}]\phantomsection\label{I-defdens}

Une fonction réelle \(f\) sur \(\mathbb{R}\) est une \emph{densité de
probabilité} (ou plus simplement une \emph{densité}) si elle est
positive, intégrable et vérifie \[\int_\mathbb{R}f(x)\, dx = 1,\]

où \(dx\) désigne la mesure de Lebesgue sur \(\mathbb{R}\).

Si \(f\) est une densité, la fonction
\[F(x) =\int_{-\infty}^x f(y)\, dy\] est la fonction de répartition
d'une probabilité sur \(\mathbb{R}\).
\end{definition}

\addcontentsline{toc}{subsection}{Variable aléatoire réelle à densité}
\begin{definition}[{Variable aléatoire réelle à
densité}]\phantomsection\label{I-va.densituxe9}

Soit \(X\) une variable aléatoire. On dit que \(X\) a une \emph{loi de
densité \(f\)} (ou par abus de language ``est de densité \(f\)''), si
\(\mathbb{P}_X\) admet la densité \(f\) et donc si pour tout réel \(x\),
\[ \mathbb{P}(X\leq x) = \int_{-\infty}^x f(y) dy.\]
\end{definition}

\addcontentsline{toc}{subsection}{Fonction de répartition d'une variable
aléatoire à densité}
\begin{proposition}[{Fonction de répartition d'une variable aléatoire à
densité}]\phantomsection\label{I-fdr-va-dens}

Soit \(X\) de loi de densité \(f\).

\begin{enumerate}
\def\labelenumi{\arabic{enumi}.}
\tightlist
\item
  Sa fonction de répartition \(F\) est continue, de sorte que
  \[\mathbb{P}(X=x) = 0,~~~\forall x \in \mathbb{R}.\]
\item
  La fonction \(F\) est dérivable en tout point \(x\) où \(f\) est
  continue et \[F'(x) = f(x).\]
\item
  A l'inverse, si la fonction de répartition \(F\) de \(X\) est
  dérivable, ou seulement continue partout et dérivable par morceaux,
  alors \(X\) admet la densité \[F'(x) = f(x).\]
\end{enumerate}
\end{proposition}

Il existe bien sûr des variables aléatoires qui n'ont pas de densité~:
c'est le cas des variables discrètes données en exemple
\protect\hypertarget{I-ex.discret}{}{ci-dessus}. Il existe aussi des cas
``mixtes''~: soient d'une part \(f\) une fonction positive, intégrable
et d'intégrale strictement positive et d'autre part une partie finie ou
dénombrable \(E\) de \(\mathbb{R}\) et des poids \(p_i>0\) indexés par
\(i \in E\), tels que~:
\[ \int_\mathbb{R}f(x)\, dx + \sum_{i\in E} p_i = 1.\] Alors la fonction
\[ F(x) = \int_{-\infty}^x f(x)\, dx + \sum_{\substack{i\in E \\ i \leq x}} p_i\]
est une fonction de répartition, et la probabilité associée
\(\mathbb{P}_X\) n'admet pas de densité et n'est pas non plus discrète.

La mesure de probabilité associée s'écrit
\(\mathbb{P}_X = f\lambda + \sum_{i\in E} p_i \delta_i\) où \(\lambda\)
est la mesure de Lebesgue.

\begin{remarque}\phantomsection\label{I-remarques-1}
\begin{enumerate}
\def\labelenumi{\arabic{enumi}.}
\item
  La fonction de répartition est entièrement déterminée par la
  probabilité \(\mathbb{P}\). Il n'en est pas de même de la densité
  lorsqu'elle existe~: si en effet on a
  \(F(x) =\int_{-\infty}^x f(y)\, dy\) et si on pose \(g(x) = f(x)\) si
  \(x \notin E\) et \(g(x) =f(x)+1\) (par exemple) si \(x\in E\), où
  \(E\) est un ensemble négligeable, alors \(g\) est encore une densité
  de \(\mathbb{P}\). Autrement dit, la densité
\item
  Une interprétation intuitive de la densité \(f\) de \(\mathbb{P}\)~:
  si \(dx\) est un petit accroissement de la variable \(x\), on a (si du
  moins \(f\) est continue en \(x\))~:
  \[ f(x) \sim \frac{\mathbb{P}([x,x+dx])}{dx}.\]
\end{enumerate}
\end{remarque}

\runinhead{Exercice -- Durée de
fonctionnement (\(\mathord{\bullet}\))}\phantomsection\addcontentsline{toc}{subsubsection}{Durée
de fonctionnement}\label{I-ex.expo}
On suppose que la durée de fonctionnement, en heures, d'un ordinateur
avant sa première panne est une variable aléatoire positive de loi
exponentielle de paramètre 1/100, de densité donnée par

\[f(x) = \left\{ \begin{array}{ll}
        \frac{1}{100}\exp\left(-\frac{x}{100}\right) & x \geq 0, \\
        0 & x < 0.
        \end{array}
        \right.\]

\begin{enumerate}
\def\labelenumi{\arabic{enumi}.}
\tightlist
\item
  Calculer la probabilité que cette durée de fonctionnement \(X\) soit
  comprise entre 50 et 150 heures.
\item
  Calculer la probabilité que l'ordinateur fonctionne moins de 100
  heures.
\item
  Calculer la probabilité que l'ordinateur fonctionne exactement 100
  heures avant sa première panne.
\end{enumerate}

\solref{I-answer-ex.expo}

\runinhead{Exercice -- Précipitation (\(\mathord{\bullet}\))}\phantomsection\addcontentsline{toc}{subsubsection}{Précipitation}\label{I-ex.pluie}
On souhaite modéliser la hauteur de précipitation (en mm) tombée sur
Paris en une journée par une variable aléatoire \(X\). On considère que
cette hauteur suit une loi exponentielle de paramètre \(1/10\) lorsqu'il
pleut et qu'il pleut en moyenne un jour sur deux.

\begin{enumerate}
\def\labelenumi{\arabic{enumi}.}
\tightlist
\item
  Donner la Fonction de répartition de \(X\).
\item
  \(X\) admet-elle une densité\,?
\end{enumerate}

\solref{I-answer-ex.pluie}

\addcontentsline{toc}{subsection}{Quelques lois à densités}
\begin{exemple}[{Quelques lois à
densités}]\phantomsection\label{I-exampledens}

\begin{enumerate}
\def\labelenumi{\arabic{enumi}.}
\item
  La \emph{loi uniforme} sur \([a,b]\), avec \(a,b \in \mathbb{R}\) tels
  que \(a < b\), notée \(\mathcal{U}_{[a,b]}\), est la probabilité
  \(\mathbb{P}\) qui admet la densité \[f(x) = \left\{\begin{array}{ll}
          \frac{1}{b-a} & \text{si } a \leq x \leq b,\\
          0 & \text{sinon}.
          \end{array}\right.\] 
          
          \begin{center}
          \includegraphics{PdfUnif.tex.pdf}
          \end{center}
  En
  vertu de la remarque 1. ci-dessus, on aurait pu choisir
  \(f(a) = f(b) = 0\). Au vu de l'interprétation 2. ci-dessus, le fait
  que \(f\) soit constante sur \([a, b]\) exprime que si on choisit un
  point selon la probabilité uniforme, on a ``autant de chances'' de
  tomber au voisinage de chaque point de l'intervalle \([a, b]\). Cela
  explique le nom ``uniforme''. On remarque aussi que
  \(\mathbb{P}(\{x\}) = 0\) pour tout \(x\) (comme pour toutes les
  probabilités à densité). On a donc une probabilité nulle de tomber
  exactement en un point \(x\) fixé à l'avance. On vérifie aisément que
  sa fonction de répartition est \[F(x) = \left\{\begin{array}{ll}
          0 & \text{si } x < a,\\
          \frac{x-a}{b-a} & \text{si } a \leq x \leq b,\\
          1 & \text{si } x > b.
          \end{array}\right.\] \begin{center}
          \includegraphics{CdfUnif.tex.pdf}
          \end{center}
\item
  La \emph{loi exponentielle} de paramètre \(\theta\), notée
  \(\mathcal{E}_\theta\), est la loi de densité
  \[f(x) = \left\{\begin{array}{ll}
          0 & \text{si } x < 0,\\
          \theta e^{-\theta x} & \text{sinon}.
          \end{array}\right.\] 
          \begin{center}
          \includegraphics{PdfExp.tex.pdf}
          \end{center}
          et de
  fonction de répartition \[F(x) = \left\{\begin{array}{ll}
          0 & \text{si } x < 0,\\
          1 - e^{-\theta x} & \text{sinon}.
          \end{array}\right.\] \begin{center}
          \includegraphics{CdfExp.tex.pdf}
          \end{center} Dans la
  pratique, on utilise fréquemment la loi exponentielle pour modéliser
  une durée de vie ou le temps d'attente avant l'arrivée d'un événement
  spécifique. Par exemple la durée de vie d'une bactérie, d'un composant
  électronique, la durée d'une conversation téléphonique ou le temps qui
  nous sépare du prochain tremblement de terre peuvent être considérées
  comme des variables aléatoires de loi exponentielle.
\item
  La \emph{loi normale} (ou de \emph{Gauss} ou encore \emph{gaussienne})
  de paramètres \(\mu\) et \(\sigma^2\), notée
  \(\mathcal{N}(\mu,\sigma^2)\), est la loi de densité

  \[f(x) = \frac{1}{\sqrt{2\pi}\sigma}\exp\left(-\frac{(x-\mu)^2}{2\sigma^2}\right)\]
  \begin{center}
  \includegraphics{PdfGauss.tex.pdf}
  \end{center}
   Lorsque \(\mu = 0\) et
  \(\sigma^2 = 1\), on l'appelle la loi normale \emph{centrée réduite}
  pour des raisons que nous expliciterons au prochain chapitre.

  La distribution normale fut introduite par De Moivre en 1733. Celui-ci
  l'utilisa pour approximer une variable aléatoire binomiale quand le
  paramètre \(n\) de celle-ci était grand. Ce résultat fut ensuite
  progressivement généralisé par Laplace et d'autres confrères pour
  devenir le théorème connu sous le nom de théorème de la limite
  centrale, qui sera démontré au Chapitre 4. Ce théorème est l'un des
  plus importants de la théorie des probabilités et prouve que de très
  nombreux phénomènes aléatoires suivent approximativement une loi
  normale. On peut citer à titre d'exemple la taille d'un individu
  choisi au hasard, les composantes de la vitesse d'une molécule de gaz
  ou l'erreur de mesure d'une quantité physique.

  Il est difficile de faire des calculs avec la loi normale car sa
  densité n'admet pas de primitive explicite. Aussi des tables
  numériques ont-elles été construites pour permettre aux utilisateurs
  d'obtenir très rapidement des valeurs numériques. Elles sont
  disponibles dans la plupart des logiciels et permettent entre autre de
  représenter sa fonction de répartition.

\begin{center}
\includegraphics{CdfGauss.tex.pdf}
\captionof{figure}{fonction de répartition de la loi normale}
\end{center} 
\end{enumerate}

Nous aurons l'occasion de voir par la suite un grand nombre d'autres
exemples de probabilités à densité.
\end{exemple}

\begin{figure}
\centering
\includegraphics{PdfUnif.tex.pdf}
\caption{Densité de probabilité de la loi
uniforme}\label{I-fig-PdfUnif.tex}
\end{figure}

\begin{figure}
\centering
\includegraphics{CdfUnif.tex.pdf}
\caption{fonction de répartition de la loi
uniforme}\label{I-fig-CdfUnif.tex}
\end{figure}

\begin{figure}
\centering
\includegraphics{PdfExp.tex.pdf}
\caption{Densité de probabilité de la loi
exponentielle}\label{I-fig-PdfExp.tex}
\end{figure}

\begin{figure}
\centering
\includegraphics{CdfExp.tex.pdf}
\caption{fonction de répartition de la loi
exponentielle}\label{I-fig-CdfExp.tex}
\end{figure}

\begin{figure}
\centering
\includegraphics{PdfGauss.tex.pdf}
\caption{Densité de probabilité de la loi
normale}\label{I-fig-PdfGauss.tex}
\end{figure}

\begin{figure}
\centering
\includegraphics{CdfGauss.tex.pdf}
\caption{fonction de répartition de la loi
normale}\label{I-fig-CdfGauss.tex}
\end{figure}

\chapter{Exercices
complémentaires}\label{I-exercices-compluxe9mentaires}

\section{Indépendance et
conditionnement}\label{I-induxe9pendance-et-conditionnement}

\runinhead{Question
1}\phantomsection\addcontentsline{toc}{subsubsection}{Question
1}\label{I-ic-1}
Dans quelles circonstances particulières deux événements \(A\) et \(B\)
tels que \(A\cup B\) soit presque sûr sont-ils
indépendants\,?\solref{I-answer-ic-1}

\runinhead{Question
2}\phantomsection\addcontentsline{toc}{subsubsection}{Question
2}\label{I-ic-2}
Soit \(A\), \(B\) et \(C\) trois événements indépendants \textbf{2 à 2},
avec \(\mathbb{P}(C)>0\). Donner une condition nécessaire et suffisante
sur \(A\), \(B\) et \(C\) pour que \(A\) et \(B\) soient indépendants
relativement à la probabilité conditionnelle \(\mathbb{P}(\,\,|C)\), à
la place de \(\mathbb{P}\).\solref{I-answer-ic-2}

\runinhead{Question
3}\phantomsection\addcontentsline{toc}{subsubsection}{Question
3}\label{I-ic-3}
Etablir que deux événements \(A\) et \(B\) sont indépendants ssi
\[ \mathbb{P}(A\cap B)\mathbb{P}(A^c \cap B^c) = \mathbb{P}(A \cap B^c)\mathbb{P}(A^c\cap B)\]

\solref{I-answer-ic-3}

\section{Problème de Monty Hall}\label{I-mh}

Dans un jeu télévisé, vous êtes confrontés au problème suivant~: devant
vous se trouvent 3 portes fermées. Derrière l'une d'entre elle se trouve
une voiture, derrière chacune des deux autres, une chèvre. Vous désignez
au hasard l'une d'entre elle. Le présentateur, qui sait où se trouve la
voiture, ouvre alors une porte, qui n'est ni celle que vous avez
choisie, ni celle qui cache la voiture. Il vous offre alors la
possibilité de réviser votre choix. Que choisissez-vous\,?

Calculer la probabilité de remporter la voiture selon les deux
stratégies (changer ou non son choix de porte).\solref{I-answer-mh}

\section{Mesure de Probabilité}\label{I-mesure-de-probabilituxe9}

Soit \((\Omega,\mathcal{A},\mathbb{P})\) un espace probabilisé et \(X\)
une variable aléatoire réelle définie sur
\((\Omega,\mathcal{A},\mathbb{P})\) telle que \(X \geq 0\) p.s. et
\(\mathbb{E}(X) = \int_\Omega X(\omega) d\mathbb{P}(\omega) = 1\). (on
reviendra sur cette définition dans le chapitre suivant).

On définit \(Q : \mathcal{A}\to \mathbb{R}\) par
\(Q(A) = \mathbb{E}(X\,1_A) = \int_{\Omega} X(\omega) 1_A(\omega)d\mathbb{P}(\omega)\)
pour \(A\in \mathcal{A}\).

\runinhead{Question
1}\phantomsection\addcontentsline{toc}{subsubsection}{Question
1}\label{I-mes-proba1}
Montrer que \(Q\) est une mesure de probabilité sur
\((\Omega,\mathcal{A})\).\solref{I-answer-mes-proba1}

\runinhead{Question
2}\phantomsection\addcontentsline{toc}{subsubsection}{Question
2}\label{I-mes-proba2}
Montrer que si \(\mathbb{P}(A) = 0\), alors \(Q(A) = 0\). Donner un
exemple montrant que \(Q(A) = 0\) n'implique pas en général que
\(\mathbb{P}(A) = 0\).\solref{I-answer-mes-proba2}

\runinhead{Question
3}\phantomsection\addcontentsline{toc}{subsubsection}{Question
3}\label{I-mes-proba3}
En notant \(\mathbb{E}_Q\), l'espérance par rapport à \(Q\), montrer que
\(\mathbb{E}_Q(Y) = \int_{\Omega} Y(\omega) Q(d\omega) = \mathbb{E}_\mathbb{P}(XY)\)
pour toute variable aléatoire réelle \(Y\) définie sur
\((\Omega,\mathcal{A},\mathbb{P})\) positive ou
bornée.\solref{I-answer-mes-proba3}

\runinhead{Question
4}\phantomsection\addcontentsline{toc}{subsubsection}{Question
4}\label{I-mes-proba4}
Supposons que \(\mathbb{P}(X > 0) = 1\).

\begin{enumerate}
\def\labelenumi{\roman{enumi})}
\item
  Montrer que \(\frac{1}{X}\) est intégrable pour \(Q\).
\item
  Soit \(R : \mathcal{A}\to \mathbb{R}\) telle que
  \(\forall A \in \mathcal{A}\,\, R(A) = \mathbb{E}_Q(\frac{1}{X}1_A)\).
  Montrer que \(R\) est exactement la probabilité \(\mathbb{P}\).
\item
  Montrer que \(Q(A) = 0\) implique \(\mathbb{P}(A) = 0\).
\end{enumerate}

\solref{I-answer-mes-proba4}

\section{Queues de distributions}\label{I-queues-de-distributions}

Soit une densité \(f\) sur \(\mathbb{R}\) telle que la fonction
\(h: x \in \mathbb{R}\mapsto x\,f(x)\) est intégrable. Cette hypothèse
sera appelée \((\mathcal{H})\). On note \(F\) la fonction de répartition
correspondante. Nous allons montrer que\footnote{Soient
  \(a \in [-\infty,+\infty]\) et \(f\) et \(g\) deux fonctions définies
  sur un même voisinage \(I\) de \(a\). On dit que \(f\) est
  \emph{négligeable} par rapport à \(g\) (ou que \(g\) est
  \emph{prépondérante} par rapport à \(f\)) au voisinage de \(a\) et on
  note note \(f(x) \underset{x\to a}{=} o(g(x))\) s'il existe une
  fonction \(\epsilon\) définie sur \(I\) telle que
  \(\epsilon(x) \xrightarrow[x\to a]{} 0\) et \(f = \epsilon \times g\).}
\[\left|\begin{array}{rl} 1-F(x) \underset{x\to+\infty}{=} o\left(\dfrac{1}{x}\right), & (1)\\[1em] F(x) \underset{x\to-\infty}{=} o\left(\dfrac{1}{x}\right).& (2) \end{array} \right. \]

\runinhead{Etude en
\(+\infty\)}\phantomsection\addcontentsline{toc}{subsubsection}{Etude en
\(+\infty\)}\label{I-dlfdr-pi}

\begin{enumerate}
\def\labelenumi{\arabic{enumi}.}
\item
  Vérifier que pour tout \(u > 0\),
  \[\int_{-\infty}^{+\infty} h(x)\,dx \geq \int_{-\infty}^u h(x)\,dx + u\,\left(1-F(u)\right).\]
\item
  En déduire \((1)\).
\end{enumerate}

\solref{I-answer-dlfdr-pi}

\runinhead{Etude en
\(-\infty\)}\phantomsection\addcontentsline{toc}{subsubsection}{Etude en
\(-\infty\)}\label{I-dlfdr-mi}

\begin{enumerate}
\def\labelenumi{\arabic{enumi}.}
\setcounter{enumi}{2}
\tightlist
\item
  En vous inspirant du raisonnement précédent, montrer \((2)\).
\end{enumerate}

\solref{I-answer-dlfdr-mi}

\runinhead{Interprétation}\phantomsection\addcontentsline{toc}{subsubsection}{Interprétation}\label{I-dlfdr-interpret}

\begin{enumerate}
\def\labelenumi{\arabic{enumi}.}
\setcounter{enumi}{3}
\tightlist
\item
  Proposez une interprétation de ce résultat.
\end{enumerate}

\solref{I-answer-dlfdr-interpret}

\runinhead{Etude de la
réciproque}\phantomsection\addcontentsline{toc}{subsubsection}{Etude de
la réciproque}\label{I-dlfdr-reciproque}
Considérons maintenant la fonction de répartition
\[ F: x \in \mathbb{R}\mapsto \left|\begin{array}{ll} 1 - \dfrac{e}{x\,\ln(x)} & \text{si } x \geq e,\\ 0 & \text{si } x <e, \end{array}\right. \]

\begin{enumerate}
\def\labelenumi{\arabic{enumi}.}
\setcounter{enumi}{4}
\item
  Vérifier que cette fonction possède bien une densité \(f\) et
  l'expliciter.
\item
  A-t-on les développements asymptotiques \((1)\) et \((2)\)\,?
\item
  La fonction \(h\) correspondante est-elle absolument intégrable\,?
\item
  Qu'en concluez-vous\,?
\end{enumerate}

\solref{I-answer-dlfdr-reciproque}

\runinhead{Liens entre les développements
asymptotiques}\phantomsection\addcontentsline{toc}{subsubsection}{Liens
entre les développements asymptotiques}\label{I-dlfdr-liensdl}
Soient les fonctions de répartition
\[F: x\in\mathbb{R}\mapsto \left|\begin{array}{ll} 1 - \dfrac{1}{\ln(x)} & \text{si } x \geq e,\\ 0 & \text{sinon,}\end{array}\right.\]
et \(G : x\in\mathbb{R}\mapsto 1 - F(-x)\).

\begin{enumerate}
\def\labelenumi{\arabic{enumi}.}
\setcounter{enumi}{8}
\tightlist
\item
  Que peut-on dire sur les éventuels liens entre \((1)\) et \((2)\) à
  l'aide de ces deux fonctions\,?
\end{enumerate}

\solref{I-answer-dlfdr-liensdl}

\section{Densité et fonction de répartition d'une loi
Normale}\label{I-densituxe9-et-fonction-de-ruxe9partition-dune-loi-normale}

On considère la densité d'une loi Normale centrée réduite~:
\[f : x\in \mathbb{R}\mapsto \dfrac{1}{\sqrt{2\,\pi}}\,\exp\left\{-\dfrac{x^2}{2}\right\},\]
où \(f(-\infty) := \lim\limits_{x\to-\infty} f(x)\) et
\(f(+\infty) := \lim\limits_{x\to+\infty} f(x)\).

\runinhead{Propriétés
générales}\phantomsection\addcontentsline{toc}{subsubsection}{Propriétés
générales}\label{I-fdrgaussprop}

\begin{enumerate}
\def\labelenumi{\arabic{enumi}.}
\item
  Faire l'étude de \(f\) (domaine, parité, limites aux bornes,
  dérivabilité, variations, convexité/concavité).
\item
  Donner la définition de la fonction de répartition \(F\) associée à
  \(f\). En donner une expression faisant apparaître la \emph{fonction
  d'erreur} (qui est une fonction spéciale)
  \[\text{erf} : x\in\mathbb{R}\mapsto \dfrac{2}{\sqrt{\pi}}\,\int_{0}^x e^{-u^2}\,du.\]
\end{enumerate}

\solref{I-answer-fdrgaussprop}

\runinhead{Encadrement de \(1 - F(x)\) pour tout
\(x > 0\)}\phantomsection\addcontentsline{toc}{subsubsection}{Encadrement
de \(1 - F(x)\) pour tout \(x > 0\)}\label{I-fdrgaussenca}
Nous allons maintenant démontrer la propriété suivante~:
\(\forall\,x > 0\)
\[ \dfrac{x}{1+x^2}\,f(x) \leq 1 - F(x) \leq \dfrac{f(x)}{x}. \]

\begin{enumerate}
\def\labelenumi{\arabic{enumi}.}
\setcounter{enumi}{2}
\item
  En observant que \(u\geq x >0\) implique \(\dfrac{u}{x} \geq 1\),
  montrer l'inégalité de droite.
\item
  En observant que \(u\geq x > 0\) implique \(u^{-2} \leq x^{-2}\),
  montrer que \(\forall\,x>0\),
  \[\left(1+\dfrac{1}{x^2}\right)\,\left(1-F(x)\right) \geq \dfrac{f(x)}{x}.\]
  En déduire l'inégalité de gauche.
\end{enumerate}

\solref{I-answer-fdrgaussenca}

\runinhead{Equivalent de \(1-F(x)\) quand
\(x\to+\infty\).}\phantomsection\addcontentsline{toc}{subsubsection}{Equivalent
de \(1-F(x)\) quand \(x\to+\infty\).}\label{I-fdrgaussequi}

\begin{enumerate}
\def\labelenumi{\arabic{enumi}.}
\setcounter{enumi}{4}
\tightlist
\item
  Déduire de l'encadrement prédédent un équivalent de \(1-F(x)\) lorsque
  \(x\to+\infty\).
\end{enumerate}

\solref{I-answer-fdrgaussequi}

\begin{enumerate}
\def\labelenumi{\arabic{enumi}.}
\setcounter{enumi}{5}
\tightlist
\item
  Comparer ce résultat à celui de l'exercice sur les développements
  limités pour les fonctions de répartition. En particulier,
  \(h:x\in[-\infty,+\infty] \mapsto x\,f(x)\) est-elle intégrable\,?
\end{enumerate}

\runinhead{Loi Normale
générale.}\phantomsection\addcontentsline{toc}{subsubsection}{Loi
Normale générale.}\label{I-fdrgaussgen}
Considérons maintenant la densité plus générale
\[g(x) = \dfrac{1}{\sqrt{2\,\pi}\,\sigma}\, \exp\left\{-\dfrac{(x-\mu)^2}{2\,\sigma^2} \right\},\]
où \(\mu \in \mathbb{R}\) et \(\sigma^2 > 0\). On note \(G\) sa fonction
de répartition associée.

\begin{enumerate}
\def\labelenumi{\arabic{enumi}.}
\setcounter{enumi}{6}
\tightlist
\item
  En réécrivant \(g\) en fonction de \(f\), déduire des questions
  précédentes un équivalent de \(1-G(x)\) lorsque \(x\to+\infty\).
\end{enumerate}

\solref{I-answer-fdrgaussgen}

\section{Loi de vie et de mort}\label{I-loi-de-vie-et-de-mort}

La durée de vie d'un être vivant ou d'un matériel peut-être assimilée à
une variable aléatoire strictement positive \(T\). Dans ce cadre, on
peut définir les notions suivantes~:

\begin{itemize}
\tightlist
\item
  Loi de vie a priori~: il s'agit de la loi du temps \(T\) caractérisée
  par fonction de répartition complémentaire \[G(t) = \mathbb{P}(T>t)\]
\item
  Loi de survie après \(t_0 \geq 0\)~: il s'agit de la loi du temps
  \(T-t_0\) qu'il lui reste à vivre, sachant qu'il était encore en vie à
  \(t_0\), de fonction de répartition complémentaire
  \[G_{t_0}(t) =\mathbb{P}(T>t+t_0|T>t_0)\]
\end{itemize}

On dira que la loi de vie satisfait la propriété de non vieillissement
(ou d'absence de mémoire) si la loi de survie et la loi de vie sont
égales~:
\[\forall t \geq 0 \text{ et } \forall t_0 \geq 0, G_{t_0}(t) = G(t)\]

\runinhead{Question
1}\phantomsection\addcontentsline{toc}{subsubsection}{Question
1}\label{I-viemort1}
Exprimer la loi de survie à partir de la loi de vie a
priori.\solref{I-answer-viemort1}

\runinhead{Question
2}\phantomsection\addcontentsline{toc}{subsubsection}{Question
2}\label{I-viemort2}
Montrer qu'une variable aléatoire \(T\) satisfait la propriété de
non-vieillissement si et seulement si elle est de loi
exponentielle.\solref{I-answer-viemort2}

\runinhead{Question
3}\phantomsection\addcontentsline{toc}{subsubsection}{Question
3}\label{I-viemort3}
On suppose que \(T\) admet une densité continue sur
\(\mathbb{R}_+^\ast\). Montrer que l'on peut définir le taux de mort à
l'instant \(t\) par
\[D(t) = \lim_{\Delta t \to 0} \frac{1}{\Delta t}\mathbb{P}(t < T \leq t + \Delta t | T >t)\]
et exprimer \(G\) en fonction de \(D\).

Quelles lois correspondent-elles à \(D\) constant\,?
\solref{I-answer-viemort3}

\begin{center}\rule{0.5\linewidth}{0.5pt}\end{center}



\chapter{Références}\label{I-ruxe9fuxe9rences}

\hypertarget{refs}{}
\begin{cslreferences}
\leavevmode\hypertarget{refI-polygarnier}{}%
Garnier, Josselin. 2018. \emph{Aléatoire}.
\url{http://josselin-garnier.org/wp-content/uploads/2018/03/MAP311_2018.pdf}.

\leavevmode\hypertarget{refI-Jacod}{}%
Jacod, J., and P. Protter. 2003. \emph{L'essentiel En Théorie Des
Probabilités}. Cassini.
\url{https://hal.archives-ouvertes.fr/hal-00104956}.

\leavevmode\hypertarget{refI-polyponts}{}%
Jourdain, Benjamin. 2016. \emph{Probabilités et Statistique}. Ellipses.
\url{http://cermics.enpc.fr/~jourdain/probastat/poly.pdf}.

\leavevmode\hypertarget{refI-matheron}{}%
Matheron, Georges. 1978. \emph{Estimer et choisir: essai sur la pratique
des probabilités}. Fontainebleau, France: Ecole nationale supérieure des
mines de Paris.
\url{http://cg.ensmp.fr/bibliotheque/public/MATHERON_Ouvrage_00208.pdf}.

\leavevmode\hypertarget{refI-tebaldi}{}%
Tebaldi, Claudia, and Reto Knutti. 2007. ``The Use of the Multi-Model
Ensemble in Probabilistic Climate Projections.'' \emph{Philosophical
Transactions of the Royal Society A: Mathematical, Physical and
Engineering Sciences} 365 (1857): 2053--75.
\url{https://doi.org/10.1098/rsta.2007.2076}.
\end{cslreferences}


\part{Moments et vecteurs aléatoires}
\chapter{Objectifs d'apprentissage}\label{II-objectifs-dapprentissage}

\runinhead{Moments d'une variable
aléatoire}\phantomsection\label{II-moments-dune-variable-aluxe9atoire}

\begin{itemize}
\tightlist
\item
  $\mathord{\bullet}$ connaître la définition de l'espérance d'une
  variable aléatoire
\item
  $\mathord{\bullet}$ connaître l'interprétation de l'espérance
\item
  $\mathord{\bullet}$ connaître les propriétés de base de l'espérance
\item
  $\mathord{\bullet}$ connaître la définition de la variance d'une
  variable aléatoire
\item
  $\mathord{\bullet}$ savoir que la variance est toujours positive
\item
  $\mathord{\bullet}$ savoir ce que signifie variable centrée-réduite
\item
  $\mathord{\bullet}$ connaître la définition de la covariance
\item
  $\mathord{\bullet}$ savoir que la covariance est une forme bilinéaire
\item
  $\mathord{\bullet}$ connaître et savoir manipuler l'inégalité de
  Cauchy-Schwarz
\item
  $\mathord{\bullet}\mathord{\bullet}$ connaître les conditions
  d'existence de l'espérance d'une fonction (mesurable) d'une variable
  aléatoire
\item
  $\mathord{\bullet}\mathord{\bullet}$ savoir exprimer la probabilité
  d'un événement de la forme \(X \in A\) sous la forme d'une espérance
\item
  $\mathord{\bullet}\mathord{\bullet}$ connaître et savoir calculer les
  moments des lois usuelles
\end{itemize}

\runinhead{Vecteurs
aléatoires}\phantomsection\label{II-vecteurs-aluxe9atoires}

\begin{itemize}
\tightlist
\item
  $\mathord{\bullet}$ savoir que les composantes d'un vecteur aléatoire
  sont des variables aléatoires définies sur le même espace de
  probabilité
\item
  $\mathord{\bullet}$ connaître la définition d'un vecteur aléatoire à
  densité
\item
  $\mathord{\bullet}\mathord{\bullet}$ connaître les conditions
  d'existence de l'espérance d'une fonction (mesurable) d'un vecteur
  aléatoire à densité
\item
  $\mathord{\bullet}\mathord{\bullet}$ savoir calculer les densités
  marginales d'un vecteur aléatoire à densité
\item
  $\mathord{\bullet}\mathord{\bullet}$ savoir qu'un couple de variables
  aléatoires à densité n'admet pas nécessairement de densité
\item
  $\mathord{\bullet}$ connaître la définition de l'espérance d'un
  vecteur aléatoire à densité
\item
  $\mathord{\bullet}$ connaître la définition de la matrice de
  covariance d'un vecteur aléatoire à densité
\item
  $\mathord{\bullet}$ savoir que la matrice de covariance est
  semi-définie positive
\end{itemize}

\runinhead{Variables et vecteurs aléatoires
indépendantes}\phantomsection\label{II-variables-et-vecteurs-aluxe9atoires-induxe9pendantes}

\begin{itemize}
\tightlist
\item
  $\mathord{\bullet}$ connaître la définition de l'indépendance de deux
  vecteurs/variables aléatoires
\item
  $\mathord{\bullet}$ savoir caractériser l'indépendance de deux
  vecteurs/variables aléatoires à densité
\item
  $\mathord{\bullet}\mathord{\bullet}$ savoir que si \(X\) et \(Y\) sont
  indépendantes, alors \(g(X)\) et \(h(Y)\) sont aussi indépendantes
\item
  $\mathord{\bullet}\mathord{\bullet}$ savoir que si \(X\) et \(Y\) sont
  indépendantes et si \(g(X)\) et \(h(Y)\) sont intégrables alors
  \(\mathbb{E}(g(X)g(Y)) = \mathbb{E}(g(X))\mathbb{E}(g(Y))\)
\item
  $\mathord{\bullet}$ savoir que la covariance de deux variables
  indépendantes est nulle
\item
  $\mathord{\bullet}\mathord{\bullet}$ savoir que la réciproque est
  fausse et connaître un contre-exemple
\end{itemize}

\runinhead{Identification de
densité}\phantomsection\label{II-identification-de-densituxe9}

\begin{itemize}
\tightlist
\item
  $\mathord{\bullet}\mathord{\bullet}$ savoir identifier la densité
  d'une fonction d'une variable aléatoire par la méthode de la fonction
  muette
\item
  $\mathord{\bullet}\mathord{\bullet}$ savoir identifier la densité
  d'une fonction d'un vecteur aléatoire par la méthode de la fonction
  muette
\end{itemize}

\chapter{Moments d'une variable
aléatoire}\label{II-moments-dune-variable-aluxe9atoire-1}

La loi de probabilité d'une variable aléatoire définie sur l'espace
probabilisé \((\Omega,\mathcal{A},\mathbb{P})\) va nous permettre de
calculer aisément des grandeurs caractéristiques telles que sa valeur
moyenne et sa variance (lorsqu'elles existent). Elles sont définies
ci-dessous.

\addcontentsline{toc}{subsection}{Espérance d'une variable aléatoire
réelle}
\begin{definition}[{Espérance d'une variable aléatoire
réelle}]\phantomsection\label{II-defesp}

La variable aléatoire \(X : \Omega \to \mathbb{R}\) de loi
\(\mathbb{P}_X\) est dite \emph{intégrable} si l'intégrale
\(\int_\mathbb{R}|x| d\mathbb{P}_X(x)\) est finie, autrement dit si
\(x\) est \(\mathbb{P}_X\)-intégrable. On définit alors son
\emph{espérance} par
\[\mathbb{E}(X) = \int_\mathbb{R}x d\mathbb{P}_X(x) = \int_\Omega X(\omega) d\mathbb{P}(\omega)\]

Si la variable aléatoire \(X\) admet une densité \(f\) par rapport à la
mesure de Lebesgue, elle sera \emph{intégrable} si l'intégrale
\(\int_\mathbb{R}|x|f(x) dx\) est finie, autrement dit si le produit
\(x f(x)\) est intégrable par rapport à la mesure de Lebesgue. Son
\emph{espérance} vaut alors \[\mathbb{E}(X) = \int_\mathbb{R}x f(x)dx\]
\end{definition}

\addcontentsline{toc}{subsection}{Interprétation}
\begin{remarque}[{Interprétation}]\phantomsection\label{II-interpruxe9tation}

\(\mathbb{E}(X)\) est un nombre réel qui donne une valeur moyenne
résumant la variable aléatoire \(X\).

L'espérance mathématique d'une variable aléatoire est un concept
fondamental de la théorie des probabilités. La dénomination d'espérance
pour cette quantité fait référence aux problèmes de jeux et d'espérance
de gain. Cette terminologie imagée a été introduite par Pascal.

On note \(\mathcal{L}^1\) l'ensemble de toutes les variables réelles
\(X\) intégrables. Les propriétés suivantes découlent directement des
propriétés de l'intégrale.
\end{remarque}

\addcontentsline{toc}{subsection}{Propriétés de base}
\begin{proposition}[{Propriétés de
base}]\phantomsection\label{II-propl1}

\begin{itemize}
\tightlist
\item
  \(\mathcal{L}^1\) est un espace vectoriel et
  \(\forall X,Y \in \mathcal{L}^1\), \(\forall a,b \in \mathbb{R}\)
  \[\mathbb{E}(aX + bY) = a\mathbb{E}(X) + b\mathbb{E}(Y).\]
\item
  \(X \in \mathcal{L}^1 \Leftrightarrow |X| \in \mathcal{L}^1\), et dans
  ce cas \[|\mathbb{E}(X)| \leq \mathbb{E}(|X|).\]
\item
  Si \(X \geq 0\) et \(X \in \mathcal{L}^1\), alors
  \(\mathbb{E}(X) \geq 0.\)
\item
  Si \(X,Y \in \mathcal{L}^1\) sont telles que \(X \leq Y\), alors
  \[ \mathbb{E}(X) \leq \mathbb{E}(Y).\]
\item
  L'espérance d'une variable presque-sûrement constante est égale à
  cette constante~:
  \[ \text{Si } \mathbb{P}(X(\omega) = a ) = 1, \text{ alors } \mathbb{E}(X) = a.\]
\item
  Si \(\exists b \in \mathbb{R}_+\) tel que \(|X| \leq b\), alors
  \(X \in \mathcal{L}^1\) et \(\mathbb{E}(X) \leq b\).
\end{itemize}
\end{proposition}

\runinhead{Exercice -- Démontrer ces
propriétés (\(\mathord{\bullet}\))}\phantomsection\addcontentsline{toc}{subsubsection}{Démontrer
ces propriétés}\label{II-basics}
\addcontentsline{toc}{subsection}{Rappel~: cas discret}

\begin{remarque}[{Rappel~: cas
discret}]\phantomsection\label{II-rappel-cas-discret}
Dans le cas d'une variable aléatoire discrète \(X\) à valeurs dans
\(\mathbb{N}^\ast\), son espérance est définie par la quantité
\(\mathbb{E}(X) = \sum_{i\in\mathbb{N}^\ast} i\mathbb{P}(X=i)\), pourvu
que celle-ci soit finie. On voit immédiatement que les propriétés
ci-dessus sont également vérifiées.
\end{remarque}

Outre l'espace \(\mathcal{L}^1\), nous pouvons définir l'espace
\(\mathcal{L}^2\) des variables aléatoires réelles dont le carré \(X^2\)
est dans \(\mathcal{L}^1\).

\addcontentsline{toc}{subsection}{Variance}
\begin{definition}[{Variance}]\phantomsection\label{II-defvar}

La variable aléatoire \(X : \Omega \to \mathbb{R}\) de loi
\(\mathbb{P}_X\) est dite \emph{de carré intégrable} si
\(\mathbb{E}(X^2) = \int_{\mathbb{R}} x^2 \mathbb{P}_X(dx) = \int_{\Omega} X^2(\omega)\mathbb{P}(d\omega) < +\infty\).
Sa \emph{variance} est définie par
\[\mathbb{V}(X) = \mathbb{E}((X-\mathbb{E}(X))^2)\]

Si la variable aléatoire \(X\) admet une densité \(f\) par rapport à la
mesure de Lebesgue, elle sera \emph{de carré intégrable} si
\(\mathbb{E}(X^2) = \int_\mathbb{R}x^2 f(x)dx\) est définie, autrement
dit si le produit \(x^2 f(x)\) est intégrable par rapport à la mesure de
Lebesgue.
\end{definition}

\addcontentsline{toc}{subsection}{Propriétés}
\begin{proposition}[{Propriétés}]\phantomsection\label{II-propl2}

\(\mathcal{L}^2\) est un sous-espace vectoriel de \(\mathcal{L}^1\), et
si \(X \in \mathcal{L}^2\),
\[|\mathbb{E}(X)| \leq \mathbb{E}(|X|) \leq \sqrt{\mathbb{E}(X^2)}\]
\end{proposition}

\begin{preuve}
Soient \(X\) et \(Y\) deux variables aléatoires réelles de
\(\mathcal{L}^2\) et \(a \in \mathbb{R}\). Comme
\((aX+ Y)^2 \leq 2 a^2 X^2 + 2 Y^2\), alors
\(aX + Y \in \mathcal{L}^2\). Ainsi, \(\mathcal{L}^2\) est un espace
vectoriel.

L'inclusion \(\mathcal{L}^2 \subset \mathcal{L}^1\) découle de
\(|X| \leq 1 + X^2\) et de la \protect\hyperref[II-propl1]{proposition
précédente (p.~\pageref*{II-propl1})} (linéarité).

La première inégalité a déjà été vue
\protect\hyperref[II-propl1]{ci-dessus (p.~\pageref*{II-propl1})}. Pour
la seconde, nous pouvons nous limiter au cas où \(X\) est positive. Soit
alors \(a = \mathbb{E}(X)\) et \(Y = X-a\). Par linéarité, on a
\[ \mathbb{E}(Y^2) = \mathbb{E}(X^2) - 2a \mathbb{E}(X) + a^2 = \mathbb{E}(X^2)-a^2.\]
Et \(\mathbb{E}(Y^2) \geq 0\) par le troisième point de la
\protect\hyperref[II-propl1]{proposition
ci-dessus (p.~\pageref*{II-propl1})}. Par conséquent,
\(\mathbb{E}(X)^2 = a^2 \leq \mathbb{E}(X^2)\) ce qui est le résultat
recherché.\hfill$\blacksquare$
\end{preuve}

\runinhead{Exercice -- Approximation par une
constante (\(\mathord{\bullet}\))}\phantomsection\addcontentsline{toc}{subsubsection}{Approximation
par une constante}\label{II-exo-approx}
Soit \(X\) une v.a. de carré intégrable. Montrer que
\(\mathbb{E}((X - a)^2 )\) atteint son minimum lorsque
\(a = \mathbb{E}(X)\).\solref{II-answer-exo-approx}

\addcontentsline{toc}{subsection}{La variance est positive}
\begin{remarque}[{La variance est
positive}]\phantomsection\label{II-la-variance-est-positive}

En vertu de cette
\protect\hyperref[II-propl2]{proposition (p.~\pageref*{II-propl2})},
\(\mathbb{V}(X)\) est \textbf{positive} et sa racine carrée \(\sigma_X\)
s'appelle l'\emph{écart-type} de \(X\). L'écart-type est une grandeur
qui mesure la moyenne (en un certain sens) de l'écart des valeurs de
\(X\) à sa moyenne, d'où son nom.

On a également \[\mathbb{V}(X) = \mathbb{E}(X^2)-\mathbb{E}(X)^2\] que
l'on obtient en développant \((X-\mathbb{E}(X))^2\). Cette manipulation
anodine est fort utile dans la pratique. On retiendra que ``La variance
est égale à la moyenne des carrés moins le carré de la moyenne''. On
désigne le terme \(\mathbb{E}(X^2)\) par l'expression \emph{moment
d'ordre deux} tandis que la variance est parfois appelée \emph{moment
centré d'ordre deux}.
\end{remarque}

\addcontentsline{toc}{subsection}{Variable centrée réduite}
\begin{remarque}[{Variable centrée
réduite}]\phantomsection\label{II-variable-centruxe9e-ruxe9duite}

D'après ce qui précède, si \(X\) est une variable aléatoire de carré
intégrable, d'espérance \(\mathbb{E}(X)\) et d'écart-type
\(\sigma_X >0\), alors la variable aléatoire
\[\frac{X-\mathbb{E}(X)}{\sigma_X}\] est d'espérance nulle et de
variance 1. On dira qu'une telle variable aléatoire est \emph{centrée et
réduite}.

On peut remarquer que si \(X\) et \(Y\) sont dans \(\mathcal{L}^2\), la
variable aléatoire \(XY\) est dans \(\mathcal{L}^1\), puisque
\(|XY| \leq \frac{1}{2}(X^2+Y^2)\). On peut ainsi définir la
\emph{covariance} de deux variables aléatoires~:
\end{remarque}

\addcontentsline{toc}{subsection}{Covariance}
\begin{definition}[{Covariance}]\phantomsection\label{II-defcov}

Si \(X\) et \(Y\) sont dans \(\mathcal{L}^2\), la variable aléatoire
\((X-\mathbb{E}(X))(Y-\mathbb{E}(Y))\) est intégrable. On appelle la
\emph{covariance} de \(X\) et \(Y\) l'espérance de cette variable
aléatoire et on la note~:
\[\text{Cov}(X,Y) = \mathbb{E}((X-\mathbb{E}(X))(Y-\mathbb{E}(Y))).\] Le
\emph{coefficient de corrélation} des variables aléatoires \(X\) et
\(Y\) est le nombre
\[\rho(X,Y) = \frac{\text{Cov}(X,Y)}{\sqrt{\mathbb{V}(X)}\sqrt{\mathbb{V}(Y)}}\]
qui est bien défini lorsque \(\mathbb{V}(X)>0\) et \(\mathbb{V}(Y)>0\).
\end{definition}

\runinhead{Exercice -- Calcul d'une
covariance (\(\mathord{\bullet}\))}\phantomsection\addcontentsline{toc}{subsubsection}{Calcul
d'une covariance}\label{II-cov-unif}
Soit \(U \sim \mathcal{U}_{[0,1]}\) et \(V = 1-U\). Calculer
\(\text{Cov}(U,V)\) et \(\rho(U,V)\).\solref{II-answer-cov-unif}

Du fait de la linéarité de l'espérance, on a
\[\text{Cov}(X,Y) = \mathbb{E}(XY) - \mathbb{E}(X)\mathbb{E}(Y)\] et
d'ailleurs, on voit que la formule de calcul de la variance donnée plus
haut est un cas particulier de cette formulation car
\(\mathbb{V}(X) = \text{Cov}(X,X)\). La linéarité de l'espérance nous
donne encore pour \(a,a',b,b' \in \mathbb{R}\) \begin{align*} 
\mathbb{E}((aX+b)(a'Y+b')) &= aa'\mathbb{E}(XY) + ab'\mathbb{E}(X) +a'b\mathbb{E}(Y) + bb'\\
\mathbb{E}(aX+b)\mathbb{E}(a'Y+b') &= aa'\mathbb{E}(X)\mathbb{E}(Y) + ab'\mathbb{E}(X) +a'b\mathbb{E}(Y) + bb'
\end{align*} On en déduit que la covariance est une forme bilinéaire sur
l'espace vectoriel des variables aléatoires de carré intégrable, et nous
avons \[\text{Cov}(aX+b,a'Y+b') = aa'\text{Cov}(X,Y)\] En particulier,
on a \[ \mathbb{V}(aX) = a^2\mathbb{V}(X)\] On en déduit que les
coefficients de corrélation de \(X\) et \(Y\) et de \(aX+b\) et
\(a'Y+b'\) sont égaux lorsque \(aa' >0\).

\addcontentsline{toc}{subsection}{Inégalité de Cauchy-Schwarz}
\begin{proposition}[{Inégalité de
Cauchy-Schwarz}]\phantomsection\label{II-CS}

Soit \(X\) et \(Y\) deux variables aléatoires de carré intégrable, alors
on a \emph{l'inégalité de Cauchy-Schwarz}~:
\[|\mathbb{E}(XY)| \leq \mathbb{E}(|XY|) \leq \sqrt{\mathbb{E}(X^2)\mathbb{E}(Y^2)}\]
avec égalité si et seulement si \(X\) et \(Y\) sont presque-sûrement
proportionnelles.
\end{proposition}

\begin{preuve}
La première inégalité a été démontrée plus haut. Pour la seconde, on a
\(\forall x \in \mathbb{R}\) d'après la linéarité de l'espérance~:
\[x^2\mathbb{E}(X^2) + 2x\mathbb{E}(XY)+ \mathbb{E}(Y^2) = \mathbb{E}((xX+Y)^2) \geq 0.\]
Mais ceci n'est possible que si ce trinôme en \(x\) n'a au plus qu'une
seul racine réelle\,; son discriminant
\[4\left((\mathbb{E}(XY))^2 - \mathbb{E}(X^2)\mathbb{E}(Y^2)\right)\]
doit donc être négatif ou nul ce qui donne le résultat.

Le discriminant est nul si et seulement si le trinôme admet une racine
double \(x_0\) et dans ce cas, \(Y(\omega) = -x_0 X(\omega)\) pour
presque tout \(\omega\).\hfill$\blacksquare$
\end{preuve}

\runinhead{Exercice -- Coefficient de
corrélation (\(\mathord{\bullet}\))}\phantomsection\addcontentsline{toc}{subsubsection}{Coefficient
de corrélation}\label{II-CS-corr}
Montrer que le coefficient de corrélation de \(X\) et \(Y\) vérifie
\[-1\leq \rho(X,Y) \leq 1.\]

\solref{II-answer-CS-corr}

Enfin, il peut être intéressant de pouvoir calculer l'espérance d'une
fonction mesurable d'une variable aléatoire réelle qui est une variable
aléatoire en vertu de la
\href{Probabilité\%20I.pdf\%20\#composition}{proposition de
composition}.

\addcontentsline{toc}{subsection}{Espérance de \(g(X)\)}
\begin{proposition}[{Espérance de
\(g(X)\)}]\phantomsection\label{II-esperanceg}

Soit \(X\) une variable aléatoire réelle de loi \(\mathbb{P}_X\), et
\(g\) une fonction mesurable de \((\mathbb{R},\mathcal{B}(\mathbb{R}))\)
dans \((\mathbb{R},\mathcal{B}(\mathbb{R}))\). Alors \(g(X)\) est
\(\mathbb{P}_X\) - intégrable si et seulement si l'intégrale
\[\int_\mathbb{R}|g(x)| d\mathbb{P}_X(x),\] est finie et dans ce cas
\[\mathbb{E}(g(X)) = \int_\mathbb{R}g(x)d\mathbb{P}_X(x) = \int_\Omega g(X(\omega)) d\mathbb{P}(\omega)\]

Si \(\mathbb{P}_X\) admet une densité \(f\) par rapport à la mesure de
Lebesgue, alors \(g(X)\) est \(\mathbb{P}_X\) - intégrable si et
seulement si le produit \(fg\) est intégrable par rapport à la mesure de
Lebesgue. On a alors~:
\[\mathbb{E}(g(X)) =  \int_\mathbb{R}g(x)d\mathbb{P}_X(x) = \int_\mathbb{R}g(x)f(x)dx = \int_\Omega g(X(\omega)) d\mathbb{P}(\omega)\]
\end{proposition}

\begin{preuve}
On rappelle que la loi de probabilité \(\mathbb{P}_X\) est définie pour
\(B\in \mathcal{B}(\mathbb{R})\) par
\(\mathbb{P}_X(B)= \mathbb{P}(X^{-1}(B))\). Par conséquent
\[\mathbb{E}(1_B(X)) = \mathbb{P}(X^{-1}(B)) = \mathbb{P}_X(B) = \int_\mathbb{R}1_B (X) d\mathbb{P}(x).\]
Ainsi, si \(g\) est une fonction étagée positive
(\(g \in \mathcal{E}^+\)), on obtient le résultat voulu par linéarité.

Si \(g\) est positive, soit \((g_n)_{n \in \mathbb{N}^\star}\) une suite
croissante de fonctions étagées positives convergeant simplement vers
\(g\). On a alors \begin{align*}
\mathbb{E}(g(x)) & = \mathbb{E}(\lim_{n\to \infty}g_n(X))\\
           & = \lim_{n\to \infty} \mathbb{E}(g_n(X)) ~\text{ par le théorème de convergence monotone}\\
           & = \lim_{n\to \infty} \int_\mathbb{R}g_n(x) d\mathbb{P}_X(x)\\
           & = \int_\mathbb{R}\lim_{n\to \infty} g_n(x) d\mathbb{P}_X(x) ~\text{ par le théorème de convergence monotone}\\
           & = \int_\mathbb{R}g(x) d\mathbb{P}_X(x).
\end{align*}

Enfin, si \(g\) n'est pas positive, on utilise la décomposition
\(g = g^+ - g^-\) et on déduit le résultat par
soustraction.\hfill$\blacksquare$
\end{preuve}

\addcontentsline{toc}{subsection}{Cas particulier}
\begin{remarque}[{Cas
particulier}]\phantomsection\label{II-cas-particulier}

L'espérance et la variance sont des cas particuliers de ce résultat.
\end{remarque}

\addcontentsline{toc}{subsection}{Inégalité de Jensen}
\begin{proposition}[{Inégalité de
Jensen}]\phantomsection\label{II-jensen}

Soit \(X\) une variable aléatoire réelle intégrable de loi
\(\mathbb{P}_X\), et soit \(g : \mathbb{R}\to \mathbb{R}\) une fonction
convexe telle que \(g(X)\) soit \(\mathbb{P}_X\) - intégrable. On a
alors \[ g(\mathbb{E}(X)) \leq \mathbb{E}(g(X)).\]
\end{proposition}

\begin{preuve}
Puisque \(g\) est convexe, pour tout \(a \in \mathbb{R}\) il existe
\(\lambda_a \in \mathbb{R}\) tel que pour tout \(x\in\mathbb{R}\) on a
\[g(x) \geq g(a) + \lambda_a\,(x-a).\] C'est une conséquence directe de
la caractérisation de la convexité par les inégalités des pentes. C'est
vrai en particulier pour \(x =X(\omega)\), \(\omega \in \Omega\), et
\(a = \mathbb{E}(X)\)~: pour tout \(\omega \in \Omega\),
\[g\left(X(\omega)\right) \geq g\left(\mathbb{E}(X)\right) + \lambda_{\mathbb{E}(X)}\,\left(X(\omega) - \mathbb{E}(X)\right).\]
En intégrant de chaque côté de l'inégalité, on obtient bien
\begin{align*}
\int_\Omega g\left(X(\omega)\right) \,\mathbb{P}(d\omega) = \mathbb{E}\left(g(X)\right) &\geq g\left(\mathbb{E}(X)\right) + \lambda_{\mathbb{E}(X)}\,\left(\int_\Omega X(\omega)\,\mathbb{P}(d\omega) - \mathbb{E}(X)\right)\\
& = g\left(\mathbb{E}(X)\right) + \lambda_{\mathbb{E}(X)}\,\left(\mathbb{E}(X) - \mathbb{E}(X)\right)\\
& = g\left(\mathbb{E}(X)\right).
\end{align*}\hfill$\blacksquare$
\end{preuve}

\section{Exemples de loi à
densité}\label{II-exemples-de-loi-uxe0-densituxe9}

Nous donnons ici quelques exemples de densités de probabilité. Nous
reprenons en particulier les
\href{Probabilité\%20I.pdf\%20\#exampledens}{trois exemples de densités
donnés au premier cours}~:

\subsection{\emph{Loi uniforme}}\label{II-loi-uniforme}

sur \([a,b]\), où \(a < b\) et on note \(X \sim \mathcal{U}_{[a,b]}\) si
\(X\) est de densité \[ f(x) = \frac{1}{b-a} 1_{[a,b]} (x).\]

\runinhead{Exercice -- Moments d'une variable aléatoire de loi uniforme
sur
\([a,b]\) (\(\mathord{\bullet}\))}\phantomsection\addcontentsline{toc}{subsubsection}{Moments
d'une variable aléatoire de loi uniforme sur
\([a,b]\)}\label{II-moments-unif}
Soit \(X \sim \mathcal{U}_{[a,b]}\). Montrer que son espérance vaut
\[ \mathbb{E}(X) = \int_a^b \frac{x}{b-a} dx = \frac{a+b}{2}\] et que sa
variance vaut
\[ \mathbb{V}(X) = \mathbb{E}(X^2) - \mathbb{E}(X)^2 = \frac{(b-a)^2}{12}.\]

\solref{II-answer-moments-unif}

\subsection{\emph{Loi exponentielle}}\label{II-loi-exponentielle}

de paramètre \(\theta > 0\) et on note \(X \sim \mathcal{E}(\theta)\) si
\(X\) est de densité \[ f(x) = \theta e^{-\theta x} 1_{\{x>0\}} (x).\]

\runinhead{Exercice -- Moments d'une v.a.
exponentielle (\(\mathord{\bullet}\))}\phantomsection\addcontentsline{toc}{subsubsection}{Moments
d'une v.a. exponentielle}\label{II-moments-expo}
Montrer que son espérance et sa variance valent
\[ \mathbb{E}(X) = \frac{1}{\theta} \text{ et } \mathbb{V}(X) = \frac{1}{\theta^2}\]

\solref{II-answer-moments-expo}

\subsection{\emph{Loi gamma}}\label{II-loi-gamma}

On rappelle tout d'abord que la fonction gamma est définie pour
\(\alpha \in \left] 0, + \infty \right[\) par
\[\Gamma(\alpha) = \int_0^{+\infty} x^{\alpha-1}e^{-x} dx.\] En
intégrant par partie, on obtient la relation
\(\Gamma(\alpha+1) = \alpha \Gamma(\alpha)\), et on a \(\Gamma(1) = 1\).
On en déduit que \(\Gamma(n) = (n-1)!\).

Une variable aléatoire \(X\) suit une loi gamma de paramètre d'échelle
\(\theta\) et d'indice \(\alpha\) et on note
\(X \sim \Gamma(\alpha,\theta)\), si sa loi admet la densité
\[f(x) = \frac{1}{\Gamma(\alpha)}\theta^\alpha x^{\alpha -1} e^{-\theta x}.\]

\runinhead{Exercice -- Moments d'une v.a. de loi
gamma (\(\mathord{\bullet}\))}\phantomsection\addcontentsline{toc}{subsubsection}{Moments
d'une v.a. de loi gamma}\label{II-moments-gamma}
Soit \(X\sim \Gamma(\alpha,\theta)\). Montrer que son espérance et sa
variance valent~:

\[ \mathbb{E}(X) = \frac{\alpha}{\theta} \text{ et } \mathbb{V}(X) = \frac{\alpha}{\theta^2}.\]

On remarquera que \(\Gamma(1,\theta)\) est la loi exponentielle de
paramètre \(\theta\).

Lorsque \(\alpha\) est entier, la loi gamma permet de modéliser le temps
d'attente avant la \(n\)-ième occurence d'événements indépendants de loi
exponentielle de paramètre \(\theta\).\solref{II-answer-moments-gamma}

\subsection{\emph{Loi normale}}\label{II-loi-normale}

de paramètres \(\mu\) et \(\sigma^2\) et on note
\(X \sim \mathcal{N}(\mu,\sigma^2)\) si \(X\) est de densité
\[f(x) = \frac{1}{\sqrt{2\pi}\sigma}\exp\left(-\frac{(x-\mu)^2}{2\sigma^2}\right)\]
Son espérance et sa variance valent
\[\mathbb{E}(X) =  \mu \text{ et } \mathbb{V}(X) = \sigma^2\]

\runinhead{Exercice -- Moments d'une v.a.
gaussienne (\(\mathord{\bullet}\))}\phantomsection\addcontentsline{toc}{subsubsection}{Moments
d'une v.a. gaussienne}\label{II-moments-gauss}
Après avoir vérifié que \(f\) est bien une densité, calculer l'espérance
et la variance de
\(X \sim \mathcal{N}(\mu,\sigma^2)\)\solref{II-answer-moments-gauss}

\addcontentsline{toc}{subsection}{Estimation statistique}
\begin{remarque}[{Estimation
statistique}]\phantomsection\label{II-estimation-statistique}

Dans les exemples ci-dessus, on peut remarquer que les densités sont
paramétrées par des nombres réels qui sont liés directement aux valeurs
de l'espérance et de la variance de la variable. C'est très utile en
statistique où l'on cherchera à estimer la valeur de ces paramètres à
partir des observations disponibles. Dans le cas de la loi normale, la
moyenne et la variance (empiriques) des échantillons fourniront ainsi
directement des estimateurs des paramètres.

Il existe des variables aléatoires qui n'ont pas d'espérance, comme le
montre l'exercice suivant.
\end{remarque}

\runinhead{Exercice -- \emph{Loi de
Cauchy} (\(\mathord{\bullet}\mathord{\bullet}\))}\phantomsection\addcontentsline{toc}{subsubsection}{\emph{Loi
de Cauchy}}\label{II-exo-cauchy}
Un gyrophare envoie un flash lumineux dans une direction aléatoire
uniforme d'angle \(\theta\). On cherche la distribution de l'abscisse
\(X\) du point d'impact du rayon lumineux sur un écran plan infini situé
à distance 1 du gyrophare. Donner la densité de \(X\). Calculer son
espérance.\solref{II-answer-exo-cauchy}

\chapter{Rappels sur les intégrales
multiples}\label{II-rappels-sur-les-intuxe9grales-multiples}

Dans la suite, nous aurons besoin de quelques résultats à propos des
intégrales multiples que nous reformulons ici dans le contexte des
probabilités.

\addcontentsline{toc}{subsection}{Théorème de Fubini (pour les
probabilités)}
\begin{theoreme}[{Théorème de Fubini (pour les
probabilités)}]\phantomsection\label{II-fubiniproba}

Soient \(\mathbb{P}_X\) et \(\mathbb{P}_Y\) deux probabilités définies
respectivement sur \((\mathbb{R}^n, \mathcal{B}(\mathbb{R}^n))\) et
\((\mathbb{R}^m, \mathcal{B}(\mathbb{R}^m))\).

\begin{enumerate}
\def\labelenumi{\arabic{enumi}.}
\tightlist
\item
  Soient \(A \in \mathcal{B}(\mathbb{R}^n)\) et
  \(B \in \mathcal{B}(\mathbb{R}^m)\), alors
  \[\mathbb{P}(A \times B) = \mathbb{P}_X(A)\mathbb{P}_Y(B)\] définit de
  manière unique une probabilité sur
  \((\mathbb{R}^{n+m}, \mathcal{B}(\mathbb{R}^{n+m}))\) (que l'on note
  aussi \(\mathbb{P}_X \otimes \mathbb{P}_Y\)).
\item
  Soit \(f\) une fonction \(\mathcal{B}(\mathbb{R}^{n+m})\)-mesurable
  positive ou \(\mathbb{P}_X \otimes \mathbb{P}_Y\)-intégrable, alors la
  fonction \(x \mapsto \int f(x, y)\mathbb{P}_Y(dy)\) est
  \(\mathcal{B}(\mathbb{R}^{n})\)-mesurable, la fonction
  \(y \mapsto \int f (x, y) \mathbb{P}_X(dx)\) est
  \(\mathcal{B}(\mathbb{R}^m)\)-mesurable et
  \[\int f d\mathbb{P}_X\otimes\mathbb{P}_Y = \int \left(\int f(x,y) \mathbb{P}_Y(dy) \right) \mathbb{P}_X(dx) = \int \left(\int f(x,y) \mathbb{P}_X(dx) \right) \mathbb{P}_Y(dy)\]
\end{enumerate}
\end{theoreme}

\runinhead{Exercice -- Calcul de l'aire d'un
triangle (\(\mathord{\bullet}\))}\phantomsection\addcontentsline{toc}{subsubsection}{Calcul
de l'aire d'un triangle}\label{II-triangle}
Considérons le triangle \[
T = \{(x, y) \in \mathbb{R}^2 \; | \; x \geq 0, \, y \geq 0 \mbox{ et } x + y \leq  1\}.
\] En supposant l'intégrale ci-dessous bien définie, calculer~: \[
a(T) = \int_{\mathbb{R}^2} 1_T(x, y) \, dxdy.
\]

\solref{II-answer-triangle}

\runinhead{Exercice -- Contre-exemple (\(\mathord{\bullet}\mathord{\bullet}\))}\phantomsection\addcontentsline{toc}{subsubsection}{Contre-exemple}\label{II-fubini-counter-example}
Comparer les valeurs des intégrales multiples \[
\int_0^1 \left[\int_0^1 \frac{x^2 - y^2}{(x^2+y^2)^2} \, dx \right] \, dy
\; \mbox{ et } \;
\int_0^1 \left[\int_0^1 \frac{x^2 - y^2}{(x^2+y^2)^2} \, dy \right] \, dx,
\] puis expliquer le résultat. Indication~: on remarquera que \[
\frac{x^2 - y^2}{(x^2+y^2)^2} = \frac{\partial}{\partial x} \left(-\frac{x}{x^2+y^2}\right) 
\] \solref{II-answer-fubini-counter-example}

Pour pouvoir appliquer le théorème de Fubini, il faut savoir a priori
que la fonction est intégrable et pas seulement que ses intégrales
multiples sont bien définies
(cf.~\protect\hyperref[II-fubini-counter-example]{``Contre-exemple''
ci-dessus (p.~\pageref*{II-fubini-counter-example})}). Si la fonction
est à valeurs positives toutefois, l'examen de ses intégrales multiples
permet de s'assurer de l'intégrabilité\,; c'est le théorème de
Fubini-Tonelli vu dans le cours de calcul intégral.

Autrement dit, en pratique, on appliquera Fubini-Tonelli à la valeur
absolue de la fonction pour vérifier que celle-ci est intégrable, puis
Fubini, s'il est applicable, pour effectuer le calcul de l'intégrale.

\addcontentsline{toc}{subsection}{Changement de variables}
\begin{theoreme}[{Changement de
variables}]\phantomsection\label{II-theorem-changement-de-variables}

Soient \(D_1\) et \(D_2\) des ouverts de \(\mathbb{R}^n\) et
\(h: D_1 \to D_2\) un \(C^1\)-difféomorphisme de \(D_1\) sur \(D_2\),
c'est-à-dire une fonction continûment différentiable et bijective dont
l'inverse \(h^{-1}: D_2 \to D_1\) est également continûment
différentiable. La matrice jacobienne associée à la différentielle de
\(h\) étant notée \(J_h\), la fonction \(f: D_2 \to \mathbb{R}\) est
intégrable si et seulement si la fonction
\((f \circ h) |\det J_h| : D_1 \to \mathbb{R}\) est intégrable et dans
ce cas, \[
\int_{D_2} f(y) \, dy = \int_{D_1} f(h(x)) |\det J_h(x)| \, dx.
\]
\end{theoreme}

\begin{figure}
\centering
\fbox{[figure manquante : changement-de-variables.svg.pdf]}
\caption{Changement de
variables}\label{II-fig-changement-de-variables.svg}
\end{figure}

\runinhead{Exercice -- Homothétie (\(\mathord{\bullet}\))}\phantomsection\addcontentsline{toc}{subsubsection}{Homothétie}\label{II-h}
Soit \(f:\mathbb{R}^n \to \mathbb{R}\) une fonction intégrable. Montrer
que pour tout coefficient \(\alpha > 0\), l'intégrale \[
\int_{\mathbb{R}^n} f(\alpha x) dx
\] est bien définie et la calculer en fonction de l'intégrale de \(f\)
sur \(\mathbb{R}^n\).\solref{II-answer-h}

\runinhead{Exercice -- Volume et
translation (\(\mathord{\bullet}\))}\phantomsection\addcontentsline{toc}{subsubsection}{Volume
et translation}\label{II-vr}
Montrer que si \(A\) est mesurable et de volume fini dans
\(\mathbb{R}^3\) l'image de \(A\) par une translation est également
mesurable et de même volume.\solref{II-answer-vr}

\runinhead{Exercice -- Coordonnées
polaires (\(\mathord{\bullet}\mathord{\bullet}\))}\phantomsection\addcontentsline{toc}{subsubsection}{Coordonnées
polaires}\label{II-cp}
Soit \[
C = \{(x, y) \in \mathbb{R}^2 \; | \; y \neq 0 \mbox{ ou } x > 0\}
\; \mbox{ et } \;
P = \{(r,\theta) \in \mathbb{R}^2 \; | \; r>0 \mbox{ et } -\pi < \theta < \pi\}.\]
On note \(h\) la fonction de \(P\) dans \(C\) définie par
\(h(r, \theta) = (r \cos \theta, r \sin \theta)\). Montrer que pour
toute fonction \(f: C \to \mathbb{R}\) intégrable, si l'on pose
\(g(r,\theta) = f(x, y)\) où \((x, y) = h(r,\theta)\), alors \[
\int_{C} f(x, y) \, d(x,y) = \int_{P} g(r,\theta)  r \, d(r, \theta).
\]

\solref{II-answer-cp}

\runinhead{Exercice -- Absence du déterminant
jacobien (\(\mathord{\bullet}\mathord{\bullet}\))}\phantomsection\addcontentsline{toc}{subsubsection}{Absence
du déterminant jacobien}\label{II-adj}
Supposons \(D_1\), \(D_2\), \(h\) et \(f\) conformes aux hypothèses
\protect\hyperref[II-theorem-changement-de-variables]{du théorème de
changement de
variables (p.~\pageref*{II-theorem-changement-de-variables})}. On
suppose de plus que \(f \circ h\) est intégrable sur \(D_1\). Exprimer
l'intégrale \[
\int_{D_1} f (h(x)) \, dx
\] comme une intégrale sur \(D_2\).\solref{II-answer-adj}

\chapter{Vecteurs aléatoires}\label{II-vecteurs-aluxe9atoires-1}

Nous allons généraliser ici la notion de variable aléatoire en
considérant qu'elle peut prendre ses valeurs dans \(\mathbb{R}^n\). Les
vecteurs aléatoires se rencontrent naturellement lorsqu'on s'intéresse à
plusieurs quantités conjointement, par exemple dans le cas de la météo,
la température, la pluviométrie et la vitesse et la direction du vent.

\section{Définitions}\label{II-duxe9finitions}

Une variable aléatoire \(X\) à valeurs dans \(\mathbb{R}^n\) (ou vecteur
aléatoire) est simplement une collection de \(n\) variables réelles
définies sur le même espace probabilisé
\((\Omega, \mathcal{A}, \mathbb{P})\), qui sont les \emph{composantes}
de \(X\). On écrit \(X = (X_1,\ldots,X_n)\).

De même qu'en dimension 1, la loi de \(X\) est caractérisée par la
fonction de répartition multi-dimensionnelle
\(F : \mathbb{R}^n \to \mathbb{R}\) définie par
\[F(x_1,\ldots,x_n) = \mathbb{P}_X(X_1\leq x_1,\ldots,X_n\leq x_n)\]
Cependant, caractériser les fonctions de répartition sur
\(\mathbb{R}^n\) est délicat, de sorte que cette notion est rarement
utilisée. Nous allons plus particulièrement nous intéresser aux vecteurs
aléatoires à densité. On gardera tout de même à l'esprit que les lois de
probabilités définies sur \((\mathbb{R}^n,\mathcal{B}(\mathbb{R}^n))\)
n'admettent pas toutes une densité\,!

\addcontentsline{toc}{subsection}{Vecteur aléatoire à densité}
\begin{definition}[{Vecteur aléatoire à
densité}]\phantomsection\label{II-defvect}

On dit que \(X\) admet la densité \(f\) si la fonction réelle \(f\) sur
\(\mathbb{R}^n\) est positive, intégrable et vérifie
\[\int_{\mathbb{R}^n} f(x) dx = \int_{-\infty}^{+\infty} \ldots \int_{-\infty}^{+\infty} f(x_1,\ldots,x_n) dx_1 \ldots dx_n= 1\]
et si
\[\mathbb{P}_X(X_1\leq x_1,\ldots,X_n\leq x_n) = \int_{-\infty}^{x_1} \ldots \int_{-\infty}^{x_n} f(x_1,\ldots,x_n) dx_1 \ldots dx_n.\]

De la même manière que dans le \protect\hyperref[II-esperanceg]{cas
unidimensionnel (p.~\pageref*{II-esperanceg})}, on a~:
\end{definition}

\addcontentsline{toc}{subsection}{Fonction d'un vecteur aléatoire}
\begin{proposition}[{Fonction d'un vecteur
aléatoire}]\phantomsection\label{II-esperancegvect}

Soit \(X\) un vecteur aléatoire de densité \(f\), et soit \(g\) une
fonction de \(\mathbb{R}^n\) dans \(\mathbb{R}\), mesurable. On a alors
\(g(X)\in \mathcal{L}^1\) si et seulement si
\[\int_{-\infty}^{+\infty} \ldots \int_{-\infty}^{+\infty} |g(x_1,\ldots,x_n)|f(x_1,\ldots,x_n) dx_1 \ldots dx_n,\]
est finie et dans ce cas, on a
\[\mathbb{E}(g(X)) = \int_{-\infty}^{+\infty} \ldots \int_{-\infty}^{+\infty} g(x_1,\ldots,x_n)f(x_1,\ldots,x_n) dx_1 \ldots dx_n.\]

Pour revenir à la densité d'une composante d'un vecteur aléatoire, on
intègre par rapport aux autres variables. On le présente ici dans le cas
d'un couple \(Z=(X,Y)\) de variables aléatoires. On généralise aisément
à une dimension supérieure.
\end{proposition}

\addcontentsline{toc}{subsection}{Densités marginales}
\begin{proposition}[{Densités
marginales}]\phantomsection\label{II-loimarg}

Supposons que \(Z\) admette une densité \(f\). Alors \(X\) et \(Y\)
admettent les densités \(f_X\) et \(f_Y\) données par
\[f_X(x) = \int_\mathbb{R}f(x,y) dy,\,\,\, f_Y(y) = \int_\mathbb{R}f(x,y) dx.\]

Les fonctions \(f_X\) et \(f_Y\) s'appellent les \emph{densités
marginales} de \(f\).
\end{proposition}

\begin{preuve}
Pour tout \(x \in \mathbb{R}\), on a par définition
\[ \mathbb{P}(X \leq x) = \mathbb{P}(Z \in ]-\infty,x] \times \mathbb{R}) = \mathbb{E}(1_{]-\infty,x] \times \mathbb{R}}(Z)) = \int_{-\infty}^x du \left(\int_\mathbb{R}f(u,v) dv\right),\]
où l'on a utilisé le théorème de Fubini-Tonelli pour
\(1_{]-\infty,x] \times \mathbb{R}}(u,v)f(u,v)\) qui est bien positive.
Donc si \(f_X\) est définie par \(f_X(x) = \int_\mathbb{R}f(x,y) dy\),
on obtient que \(\mathbb{P}(X \leq x) = \int_{-\infty}^x f_X(u) du\), ce
qui montre que \(f_X\) est la densité de \(X\). Le raisonnement est
analogue pour \(Y\).\hfill$\blacksquare$
\end{preuve}

\addcontentsline{toc}{subsection}{Réciproque\,?}
\begin{remarque}[{Réciproque\,?}]\phantomsection\label{II-ruxe9ciproque}

La réciproque de cette proposition est fausse en revanche~: les
variables aléatoires \(X\) et \(Y\) peuvent avoir des densités sans que
le couple \(Z = (X, Y )\) en ait une.

Supposons par exemple que \(X = Y\). Si
\(\Delta = \{(x,x) ; x\in \mathbb{R}\}\) est la diagonale de
\(\mathbb{R}^2\), nous avons évidemment \(\mathbb{P}_Z(\Delta) = 1\)
mais si la \protect\hyperref[II-esperancegvect]{proposition
précédente (p.~\pageref*{II-esperancegvect})} était valide pour
\(\mathbb{P}_Z\), on aurait
\(\mathbb{P}_Z(\Delta) = \mathbb{E}(1_\Delta) = \int_{\mathbb{R}^2} 1_\Delta f_Z(z)dz = 0\)
car \(\Delta\) est de volume nul dans \(\mathbb{R}^2\).

En particulier, il faut faire attention au fait que dans le cas général,
la densité d'un couple de variables aléatoires n'est pas le produit des
densités.
\end{remarque}

\runinhead{Exercice -- Lancer de
fléchette (\(\mathord{\bullet}\mathord{\bullet}\))}\phantomsection\addcontentsline{toc}{subsubsection}{Lancer
de fléchette}\label{II-exo-flechette}
On lance une fléchette sur une cible circulaire de rayon unité. Le
joueur est suffisamment maladroit pour que le point \(M\) d'impact de la
fléchette soit supposé uniformément distribué sur la cible (On décide de
n'observer que les lancés qui atteignent la cible). Exprimer la densité
du couple formé par les coordonnées cartésiennes \((X,Y)\) de \(M\) puis
leurs densités marginales.\solref{II-answer-exo-flechette}

\section{Moments d'un vecteur
aléatoire}\label{II-moments-dun-vecteur-aluxe9atoire}

\addcontentsline{toc}{subsection}{Vecteur espérance et covariance}
\begin{definition}[{Vecteur espérance et
covariance}]\phantomsection\label{II-vecteur-espuxe9rance-et-covariance}

Si les composantes \(X_i\) du vecteur aléatoire \(X = (X_1,\ldots,X_n)\)
sont intégrables, on peut définir le \emph{vecteur espérance}
\[\mathbb{E}(X) = (\mathbb{E}(X_1),\ldots,\mathbb{E}(X_n))\] Si les
composantes \(X_i\) du vecteur aléatoire \(X = (X_1,\ldots,X_n)\) sont
de carré intégrable, la \emph{matrice de covariance} de \(X\) est la
matrice \(C_X = (c_{i,j})_{1 \leq i \leq n , 1 \leq j \leq n}\) de
taille \(n \times n\) et dont les éléments valent
\[c_{i,j} = \text{Cov}(X_i,X_j)\]
\end{definition}

\addcontentsline{toc}{subsection}{Matrice de covariance}
\begin{proposition}[{Matrice de
covariance}]\phantomsection\label{II-matrice-de-covariance}

La matrice de covariance est symétrique non-négative (ou encore
semi-définie positive).
\end{proposition}

\begin{preuve}
La symétrie est évidente. Non-négative signifie que pour tous réels
\(a_1,\ldots,a_n\), on a
\(\sum_{i=1}^n \sum_{j=1}^n a_i a_j c_{i,j} \geq 0\). Un calcul simple
montre que
\[ \sum_{i=1}^n \sum_{j=1}^n a_i a_j c_{i,j} = \mathbb{V}(\sum_{i=1}^n a_i X_i) \geq 0.\]

\hfill$\blacksquare$
\end{preuve}

\addcontentsline{toc}{subsection}{Vecteur Gaussien \(n\)-dimensionel}
\begin{exemple}[{Vecteur Gaussien
\(n\)-dimensionel}]\phantomsection\label{II-vecteur-gaussien-n-dimensionel}

Un exemple important de vecteurs aléatoires est celui des vecteurs
gaussiens, que nous étudierons plus en détail au chapitre suivant.
Soient \(m \in \mathbb{R}^n\) et \(C\) une matrice symétrique définie
positive (c'est-à-dire telle que pour tout \(x \in \mathbb{R}^n\) non
identiquement nul \(x^tCx > 0\) où \(^t\) désigne la transposée). Le
vecteur \(X \in \mathbb{R}^n\) est un vecteur aléatoire gaussien
d'espérance \(m\) et de matrice de covariance \(C\) si sa densité
s'écrit
\[ f(x) = \frac{1}{(2\pi)^{n/2}\sqrt{\det (C)}}\exp (-\frac{1}{2}(x-m)^tC^{-1}(x-m)) \]
On a alors \(\mathbb{E}(X) = m\) et \(C_X =C\).
\end{exemple}

\begin{figure}
\centering
\includegraphics{PdfGauss3D.tex.pdf}
\caption{Densité de probabilité de la loi normale bivariée centrée,
réduite, de coefficient de corrélation 1/2}\label{II-fig-PdfGauss3D.tex}
\end{figure}

\chapter{Variables aléatoires
indépendantes}\label{II-variables-aluxe9atoires-induxe9pendantes}

Lorsqu'on modélise plusieurs variables conjointement, une hypothèse
importante est celle de l'indépendance. Ce caractère traduit l'absence
de lien de causalité entre les variables. Par exemple, on fait
naturellement l'hypothèse d'indépendance lorsque l'on considère une
répétition d'une même expérience dans les mêmes conditions.

Dans ce paragraphe, on considère un couple \((X,Y)\) de vecteurs
aléatoires respectivement à valeurs dans \(\mathbb{R}^m\) et
\(\mathbb{R}^n\). Les résultats s'étendent sans peine à une famille
finie quelconque.

On peut se ramener aux évènements pour caractériser l'indépendance de
deux vecteurs aléatoires. En effet, considérons le vecteur aléatoire
\(Z = (X,Y)\), \(A\) et \(B\) deux ensembles dans
\(\mathcal{B}({\mathbb{R}^m})\) et \(\mathcal{B}({\mathbb{R}^n})\). On a
vu que les évènements \(X\in A\) et \(Y \in B\) sont indépendants si et
seulement si
\(\mathbb{P}_Z(X \in A, Y \in B) = \mathbb{P}(X^{-1}(A) \cap Y^{-1}(B)) = \mathbb{P}(X^{-1}(A))\mathbb{P}(Y^{-1}(B)) = \mathbb{P}_X(X \in A)\mathbb{P}_Y(Y \in B)\).
Pour que deux vecteurs aléatoires soient indépendants, on va donc
demander que ceci soit valable quels que soient \(A\) et \(B\).

\addcontentsline{toc}{subsection}{Vecteurs aléatoires indépendants}
\begin{definition}[{Vecteurs aléatoires
indépendants}]\phantomsection\label{II-defvai}

Les vecteurs aléatoires \(X\) et \(Y\) sont \emph{indépendants} si pour
tous ensembles \(A\) et \(B\) des tribus correspondantes,
\[\mathbb{P}(X \in A, Y \in B) = \mathbb{P}(X \in A)\mathbb{P}(Y \in B)\]
\end{definition}

En termes de mesures de probabilité, cette égalité s'écrit
\[\mathbb{P}_{(X,Y)}(A \times B) = \mathbb{P}_X(A)\mathbb{P}_Y(B) = (\mathbb{P}_X \otimes \mathbb{P}_Y) (A \times B).\]
Autrement dit, \(X\) et \(Y\) sont indépendantes si et seulement si la
mesure de probabilité du couple \((X,Y)\) est la mesure produit
\(P_X \otimes P_Y\) dont l'existence et l'unicité sont garanties par le
théorème de
\protect\hyperref[II-fubiniproba]{Fubini (p.~\pageref*{II-fubiniproba})}.

Cette définition se traduit dans le cas à densité dans la proposition
suivante que l'on énonce sans perte de généralité pour un couple de
variables aléatoires.

\addcontentsline{toc}{subsection}{Caractérisation de l'indépendance --
cas à densité}
\begin{proposition}[{Caractérisation de l'indépendance -- cas à
densité}]\phantomsection\label{II-caractuxe9risation-de-linduxe9pendance-cas-uxe0-densituxe9}

Soient \(X\) et \(Y\) deux variables aléatoires réelles de densités
\(f_X\) et \(f_Y\). \(X\) et \(Y\) sont indépendantes si et seulement si
le couple \(Z = (X,Y)\) a pour densité (sur \(\mathbb{R}^2\))~:
\[f_Z(x,y) = f_X(x)f_Y(y).\]
\end{proposition}

\begin{preuve}
S'il y a indépendance, la
\protect\hyperref[II-defvai]{définition (p.~\pageref*{II-defvai})}
implique
\[P(X\leq x, Y\leq y) = P(X\leq x) \mathbb{P}(Y\leq y) = \int_{-\infty}^x f_X(u) du \int_{-\infty}^y f_Y(v) dv\]
ce qui montre que \(\mathbb{P}_Z\) vérifie la
\protect\hyperref[II-defvect]{définition d'un vecteur aléatoire à
densité (p.~\pageref*{II-defvect})} avec \(f_Z=f_X f_Y\).

Inversement, si \(f_Z=f_X f_Y\), on a pour tous \(A\), \(B\) de
\(\mathcal{B}({\mathbb{R}})\)
\[\mathbb{P}(X\in A, Y\in B) = \int_A \int_B f_X(x) f_Y(y) dxdy = \mathbb{P}(X \in A)\mathbb{P}(Y \in B)\]
où on a utilisé le théorème de Fubini-Tonelli.\hfill$\blacksquare$
\end{preuve}

\runinhead{Exercice -- Lancer de fléchette
(suite) (\(\mathord{\bullet}\))}\phantomsection\addcontentsline{toc}{subsubsection}{Lancer
de fléchette (suite)}\label{II-exo-flechettebis}
Les coordonnées cartésiennes \((X,Y)\) de \(M\) sont elles
indépendantes\,?\solref{II-answer-exo-flechettebis}

Considérons maintenant deux fonctions mesurables \(g\) et \(h\) définies
sur \(\mathbb{R}^m\) et \(\mathbb{R}^n\).

\addcontentsline{toc}{subsection}{Fonctions de vecteurs aléatoires
indépendants}
\begin{proposition}[{Fonctions de vecteurs aléatoires
indépendants}]\phantomsection\label{II-indep_fct}

Avec les notations précédentes, si \(X\) et \(Y\) sont indépendantes de
densités respectives \(f_X\) et \(f_Y\), les variables aléatoires
\(g(X)\) et \(h(Y)\) sont aussi indépendantes. Si de plus \(g(X)\) et
\(h(Y)\) sont intégrables, alors le produit \(g(X)h(Y)\) est aussi
intégrable, et on a
\[ \mathbb{E}(g(X)h(Y)) = \mathbb{E}(g(X))\mathbb{E}(h(Y))\]
\end{proposition}

\begin{preuve}
La première assertion est évidente par définition de l'indépendance. Par
ailleurs, si \(g(X)\) et \(h(Y)\) sont intégrables, en notant
\(f_{(X,Y)}\) la densité du couple \((X,Y)\), et en utilisant le
théorème de Fubini, on a \begin{align*}
\mathbb{E}(g(X)h(Y))  & = \int_{\mathbb{R}^{m+n}} g(x)h(y) f_{(X,Y)}(x,y) dx dy \\
                & = \int_{\mathbb{R}^m} \int_{\mathbb{R}^n} g(x)h(y)f_X(x) f_Y(y) dx dy \\
                & = \left(\int_{\mathbb{R}^m}  g(x) f_X(x) dx \right) \left(\int_{\mathbb{R}^n} h(y) f_Y(y) dy\right)\\
                & = \mathbb{E}(g(X))\mathbb{E}(h(Y))
\end{align*}\hfill$\blacksquare$
\end{preuve}

\addcontentsline{toc}{subsection}{Cas général}
\begin{remarque}[{Cas
général}]\phantomsection\label{II-cas-guxe9nuxe9ral}

Ce résultat est encore valable dans le cas général (sans densité). Il
suffit de remplacer \(f_{(X,Y)}(x,y) dx dy\) par
\(d\mathbb{P}_{X,Y}(x,y) = d\mathbb{P}_X(x)d\mathbb{P}_Y(y)\) dans la
démonstration précédente.

On déduit de ce résultat et de la
\protect\hyperref[II-defcov]{définition de la
covariance (p.~\pageref*{II-defcov})} que~:
\end{remarque}

\addcontentsline{toc}{subsection}{Covariance de variables aléatoires
indépendantes}
\begin{corollaire}[{Covariance de variables aléatoires
indépendantes}]\phantomsection\label{II-covariance-de-variables-aluxe9atoires-induxe9pendantes}

Si les variables aléatoires réelles \(X\) et \(Y\) sont indépendantes et
de carré intégrable, alors \(\text{Cov}(X,Y) = 0\) et \(\rho(X,Y) = 0\).
\end{corollaire}

\runinhead{Exercice -- Démonstration (\(\mathord{\bullet}\))}\phantomsection\addcontentsline{toc}{subsubsection}{Démonstration}\label{II-demo-cov-indep}
Démontrer ce résultat.\solref{II-answer-demo-cov-indep}

\subsection{Remarque}\label{II-remarque}

Attention, la réciproque est fausse. Par exemple, si
\(X \sim \mathcal{U}_{[-1,1]}\) et \(Y = X^2\). \(X\) et \(Y\) ne sont
clairement pas indépendantes mais on a

\[\text{Cov}(X,Y) = \text{Cov}(X,X^2) = \mathbb{E}(X^3) - \mathbb{E}(X)\mathbb{E}(X^2) = 0\]

\chapter{Identification de
densité}\label{II-identification-de-densituxe9-1}

Un problème important est le suivant. Soit \(X\) une variable aléatoire
réelle, admettant la densité \(f_X\). Soit \(g\) une fonction mesurable,
de sorte que \(Y = g(X)\) est aussi une variable aléatoire. Est-ce que
\(Y\) admet une densité, et si oui, comment la calculer\,?

On peut déjà remarquer que cette densité n'existe pas toujours. Si par
exemple \(g(x) = a\) pour tout \(x\), la loi de \(Y\) est la mesure de
Dirac en \(a\), qui n'a pas de densité par rapport à la mesure de
Lebesgue.

Pour résoudre ce problème, l'idée consiste à essayer de mettre
\(E(h(Y )) = E(h \circ g(X))\) sous la forme \(\int h(y)f_Y (y)dy\) pour
une fonction convenable \(f_Y\), et une classe de fonctions \(h\)
suffisamment grande. La fonction \(f_Y\) sera alors la densité cherchée.

La \protect\hyperref[II-esperanceg]{proposition qui assure que \(g(X)\)
est intégrable pour une variable aléatoire réelle
\(X\) (p.~\pageref*{II-esperanceg})} implique
\[ \mathbb{E}(h(Y)) = \mathbb{E}(h \circ g (X)) = \int_{\mathbb{R}} h \circ g(x) f_X(x) dx\]

et on fait le changement de variable \(y = g(x)\) dans cette intégrale.
Cela nécessite que \(g\) soit dérivable et bijective ``par morceaux'',
et il faut faire très attention aux domaines où \(g\) est croissante ou
décroissante. Puisque la fonction \(h\) est arbitraire, on appelle
couramment cette technique la \emph{méthode de la fonction muette}.
Cette approche résulte en fait de la proposition suivante~:

\addcontentsline{toc}{subsection}{Méthode de la fonction muette}
\begin{proposition}[{Méthode de la fonction
muette}]\phantomsection\label{II-muxe9thode-de-la-fonction-muette}

Si il existe une fonction \(f\) positive telle que pour toute fonction
mesurable \(h\) telle que \(h(x) f(x)\) soit intégrable,
\[\mathbb{E}(h(X)) = \int_\mathbb{R}h(x) f(x) dx\] alors la loi de \(X\)
admet la densité \(f\) par rapport à la mesure de Lebesgue.
\end{proposition}

\begin{preuve}
Supposons qu'un tel \(f\) existe. Soit
\(B \in \mathcal{B}(\mathbb{R})\), alors \(1_B\) est mesurable. En
particulier, si \(B = ]-\infty,x], ~x\in\mathbb{R}\), on a
\[\mathbb{P}(X\leq x) = \mathbb{P}(X \in B) = \mathbb{P}_X(B) = \mathbb{E}(1_B(X)) = \int_{- \infty}^x f(y)dy\]
Par ailleurs, prenant \(B = \mathbb{R}\), on a
\(1 = \mathbb{P}(X\in \mathbb{R}) = \int_\mathbb{R}f(y) dy\). Par
conséquent, \(\mathbb{P}_X\) admet la densité \(f\).\hfill$\blacksquare$
\end{preuve}

Nous donnons ici quelques exemples d'application de cette méthode~:

\begin{itemize}
\tightlist
\item
  Soit \(Y = aX + b\), où \(a\) et \(b\) sont des constantes. Si
  \(a=0\), alors \(Y = b\) et la loi de \(Y\) est la masse de Dirac en
  \(b\) (sans densité). Si \(a \neq 0\), on fait le changement de
  variable \(y = ax+b\), ce qui donne
  \[ \mathbb{E}(h(Y)) = \int_\mathbb{R}h(ax+b) f_X(x) dx = \int_\mathbb{R}h(y) f_X(\frac{y-b}{a})\frac{1}{|a|} dy \]
  Donc \[ f_Y(y) = f_X(\frac{y-b}{a})\frac{1}{|a|} \]
\item
  Soit \(Y =X^2\). La fonction \(g\) est décroissante sur
  \(\mathbb{R}_-\) et croissante sur \(\mathbb{R}_+\). Le changement de
  variable \(y = x^2\) donne alors \begin{align*}
       \mathbb{E}(h(Y)) &= \int_{-\infty}^0 h(x^2) f_X(x) dx + \int_0^{+\infty} h(x^2) f_X(x) dx\\
                  &= \int_0^{+\infty} h(y) f_X(-\sqrt{y})\frac{1}{2\sqrt{y}} dy + \int_0^{+\infty} h(y) f_X(\sqrt{y})\frac{1}{2\sqrt{y}}dy
       \end{align*} et on en déduit
  \[f_Y(y) = (f_X(-\sqrt{y})+f_X(\sqrt{y}))\frac{1}{2\sqrt{y}}1_{\left]0,+\infty\right[}\]
\end{itemize}

Dans le cas des vecteurs aléatoires, l'idée est la même. Soit
\(X = (X_1,\ldots,X_n)\), un vecteur aléatoire de densité \(f_X\) sur
\(\mathbb{R}^n\), \(g\) une fonction de \(\mathbb{R}^n\) dans
\(\mathbb{R}^m\) et \(Y = g(X)\). Plusieurs cas sont à considérer~:

\begin{enumerate}
\def\labelenumi{\arabic{enumi}.}
\tightlist
\item
  \(m > n\), le vecteur \(Y\) n'admet pas de densité.
\item
  \(m=n\), on utilise comme dans le cas unidimensionel le changement de
  variable \(y = g(x)\) dans
  \[ \mathbb{E}(h(Y)) = \mathbb{E}(h \circ g (X)) = \int_{\mathbb{R}^n} h \circ g(x) f_X(x) dx \]
  Supposons d'abord que \(g\) soit une bijection continûment
  différentiable de \(A\) dans \(B\), ouverts de \(\mathbb{R}^n\). Sous
  l'hypothèse que \(h \circ g(x) f_X(x)\) soit intégrable, le
  \href{Calcul\%20Intégral\%20III.pdf}{théorème de changement de
  variable} nous assure~:
  \[ \int_A h\circ g (x) f_X(x) dx = \int_B h(y) f_X \circ g^{-1}(y) \frac{1}{|\det J_g (g^{-1}(y))|} dy,\]
  où \(J_g\) désigne la matrice de Jacobi associée à la différentielle
  de \(g\). Dans le cas où \(f_X(x) = 0\) en dehors de \(A\), on obtient
  que \(Y\) admet la densité
  \[ f_Y(y) = 1_B(y)f_X \circ g^{-1}(y)\frac{1}{|\det J_g (g^{-1}(y))|}.\]
  Lorsque \(g\) est simplement continûment différentiable, il existe
  souvent une partition finie \((A_i)_{1\leq i \leq n}\) de l'ensemble
  \(\{x ; f(x) >0\}\), telle que \(g\) soit bijective sur chaque
  \(A_i\). On note alors \(B_i = g(A_i)\) l'image de \(A_i\) par \(g\).
  On découpe alors l'intégrale selon les \(A_i\), on applique la formule
  précédente à chaque morceau et on somme pour obtenir~:
  \[ f_Y(y) = \sum_{i=1}^n 1_{B_i}(y)f_X \circ g^{-1}(y) \frac{1}{|\det J_g (g^{-1}(y))|},\]
  où \(g^{-1}\) est bien définie sur chaque \(B_i\) comme image
  réciproque de la restriction de g à \(A_i\).
\item
  \(m < n\), on commence par ``compléter'' \(Y\), en essayant de
  construire une application \(g'\) de \(\mathbb{R}^n\) dans
  \(\mathbb{R}^n\) dont les \(m\) premières composantes coïncident avec
  les composantes de \(g\) et pour laquelle on peut appliquer l'une des
  deux formules précédentes. On obtient ainsi la densité \(f_Y'\) de
  \(Y' = g'(X)\) puis on obtient la densité de \(Y\) en calculant sa loi
  marginale~:
  \[f_Y(y_1,\ldots,y_m) = \int_{\mathbb{R}^{n-m}} f_{Y'}(y_1,\ldots,y_m,y_{m+1},\ldots y_n) dy_{m+1}\ldots dy_n.\]
\end{enumerate}

\addcontentsline{toc}{subsection}{\textbf{Coordonnées polaires}}
\begin{exemple}[{\textbf{Coordonnées
polaires}}]\phantomsection\label{II-polar}

Soit \(X = (U,V)\) un vecteur aléatoire de \(\mathbb{R}^2\), et
\(Y = (R,\Theta)\) ses coordonnées polaires. La transformation \(g\) est
un difféomorphisme de \(A = \mathbb{R}^2\setminus\{0\}\) dans
\(B = \left]0,+\infty\right[\times \left]0,2\pi\right]\), et son inverse
\(g^{-1}\) s'écrit~: \(u = r\cos \theta,\, v= r\sin \theta\). Le
Jacobien de \(g^{-1}\) au point \((r,\theta)\) vaut \(r\), donc le point
2. ci-dessus entraîne que
\[ f_Y(r,\theta) = r f_X(r\cos\theta,r\sin\theta)1_B(r,\theta)\] Par
exemple, si \(U\) et \(V\) sont indépendantes et de loi
\(\mathcal{N}(0,1)\), on a
\(f_X(u,v) = \frac{1}{2\pi}\exp\left(-\frac{u^2+v^2}{2}\right)\) et donc
\[ f_Y(r,\theta) = \frac{1}{2\pi} r\exp\left(-\frac{r^2}{2}\right)1_{\left]0,+\infty\right[ }(r)1_{\left]0,2\pi\right]}(\theta).\]
En particulier, on remarque que \(R\) et \(\Theta\) sont indépendantes
de densités respectives
\(r\exp\left(-\frac{r^2}{2}\right)1_{\left]0,+\infty\right[}(r)\) et
\(1_{\left]0,2\pi\right]}(\theta)\).
\end{exemple}

\runinhead{Exercice -- Lancer de fléchette
(fin) (\(\mathord{\bullet}\))}\phantomsection\addcontentsline{toc}{subsubsection}{Lancer
de fléchette (fin)}\label{II-exo-flechetteter}
Donner la densité du couple \((R,\Theta)\) des coordonnées polaires de
\(M\). Sont-elles indépendantes\,?\solref{II-answer-exo-flechetteter}

\addcontentsline{toc}{subsection}{\textbf{Somme de deux variables
aléatoires indépendantes}}
\begin{exemple}[{\textbf{Somme de deux variables aléatoires
indépendantes}}]\phantomsection\label{II-sum}

Soient \(U\) et \(V\) indépendantes et admettant les densités \(f_U\) et
\(f_V\), on cherche la densité de la somme \(Z = U+V\). On commence par
compléter \(Z\) en le couple \(T = (U,Z)\) (par exemple), correspondant
à la bijection \(g(u,v) = (u, u+v)\) sur \(\mathbb{R}^2\) dont le
jacobien est 1 et d'inverse \(g^{-1}(x,y) = (x,x-y)\). Appliquant le
point 3. ci-dessus, on obtient~:
\[f_Z(z) = \int_{\mathbb{R}}f_U(u)f_V(z-u) du = \int_{\mathbb{R}}f_U(z-v)f_V(v)dv.\]
La fonction \(f_Z\) est appelée \emph{le produit de convolution} des des
fonctions \(f_U\) et \(f_V\).
\end{exemple}

\runinhead{Exercice -- Loi
bêta (\(\mathord{\bullet}\mathord{\bullet}\mathord{\bullet}\))}\phantomsection\addcontentsline{toc}{subsubsection}{Loi
bêta}\label{II-exo-loibeta}
Soit \(X = (U,V)\) un vecteur aléatoire de \(\mathbb{R}^2\), avec \(U\)
et \(V\) indépendantes de lois \(\Gamma(\alpha,\theta)\) et
\(\Gamma(\beta,\theta)\). Identifier la densité de
\(Y = \frac{U}{U+V}\), puis de
\(Z = U+V\).\solref{II-answer-exo-loibeta}

\chapter{Modélisation statistique}\label{II-moduxe9lisation-statistique}

On a vu un certain nombre d'exemples de lois de probabilité à densité.
Celles-ci font intervenir des paramètres qui permettent d'ajuster la
forme de la densité aux données observées. Par exemple, la loi normale
fait intervenir deux paramètres~: l'espérance et la variance. Dans cette
section, nous formalisons l'approche statistique qui consiste à choisir
une loi de probabilité pour modéliser un phénomène aléatoire, puis à
estimer les paramètres de cette loi à partir d'observations.

La première hypothèse que nous faisons est que le phénomène étudié peut
être modélisé par une variable aléatoire \(X\) de densité \(f\)
inconnue. La première étape de la modélisation consiste à choisir une
famille de densités \((f_\theta)_{\theta \in \Theta}\), où \(\Theta\)
est un ensemble de paramètres. Par exemple, pour la loi normale, on peut
choisir \(\Theta = \mathbb{R}\times \mathbb{R}_+^*\) et
\(f_{(m,\sigma)}\) la densité de la loi \(\mathcal{N}(m,\sigma^2)\).

La seconde étape consiste à estimer le paramètre \(\theta\) de manière à
ce que la densité \(f_\theta\) soit la plus proche possible de la
densité \(f\) inconnue. On a donc besoin d'une notion de proximité entre
deux densités. Plusieurs choix sont possibles, mais le plus courant est
celui de la divergence de Kullback-Leibler.

\addcontentsline{toc}{subsection}{Divergence de Kullback-Leibler}
\begin{definition}[{Divergence de
Kullback-Leibler}]\phantomsection\label{II-defKL}

Soient \(f\) et \(g\) deux densités de probabilité sur \(\mathbb{R}\),
telles que \(g(x) > 0\) pour tout \(x\) tel que \(f(x) > 0\). La
\emph{divergence de Kullback-Leibler} de \(g\) par rapport à \(f\) est
définie par
\[ D_{KL}(f||g) = \int f(x) \log\left(\frac{f(x)}{g(x)}\right) dx.\]
\end{definition}

\begin{remarque}\phantomsection\label{II-remarque-1}
La divergence de Kullback-Leibler n'est pas une distance au sens
mathématique du terme. En effet, elle n'est pas symétrique et ne
satisfait pas l'inégalité triangulaire. En revanche, on a toujours
\(D_{KL}(f||g) \geq 0\), avec égalité si et seulement si \(f = g\)
presque partout.
\end{remarque}

\addcontentsline{toc}{subsection}{Positivité de la divergence de
Kullback-Leibler}
\begin{proposition}[{Positivité de la divergence de
Kullback-Leibler}]\phantomsection\label{II-corKL}

Soient \(f\) et \(g\) deux densités de probabilité sur \(\mathbb{R}\),
telles que \(g(x) > 0\) pour tout \(x\) tel que \(f(x) > 0\). Alors
\[ D_{KL}(f||g) \geq 0,\] avec égalité si et seulement si \(f = g\)
presque partout.
\end{proposition}

\begin{preuve}
La fonction \(\log\) est strictement concave, donc \(-\log\) est
strictement convexe. On peut donc appliquer
\protect\hyperref[II-jensen]{l'inégalité de
Jensen (p.~\pageref*{II-jensen})}~: \begin{align*}
D_{KL}(f||g) & = \int_\mathbb{R}f(x) \log\left(\frac{f(x)}{g(x)}\right) dx = - \int_\mathbb{R}f(x) \left(-\log\left(\frac{g(x)}{f(x)}\right)\right) dx \\
& = - \mathbb{E}\left(-\log\left(\frac{g(X)}{f(X)}\right)\right) \geq - \left(-\log\left(\mathbb{E}\left(\frac{g(X)}{f(X)}\right)\right)\right) = -\log(1) = 0.
\end{align*} On a égalité si et seulement si \(\frac{g(x)}{f(x)}\) est
presque sûrement constant, ce qui revient à dire que \(f = g\) presque
partout.\hfill$\blacksquare$
\end{preuve}

\runinhead{Exercice -- Divergence de Kullback-Leibler entre deux
gaussiennes}\phantomsection\addcontentsline{toc}{subsubsection}{Divergence
de Kullback-Leibler entre deux gaussiennes}\label{II-exoklgauss}
Soient \(f\) et \(g\) les densités de deux lois normales
\(\mathcal{N}(m_1,\sigma_1^2)\) et \(\mathcal{N}(m_2,\sigma_2^2)\).
Montrer que
\(D_{KL}(f||g) = \log \left(\frac{\sigma_2}{\sigma_1}\right) + \frac{\sigma_1^2 + (\mu_1 - \mu_2)^2}{2\sigma_2^2} - \frac{1}{2}\).\solref{II-answer-exoklgauss}

Revenons à notre problème d'estimation de paramètres. Soit \(X\) une
variable aléatoire de densité inconnue \(f\). On choisit une famille de
densités \((f_\theta)_{\theta \in \Theta}\), et on cherche à minimiser
la divergence de Kullback-Leibler entre \(f\) et \(f_\theta\)~:
\[ \hat{\theta} = \arg\min_{\theta \in \Theta} D_{KL}(f||f_\theta).\] On
peut réécrire la divergence de Kullback-Leibler sous la forme suivante~:
\begin{align*}
D_{KL}(f||f_\theta) & = \int f(x) \log\left(\frac{f(x)}{f_\theta(x)}\right) dx \\
& = \int f(x) \log(f(x)) dx - \int f(x) \log(f_\theta(x)) dx \\
& = -H(f) - \mathbb{E}_f(\log f_\theta(X)),
\end{align*} où \(H(f) = -\int f(x) \log f(x) dx\) est
l'\textbf{entropie} de la densité \(f\). Celle-ci s'interprète comme Ce
terme ne dépend pas de \(\theta\), donc minimiser
\(D_{KL}(f||f_\theta)\) revient à maximiser
\(\mathbb{E}_f(\log f_\theta(X))\). En statistique, \(f_\theta(X)\) est
appelée la \emph{vraisemblance} de l'observation \(X\) pour le paramètre
\(\theta\), et \(\log f_\theta(X)\) la \emph{log-vraisemblance}.
L'estimateur \(\hat{\theta}\) est ainsi appelé \emph{estimateur du
maximum de vraisemblance}.

En pratique, on dispose d'un échantillon \(X_1,\ldots,X_n\) de variables
aléatoires indépendantes et identiquement distribuées selon la loi de
densité \(f\). Le choix de la forme paramétrique est généralement guidé
par des considérations théoriques, portant sur la physique du phénomène
étudié, ou empiriques, en regardant par exemple l'histogramme des
données. On peut ensuite estimer \(\mathbb{E}_f(\log f_\theta(X))\) par
la moyenne empirique \[ \frac{1}{n} \sum_{i=1}^n \log f_\theta(X_i),\]
et l'estimateur du maximum de vraisemblance s'écrit
\[ \hat{\theta}_n = \arg\max_{\theta \in \Theta} \frac{1}{n} \sum_{i=1}^n \log f_\theta(X_i).\]
Cela revient à remplacer \(f\) par la mesure empirique de l'échantillon
\(\tilde{p} = \frac{1}{n} \sum_{i=1}^n \delta_{X_i}\) dans la définition
de la divergence de Kullback-Leibler.

\chapter{Exercices
complémentaires}\label{II-exercices-compluxe9mentaires}

\section{Durée de vie}\label{II-duruxe9e-de-vie}

La durée de vie, exprimée en années, d'un circuit électronique est une
variable aléatoire \(T\) dont la fonction de répartition \(F\) est
définie par~: \[F (t) = (1 - e^{-t^2 /2} )1_{t\geq 0}\]

\runinhead{Question
1}\phantomsection\addcontentsline{toc}{subsubsection}{Question
1}\label{II-dureevie1}
Donner la densité de probabilité \(f\) de \(T\) . Calculer
\(\mathbb{E}(T)\).\solref{II-answer-dureevie1}

\runinhead{Question
2}\phantomsection\addcontentsline{toc}{subsubsection}{Question
2}\label{II-dureevie2}
Sachant que le circuit a déjà fonctionné durant 1 an, quelle est la
probabilité qu'il continue à fonctionner encore durant au moins 2 ans\,?
La loi est-elle sans mémoire\,?

Un équipement électronique E est composé de 10 circuits identiques et
indépendants. Au circuit \(i (1 \leq i \leq 10)\) est associée la
variable aléatoire \(X_i\) , avec \(X_i = 1\) si la durée de vie du
circuit \(i\) est inférieure à un an et \(X_i = 0\)
sinon.\solref{II-answer-dureevie2}

\runinhead{Question
3}\phantomsection\addcontentsline{toc}{subsubsection}{Question
3}\label{II-dureevie3}
Quelle est la loi du nombre \(N\) de circuits de E dont la durée de vie
est inférieure à 1 an\,?\solref{II-answer-dureevie3}

\runinhead{Question
4}\phantomsection\addcontentsline{toc}{subsubsection}{Question
4}\label{II-dureevie4}
L'équipement E est dit en série si la défaillance de l'un de ses
circuits entraîne sa défaillance. Quelle est alors la probabilité qu'il
soit défaillant avant 1 an\,?\solref{II-answer-dureevie4}

\runinhead{Question
5}\phantomsection\addcontentsline{toc}{subsubsection}{Question
5}\label{II-dureevie5}
L'équipement E est dit en parallèle si sa défaillance ne peut se
produire que si tous ses circuits sont défaillants. Quelle est alors la
probabilité qu'il soit défaillant avant 1
an\,?\solref{II-answer-dureevie5}

\section{Crues centennales de la
Seine}\label{II-crues-centennales-de-la-seine}

On suppose que la hauteur d'eau maximale au cours de l'année \(n\) est
décrite par une variable aléatoire \(X_n\) de densité \(f(x)\) et de
fonction de répartition \(F(x)\), identique pour toutes les \(X_n\), que
l'on suppose également indépendantes.

\runinhead{Question
1}\phantomsection\addcontentsline{toc}{subsubsection}{Question
1}\label{II-crue1}
Calculer la fonction de répartition du maximum sur \(n\) années
\(Y = \max(X_1,\ldots,X_n)\).\solref{II-answer-crue1}

\runinhead{Question
2}\phantomsection\addcontentsline{toc}{subsubsection}{Question
2}\label{II-crue2}
En déduire la densité de \(Y\).\solref{II-answer-crue2}

\runinhead{Question
3}\phantomsection\addcontentsline{toc}{subsubsection}{Question
3}\label{II-crue3}
Expliciter la densité de \(Y\) lorsque les \(X_i\) sont de loi uniforme
sur \([0,1]\).\solref{II-answer-crue3}

\runinhead{Question
4}\phantomsection\addcontentsline{toc}{subsubsection}{Question
4}\label{II-crue4}
Selon mêmes hypothèses, calculer la hauteur minimale des quais pour que
la probabilité de voir la Seine déborder sur une période de \(n\) années
soit inférieure à 1/1000.\solref{II-answer-crue4}

\runinhead{Question
5}\phantomsection\addcontentsline{toc}{subsubsection}{Question
5}\label{II-crue5}
Toujours selon les mêmes hypothèses, calculer la médiane de \(Y\). La
comparer avec celle des \(X_i\). Commenter.\solref{II-answer-crue5}

\section{Indépendance et vecteurs
gaussiens}\label{II-induxe9pendance-et-vecteurs-gaussiens}

Soit \(Z := (X,Y)\) un vecteur gaussien (i.e.~qui suit une loi Normale
bi-variée) d'espérance \(m\) et de matrice de covariance définie
positive \(C\). Notons \(f_Z\) sa densité.

\runinhead{Question
1}\phantomsection\addcontentsline{toc}{subsubsection}{Question
1}\label{II-covindepgauss-joint}
Rappeler la forme de \(f_Z\) et en donner une expression non matricielle
faisant apparaître le coefficient de corrélation noté
\(\rho\).\solref{II-answer-covindepgauss-joint}

\runinhead{Question
2}\phantomsection\addcontentsline{toc}{subsubsection}{Question
2}\label{II-covindepgauss-marg}
Expliciter les lois marginales de \(X\) et de
\(Y\).\solref{II-answer-covindepgauss-marg}

\runinhead{Question
3}\phantomsection\addcontentsline{toc}{subsubsection}{Question
3}\label{II-covindepgauss-indep}
Montrer que \(X\) et \(Y\) sont indépendantes ssi
\(\text{Cov}\left(X,Y\right) = 0\).\solref{II-answer-covindepgauss-indep}

\section{Symétrie de la
gaussienne}\label{II-symuxe9trie-de-la-gaussienne}

Soit \(X\) une variable aléatoire réelle de loi Normale centrée réduite,
dont la densité est notée \(f\) et la fonction de répartition \(F\).
Pour tout \(c>0\) on pose
\[X_c := \left|\begin{array}{ll} X & \text{si } |X| > c,\\ -X & \text{si } |X| \leq c.\end{array}\right.\]

\runinhead{Question
1}\phantomsection\addcontentsline{toc}{subsubsection}{Question
1}\label{II-gausssym-loi}
Déterminer la loi de \(X_c\).\solref{II-answer-gausssym-loi}

\runinhead{Question
2}\phantomsection\addcontentsline{toc}{subsubsection}{Question
2}\label{II-gausssym-cov}
Calculer la covariance de \(X\) et de \(X_c\) et montrer qu'il existe
une valeur \(c_0\) de \(c\) qui l'annule.\solref{II-answer-gausssym-cov}

\runinhead{Question
3}\phantomsection\addcontentsline{toc}{subsubsection}{Question
3}\label{II-gausssym-indep}
Les deux variables \(X\) et \(X_{c_0}\) sont-elles
indépendantes\,?\solref{II-answer-gausssym-indep}

\runinhead{Question
4}\phantomsection\addcontentsline{toc}{subsubsection}{Question
4}\label{II-gausssym-vecgauss}
Le vecteur aléatoire \((X,X_{c_0})\) est-il
gaussien\,?\solref{II-answer-gausssym-vecgauss}

\section{Loi du \(\chi^2\)}\label{II-loi-du-chi2}

On considère une variable aléatoire \(X\) gaussienne centrée réduite. On
note \(f\) sa densité et \(F\) sa fonction de répartition.

\runinhead{A 1 degré de
liberté}\phantomsection\addcontentsline{toc}{subsubsection}{A 1 degré de
liberté}\label{II-khi2-1dl}

\begin{enumerate}
\def\labelenumi{\arabic{enumi}.}
\tightlist
\item
  Calculer la densité de la variable aléatoire \(Y := X^2\). En donner
  une expression faisant apparaître la fonction gamma (on rappelle que
  \(\Gamma\left(\frac{1}{2}\right)=\sqrt{\pi}\)).
\end{enumerate}

On dit que \(Y\) suit une \emph{loi du \(\chi^2\) à \(1\) degré de
liberté}, et on note \(Y \sim \chi^2\).\solref{II-answer-khi2-1dl}

\runinhead{A n degrés de
liberté}\phantomsection\addcontentsline{toc}{subsubsection}{A n degrés
de liberté}\label{II-khi2-ndl}
Soient maintenant \(n\in\mathbb{N}^\ast\) et \(X_1,\dots,X_n\) des
copies indépendantes de \(X\). Pour tout \(i \in \{1,\dots,n\}\) on pose
\(Y_i := X_i^2\).

\begin{enumerate}
\def\labelenumi{\arabic{enumi}.}
\setcounter{enumi}{1}
\tightlist
\item
  Montrer que la variable aléatoire \(Y := \sum_{i = 1}^n Y_i\) admet
  pour densité
  \[ f_Y : x\in\mathbb{R}\mapsto \left|\begin{array}{ll} \dfrac{x^{\frac{n}{2} - 1}}{2^{\frac{n}{2}}\,\Gamma\left(\frac{n}{2}\right)}\,\exp\left\{ -\dfrac{x}{2} \right\} & \text{si } x> 0,\\ 0 & \text{si } x \leq 0. \end{array}\right.\]
  On pourra procéder par récurrence sur \(n\).
\end{enumerate}

On dit que \(Y\) suit une \emph{loi du \(\chi^2\) à \(n\) degrés de
liberté} et on note \(Y \sim \chi^2_n\).\solref{II-answer-khi2-ndl}

\section{Combinaisons linéaires de variables Gaussiennes
indépendantes}\label{II-combinaisons-linuxe9aires-de-variables-gaussiennes-induxe9pendantes}

Soit \(X\) une variable aléatoire de loi Normale centrée réduite, dont
on note \(f\) la densité et \(F\) la fonction de répartition.

\runinhead{Préliminaires}\phantomsection\addcontentsline{toc}{subsubsection}{Préliminaires}\label{II-CLIGauss-pre}

\begin{enumerate}
\def\labelenumi{\arabic{enumi}.}
\tightlist
\item
  Rappeler la loi de la variable aléatoire \(s\,X + m\), où
  \(s, m \in \mathbb{R}\).
\end{enumerate}

\solref{II-answer-CLIGauss-pre}

\runinhead{Combinaisons
linéaires}\phantomsection\addcontentsline{toc}{subsubsection}{Combinaisons
linéaires}\label{II-CLIGauss-cl}
Soient maintenant \(n\in\mathbb{N}^\ast\) variables aléatoires
\(X_1,\dots,X_n\) indépendantes et de même loi que \(X\) (on dit que ce
sont des \emph{copies indépendantes} de \(X\)). Pour tout vecteur
\(a = (a_1,\dots,a_n) \in (\mathbb{R}^{*})^n\) on pose
\(S_n^a := \sum_{i = 1}^n a_i\,X_i\).

\begin{enumerate}
\def\labelenumi{\arabic{enumi}.}
\setcounter{enumi}{1}
\item
  Montrer que pour \(S_n^a\) suit une loi Normale d'espérance nulle et
  de variance \(\sum_{i = 1}^n a_i^2\) (on pourra raisonner par
  récurrence sur \(n\)).
\item
  Soient \(a,b\in(\mathbb{R}^{*})^n\). Sous quelle condition a-t-on
  \(\text{Cov}\left(S_n^a, S_n^b\right) = 0\)\,?
\end{enumerate}

\solref{II-answer-CLIGauss-cl}

\section{Loi de Maxwell}\label{II-loi-de-maxwell}

On désire déterminer la distribution des vitesses des molécules d'un gaz
monoatomique parfait à l'équilibre (loi de Maxwell (1859)).

On représente la vitesse d'une molécule d'un gaz monoatomique parfait à
l'équilibre dans un repère orthonormal par un vecteur aléatoire
\(V = (V_1 , V_2 , V_3 )\). Le choix du repère étant arbitraire, il est
naturel de supposer que la loi de \(V\) est invariante par rotation
(autour de l'origine) et que les composantes de \(V\) sont
indépendantes.

\subsection{Partie 1}\label{II-partie-1}

Soit \((X, Y, Z)\) un vecteur aléatoire à valeurs dans \(\mathbb{R}^3\)
de densité \(f_{(X, Y, Z)}\). Ce vecteur aléatoire est supposé invariant
par rotation~: il existe une fonction
\(\phi: \left[0, +\infty\right[ \to \left]0, +\infty\right[\)
(strictement positive) telle que \[
f_{(X, Y, Z)}(x,y,z) = \phi(x^2+y^2+z^2) \;  \mbox{ pour tout $(x, y, z) \in \mathbb{R}^3$.}
\]

\runinhead{Question
1}\phantomsection\addcontentsline{toc}{subsubsection}{Question
1}\label{II-Maxwell-1}
On suppose de plus les variables aléatoires \(X\), \(Y\) et \(Z\)
indépendantes. Montrer qu'il existe une fonction
\(f: \mathbb{R}\to \left[0, +\infty\right[\) qui soit une densité telle
que \[
f_{(X, Y, Z)}(x, y, z) = f(x) f(y) f(z) \; \mbox{ pour tout $(x, y, z) \in \mathbb{R}^3$.} 
\] Montrer finalement que \(f(x) > 0\) pour tout
\(x \in \mathbb{R}\).\solref{II-answer-Maxwell-1}

\runinhead{Question
2}\phantomsection\addcontentsline{toc}{subsubsection}{Question
2}\label{II-Maxwell-2}
On suppose de plus que les densités marginales de \(X, Y\) et \(Z\)
ainsi que la fonction \(\phi\) sont continûment différentiables
(c'est-à-dire de classe \(C^1\)). Montrer que la densité de chacune des
composantes s'écrit sous la forme \(f (x) = a e^{cx^2/2}\), avec \(a\)
et \(c\) deux constantes réelles.\solref{II-answer-Maxwell-2}

\runinhead{Question
3}\phantomsection\addcontentsline{toc}{subsubsection}{Question
3}\label{II-Maxwell-3}
En déduire que le vecteur \((X, Y, Z)\) suit une loi gaussienne
d'espérance l'origine dont on précisera la matrice de covariance en
fonction de la constante
\(\sigma = 1/\sqrt{|c|}\).\solref{II-answer-Maxwell-3}

\subsection{Partie 2}\label{II-partie-2}

On suppose que le vecteur aléatoire \((X, Y, Z) = V\) vérifie les
hypothèses des questions précédentes.

\runinhead{Question
4}\phantomsection\addcontentsline{toc}{subsubsection}{Question
4}\label{II-Maxwell-4}
Calculer l'énergie cinétique moyenne d'un atome du gaz, c'est-à-dire
l'espérance \[
E_c := \mathbb{E}\left(\frac{1}{2}m\|V\|^2\right)
\] où \(m\) est la masse d'un atome du gaz. L'énergie cinétique moyenne
d'un atome du gaz de masse \(m\) étant égale à \(\frac{3}{2} kT\) où
\(k\) est la constante de Boltzmann et \(T\) la température du gaz, en
déduire la valeur de \(\sigma^2\) en fonction de \(k\), \(T\) et
\(m\).\solref{II-answer-Maxwell-4}

\runinhead{Question
5}\phantomsection\addcontentsline{toc}{subsubsection}{Question
5}\label{II-Maxwell-5}
On rappelle que si \(X\) et \(Y\) sont deux variables aléatoires
indépendantes de loi respective \(\Gamma (a,\lambda)\) et
\(\Gamma (b,\lambda)\), alors la loi de \(X + Y\) est la loi
\(\Gamma (a+b,\lambda)\). On rappelle également que la densité \(g\) de
la loi \(\Gamma (a,\lambda)\) vérifie \[
g(x) = \frac{1}{\Gamma(a)}\lambda^a x^{a -1} e^{-\lambda x} 1_{]0,+\infty[}(x)
\; \mbox{ pour tout $x\in\mathbb{R}$},
\] où \(\Gamma\) est la fonction définie par
\[\Gamma(a) = \int_0^{+\infty} x^{a-1}e^{-x} dx \text{ pour } a \in \left] 0, + \infty \right[.\]
On remarquera que \(\Gamma(1/2) = \sqrt{\pi}\) et que
\(\Gamma(x+1) = x \Gamma(x).\)

Calculer la loi de \(V_1^2\). En déduire la loi de \(\|V \|^2\) puis la
densité de \(\|V \| = \sqrt{V_1^2 + V_2^2 + V_3^2}\). La probabilité
associée est appelée loi de Maxwell.

\begin{center}\rule{0.5\linewidth}{0.5pt}\end{center}

\solref{II-answer-Maxwell-5}


\part{Conditionnement}
\chapter{Objectifs d'apprentissage}\label{III-objectifs-dapprentissage}

Cette section s'efforce d'expliciter et de hiérarchiser les acquis
d'apprentissages associés au chapitre. Ces objectifs sont organisés en
paliers~:

($\mathord{\boldsymbol{\circ}}$) Prérequis ($\mathord{\bullet}$)
Fondamental ($\mathord{\bullet}\mathord{\bullet}$) Standard
($\mathord{\bullet}\mathord{\bullet}\mathord{\bullet}$) Avancé
($\mathord{\bullet}\mathord{\bullet}\mathord{\bullet}\mathord{\bullet}$)
Expert

Sauf mention particulière, la connaissance des démonstrations du
document n'est pas exigible\footnote{l'étude des démonstrations du cours
  peut toutefois contribuer à votre apprentissage, au même titre que la
  résolution d'exercices.}

\runinhead{Lois
conditionnelles}\phantomsection\label{III-lois-conditionnelles}

\begin{itemize}
\tightlist
\item
  $\mathord{\bullet}\mathord{\bullet}$ connaître le théorème de Fubini
  conditionnel
\item
  $\mathord{\bullet}\mathord{\bullet}$ connaître les formules de
  balayage conditionnel
\item
  $\mathord{\bullet}\mathord{\bullet}$ savoir appliquer ces résultats
  pour différents types de loi de probabilité (à densité ou non)
\item
  $\mathord{\bullet}$ connaître le critère d'indépendance qui résulte du
  théorème de Fubini conditionnel
\end{itemize}

\runinhead{Espérance
conditionnelle}\phantomsection\label{III-espuxe9rance-conditionnelle}

\begin{itemize}
\tightlist
\item
  $\mathord{\bullet}$ connaître les deux points de la définition de
  l'espérance conditionnelle
\item
  $\mathord{\bullet}$ connaître les deux points de la définition de
  l'espérance conditionnelle d'une fonction de variables aléatoires
\item
  $\mathord{\bullet}$ connaître et savoir utiliser la formule de
  l'espérance totale
\item
  $\mathord{\bullet}\mathord{\bullet}$ savoir calculer la densité
  conditionnelle d'un certain nombre de composantes sachant les autres
  dans un vecteur gaussien à densité
\end{itemize}

\runinhead{Cas
\(\mathcal{L}^2\)}\phantomsection\label{III-cas-mathcall2}

\begin{itemize}
\tightlist
\item
  $\mathord{\bullet}\mathord{\bullet}$ savoir que la régression linéaire
  est la meilleure approximation linéaire (au sens des moindres carrés)
  d'une variable aléatoire par une autre
\item
  $\mathord{\bullet}\mathord{\bullet}\mathord{\bullet}$ savoir retrouver
  ce résultat
\item
  $\mathord{\bullet}\mathord{\bullet}\mathord{\bullet}$ connaître
  l'interprétation géométrique de l'espérance conditionnelle dans le cas
  \(\mathcal{L}^2\)
\item
  $\mathord{\bullet}$ connaître et savoir utiliser la formule de la
  variance totale
\end{itemize}

\chapter{Lois conditionnelles}\label{III-lois-conditionnelles-1}

\section{Introduction}\label{III-introduction}

On s'est consacré jusqu'à présent à l'étude de variables aléatoires
indépendantes. En pratique cependant, on rencontre souvent des variables
dépendant les unes des autres. Dans le cas de la météo, les variables
température, vitesse du vent et pression en fournissent un exemple. Dans
les approches bayésiennes, on résume l'information disponible sur l'état
du système étudié par la \textbf{loi a priori} et on met à jour notre
connaissance du système en incorporant de l'information supplémentaire
(par exemple des observations). On cherche alors à caractériser la
\textbf{loi a posteriori} de l'état du système, qui est la loi de l'état
sachant l'information supplémentaire. On va ainsi s'attacher dans ce
chapitre à décrire les \textbf{lois conditionnelles} qui vont permettre
de résumer l'information apportée par une variable (ou un vecteur) sur
une autre et s'intéresser en particulier à l'\textbf{espérance
conditionnelle} qui indiquera le comportement moyen d'une variable
conditionnellement à une autre. Ce dernier cas pose le cadre
probabiliste d'un des problèmes fondamentaux en apprentissage
statistique~: l'apprentissage supervisé, où on dispose d'un ensemble de
réalisations d'une variable dont on cherche à prédire le comportement à
partir d'un ensemble de variables dites explicatives (ou prédicteurs).

\section{Lois conditionnelles dans un
couple}\label{III-lois-conditionnelles-dans-un-couple}

Soient deux variables aléatoire \(X\) et \(Y\) définies sur le même
espace probabilisé \((\Omega, \mathcal{A}, \mathbb{P})\). Dans le cas où
\(X\) et \(Y\) sont indépendantes, on a vu que pour tous boréliens
\(B_1\) et \(B_2\) de \(\mathcal{B}(\mathbb{R})\), on a
\[\mathbb{P}(X\in B_1, Y\in B_2)= \mathbb{P}(X\in B_1)\mathbb{P}(Y\in B_2) = \mathbb{P}_X(B_1)\mathbb{P}_Y(B_2) = \int_{B_1}\mathbb{P}_Y(B_2)d\mathbb{P}_X(x),\]
où on a utilisé le
\href{Probabilité\%20II.pdf\%20\#fubiniproba}{théorème de Fubini}.

Du fait de l'indépendance, on a aussi
\(\mathbb{P}_Y(B_2) = \mathbb{P}(Y\in B_2) = \mathbb{P}(Y \in B_2 | X \in B_1) = \mathbb{P}_Y(B_2|X \in B_1)\)
ce qui exprime que pour tout borélien \(B_1\), la loi conditionnelle de
\(Y\) sachant \(X\in B_1\) est identique à la loi de \(Y\).

Lorsque \(X\) et \(Y\) en sont pas indépendantes, on va chercher à
établir une égalité de la forme
\[\mathbb{P}(X\in B_1, Y\in B_2) = \mathbb{P}_X(B_1)\mathbb{P}_Y(B_2 |X\in B_1) = \int_{B_1}\mathbb{P}_{Y|X=x}(B_2)d\mathbb{P}_X(x)\]
et s'intéresser à caractériser la \emph{loi conditionnelle de \(Y\)
sachant \(X=x\)}, que l'on notera donc \(\mathbb{P}_{Y|X=x}\).

De même, pour toute application \(g : \mathbb{R}^2 \to \mathbb{R}\)
mesurable telle que \(g(X,Y)\) admette une espérance (relativement à la
loi du couple \(\mathbb{P}_{X,Y}\)), on voudrait écrire~:
\[\mathbb{E}(g(X,Y)) = \int_{\mathbb{R}} \left( \int_{\mathbb{R}} g(x,y) d\mathbb{P}_{Y|X=x}(y)\right) d\mathbb{P}_X(x)\]

Pour bien fixer les idées, on va décrire spécifiquement les cas où \(X\)
est discrète puis où le couple \((X,Y)\) admet une densité avant
d'aborder le cas général.

\section{Cas où \(X\) est
discrète}\label{III-cas-ouxf9-x-est-discruxe8te}

Dans ce paragraphe, on suppose que la variable aléatoire réelle \(X\)
est discrète, c'est-à-dire que l'ensemble
\(X(\Omega) \subset \mathbb{R}\) des valeurs \(x_k\) prises par \(X\)
est au plus dénombrable.

On peut imposer que \(\forall x \in X(\Omega)\) on ait
\(\mathbb{P}(X=x) > 0\), quitte à modifier \(X\) sur un ensemble de
probabilité nulle. On va ainsi pouvoir utiliser la définition de la
probabilité conditionnelle pour des événements de la forme \(\{X =x\}\).
Ceci permet d'écrire pour tous boréliens \(B_1\) et \(B_2\) de
\(\mathbb{R}\)~: \begin{align*}
\mathbb{P}(X \in B_1, Y \in B_2) &= \sum_{x \in X(\Omega)\cap B_1} \mathbb{P}(X=x, Y\in B_2)\\
                         &= \sum_{x \in X(\Omega)\cap B_1} \mathbb{P}(X=x) \mathbb{P}(Y \in B_2 | X=x)\\
                         &= \int_{B_1} \mathbb{P}(Y \in B_2 | X=x) d\mathbb{P}_X(x)
\end{align*} puisque
\(\mathbb{P}_X = \sum_{x \in X(\Omega)} \mathbb{P}(X=x)\delta_x\). On
obtient ainsi l'écriture souhaitée en posant
\[\mathbb{P}_{Y|X=x}(B_2) = \mathbb{P}(Y \in B_2 | X=x),\,\,\,\forall x \in X(\Omega), \forall B_2\in\mathcal{B}(\mathbb{R}).\]

\addcontentsline{toc}{subsection}{Image de la probabilité
conditionnelle}
\begin{remarque}[{Image de la probabilité
conditionnelle}]\phantomsection\label{III-image-de-la-probabilituxe9-conditionnelle}

\(\mathbb{P}_{Y|X=x}\) ainsi définie est simplement la probabilité sur
\((\mathbb{R},\mathcal{B}(\mathbb{R}))\) image par \(Y\) de la
probabilité conditionnelle \(\mathbb{P}(\cdot|X=x)\) définie sur
\((\Omega,\mathcal{A})\), autrement dit, la \textbf{loi de \(Y\)
relative à \(\mathbb{P}(\cdot|X=x)\)} et non à \(\mathbb{P}\).

La formule ci-dessus s'écrit
\(\mathbb{P}_{X,Y}(B_1 \times B_2) = \int_{B_1} \mathbb{P}(Y \in B_2 | X=x) d\mathbb{P}_X(x)\),
où \(\mathbb{P}_{X,Y}\) est la loi du couple. Elle se généralise à tout
borélien \(B\) de \(\mathbb{R}^2\) de la manière suivante~: \[
\begin{aligned}
\mathbb{P}_{X,Y}(B) &= \mathbb{P}((X,Y)\in B) = \sum_{x \in X(\Omega)} \mathbb{P}(X=x, (x,Y) \in B) \\
      &= \sum_{x \in X(\Omega)} \mathbb{P}(X=x) \mathbb{P}((x,Y) \in B | X=x) \\
      &= \sum_{x \in X(\Omega)} \mathbb{P}(X=x) \mathbb{P}_{Y|X=x}(B_x),
\end{aligned}
\] où \(B_x = \{y\in \mathbb{R}, (x,y) \in B\}\). Ainsi, pour tout \(B\)
borélien de \(\mathbb{R}^2\),

\[\mathbb{E}(1_B(X,Y)) = \int_{\mathbb{R}^2} 1_B(x,y)d\mathbb{P}_{X,Y}(x y) = \int_\mathbb{R}\left(\int_\mathbb{R}1_B(x,y) d\mathbb{P}_{Y|X=x}(y)\right)  d\mathbb{P}_X(x)\]

Par linéarité de l'espérance, on peut ainsi exprimer l'espérance d'une
fonction étagée. Pour avoir le résultat pour une fonction mesurable
positive, on exprime celle-ci comme limite simple d'une suite croissante
de fonctions étagées, et on applique le théorème de convergence
monotone. Enfin, on applique cette construction à \(g_+\) et \(g_-\)
pour une fonction \(g\) de signe quelconque
\(\mathbb{P}_{X,Y}\)-intégrable. En d'autres termes, on reprend le
procédé de construction de l'intégrale de Lebesgue. On obtient ainsi la
formule souhaitée~:
\[\mathbb{E}(g(X,Y)) = \int_{\mathbb{R}} \left( \int_{\mathbb{R}} g(x,y) d\mathbb{P}_{Y|X=x}(y)\right) d\mathbb{P}_X(x).\]
\end{remarque}

\addcontentsline{toc}{subsection}{Pour fixer les idées (1)}
\begin{exemple}[{Pour fixer les idées
(1)}]\phantomsection\label{III-ex1}

Soit \(X \geq 0\) une variable aléatoire à valeurs dans \(\mathbb{N}\)
et \(Y\) une variable aléatoire réelle positive telle que la loi du
couple \(\mathbb{P}_{X,Y}\) vérifie pour tout \(n\in\mathbb{N}\) et tout
borélien \(B_2\) de \(\mathbb{R}\)~:
\[\mathbb{P}_{X,Y} (\{n\}\times B_2) = (1-\alpha)\alpha^n \int_{B_2 \cap \mathbb{R}_+^\ast}e^{-t}\frac{t^n}{n!}dt,\,\,\, 0 < \alpha <1\]
\(\mathbb{P}_{X,Y}\) est bien une probabilité sur \(\mathbb{R}^2\)
puisque par convergence monotone~: \begin{align*}
\mathbb{P}_{X,Y} (\mathbb{R}^2) &= \mathbb{P}_{X,Y}(\mathbb{N}\times \mathbb{R}) \\
                &= \sum_{n\in\mathbb{N}} \mathbb{P}_{X,Y} (\{n\}\times \mathbb{R})\\
                &= \sum_{n\in\mathbb{N}} (1-\alpha)\alpha^n \int_{\mathbb{R}_+^\ast}e^{-t}\frac{t^n}{n!}dt \\
                &= (1-\alpha)\int_{\mathbb{R}_+^\ast}e^{-t} \sum_{n\in\mathbb{N}} \frac{(\alpha t)^n}{n!}dt \\
                &= (1-\alpha)\int_{\mathbb{R}_+^\ast}e^{-(1-\alpha)t} dt = 1
\end{align*}\\
où on aura reconnu la loi exponentielle de paramètre \((1-\alpha)\).
\(\forall n \in \mathbb{N}\),
\[ \int_{\mathbb{R}_+^\ast}e^{-t}\frac{t^n}{n!}dt = \int_{\mathbb{R}_+^\ast}e^{-t}\frac{t^{(n-1)}}{(n-1)!}dt = \ldots = \int_{\mathbb{R}_+^\ast}e^{-t}dt =1\]
par intégrations par parties itérées. La loi marginale de \(X\) s'écrit
donc~:
\[\forall n \in \mathbb{N}, \mathbb{P}(X=n) = \mathbb{P}_{X,Y} (\{n\}\times \mathbb{R}_+^\ast) = (1-\alpha)\alpha^n,\]
loi géométrique de paramètre \((1-\alpha)\). On en déduit la loi
conditionnelle de \(Y\) sachant \(X = n\)~:
\[\mathbb{P}_{Y|X=n}(B_2) = \mathbb{P}(Y \in B_2 | X=n) = \frac{\mathbb{P}_{X,Y} (\{n\}\times B_2)}{\mathbb{P}(X=n)} = \int_{B_2 \cap \mathbb{R}_+^\ast} e^{-t}\frac{t^n}{n!}dt\]
et \(\mathbb{P}_{Y|X=n}\) est la donc la loi gamma de paramètre
\((n+1,1)\).
\end{exemple}

\section{Densités
conditionnelles}\label{III-densituxe9s-conditionnelles}

On suppose maintenant que le couple \((X,Y)\) admet une densité
\(f_{X,Y}\) (par rapport à la mesure de Lebesgue). On note
\(f_X(x) = \int_\mathbb{R}f_{X,Y}(x,y)dy\) (respectivement
\(f_Y(y) = \int_\mathbb{R}f_{X,Y}(x,y)dx\)) la loi marginale de \(X\)
(resp. de \(Y\)). On s'intéresse à caractériser la densité de la
variable \(Y\) connaissant la valeur prise par la variable \(X\), c'est
la \emph{densité conditionnelle} de \(Y\) sachant \(\{X = x\}\)~:

\addcontentsline{toc}{subsection}{Densité conditionnelle}
\begin{proposition}[{Densité
conditionnelle}]\phantomsection\label{III-defdenscond}

La formule suivante définit une densité sur \(\mathbb{R}\), pour tout
\(x \in \mathbb{R}\) tel que \(f_X(x) > 0\).
\[ f_{Y|X=x}(y) = \frac{f_{X,Y}(x,y)}{f_X(x)}.\] Cette fonction
s'appelle la \emph{densité conditionnelle de \(Y\) sachant
\(\{X = x\}\)}. La probabilité conditionnelle de \(Y\) sachant
\(\{X = x\}\) s'écrit ainsi \(\mathbb{P}_{Y|X=x} = f_{Y|X=x}\lambda\),
où \(\lambda\) représente la mesure de Lebesgue.
\end{proposition}

\begin{preuve}
La preuve est immédiate puisque \(f_{Y|X=x}\) est une fonction positive
d'intégrale 1.\hfill$\blacksquare$
\end{preuve}

\runinhead{Exercice -- Dans un triangle
(1) (\(\mathord{\bullet}\))}\phantomsection\addcontentsline{toc}{subsubsection}{Dans
un triangle (1)}\label{III-etb1}
Soient \(X\) et \(Y\) de densité jointe
\(f_{X,Y}(x,y)= \frac{1}{x}1_T (x,y)\) où \(T\) est le triangle
\(T = \{0< y< x < 1\}\).

\begin{enumerate}
\def\labelenumi{\arabic{enumi}.}
\tightlist
\item
  Calculer la densité marginale de \(X\)
\item
  Calculer la densité conditionnelle de \(Y\) sachant \(X=x\).
\end{enumerate}

\solref{III-answer-etb1}

L'interprétation de cette définition est la suivante~: la fonction
\(f_{Y|X=x}\) est la densité de la ``loi conditionnelle de \(Y\) sachant
que \(X = x\)''. Bien sûr, nous avons \(\mathbb{P}(X = x) = 0\) puisque
\(X\) admet une densité, donc la phrase ci-dessus n'a pas réellement de
sens, mais elle se justifie heuristiquement ainsi~: \(dx\) et \(dy\)
étant de ``petits'' accroissements des variables \(x\) et \(y\) et
lorsque \(f\) et \(f_X\) sont continues et strictement positives
respectivement en \((x,y)\) et \(x\)~: \begin{align*}
f_X(x) dx & \approx \mathbb{P}(X \in [x, x+dx])\\
f_{X,Y}(x,y) dx dy & \approx \mathbb{P}(X \in [x, x+dx], Y \in [y, y+dy])\\
\end{align*} Par suite \begin{align*}
f_{Y|X=x} (y) dy & \approx \frac{\mathbb{P}(X \in [x, x+dx], Y \in [y, y+dy])}{\mathbb{P}(X \in [x, x+dx])}\\
                 & \approx \mathbb{P}(Y \in [y, y+dy]|X \in [x, x+dx])\\
\end{align*}

On a alors le résultat suivant qui résout le problème posé en
introduction~:

\begin{proposition}\phantomsection\label{III-proposition}
Pour toute fonction \(g : \mathbb{R}^2 \to \mathbb{R}\) telle que
\(g(X,Y)\) admette une espérance, on a~:
\[\mathbb{E}(g(X,Y)) = \int_\mathbb{R}\left( \int_\mathbb{R}g(x,y)f_{Y|X=x}(y) dy \right) f_X(x) dx,\]
dont on déduit, en prenant \(g=1_{B_1 \times B_2}\), que~:
\[\mathbb{P}(X\in B_1, Y\in B_2) = \int_{B_1} \left( \int_{B_2}f_{Y|X=x}(y) dy \right) f_X(x) dx.\]
\end{proposition}

\begin{preuve}
On a \begin{align*}
\mathbb{E}(g(X,Y)) &= \int_{\mathbb{R}^2} g(x,y) f_{X,Y}(x,y) dy dx \\
             &= \int_{\mathbb{R}^2} g(x,y) f_{Y|X=x}(y)f_X(x) dy dx \\
             &= \int_\mathbb{R}\left( \int_\mathbb{R}g(x,y)f_{Y|X=x}(y) dy \right) f_X(x)dx,
\end{align*} les calculs étant licites par application du
\href{Probabilité\%20II.pdf\%20\#fubiniproba}{théorème de Fubini} et du
fait que l'application
\(x \mapsto \int_\mathbb{R}g(x,y)f_{Y|X=x}(y) dy\) est définie pour
\(f_X(x) >0\), soit presque partout relativement à la mesure
\(\mathbb{P}_X = f_X \lambda\).\hfill$\blacksquare$
\end{preuve}

\section{Cas général}\label{III-cas-guxe9nuxe9ral}

On peut établir le résultat suivant, qui complète le
\href{Probabilité\%20II.pdf\%20\#fubiniproba}{théorème de Fubini} et le
résultat d'existence et d'unicité des mesures produits, et que l'on
admettra.

\addcontentsline{toc}{subsection}{Fubini conditionnel}
\begin{theoreme}[{Fubini
conditionnel}]\phantomsection\label{III-fubinicond}

Soit un couple \((X,Y)\) de variables aléatoires réelles de loi jointe
\(\mathbb{P}_{X,Y}\), il existe une famille
\(\left(\mathbb{P}_{Y|X=x}\right)_{x\in\mathbb{R}}\) de probabilités sur
\((\mathbb{R},\mathcal{B}(\mathbb{R}))\), unique à une égalité
\(\mathbb{P}_X\)-presque partout près\footnote{c'est-à-dire qu'on peut
  définir ces probabilités de la manière qu'on souhaite pour les
  boréliens \(B\) tels que \(\mathbb{P}_X(B)=0\).}, qui vérifie pour
tous \(B_1, B_2\) boréliens de \(\mathbb{R}\)~:
\[ \mathbb{P}_{X,Y}(B_1 \times B_2) = \int_{B_1} \left( \int_{B_2} d\mathbb{P}_{Y|X=x}(y) \right) d\mathbb{P}_X(x).\]
Ces probabilités sont appelées \emph{lois conditionnelles} de \(Y\)
sachant \(X =x\). On a de plus pour toute application
\(g : \mathbb{R}^2 \to \mathbb{R}\) telle que \(g(X,Y)\) admette une
espérance~:
\[\mathbb{E}(g(X,Y)) = \int_\mathbb{R}\left( \int_\mathbb{R}g(x,y)d\mathbb{P}_{Y|X=x}(y) \right) d\mathbb{P}_X(x).\]
\end{theoreme}

\addcontentsline{toc}{subsection}{A noter}
\begin{remarque}[{A noter}]\phantomsection\label{III-a-noter}

\begin{itemize}
\tightlist
\item
  Ce résultat peut être interprété comme un \textbf{théorème de Fubini
  conditionnel}, dans le sens où il permet une intégration séquentielle,
  mais ici la mesure de probabilité du couple \((X,Y)\) s'exprime comme
  un produit de mesures dont l'un des termes dépend de la variable
  d'intégration de l'autre. En particulier, si on change l'ordre
  d'intégration, on change les mesures qui interviennent.
\item
  Fréquemment, dans les applications, la famille des lois
  conditionnelles est une donnée du modèle considéré, et leur existence
  ne pose donc pas de problème\,!
\item
  On retrouve les cas vus précédemment en notant que pour tout borélien
  \(B_1\) de \(\mathbb{R}\) on a
  \(\mathbb{P}_X(B_1) = \int_{B_1}d\mathbb{P}_X(x) = \sum_{x \in B_1} \mathbb{P}(X=x)\)
  lorsque \(X\) est discrète, et que pour tous boréliens \(B_1\) et
  \(B_2\) de \(\mathbb{R}\) on a
  \(\mathbb{P}_X(B_1) = \int_{B_1} f_X(x)dx\) et
  \(\mathbb{P}_{X,Y}(B_1 \times B_2) = \int_{B_1 \times B_2}f_{X,Y}(x,y) dx dy\).
\item
  Dans tout ce qui précède, les rôles de \(X\) et \(Y\) peuvent
  évidemment être inversés.
\end{itemize}
\end{remarque}

\section{Conséquences}\label{III-consuxe9quences}

Le \protect\hyperref[III-fubinicond]{théorème
précédent (p.~\pageref*{III-fubinicond})} a deux conséquences majeures.
Il fournit d'une part un moyen efficace d'identifier la loi marginale de
\(Y\) connaissant la loi marginale de \(X\) et la loi de \(Y\) sachant
\(X = x\). En effet, en notant que pour tout borélien \(B\) de
\(\mathbb{R}\),
\(\mathbb{P}_Y(B) = \mathbb{P}_{X,Y}(\mathbb{R}\times B)\) et en
appliquant ce théorème, on a la proposition suivante~:

\addcontentsline{toc}{subsection}{Formule de balayage conditionnel}
\begin{proposition}[{Formule de balayage
conditionnel}]\phantomsection\label{III-balcond}

\begin{itemize}
\tightlist
\item
  La loi marginale \(\mathbb{P}_Y\) de \(Y\) s'exprime comme la moyenne
  des lois conditionnelles \(\mathbb{P}_{Y|X=x}\) pondérée par la loi de
  \(X\). Pour tout \(B\) borélien de \(\mathbb{R}\)
  \[\mathbb{P}_Y(B) = \int_\mathbb{R}\left( \int_{B} d\mathbb{P}_{Y|X=x}(y) \right) d\mathbb{P}_X(x) = \int_\mathbb{R}\mathbb{P}_{Y|X=x}(B) d\mathbb{P}_X(x)\]
\item
  Dans le cas où \(X\) est discrète (à valeurs dans \(I\) dénombrable),
  on retrouve une expression de la formule des probabilités totales et
  composées~:
  \[\mathbb{P}_Y(B) = \mathbb{P}(Y\in B) = \sum_{x \in I} \mathbb{P}(Y \in B | X = x)\mathbb{P}(X=x)\]
\item
  Dans le cas où le couple \((X,Y)\) admet une densité, puisqu'on a
  \(f_{X,Y}(x,y) = f_{Y | X=x}(y)f_X(x)\), on obtient l'expression
  suivante pour la densité marginale~:
  \[f_Y(y) = \int_\mathbb{R}f_{X,Y}(x,y)dx = \int_\mathbb{R}f_{Y | X=x}(y)f_X(x) dx.\]
  On a en particulier la \emph{formule de Bayes pour les densités}~:
  pour tout \(x\) tel que \(f_X(x) > 0\) et tout \(y\) tel que
  \(f_Y(y) > 0\)~:
  \[ f_{X|Y=y}(x) = \frac{f_{X,Y}(x,y)}{f_Y(y)} = \frac{f_{Y|X=x}(y)f_X(x)}{f_Y(y)} .\]
\end{itemize}
\end{proposition}

\addcontentsline{toc}{subsection}{Pour fixer les idées (2)}
\begin{exemple}[{Pour fixer les idées
(2)}]\phantomsection\label{III-pour-fixer-les-iduxe9es-2}

Poursuivons \protect\hyperref[III-ex1]{l'exemple vu plus
haut (p.~\pageref*{III-ex1})}. On rappelle qu'on a déjà identifié la loi
marginale de \(X\) ainsi que la loi conditionnelle de \(Y\) sachant
\(X=n\) pour \(n\in\mathbb{N}\) que l'on rappelle ici~:
\[\mathbb{P}(X=n) = (1-\alpha)\alpha^n,\,n\in\mathbb{N}\text{ et }\forall  B \in \mathcal{B}(\mathbb{R}), \, \mathbb{P}_{Y|X=n}(B) = \int_{B \cap \mathbb{R}_+^\ast} e^{-t}\frac{t^n}{n!}dt\]
On peut en déduire la loi marginale de \(Y\) en utilisant la
\protect\hyperref[III-balcond]{formule de balayage
conditionnel (p.~\pageref*{III-balcond})} et le théorème de convergence
monotone~: \[\begin{aligned}
\mathbb{P}_Y(B)  &= \sum_{n \in \mathbb{N}} (1-\alpha)\alpha^n \int_{B \cap \mathbb{R}_+^\ast} e^{-t}\frac{t^n}{n!} dt\\
         &= (1-\alpha) \int_{B \cap \mathbb{R}_+^\ast} e^{-t} \sum_{n \in \mathbb{N}} \frac{(\alpha t)^n}{n!} dt \\
         &= \int_B 1_{\mathbb{R}_+}(t)(1-\alpha) e^{-(1-\alpha)t} dt,
\end{aligned}\] de sorte que \(Y\) suit une loi exponentielle de
paramètre \((1-\alpha)\).

En inversant les rôles, on va pouvoir identifier la loi de \(X\) sachant
\(Y \in B\) en notant que \begin{align*}
\mathbb{P}_{X,Y}(\{n\} \times B) &= \mathbb{P}_X(\{n\})\mathbb{P}_{Y|X=n}(B) \\
                         &= \int_B \frac{(\alpha t)^n}{n!} e^{-\alpha t} d\mathbb{P}_Y(t)\\
                         &= \int_B \mathbb{P}_{X | Y =t}(\{n\}) d\mathbb{P}_Y(t)
\end{align*} où l'on reconnaît que
\(\mathbb{P}_{X =n | Y =t}(\{n\}) = \frac{(\alpha t)^n}{n!} e^{-\alpha t}\),
c'est-à-dire que \(X\) sachant \(Y = t\) suit une loi de Poisson de
paramètre \(\alpha t\) pour \(\mathbb{P}_Y\)-presque tout \(t\).
\end{exemple}

En utilisant, le \protect\hyperref[III-fubinicond]{théorème de Fubini
conditionnel (p.~\pageref*{III-fubinicond})}, on obtient également une
nouvelle caractérisation de l'indépendance de deux variables aléatoires
faisant intervenir les lois conditionnelles.

\addcontentsline{toc}{subsection}{Critère d'indépendance}
\begin{proposition}[{Critère
d'indépendance}]\phantomsection\label{III-crituxe8re-dinduxe9pendance}

\begin{enumerate}
\def\labelenumi{\arabic{enumi}.}
\item
  \(X\) et \(Y\) sont indépendantes si et seulement si, pour
  \(\mathbb{P}_X\)-presque tout \(x\), \(\mathbb{P}_{Y |X = x}\) ne
  dépend pas de \(x\) et dans ce cas, on a
  \(\mathbb{P}_{Y |X = x} = \mathbb{P}_Y\), c'est-à-dire que la loi
  conditionnelle est identique à la loi marginale.
\item
  Dans le cas où \((X,Y)\) admet une densité, \(X\) et \(Y\) sont
  indépendantes si et seulement si la densité conditionnelle de \(Y\)
  sachant \(\{X = x\}\) ne dépend pas de \(x\).
\end{enumerate}
\end{proposition}

\begin{preuve}
\begin{enumerate}
\def\labelenumi{\arabic{enumi}.}
\item
  Si \(X\) et \(Y\) sont indépendantes, pour tous \(B_1\), \(B_2\)
  boréliens de \(\mathbb{R}\),
  \(\mathbb{P}_{X,Y} (B_1 \times B_2) = \mathbb{P}_X(B_1)\mathbb{P}_Y(B_2) = \int_{B_1}\mathbb{P}_Y(B_2)d\mathbb{P}_X(x) = \int_{B_2}\mathbb{P}_X(B_1) d\mathbb{P}_Y(y)\).
  Le résultat d'unicité du \protect\hyperref[III-fubinicond]{théorème de
  Fubini conditionnel (p.~\pageref*{III-fubinicond})} (à une égalité
  \(\mathbb{P}_X\)-presque sûre près), nous indique alors que
  \(\mathbb{P}_{Y|X=x}(B_2) = \mathbb{P}_Y(B_2)\).

  Inversement, si \(\mathbb{P}_{Y |X = x} = \mathbb{P}_Y\), alors
  \(\mathbb{P}_{X,Y} (B_1 \times B_2) = \int_{B_1}\mathbb{P}_{Y|X=x}(B_2)d\mathbb{P}_X(x) = \int_{B_1}\mathbb{P}_Y(B_2)d\mathbb{P}_X(x) = \mathbb{P}_X(B_1)\mathbb{P}_Y(B_2)\).
\item
  Si \(X\) et \(Y\) sont indépendantes,
  \(f_{X,Y} (x,y) = f_X(x) f_Y(y)\), d'où \(f_{Y|X=x}(y) = f_Y(y)\).

  Inversement, si \(f_{Y|X=x}(y) = f_Y(y)\) alors
  \(f_{X,Y}(x,y) = f_{Y|X=x}(y) f_X(x) = f_Y(y)f_X(x)\) et \(X\) et
  \(Y\) sont indépendantes.
\end{enumerate}

\hfill$\blacksquare$
\end{preuve}

\chapter{Espérance
conditionnelle}\label{III-espuxe9rance-conditionnelle-1}

Puisque \(\mathbb{P}_{Y|X=x}\) est la loi d'une variable aléatoire, on
peut définir l'espérance qui lui est associée et introduire la notion
d'espérance conditionnelle dans le cas où \(Y\) est intégrable.

\addcontentsline{toc}{subsection}{Espérance conditionnelle}
\begin{definition}[{Espérance
conditionnelle}]\phantomsection\label{III-defespcond}

Soit \(Y \in \mathcal{L}^1\).

\begin{enumerate}
\def\labelenumi{\arabic{enumi}.}
\tightlist
\item
  L'\emph{espérance conditionnelle de \(Y\) sachant \(\{X=x\}\)} est
  définie par
  \[\mathbb{E}(Y|X=x) = \int_\mathbb{R}y d\mathbb{P}_{Y|X=x} (y).\]
\item
  L'\emph{espérance conditionnelle de \(Y\) sachant \(X\)} est la
  \textbf{variable aléatoire} définie par~:
  \[\mathbb{E}(Y|X) = \psi(X), \text{ avec } \psi(x) = \mathbb{E}(Y|X=x).\]
\end{enumerate}
\end{definition}

\runinhead{Exercice -- Dans un triangle
(2) (\(\mathord{\bullet}\))}\phantomsection\addcontentsline{toc}{subsubsection}{Dans
un triangle (2)}\label{III-etb2}
Soient \(X\) et \(Y\) de densité jointe
\(f_{X,Y}(x,y)= \frac{1}{x}1_T (x,y)\) où \(T\) est le triangle
\(T = \{0< y< x < 1\}\). Calculer l'espérance conditionnelle de \(Y\)
sachant \(X\).\solref{III-answer-etb2}

\addcontentsline{toc}{subsection}{Conséquences}
\begin{remarque}[{Conséquences}]\phantomsection\label{III-consuxe9quences-1}

\begin{enumerate}
\def\labelenumi{\arabic{enumi}.}
\tightlist
\item
  \(\psi(x)\) n'est définie que pour \(x \notin N\), avec
  \(\mathbb{P}(X\in N) = 0\). Par conséquent, la
  \protect\hyperref[III-defespcond]{définition (p.~\pageref*{III-defespcond})}
  définit bien l'espérance conditionnelle \(\psi(X) = \mathbb{E}(Y|X)\)
  \(\mathbb{P}_X\)-presque partout, autrement dit avec probabilité 1, ou
  encore presque sûrement.
\item
  \(\mathbb{E}(\mathbb{E}(|Y||X)) = \mathbb{E}(|Y|)\) comme conséquence
  directe du \protect\hyperref[III-fubinicond]{théorème de Fubini
  conditionnel (p.~\pageref*{III-fubinicond})}. L'espérance
  conditionnelle de \(Y\) sachant \(X\) est bien définie dès que \(Y\)
  est intégrable.
\item
  Lorsque \((X,Y)\) admet une densité, l'espérance conditionnelle de
  \(Y\) sachant \(\{X=x\}\) s'écrit
  \[\mathbb{E}(Y|X=x) = \int_\mathbb{R}y f_{Y|X=x}(y) dy.\]
\end{enumerate}
\end{remarque}

\runinhead{Exercice -- Auto-conditionnement (\(\mathord{\bullet}\))}\phantomsection\addcontentsline{toc}{subsubsection}{Auto-conditionnement}\label{III-autocond}
Montrer que \(\mathbb{E}(Y|Y) = Y\).\solref{III-answer-autocond}

On peut étendre cette définition aux variables de la forme \(g(X,Y)\).

\addcontentsline{toc}{subsection}{Espérance conditionelle d'une fonction
de variables aléatoires}
\begin{definition}[{Espérance conditionelle d'une fonction de variables
aléatoires}]\phantomsection\label{III-defespcondg}

Soit \((X,Y)\) un couple de variables aléatoires réelles et \(g\) une
fonction mesurable positive ou \(\mathbb{P}_{X,Y}\)-intégrable sur
\(\mathbb{R}^2\).

\begin{enumerate}
\def\labelenumi{\arabic{enumi}.}
\tightlist
\item
  L'\emph{espérance conditionnelle de \(g(X,Y)\) sachant \(\{X=x\}\)}
  est définie par
  \[\mathbb{E}(g(X,Y)|X=x) = \int_\mathbb{R}g(x,y) d\mathbb{P}_{Y|X=x} (y).\]
\item
  L'\emph{espérance conditionnelle de \(g(X,Y)\) sachant \(X\)} est la
  \textbf{variable aléatoire} définie par~:
  \[\mathbb{E}(g(X,Y)|X) = \psi(X), \text{ avec } \psi(x) = \mathbb{E}(g(X,Y)|X=x).\]
\end{enumerate}
\end{definition}

\addcontentsline{toc}{subsection}{Espérance totale}
\begin{theoreme}[{Espérance
totale}]\phantomsection\label{III-espuxe9rance-totale}

Si \(Y\) est intégrable, alors \(\psi(X) = \mathbb{E}(Y | X)\) est
intégrable, et \[\mathbb{E}( \psi(X)) = E( Y ) .\]
\end{theoreme}

\begin{preuve}
C'est une conséquence directe du
\protect\hyperref[III-fubinicond]{théorème de Fubini
conditionnel (p.~\pageref*{III-fubinicond})}.\hfill$\blacksquare$
\end{preuve}

Ce résultat permet de calculer \(\mathbb{E}( Y )\) en conditionnant par
une variable auxiliaire \(X\)~:
\[\mathbb{E}( Y ) = \int_\mathbb{R}\mathbb{E}(Y | X = x) d\mathbb{P}_X(x)\]

Il généralise la \href{Probabilité\%20I.pdf\%20\#formprobatot}{formule
des probabilités totales}, qui correspond ici à \(Y = 1_A\) , et
\(B_x = \{X = x\}\) où les \(B_x\) forment cette fois une partition non
dénombrable de \(\mathbb{R}\). On l'écrit souvent sous la forme
\[ \mathbb{E}\left( \mathbb{E}(Y | X) \right) = \mathbb{E}( Y )\] et on
l'appelle la \emph{formule de l'espérance totale}.

L'espérance conditionnelle étant définie comme l'espérance selon la loi
conditionnelle, elle hérite des propriétés usuelles de l'espérance~:

\begin{enumerate}
\def\labelenumi{\arabic{enumi}.}
\tightlist
\item
  si \(Y\) et \(Z\) sont intégrables,
  \(\mathbb{E}(aY + bZ | X) = a \mathbb{E}(Y | X) + b \mathbb{E}(Z | X)\),
\item
  \(\mathbb{E}(Y | X) \geq 0\) si \(Y \geq 0\),
\item
  \(\mathbb{E}(1 | X) = 1\).
\end{enumerate}

De plus, si \(g\) est mesurable positive ou \(\mathbb{P}_X\)-intégrable,
\[ \mathbb{E}(Y g(X) | X) = g(X) \mathbb{E}(Y | X) \] est une
généralisation de l'égalité 1. ci-dessus, au cas où \(a = g(X)\), qui
doit être considéré ``comme une constante'' dans le calcul de
l'espérance conditionnelle sachant \(X\) (\(X\) est fixée comme une
donnée connue a priori). En effet, on a alors
\(\mathbb{E}(g(x)Y|X=x) = g(x)\psi(x)\). Enfin, on déduit directement du
\protect\hyperref[III-fubinicond]{théorème de Fubini
conditionnel (p.~\pageref*{III-fubinicond})} la proposition suivante.

\addcontentsline{toc}{subsection}{Transfert conditionnel}
\begin{proposition}[{Transfert
conditionnel}]\phantomsection\label{III-transfcond}

Soient un couple \((X,Y)\) de variables aléatoires réelles de loi jointe
\(\mathbb{P}_{X,Y}\) et \(g\) une fonction mesurable positive ou
\(\mathbb{P}_{X,Y}\)-intégrable sur \(\mathbb{R}^2\). On a pour
\(\mathbb{P}_X\)-presque tout \(x\) dans \(\mathbb{R}\)
\[\mathbb{E}(g(X,Y)|X=x) = \mathbb{E}(g(x,Y)|X=x) = \int_{\mathbb{R}}g(x,y) d\mathbb{P}_{Y|X=x} (y)\]
Si de plus \(X\) et \(Y\) sont indépendantes, on a~:
\[\mathbb{E}(g(X,Y)|X=x) = \mathbb{E}(g(x,Y)|X=x) = \int_{\mathbb{R}}g(x,y) d\mathbb{P}_Y(y).\]

Autrement dit, lorsqu'on conditionne par l'événement \(\{X=x\}\), cela
revient à fixer la valeur de la variable aléatoire \(X\) à la constante
\(x\).
\end{proposition}

\runinhead{Exercice -- Espérance conditionnelle d'un produit de
variables (\(\mathord{\bullet}\))}\phantomsection\addcontentsline{toc}{subsubsection}{Espérance
conditionnelle d'un produit de variables}\label{III-prod}
Calculer \(\mathbb{E}(XY|X=x)\) puis
\(\mathbb{E}(XY|X)\).\solref{III-answer-prod}

\addcontentsline{toc}{subsection}{Vecteurs gaussiens à densité}
\subsection*{Exemple~: vecteurs gaussiens à
densité}\label{III-exemple-vecteurs-gaussiens-uxe0-densituxe9}

Dans ce qui précède, on a décrit les lois et les espérances
conditionnelles dans le cas d'un couple de variables aléatoires à
valeurs dans \(\mathbb{R}^2\). Ces résultats sont aussi valables pour
des couples de vecteurs, dont on décrit ici un cas particulier.

Dans le cas des vecteurs gaussiens à densité, c'est-à-dire dont la
matrice de covariance est définie positive et donc inversible, le calcul
des lois conditionnelles de certaines composantes par rapport aux autres
est particulièrement aisé. On va voir en particulier que les lois
conditionnelles ont le bon goût d'être elles-mêmes gaussiennes, ce qui
explique (en partie) le succès de ces modèles dans les applications.

On considère un vecteur gaussien \(X = (X_1,\ldots,X_n)\) à valeurs dans
\(\mathbb{R}^n\) d'espérance \(m\) et de matrice de covariance \(C\)
définie positive. On a vu au chapitre 2 que la densité du vecteur \(X\)
s'écrit pour \(x\in\mathbb{R}^d\)~:
\[f_X(x) = \frac{1}{(2\pi)^{n/2}\sqrt{\det (C)}}\exp \left(-\frac{1}{2}(x-m)^t C^{-1}(x-m)\right)\]

Soit \(1 \leq k < n\) un entier. On souhaite exprimer \(f_{Y|Z=z}\), la
densité conditionnelle de \(Y = (X_1,\ldots,X_k)\) sachant
\(Z = (X_{k+1},\ldots,X_n) = (x_k,\ldots,x_n) = z\) (si \(k+1 = n\), ce
vecteur se réduit à une seul valeur). On a vu que
\[f_X = f_{Y|Z=z} f_Z,\] où \(f_Z\) est la densité marginale de \(Z\).
On cherche donc à décomposer \(f_X\) de la sorte. On note
\(m = (m_Y,m_Z)\) et on remarque que \(C\) peut se décomposer en blocs~:
\begin{equation*}
C = \left(\begin{array}{cc} C_Y & C_{Y,Z} \\ C_{Z,Y} & C_Z \end{array}\right)
\end{equation*} où \(C_Y = \text{Cov}(Y,Y)\), \(C_Z = \text{Cov}(Z,Z)\)
et \(C_{Y,Z} = \text{Cov}(Y,Z)\). Le \emph{complément de
Schur}\footnote{voir par exemple l'excellent
  \href{https://www2.compute.dtu.dk/pubdb/edoc/imm3274.pdf}{matrix
  cookbook}.} du bloc \(C_Y\) est la matrice
\[CS_Y = C_Y - C_{Y,Z}C_Z^{-1}C_{Z,Y}\] et permet d'exprimer l'inverse
de \(C\) comme~: \begin{equation*}
C^{-1} = \left(\begin{array}{cc} CS_Y^{-1} & -CS_Y^{-1}C_{Y,Z}C_Z^{-1} \\ -C_Z^{-1}C_{Z,Y}CS_Y^{-1}  & C_Z^{-1} + C_Z^{-1}C_{Z,Y}CS_Y^{-1}C_{Y,Z}C_Z^{-1} \end{array}\right)
\end{equation*} On peut alors réarranger les termes de la forme
quadratique dans \(f_X\) et on obtient~: \begin{align*}
(x-m)^t C^{-1}(x-m) =& \left(y - (m_Y + C_{Y,Z}C_Z^{-1}(z-m_Z))\right)^t CS_Y^{-1}\\
&.\left(y - (m_Y + C_{Y,Z}C_Z^{-1}(z-m_Z))\right)\\
 &+ (z-m_Z)^t C_Z^{-1}(z-m_Z)
\end{align*} Pour la constante, on peut remarquer que~:
\[ \det (C) = \det(CS_Y)\det(C_Z).\] On en déduit ainsi que
\[f_{Y|Z=z}(y) = \frac{1}{(2\pi)^{k/2}\sqrt{\det (CS_Y)}}\exp \left(-\frac{1}{2}\left(y - \psi(z)\right)^t CS_Y^{-1}\left(y - \psi(z))\right)\right)\]

C'est-à-dire que la variable aléatoire \(Y|Z=z\) est gaussienne
d'espérance \(m_{Y|Z=z} = \psi(z) = m_Y + C_{Y,Z}C_Z^{-1}(z-m_Z)\) et de
matrice de covariance \(CS_Y = C_Y - C_{Y,Z}C_Z^{-1}C_{Z,Y}\). Autrement
dit, l'espérance conditionnelle de \(Y\) sachant \(Z\) est la variable
aléatoire \(\mathbb{E}(Y|Z) = \psi(Z) =(m_Y + C_{Y,Z}C_Z^{-1}(Z-m_Z))\).
On notera que la covariance conditionnelle donnée par \(CS_Y\) ne dépend
pas de la valeur prise par \(Z\).

\chapter{Régression et espérance conditionnelle des variables de carré
intégrable}\label{III-ruxe9gression-et-espuxe9rance-conditionnelle-des-variables-de-carruxe9-intuxe9grable}

La régression est un ensemble de méthodes (d'apprentissage) statistiques
très utilisées pour analyser la relation d'une variable par rapport à
une ou plusieurs autres. Ces méthodes visent notamment à décrire les
liens de dépendance entre variables mais aussi à prédire au mieux la
valeur d'une quantité non observée en fonction d'une ou plusieurs autres
variables.

Pour fixer les idées, supposons que l'on s'intéresse au lien entre une
variable d'intérêt \(Y\) et une variable explicative \(X\) dans le but
de prédire la valeur de \(Y\) étant donnée l'observation de \(X\). Le
but est alors de proposer une certaine fonction
\(f : \mathbb{R}\rightarrow \mathbb{R}\) qui permette d'approximer \(Y\)
par \(f(X)\). Une telle fonction \(f\) est appelée un \emph{prédicteur}
de \(Y\) en fonction de \(X\).

On va montrer dans cette section que, dans le cadre des variables de
carré intégrable, une solution naturelle à ce problème est précisément
donnée par l'espérance conditionnelle de \(Y\) sachant \(X\). Nous nous
appuierons pour ce faire sur une interprétation géométrique de
l'espérance conditionnelle en tant que projection orthogonale dans un
espace de Hilbert. Cette approche justifie en particulier l'usage des
méthodes de moindres carrés dans les problèmes de régression.

\section{Rappels sur les espaces de Hilbert et application en
probabilités}\label{III-rappels-sur-les-espaces-de-hilbert-et-application-en-probabilituxe9s}

Soit \((\Omega, \mathcal{A},\mathbb{P})\) un espace probabilisé. On
rappelle que par définition, une variable aléatoire \(Y\) sur cet espace
est une fonction \(\mathcal{A}\)-mesurable
\(Y : \Omega \rightarrow \mathbb{R}\). On considère alors l'espace
fonctionnel \(L^2(\Omega)\) introduit dans le cours de calcul intégral.
Pour rappel, cet espace contient l'ensemble des fonctions
\(\mathcal{A}\)-mesurables (i.e.~des variables aléatoires) \(Y\) telles
que \(\vert Y\vert^2\) soit intégrable c'est-à-dire telles que \[
\Vert Y\Vert_2 = \bigg(\int_{\Omega} \vert Y(\omega)\vert^2 d\mathbb{P}(\omega)\bigg)^{1/2}=\mathbb{E}(Y^2)^{1/2} < \infty.
\] En particulier, contrairement à l'espace \(\mathcal{L}^2\) introduit
au chapitre précédent, dans l'espace \(L^2(\Omega)\) les fonctions
égales \(\mathbb{P}\)-presque partout sont confondues. C'est-à-dire que
deux fonctions \(X\) et \(Y\) égales \(\mathbb{P}\)-presque partout
correspondent à un même élément de \(L^2(\Omega)\) appelé classe
d'équivalence. Ainsi, pour \(X, Y \in L^2(\Omega)\), l'égalité \(X=Y\)
signifiera donc que \(X\) et \(Y\) sont dans la même classe
d'équivalence, autrement dit que les variables \(X\) et \(Y\) sont
égales presque sûrement.

On a vu dans le cours de calcul intégral que l'espace \(L^2(\Omega)\),
muni de la norme \(\Vert \cdot\Vert_2\), était un espace de Banach,
c'est-à-dire un espace vectoriel normé et complet pour cette norme. On
va maintenant introduire un produit scalaire sur cet espace, dont la
norme associée sera précisément la norme \(\Vert \cdot\Vert_2\), et
montrer que \(L^2(\Omega)\) est ainsi un espace de Hilbert.

On considère à cet effet l'application bilinéaire
\(\langle \cdot \vert \cdot\rangle\) définie par
\[\langle X \vert Y\rangle = \int_{\Omega} X(\omega) Y(\omega) d\mathbb{P}(\omega) = \mathbb{E}(XY), \quad X,Y\in L^2(\Omega).\]
L'application \(\langle \cdot \vert \cdot\rangle\) est bien définie
puisque, d'après l'inégalité de Cauchy-Schwartz, on a pour tous
\(X,Y\in L^2(\Omega)\):
\[\langle X \vert Y\rangle^2=\mathbb{E}(XY)^2 \leq \mathbb{E}(X^2)\mathbb{E}(Y^2)=\Vert X\Vert_2^2\Vert Y\Vert_2^2<\infty.\]
De plus, l'application \(\langle \cdot \vert \cdot\rangle\) définit un
produit scalaire sur \(L^2(\Omega)\) car elle vérifie les propriétés
suivantes pour tous \(X,Y\in L^2(\Omega)\):

\begin{itemize}
\tightlist
\item
  Symétrique:
  \(\langle X \vert Y\rangle= \mathbb{E}(XY)=\mathbb{E}(YX)=\langle X \vert Y\rangle\),
\item
  Positive: \(\langle X \vert X\rangle= \mathbb{E}(X^2) \ge 0\),
\item
  Définie:
  \(\langle X \vert X\rangle=0 \Leftrightarrow \mathbb{E}(X^2)=0 \Leftrightarrow X=0\)
  .
\end{itemize}

On remarque enfin que la norme associée au produit scalaire
\(\langle \cdot \vert \cdot\rangle\) est bien la norme
\(\Vert\cdot\Vert_2\) introduite précédemment (car pour tout
\(X\in L^2(\Omega)\),
\(\langle X \vert X\rangle^{1/2}=\mathbb{E}(X^2)^{1/2}=\Vert X\Vert_2\)).
On en déduit que l'espace \(L^2(\Omega)\) muni du produit scalaire
\(\langle \cdot \vert \cdot\rangle\) est un espace de Hilbert.

L'intérêt de travailler avec un espace de Hilbert est qu'il est possible
d'y définir de manière univoque la notion de projection orthogonale. En
effet, la projection orthogonale d'un élément \(Y\in L^2(\Omega)\) sur
un sous-espace convexe, fermé et non-vide \(K \subset L^2(\Omega)\) est
définie comme l'unique élément \(P_K(Y)\in K\) qui minimise la distance
\(\Vert Y - P_K(Y)\Vert_2\), autrement dit l'erreur des moindres carrés
\(\mathbb{E}((Y-P_K(Y))^2)\) entre \(Y\) et les éléments de \(K\). On
retrouve en particulier l'objectif de prédiction décrit en introduction,
ce qui permet de conclure que la projection orthogonale \(P_K(Y)\) est
le prédicteur optimal de \(Y\) au sens des moindres carrés parmi les
éléments de \(K\). On en déduit la propriété suivante qui permet de
caractériser ce prédicteur optimal d'après les propriétés de la
projection orthogonale vues dans le cours de calcul intégral.

\addcontentsline{toc}{subsection}{Prédicteur optimal sur un convexe
fermé}
\begin{proposition}[{Prédicteur optimal sur un convexe
fermé}]\phantomsection\label{III-identEC}

Soit \(Y\in L^2(\Omega)\) et soit \(K \subset L^2(\Omega)\) convexe,
fermé et non-vide. Le prédicteur optimal, au sens des moindres carrés,
de \(Y\) sur \(K\) est la projection orthogonale de \(Y\) sur \(K\). Il
est donc caractérisé comme étant l'unique élément de \(K\) vérifiant la
propriété
\[ \langle Y - P_K(Y) \vert U - P_K(Y)\rangle \le 0 \quad  \forall ~U\in K.\]
Si de plus \(K\) est un sous-espace vectoriel, cette caractérisation
prend une forme plus simple, à savoir:
\[ \langle Y - P_K(Y) \vert U \rangle = 0 \quad  \forall ~U\in K.\]
\end{proposition}

Pour revenir à notre problème initial de prédiction optimale de \(Y\)
par une fonction de \(X\), il s'agit donc de choisir le sous-espace
\(K\) de projection de manière à ce qu'il ne contienne que des éléments
s'exprimant comme des fonctions de \(X\), et d'identifier parmi les
éléments de \(K\) celui qui satisfait la
\protect\hyperref[III-identEC]{propriété
ci-dessus (p.~\pageref*{III-identEC})}.

\section{Cas particulier: la régression
linéaire}\label{III-cas-particulier-la-ruxe9gression-linuxe9aire}

Reprenons les variables \(X,Y\in L^2(\Omega)\) introduites en début de
section, et le problème initial de prédiction de \(Y\) à partir de
\(X\). Dans cette sous-section, nous nous restreignons à des prédicteurs
affines, c'est-à-dire que nous souhaitons trouver la meilleure
approximation, au sens des moindres carrés, de \(Y\) par une fonction
affine de \(X\) de la forme \(f(X)=aX + b\). Un tel prédicteur est
appelé \emph{régression linéaire de \(Y\) par \(X\)}.

Tout prédicteur affine \(f(X)=aX+b\) peut être vu comme une combinaison
linéaire de deux variables aléatoires de \(L^2(\Omega)\)~: \(X\), et la
variable aléatoire \(1\) égale à \(1\) presque sûrement. Par conséquent,
\(f(X)\in K\) où \(K\) désigne le sous-espace vectoriel de
\(L^2(\Omega)\) engendré par \(X\in L^2(\Omega)\) et
\(1\in L^2(\Omega)\). Ainsi, trouver un prédicteur affine optimal de
\(Y\) revient à trouver un prédicteur optimal de \(Y\) dans le
sous-espace vectoriel \(K\).

\(K\) étant un sous-espace vectoriel fermé de \(L^2(\Omega)\) (car de
dimension finie), la \protect\hyperref[III-identEC]{propriété de
caractérisation (p.~\pageref*{III-identEC})} du prédicteur optimal
permet de conclure que le prédicteur affine optimal de \(Y\) par \(X\)
est la projection orthogonale \(P_K(Y)\) de \(Y\) sur \(K\). En
particulier, puisque tout élément de \(K\) est une combinaison linéaire
de \(X\) et \(1\), on peut écrire que \(P_K(Y) = a^*X + b^* 1\), où les
coefficients \(a^*,b^*\) vérifient :
\[ \langle Y - (a^*  X + b^*  1) \vert a  X + b  1 \rangle = 0 \quad  \forall a,b\in\mathbb{R},\]
Cette égalité devant être vérifiée pour tous \(a,b\in\mathbb{R}\), elle
implique, par bilinéarité du produit scalaire, que les coefficients
\(a^*\) et \(b^*\) soient solutions du système : \[
\left\lbrace\begin{aligned}
   &\langle Y \vert X \rangle - a^*\langle X \vert X \rangle-b^*\langle 1 \vert X \rangle  = 0 \\
   &\langle Y \vert 1 \rangle - a^*\langle X \vert 1 \rangle-b^*\langle 1 \vert 1 \rangle = 0
\end{aligned}\right.
\] Par définition du produit scalaire \(\langle\cdot\vert\cdot\rangle\),
ce système se réécrit sous la forme : \[
\left\lbrace\begin{aligned}
   &\mathbb{E}(XY) = a^*\mathbb{E}(X^2)+b^*\mathbb{E}(X)   \\
   &\mathbb{E}(Y) = a^*\mathbb{E}(X)+b^* 
\end{aligned}\right.
\] et a pour solution \[\begin{aligned}
a^* & = \frac{\mathbb{E}(XY)-\mathbb{E}(X)\mathbb{E}(Y)}{\mathbb{E}(X^2)-\mathbb{E}(X)^2}=\frac{\text{Cov}(X,Y)}{\mathbb{V}(X)} = \rho(X,Y)\frac{\sigma_Y}{\sigma_X} \\
b^* & = \mathbb{E}(Y)  - a^*~\mathbb{E}(X)
\end{aligned}\]

En conclusion, le meilleur prédicteur affine, au sens des moindres
carrés, de \(Y\) par \(X\) est la fonction \(f^*(X)\) définie par
\[ f^*(X)=P_K(Y)=a^* X + b^*=\mathbb{E}(Y) + \rho(X,Y)\frac{\sigma_Y}{\sigma_X} (X -\mathbb{E}(X)),\]
et l'erreur quadratique moyenne associée à ce prédicteur vaut :
\[\begin{aligned}
\mathbb{E}\left(\left(Y - f^*(X)\right)^2\right) & = \sigma_Y^2 + \rho^2(X,Y)\sigma^2_Y - 2\rho^2(X,Y)\sigma^2_Y\\
                        & = \sigma^2_Y(1-\rho^2(X,Y)).
\end{aligned}\] On voit ainsi que cette erreur est proche de 0 lorsque
\(|\rho(X,Y)| \approx 1\) tandis qu'elle est proche de
\(\mathbb{V}(Y) = \sigma^2_Y\) lorsque \(\rho(X,Y) \approx 0\). On
notera au passage qu'on obtient que la meilleure approximation de \(Y\)
par une constante (ce qui revient à prendre \(X=0\)) est son espérance.

\runinhead{Exercice -- Une deuxième
approche (\(\mathord{\bullet}\))}\phantomsection\addcontentsline{toc}{subsubsection}{Une
deuxième approche}\label{III-detail}
Montrer que les coefficients \(a^*\) et \(b^*\) sont les uniques
minimiseurs de l'erreur \(\ell(a,b)=\mathbb{E}((Y - (aX + b))^2)\), pour
\(a,b\in\mathbb{R}\).\solref{III-answer-detail}

\addcontentsline{toc}{subsection}{Limites de la régression linéaire}
\begin{remarque}[{Limites de la régression
linéaire}]\phantomsection\label{III-limites-de-la-ruxe9gression-linuxe9aire}

L'hypothèse d'une relation linéaire est très forte et pas nécessairement
toujours adaptée pour expliquer des relations de dépendances entre
variables. Soit en effet une variable aléatoire réelle \(X\) de
\(\mathcal{L}^3\) (i.e.~\(X^3\) est \(\mathbb{P}_X\) intégrable)
symétrique, c'est-à-dire telle que \(X\) et \(-X\) sont de même loi. On
a alors \(\mathbb{E}(X) = -\mathbb{E}(X) = 0\). Les variables \(X\) et
\(X^2\) ne sont clairement pas indépendantes. Pour autant, on a
\(\text{Cov}(X,X^2) = \mathbb{E}(X^3) = -\mathbb{E}(X^3) = 0\) et le
coefficient de régression \(a^*\) est nul.
\end{remarque}

\section{Cas général}\label{III-cas-guxe9nuxe9ral-1}

Dans le paragraphe précédent, on s'est intéressé à approximer une
variable aléatoire \(Y\in L^2(\Omega)\) par une fonction affine d'une
autre variable \(X\in L^2(\Omega)\). On va montrer ici que si on
généralise la recherche du prédicteur optimal à l'ensemble des fonctions
\(f\) telles que \(f(X)\) est de carré intégrable, la solution est
donnée par l'espérance conditionnelle \(\psi(X) = \mathbb{E}(Y|X)\).

En effet, considérons le sous-espace
\(L^2_X(\Omega)\subset L^2(\Omega)\) constitué des (classes
d'équivalence) des variables aléatoires de la forme \(f(X)\), avec
\(f : \mathbb{R}\rightarrow\mathbb{R}\) borélienne et telle que
\(f(X) \in L^2(\Omega)\). On remarque que \(L^2_X(\Omega)\) est
sous-espace vectoriel de \(L^2(\Omega)\). On admettra par ailleurs que
\(L^2_X(\Omega)\) est fermé. Par conséquent, le prédicteur optimal de
\(Y\) dans \(L^2_X(\Omega)\) est donné par la projection orthogonale de
\(Y\) sur \(K=L^2_X(\Omega)\), et on peut une nouvelle fois appliquer la
\protect\hyperref[III-identEC]{propriété de
caractérisation (p.~\pageref*{III-identEC})} associée pour l'identifier.

On note par ailleurs que l'espérance conditionnelle
\(\psi(X) = \mathbb{E}(Y|X) \in L^2_X(\Omega)\) et que pour tout
\(f(X)\in L^2_X(\Omega)\), \[ 
\begin{aligned}
  \langle Y - \psi(X) \vert f(X) \rangle =  \mathbb{E}\big((Y - \psi(X))f(X)\big)
&= \mathbb{E}\big(Yf(X)\big) - \mathbb{E}\big(\mathbb{E}(Y|X)f(X)\big) \\
&= \mathbb{E}\big(Yf(X)\big) - \mathbb{E}\big(\mathbb{E}(Yf(X)|X)\big) \\
&= \mathbb{E}\big(Yf(X)\big) - \mathbb{E}\big(Yf(X)\big)=0
\end{aligned}
\] où la seconde égalité se déduit du fait que \(f(X)\) est
\(\mathbb{P}\)-intégrable, et la troisième découle du théorème de
l'espérance totale. Ainsi l'espérance conditionnelle vérifie la
\protect\hyperref[III-identEC]{propriété de
caractérisation (p.~\pageref*{III-identEC})} et par unicité de la
projection orthogonale, on en déduit que
\(P_{L^2_X(\Omega)}(Y)=\psi(X)\).

\begin{figure}
\centering
\includegraphics{Proj.tex.pdf}
\caption{Illustration géométrique de l'espérance conditionnelle comme
projection orthogonale sur un sous-espace de
\(L^2(\Omega)\)}\label{III-fig_proj}
\end{figure}

En conclusion, l'espérance conditionnelle \(\mathbb{E}(Y|X)\) est le
meilleur prédicteur de \(Y\) au sens des moindres carrés parmi les
fonctions de \(X\). On notera qu'en pratique, dans les problèmes de
régression, on se restreint souvent à des classes de fonctions
paramétrées (par exemple les fonctions affines vues précédemment, ou les
réseaux de neurones dans l'apprentissage profond) car on ne dispose que
rarement de l'expression explicite de l'espérance conditionnelle.

Il est ensuite immédiat que le ``résidu'' \(Y-\mathbb{E}(Y|X)\) est non
corrélé avec \(X\) du fait de l'orthogonalité. On en déduit ainsi la
\emph{formule de la variance totale}~: \[
\begin{aligned}
\mathbb{V}(Y) = \|Y - \mathbb{E}(Y)\|^2 &=  \|Y - \mathbb{E}(Y|X) + \mathbb{E}(Y|X) - \mathbb{E}(Y)\|^2 \\
                          &=  \|Y - \mathbb{E}(Y|X)\|^2 + \|\mathbb{E}(Y|X) - \mathbb{E}(Y)\|^2 \\
                          &=  \mathbb{E}((Y - \mathbb{E}(Y|X))^2) + \mathbb{E}((\mathbb{E}(Y|X) - \mathbb{E}(Y))^2)\\
                          &= \mathbb{E}(\mathbb{E}((Y - \mathbb{E}(Y|X))^2|X)) + \mathbb{V}(\mathbb{E}(Y|X))\\
                          &= \mathbb{E}(\mathbb{V}(Y|X)) + \mathbb{V}(\mathbb{E}(Y|X)).
\end{aligned}
\] où on a utilisé la formule de l'espérance totale et introduit la
variable aléatoire variance conditionnelle
\(\mathbb{V}(Y|X) = \mathbb{E}((Y - \mathbb{E}(Y|X))^2|X)\) comme cas
particulier de la \protect\hyperref[III-defespcondg]{définition vue plus
haut (p.~\pageref*{III-defespcondg})}. Cette décomposition est illustrée
dans la \protect\hyperref[III-fig_proj2]{figure
ci-après (figure~\ref*{III-fig_proj2})}.

\begin{figure}
\centering
\includegraphics{Proj2.tex.pdf}
\caption{Illustration géométrique de la formule de la variance totale.
L'espace \(\mathrm{Vect}\lbrace 1\rbrace\) désigne le sous-espace de
\(L^2(\Omega)\) composé des variables aléatoires p.s.
constantes.}\label{III-fig_proj2}
\end{figure}

\runinhead{Exercice -- Variance
totale (\(\mathord{\bullet}\))}\phantomsection\addcontentsline{toc}{subsubsection}{Variance
totale}\label{III-vartot}
Démontrer ce résultat sans utiliser la notion
d'orthogonalité.\solref{III-answer-vartot}

\chapter{Exercices}\label{III-exercices}

\section{Couple de variables}\label{III-couple-de-variables}

Soient \(X\) et \(Y\) deux v.a. réelles. On suppose que la densité
conditionnelle de \(X\) sachant \(Y = y\) est la densité
\(1_{\mathbb{R}_+}(x)y^2 xe^{-xy}\) et que la loi de \(Y\) est de
densité \(\frac{1}{y^2} 1_{[1,+\infty[} (y)\). On pose \(T = XY\).

\runinhead{Question
1}\phantomsection\addcontentsline{toc}{subsubsection}{Question
1}\label{III-ex1-1}
Trouver la loi du couple \((T,Y)\). Qu'en
déduit-on\,?\solref{III-answer-ex1-1}

\runinhead{Question
2}\phantomsection\addcontentsline{toc}{subsubsection}{Question
2}\label{III-ex1-2}
Trouver la loi conditionnelle de \(Y\) sachant
\(X = x\).\solref{III-answer-ex1-2}

\runinhead{Question
3}\phantomsection\addcontentsline{toc}{subsubsection}{Question
3}\label{III-ex1-3}
Calculer \(E(Y |X)\).\solref{III-answer-ex1-3}

\section{Mélanges de lois}\label{III-muxe9langes-de-lois}

\emph{Adapté du cours de probabilités de S. Bonnabel et M. Schmidt
(MINES ParisTech).}

Pour modéliser un phénomène multimodal, on utilise souvent des mélanges
de gaussiennes. C'est le cas notamment en classification non-supervisée,
où on fait l'hypothèse que chacune des classes suit une loi gaussienne.
Soient \(n\in\mathbb{N}^\ast\) et \(K\) une variable aléatoire prenant
les valeurs \(1,\dots,n\) avec les probabilités non nulles
\(p_1,\dots,p_n\) telles que \(\sum_{i = 1}^n p_i = 1\). Soient
\(X_1,\dots,X_n\) des variables aléatoires gaussiennes mutuellement
indépendantes, d'espérances respectives \(m_1,\dots,m_n \in \mathbb{R}\)
et de variances respectives
\(\sigma_1^2,\dots,\sigma_n^2 \in\mathbb{R}_+^\ast\), toutes
indépendantes de \(K\). On appelle mélange de gaussiennes la loi de la
variable aléatoire \(X = X_K\). Pour tout \(i\in\{1,\dots,n\}\), on
notera \(f_i\) la densité de la variable aléatoire \(X_i\).

\runinhead{Question
1}\phantomsection\addcontentsline{toc}{subsubsection}{Question
1}\label{III-melloi1}
Soit \(i \in \{1,\dots,n\}\). Quelle est la densité \(f_{X\mid K = i}\)
de \(X\) conditionnellement à l'événement
\(\{K = i\}\)\,?\solref{III-answer-melloi1}

\runinhead{Question
2}\phantomsection\addcontentsline{toc}{subsubsection}{Question
2}\label{III-melloi2}
Calculer la densité de probabilité de la variable
\(X\).\solref{III-answer-melloi2}

\runinhead{Question
3}\phantomsection\addcontentsline{toc}{subsubsection}{Question
3}\label{III-melloi3}
Calculer \(\mathbb{E}(X)\). Montrer que
\(\mathbb{V}(X) = \sum_{i = 1}^n p_i\sigma_i^2 + \bar{\sigma}^2\), où ce
dernier terme peut être interprété comme la dispersion des
espérances.\solref{III-answer-melloi3}

\runinhead{Question
4}\phantomsection\addcontentsline{toc}{subsubsection}{Question
4}\label{III-melloi4}
Comment approximeriez-vous le mélange par une unique gaussienne\,? Faire
un schéma dans le cas \(m = 2\).\solref{III-answer-melloi4}

\section{Lois conjuguées}\label{III-lois-conjuguuxe9es}

Soit un vecteur aléatoire \((X,Y)\) de loi jointe \(\mathbb{P}_{X,Y}\).
Expliciter la loi conditionnelle de \(Y\) sachant \(\{X=x\}\) dans les
situations suivantes, en prenant soin d'expliciter pour quelles valeurs
de \(x\) ces dernières ont du sens.

\runinhead{Question
1}\phantomsection\addcontentsline{toc}{subsubsection}{Question
1}\label{III-loiconj-expexp}
\(Y\) suit une loi Exponentielle de paramètre
\(\lambda \in\mathbb{R}_+^\ast\) et pour tout \(y\in\mathbb{R}_+^\ast\),
la variable aléatoire \(X\) sachant \(\{Y=y\}\) suit une loi
Exponentielle de paramètre \(y\).\solref{III-answer-loiconj-expexp}

\runinhead{Question
2}\phantomsection\addcontentsline{toc}{subsubsection}{Question
2}\label{III-loiconj-gampoi}
\(Y\) suit une loi Gamma de paramètres
\(\alpha,\theta \in \mathbb{R}_+^\ast\) et pour tout
\(y\in\mathbb{R}_+^\ast\), la variable aléatoire \(X\) sachant
\(\{Y=y\}\) suit une loi de Poisson de paramètre
\(y\).\solref{III-answer-loiconj-gampoi}

\section{Randomisation}\label{III-randomize}

\emph{Extrait du cours de probabilités de S. Bonnabel et M. Schmidt
(MINES ParisTech).}

Des clients arrivent à la boutique SNCF du boulevard Saint-Michel à des
instants aléatoires. On note \(T_0\) l'heure d'ouverture puis
\(T_1, T_2, \dots\) les temps successifs d'arrivée des clients jusqu'à
l'heure de fermeture. Les études statistiques montrent qu'on peut, dans
une tranche horaire donnée, supposer que les temps d'attente
\(X_1 = T_1 -T_0, X_2 = T_2 -T_1,\dots\) peuvent être modélisés par des
variables aléatoires indépendantes et de même loi qu'une variable
aléatoire positive \(X\). Par ailleurs, une loterie interne décide que
chaque jour dans la tranche horaire considérée, le \(N^{\text{ème}}\)
client sera l'heureux gagnant d'un trajet gratuit Paris-La Ciotat, où
\(N\) est une variable aléatoire bornée dont la loi dépend du processus
de loterie (e.g.~tous les clients entre le premier et le
\(30^{\text{ème}}\) ont une chance \(1/30\) d'être tirés au sort, en
supposant qu'on est sûr d'avoir au moins \(30\) clients dans la tranche
horaire).

On se demande alors~: quel est le temps d'attente moyen avant d'obtenir
un gagnant\,?\solref{III-answer-randomize}

\section{Etats cachés --- indépendance
conditionnelle}\label{III-etats-cachuxe9s-induxe9pendance-conditionnelle}

Soucieux de l'évolution du potager de l'école, des élèves à la main
verte s'intéressent à l'évolution de la température dans le jardin côté
Luxembourg. Ils récupèrent pour cela un thermomètre dans un laboratoire,
l'installent près du potager, et en relèvent les mesures à intervalles
de temps réguliers. Les résultats les surprennent rapidement~: les
températures affichées ne correspondent pas à celles prévues par
météo-France. Leur thermomètre est sans doute déréglé.

On se propose de les aider à comprendre le phénomène dont ils sont
témoins à l'aide d'un modèle probabiliste particulier, nommé
\emph{modèle de Markov caché}. Précisément, on considère la suite des
vraies températures que l'on aurait souhaité relever comme une suite de
v.a.r. non indépendantes \((X_n)_{n\in\mathbb{N}^\ast}\), dite
\emph{d'états cachés} (on ne les observe pas directement). Les erreurs
commises par le thermomètre sont quant à elles modélisées par une suite
de v.a.r. \((\epsilon_n)_{n\in\mathbb{N}^\ast}\), toutes indépendantes
et de même loi admettant une densité \(f_\epsilon\). Elles sont
supposées indépendantes de la suite \((X_n)_{n\in\mathbb{N}^\ast}\)
(l'erreur du thermomètre lui est propre et ne dépend pas de la
température réelle). A chaque instant \(n\in\mathbb{N}^\ast\), on
suppose que la mesure du thermomètre est la variable aléatoire
\[Y_n = X_n + \epsilon_n,\] et que le vecteur aléatoire
\((X_1,\dots,X_n)\) possède une densité jointe notée \(f_{1:n}\).

\runinhead{Question
1}\phantomsection\addcontentsline{toc}{subsubsection}{Question
1}\label{III-ec1}
Montrer que pour tout \(n \in \mathbb{N}^\ast\) et tout
\(x\in\mathbb{R}\), la loi de \(Y_n\) sachant \(\{X_n = x\}\) admet une
densité, que l'on explicitera.\solref{III-answer-ec1}

\runinhead{Question
2}\phantomsection\addcontentsline{toc}{subsubsection}{Question
2}\label{III-ec2}
Montrer que les \(n\in\mathbb{N}^\ast\) relevés de température
\(Y_1,\dots,Y_n\) sont indépendants \textbf{conditionnellement} aux
états cachés \(X_1,\dots,X_n\).\solref{III-answer-ec2}

\section{Covariance totale}\label{III-covtot}

Soient \(X\), \(Y\) et \(Z\) trois variables aléatoires réelles de carré
intégrable. La covariance conditionnelle de \(X\) et \(Y\) sachant \(Z\)
est définie comme la variable aléatoire
\[\text{Cov}(X,Y \mid Z) = \mathbb{E}\Bigl( \bigl( X - \mathbb{E}(X\mid Z) \bigr)\bigl( Y - \mathbb{E}(Y\mid Z) \bigr) \Bigm| Z  \Bigr).\]

Etablir la formule de la covariance totale~:
\[\text{Cov}(X,Y) = \mathbb{E}\bigl(\text{Cov}(X,Y\mid Z)\bigr) + \text{Cov}\bigl( \mathbb{E}(X\mid Z), \mathbb{E}(Y\mid Z) \bigr).\]

\solref{III-answer-covtot}

\section{Débruitage (formule de
Tweedie)}\label{III-duxe9bruitage-formule-de-tweedie}

Les photographies sont souvent entâchées d'un bruit qui dégrade leur
qualité, notamment lorsque la lumière est faible. Pour améliorer le
rendu, des algorithmes de débruitage sont mis en oeuvre. Dans cet
exercice, on se propose de donner un éclairage probabiliste sur ce
problème.

On considère que l'image d'intérêt \(X\) est un vecteur aléatoire de
\(\mathbb{R}^d\) de densité \(f_X\) (quitte à réarranger les pixels et
les canaux RGB) et que le bruit \(Z\) est de densité \(f_Z\),
gaussienne, centrée, de covariance \(\Sigma\). L'image bruitée peut donc
être représentée par \(Y = X + Z\).

\runinhead{Question
1}\phantomsection\addcontentsline{toc}{subsubsection}{Question
1}\label{III-tweedie1}
Exprimer la densité \(f_Y\) de \(Y\) en fonction de \(f_X\) et
\(f_Z\).\solref{III-answer-tweedie1}

\runinhead{Question
2}\phantomsection\addcontentsline{toc}{subsubsection}{Question
2}\label{III-tweedie2}
Calculer le gradient de \(\log f_Y\), que l'on appelle aussi le score de
la loi de \(Y\).\solref{III-answer-tweedie2}

\runinhead{Question
3}\phantomsection\addcontentsline{toc}{subsubsection}{Question
3}\label{III-tweedie3}
En déduire une expression de \(\mathbb{E}(X|Y)\) ne faisant pas
intervenir la densité \(f_X\).

Cette propriété est particulièrement intéressante en pratique puisque
pour ce type de problèmes, on ne dispose généralement que d'obervations
bruitées qui vont nous permettre d'approximer les quantités nécessaires
pour effectuer ce calcul selon une approche statistique. Elle est par
ailleurs au coeur des algorithmes modernes de génération
d'image.\solref{III-answer-tweedie3}

\section{Non-réponse}\label{III-non-ruxe9ponse}

\emph{Inspiré du cours de probabilité de M. Christine (ENSAE
ParisTech).}

Un questionnaire est diffusé aux \(n\in\mathbb{N}^\ast\) étudiants de
l'école pour savoir combien de temps ils ont consacré à l'étude des
probabilités ce semestre. On note \(Y_i\) le temps de travail de
l'étudiant \(i \in \{1,\dots,n\}\) et \(X_i\) la variable valant \(1\)
s'il a répondu au questionnaire et \(0\) sinon. On suppose que les
\((X_1,Y_1),\dots,(X_n,Y_n)\) sont des vecteurs aléatoires indépendants
de même distribution qu'un vecteur générique \((X,Y)\) tel que

\begin{itemize}
\tightlist
\item
  \(X\) est une variable de Bernoulli de paramètre \(p\in\,]0,1[\)
  indiquant la probabilité de réponse,
\item
  \(Y\) est positive, de carré intégrable, d'espérance
  \(m\in\mathbb{R}_+\) et de variance \(\sigma^2 \in\mathbb{R}_+^\ast\).
  Le coefficient de corrélation entre \(X\) et \(Y\) est enfin noté
  \(\rho \in [-1,1]\).
\end{itemize}

\runinhead{Question
1}\phantomsection\addcontentsline{toc}{subsubsection}{Question
1}\label{III-nonrep1}
En reprenant la définition de l'espérance conditionnelle
\(\mathbb{E}(Y\mid X)\) comme meilleure approximation au sens des
moindres carrés de \(Y\) par une fonction de \(X\), montrer qu'elle
coïncide ici avec l'approximation affine de \(Y\) par \(X\) puis
l'écrire en fonction de \(m\), \(\rho\), \(\sigma\) et
\(p\).\solref{III-answer-nonrep1}

\runinhead{Question
2}\phantomsection\addcontentsline{toc}{subsubsection}{Question
2}\label{III-nonrep2}
On pose \(m_0 := \mathbb{E}(Y\mid X = 0)\) et
\(m_1 = \mathbb{E}(Y\mid X = 1)\). Calculer \(m_0\) et \(m_1\) en
fonction de \(m\), \(\rho\), \(\sigma\) et
\(p\).\solref{III-answer-nonrep2}

\runinhead{Question
3}\phantomsection\addcontentsline{toc}{subsubsection}{Question
3}\label{III-nonrep3}
On pose \(\sigma^2_0 := \mathbb{V}\left(Y\mid X = 0\right)\) et
\(\sigma^2_1 := \mathbb{V}\left(Y\mid X = 1\right)\). Vérifier l'égalité
\[\sigma^2 = \dfrac{(1-p)\,\sigma^2_0 + p\,\sigma^2_1}{1-\rho^2}.\]

\solref{III-answer-nonrep3}

\runinhead{Question
4}\phantomsection\addcontentsline{toc}{subsubsection}{Question
4}\label{III-nonrep4}
Que dire des résultats obtenus aux questions 2 et 3 lorsque~:

\begin{itemize}
\tightlist
\item
  \(X\) et \(Y\) sont non corrélées,
\item
  \(X\) et \(Y\) sont indépendantes\,?
\end{itemize}
\solref{III-answer-nonrep4}

\begin{center}\rule{0.5\linewidth}{0.5pt}\end{center}



\part{Convergences et théorèmes limites}
\chapter{Objectifs d'apprentissage}\label{IV-objectifs-dapprentissage}

Cette section s'efforce d'expliciter et de hiérarchiser les acquis
d'apprentissages associés au chapitre. Ces objectifs sont organisés en
paliers~:

($\mathord{\boldsymbol{\circ}}$) Prérequis ($\mathord{\bullet}$)
Fondamental ($\mathord{\bullet}\mathord{\bullet}$) Standard
($\mathord{\bullet}\mathord{\bullet}\mathord{\bullet}$) Avancé
($\mathord{\bullet}\mathord{\bullet}\mathord{\bullet}\mathord{\bullet}$)
Expert

Sauf mention particulière, la connaissance des démonstrations du
document n'est pas exigible\footnote{l'étude des démonstrations du cours
  peut toutefois contribuer à votre apprentissage, au même titre que la
  résolution d'exercices.}

\runinhead{Préambule}\phantomsection\label{IV-pruxe9ambule}

\begin{itemize}
\tightlist
\item
  $\mathord{\bullet}$ connaître la définition des espaces
  \(\mathcal{L}^p\)
\item
  $\mathord{\bullet}\mathord{\bullet}$ connaître la propriété
  d'inclusion dans les espaces \(\mathcal{L}^p\)
\item
  $\mathord{\bullet}$ connaître l'inégalité de Markov
\item
  $\mathord{\bullet}$ connaître l'inégalité de Bienaymé-Chebyshev
\end{itemize}

\runinhead{Suites de v.a.
indépendantes}\phantomsection\label{IV-suites-de-v.a.-induxe9pendantes}

\begin{itemize}
\tightlist
\item
  $\mathord{\bullet}\mathord{\bullet}\mathord{\bullet}\mathord{\bullet}$ connaître
  le théorème d'existence et d'unicité de la mesure produit sur les
  espaces produit de dimension infinie
\item
  $\mathord{\bullet}\mathord{\bullet}$ connaître la définition d'une
  suite de variables aléatoires indépendantes
\item
  $\mathord{\bullet}\mathord{\bullet}$ connaître et savoir utiliser le
  lemme de Borel-Cantelli
\end{itemize}

\runinhead{Convergence et loi des grands
nombres}\phantomsection\label{IV-convergence-et-loi-des-grands-nombres}

\begin{itemize}
\tightlist
\item
  $\mathord{\bullet}\mathord{\bullet}$ connaître les différents modes de
  convergence des suites de v.a.
\item
  $\mathord{\bullet}\mathord{\bullet}\mathord{\bullet}$ savoir
  caractériser le(s) mode(s) de convergence d'une suite de v.a.
\item
  $\mathord{\bullet}\mathord{\bullet}$ connaître les liens d'implication
  entre les différents modes de convergence
\item
  $\mathord{\bullet}\mathord{\bullet}$ connaître la propriété de
  continuité
\item
  $\mathord{\bullet}$ connaître la loi faible des grands nombre
\item
  $\mathord{\bullet}$ connaître la loi forte des grands nombre (sous les
  hypothèses minimales)
\end{itemize}

\runinhead{Convergence en loi et théorème central
limite}\phantomsection\label{IV-convergence-en-loi-et-thuxe9oruxe8me-central-limite}

\begin{itemize}
\tightlist
\item
  $\mathord{\bullet}$ connaître la définition de la convergence en loi
\item
  $\mathord{\bullet}$ savoir que la convergence en probabilité implique
  la convergence en loi
\item
  $\mathord{\bullet}\mathord{\bullet}$ savoir que la convergence simple
  des fonctions de répartition équivaut à la convergence en loi
\item
  $\mathord{\bullet}$ connaître et savoir utiliser le théorème central
  limite
\end{itemize}

\chapter{Préambule}\label{IV-pruxe9ambule-1}

Dans ce chapitre, on va s'intéresser à l'étude des suites de variables
aléatoires. On introduira notamment différents modes de convergence. En
particulier, on établira la \emph{loi des grands nombres} qui montre
rigoureusement que, quand le nombre de répétitions de l'expérience tend
vers l'infini, la fréquence de réalisation d'un événement converge vers
la probabilité de réalisation de cet événement. Ce résultat justifie
également la démarche employée en statistique où l'on cherche par
exemple à estimer une espérance par une moyenne empirique~: cette
moyenne empirique va ainsi converger vers la quantité ciblée. Enfin, on
présentera le théorème central limite qui nous indiquera la vitesse à
laquelle cette convergence a lieu ainsi qu'un moyen de contrôler
l'erreur commise en remplaçant une espérance par sa contrepartie
empirique.

Notons d'abord qu'on peut étendre les définitions des espaces
\(\mathcal{L}^1\) et \(\mathcal{L}^2\) pour un \(p \in \mathbb{N}^\ast\)
quelconque.

\addcontentsline{toc}{subsection}{Espace \(\mathcal{L}^p\)}
\begin{definition}[{Espace
\(\mathcal{L}^p\)}]\phantomsection\label{IV-espace-mathcallp}

Soit \(X\) une variable aléatoire. On note \(X \in \mathcal{L}^p\), ou
\(\mathcal{L}^p(\Omega,\mathcal{A},\mathbb{P})\), si et seulement si
\(\mathbb{E}(|X|^p) = \int_\mathbb{R}|x|^p \mathbb{P}_X(dx) = \int_{\Omega} |X|^p(\omega)\mathbb{P}(d\omega) < +\infty\).

Si \(X \in \mathcal{L}^p\), on dit qu'elle admet un moment d'ordre
\(p\). Du fait que \(\mathbb{P}\) est une mesure finie, on a la
stabilité par inclusion suivante~:
\end{definition}

\addcontentsline{toc}{subsection}{Inclusion}
\begin{proposition}[{Inclusion}]\phantomsection\label{IV-inclusion}

Soit \(p \in \mathbb{N}^\ast\), on a l'inclusion~:
\[ \mathcal{L}^{p+1}(\Omega,\mathcal{A},\mathbb{P}) \subset \mathcal{L}^p(\Omega,\mathcal{A},\mathbb{P})\]
\end{proposition}

\begin{preuve}
Supposons \(X \in \mathcal{L}^{p+1}(\Omega,\mathcal{A},\mathbb{P})\). On
a
\[|X|^p \leq \max(1,|X|^{p+1}) = 1_{|X| < 1} + 1_{|X| \geq 1} |X|^{p+1}.\]
Le terme de droite est intégrable, en effet~:
\[\mathbb{E}\left(1_{|X| < 1} + 1_{|X| \geq 1} |X|^{p+1}\right) \leq \int_\Omega \mathbb{P}(d\omega) + \int_\Omega |X(\omega)|^{p+1} \mathbb{P}(d\omega) = 1 + \mathbb{E}(|X|^{p+1}).\]
donc \(|X|^p\) est intégrable.\hfill$\blacksquare$
\end{preuve}

On donne ici dans le cas général les inégalités de Markov et de
Bienaymé-Chebyshev, déjà vues en CPGE.

\addcontentsline{toc}{subsection}{Inégalité de Markov}
\begin{theoreme}[{Inégalité de
Markov}]\phantomsection\label{IV-inegmarkov}

Soient \(p \in \mathbb{N}^\ast\) et \(a \in \mathbb{R}_+^\ast\). Soit
\(X \in \mathcal{L}^p(\Omega,\mathcal{A},\mathbb{P})\), on a
\[\mathbb{P}(|X|\geq a) \leq \frac{\mathbb{E}(|X|^p)}{a^p}\]
\end{theoreme}

\begin{preuve}
On a \[|X|^p \geq a^p 1_{[a, +\infty[}(|X|)\] et en prenant l'espérance,
on obtient ainsi
\[\mathbb{E}(|X|^p) \geq a^p \mathbb{E}(1_{[a, +\infty[}(|X|)) = a^p \mathbb{P}(|X|\geq a).\]

\hfill$\blacksquare$
\end{preuve}

\addcontentsline{toc}{subsection}{Inégalité de Bienaymé-Chebyshev}
\begin{corollaire}[{Inégalité de
Bienaymé-Chebyshev}]\phantomsection\label{IV-inegbc}

Soit \(X \in \mathcal{L}^2\), on a
\[\mathbb{P}(|X-\mathbb{E}(X)| \geq a) \leq \frac{\mathbb{V}(X)}{a^2}\]
\end{corollaire}

\runinhead{Exercice -- Démonstration (\(\mathord{\bullet}\))}\phantomsection\addcontentsline{toc}{subsubsection}{Démonstration}\label{IV-bienaymuxe9cheby}
\solref{IV-answer-bienaymuxe9cheby}

\addcontentsline{toc}{subsection}{Interprétation}
\begin{remarque}[{Interprétation}]\phantomsection\label{IV-interpruxe9tation}

L'inégalité de Bienaymé-Chebyshev est très utile en pratique. Elle
permet de mesurer la probabilité des grands écarts entre \(X\) et sa
moyenne. Par exemple, avec \(a = 10 \sigma\), il en résulte qu'il est
improbable qu'une variable aléatoire \(X\) dévie de son espérance
\(\mathbb{E}(X)\) de plus de 10 fois son écart-type (la probabilité est
inférieure à 0.01). Cette inégalité, tout à fait générale, n'est
cependant pas très précise, et surestime très souvent en pratique le
membre de gauche. On préférera, quand c'est possible, calculer
directement ces probabilités à partir de la loi de \(X\).
\end{remarque}

\chapter{Suites de variables
indépendantes}\label{IV-suites-de-variables-induxe9pendantes}

Au début du chapitre II, nous avons vu comment construire des couples
(et même des \(n\)-uplets en itérant) de variables aléatoires réelles
indépendantes en considérant l'espace produit, munis des tribus produit
et des (mesures de) probabilités produit. Il est malheureusement
beaucoup plus délicat, mais indispensable pour les applications (en
particulier la loi des grands nombres), de construire une \textbf{suite
infinie de variables indépendantes} de lois données.

Plus précisément, pour chaque entier \(n\) on se donne une v.a.r.
\(X_n\) définie sur un espace de probabilité
\((\Omega_n,\mathcal{A}_n,\mathbb{P}_n)\), à valeurs dans
\((\mathbb{R},\mathcal{B}(\mathbb{R}))\) et de loi \(\mathbb{P}_{X_n}\)
(pour construire chaque \(X_n\), on peut procéder comme ci-dessus).
Ensuite, on pose \begin{align*}
\Omega &= \prod_{n=1}^{\infty}\Omega_n \text{        (produit cartésien dénombrable),}\\
\mathcal{A}&= \bigotimes_{n=1}^{\infty}\mathcal{A}_n,
\end{align*} où \(\bigotimes_{n=1}^{\infty}\mathcal{A}_n\) désigne la
plus petite tribu de \(\Omega\) à laquelle appartienne tous les
ensembles de la forme
\[A_1 \times A_2 \times \ldots\times A_k \times \Omega_{k+1} \times \Omega_{k+2} \times \ldots, \,\,\, A_i \in \mathcal{A}_i, \,\,\, k \in \mathbb{N}^\ast\]

On a alors le théorème suivant, que l'on admettra, qui constitue un
résultat non trivial de la théorie de la mesure et qui généralise le
\href{Probabilité\%20III.pdf\%20\#fubiniproba}{théorème de Fubini}.

\addcontentsline{toc}{subsection}{Mesure produit sur un espace de
dimension infinie}
\begin{theoreme}[{Mesure produit sur un espace de dimension
infinie}]\phantomsection\label{IV-probaespaceden}

Avec les notations ci-dessus, il existe une unique probabilité
\(\mathbb{P}\) sur \((\Omega,\mathcal{A})\), telle que
\[\mathbb{P}(A_1 \times \ldots\times A_k \times \Omega_{k+1} \times \ldots) =\prod_{i=1}^k \mathbb{P}_i(A_i)\]
pour tous \(k \in \mathbb{N}^{\ast}\) et \(A_i \in \mathcal{A}_i\),
\(i \in \{1,\ldots,k\}\).

Nous sommes maintenant en mesure de définir une suite de variables
aléatoires toutes définies sur le même espace de probabilité
\((\Omega, \mathcal{A}, \mathbb{P})\).
\end{theoreme}

\addcontentsline{toc}{subsection}{Suites de variables aléatoires
indépendantes}
\begin{definition}[{Suites de variables aléatoires
indépendantes}]\phantomsection\label{IV-suites-de-variables-aluxe9atoires-induxe9pendantes}

La suite \((X_n)_{n\in \mathbb{N}^\ast}\) de variables aléatoires est
dite \emph{indépendante} si pour tout \(n \in \mathbb{N}^\ast\), la
famille finie \(X_1,\ldots,X_n\) est indépendante.

Il est facile de vérifier que l'indépendance est préservée par certaines
transformations.
\end{definition}

\addcontentsline{toc}{subsection}{Conséquences}
\begin{proposition}[{Conséquences}]\phantomsection\label{IV-consuxe9quences}

L'indépendance de la suite \((X_n)_{n\in \mathbb{N}^\ast}\) entraîne
celle de

\begin{enumerate}
\def\labelenumi{\arabic{enumi}.}
\tightlist
\item
  toute sous-suite \((X_{i_k})_{k\in \mathbb{N}^\ast}\),
\item
  toute suite de vecteurs issus de \(X_n\),
\item
  toute suite de la forme \((f_n(X_n))_{n\in \mathbb{N}^\ast}\), où les
  fonctions \(f_n\) sont des fonctions mesurables.
\end{enumerate}
\end{proposition}

\addcontentsline{toc}{subsection}{Développement dyadique}
\begin{exemple}[{Développement
dyadique}]\phantomsection\label{IV-duxe9veloppement-dyadique}

Nous considérons l'ensemble \(\Omega = [0, 1[\) muni de la tribu
borélienne restreinte à cet ensemble, et de la mesure de Lebesgue. A
chaque réel \(\omega\), nous associons son développement dyadique
(unique si l'on impose que les \(\omega_i\) ne soient pas tous égaux à 1
à partir d'un certain rang)~:
\[ \omega = \sum_{i\in \mathbb{N}^\ast} \frac{\omega_i}{2^i},\,\,\, \omega_i \in \{0,1\}.\]
L'application \(X_i : \Omega \to \{0,1\}\), qui à \(\omega\) associe
\(X_i(\omega) = \omega_i\) est une variable aléatoire sur \(\Omega\). En
effet, pour \(x_i \in \{0,1\}\), \(1\leq i \leq n\),
\[ \{X_i = x_i\} = \bigcup_{x_1,\ldots, x_{i-1} \in \{0,1\}} \left[ \sum_{j=1}^i \frac{x_j}{2^j}, \sum_{j=1}^i \frac{x_j}{2^j} + \frac{1}{2^i}\right[, \]
qui est bien un élément de la tribu borélienne de \(\Omega = [0,1[\), et
\[\mathbb{P}(\{X_i = x_i\}) = \frac{1}{2^i} \sum_{x_1,\ldots, x_{i-1} \in \{0,1\}} 1 = \frac{1}{2}.\]
Montrons l'indépendance des variables aléatoires
\((X_i)_{1\leq i \leq n}\). Nous avons
\[ \bigcap_{1\leq i \leq n} \{X_i = x_i\} = \left[ \sum_{i=1}^n \frac{x_i}{2^i}, \sum_{i=1}^n \frac{x_i}{2^i} + \frac{1}{2^n}\right[ \]
si bien que
\[ \mathbb{P}\left(\bigcap_{1\leq i \leq n} \{X_i = x_i\}\right) = \frac{1}{2^n} = \prod_{1\leq i \leq n} \mathbb{P}(\{X_i = x_i\}).\]
Cela démontre que les variables aléatoires \(X_i\) sont indépendantes et
de loi de Bernoulli de paramètre \(\frac{1}{2}\) . Ainsi, nous venons de
construire sur \(\Omega\) une suite de variables aléatoires telle que
toute sous-suite finie est constituée de variables aléatoires
indépendantes. C'est donc une suite de variables aléatoires
indépendantes de loi de Bernoulli de paramètre \(\frac{1}{2}\) .

On note enfin le résultat très utile suivant portant sur les suites
d'événements indépendants.
\end{exemple}

\addcontentsline{toc}{subsection}{Borel-Cantelli}
\begin{lemme}[{Borel-Cantelli}]\phantomsection\label{IV-BC}

Soit \((A_n)_{n\in\mathbb{N}^\ast}\) une suite d'événements sur l'espace
probabilisé \((\Omega,\mathcal{A}, \mathbb{P})\).

\begin{enumerate}
\def\labelenumi{\arabic{enumi}.}
\item
  Si \(\sum_{n=1}^{\infty} \mathbb{P}(A_n) < \infty\), alors
  \(\mathbb{P}(\limsup\limits_{n \to \infty} A_n) = \mathbb{P}\left(\bigcap_{n\geq 1} \bigcup_{k \geq n} A_k \right) =0\).
\item
  Si \(\sum_{n=1}^{\infty} \mathbb{P}(A_n) = \infty\) et si les
  événements \(A_n\) sont mutuellement indépendants, alors on a
  \(\mathbb{P}(\limsup\limits_{n \to \infty} A_n) = 1\).
\end{enumerate}
\end{lemme}

\begin{preuve}
Exercice.\hfill$\blacksquare$
\end{preuve}

\runinhead{Exercice -- Pile ou face
infini (\(\mathord{\bullet}\))}\phantomsection\addcontentsline{toc}{subsubsection}{Pile
ou face infini}\label{IV-pf}
On considère un jeu de pile ou face infini où on a une probabilité
\(p, 0< p < 1\), d'obtenir face. Montrer que l'événement
\(A = \{\omega : \text{ il n'y a qu'un nombre fini de faces}\}\) est de
probabilité nulle. (Considérer les événements
\(A_n = \{\text{on a face au }n\text{-ième tirage}\}\), puis montrer que
\(\sum_{n=1}^\infty \mathbb{P}(A_n) = \infty\)).\solref{IV-answer-pf}

\chapter{Convergences et loi des grands
nombres}\label{IV-convergences-et-loi-des-grands-nombres}

\section{Convergences des variables
aléatoires}\label{IV-convergences-des-variables-aluxe9atoires}

Dans ce paragraphe, on va décrire les notions de convergence de
variables (on y inclut les vecteurs) aléatoires. On verra que plusieurs
notions sont possibles, non équivalentes, ce qui enrichit mais complique
aussi la description des comportements asymptotiques.

En calcul intégral, on a beaucoup étudié le cas de suites de fonctions
convergeant simplement. Une variable aléatoire étant une fonction, on a
donc la même notion qui revient à écrire qu'une suite
\((X_n)_{n\in \mathbb{N}^\ast}\) de variables aléatoires converge
simplement vers \(X\) si \(\lim_{n \to \infty} X_n(\omega) = X(\omega)\)
pour tout \(\omega \in \Omega\). Cette définition naturelle est
malheureusement à peu près inutile en probabilités, comme l'illustre
l'exemple suivant.

\addcontentsline{toc}{subsection}{Pile ou face infini (bis)}
\begin{exemple}[{Pile ou face infini
(bis)}]\phantomsection\label{IV-pfinf}

Soit \((X_n)_{n\in \mathbb{N}^\ast}\) une suite de v.a. réelles qui sont
indépendantes et de même loi (on notera \emph{i.i.d} pour
\emph{indépendantes et identiquement distribuées}), avec
\(\mathbb{P}(X_n = 1) = p\) et \(\mathbb{P}(X_n=0) = 1-p\). Ce sont donc
des v.a. de loi de Bernoulli de paramètre \(p\) qui modélisent par
exemple les résultats d'un jeu de pile ou face. Lorsque \(n\) est grand,
on s'attend à ce que la proportion de faces (\(X_n=1\)) soit à peu près
égale à \(p\) (c'est l'essence de la conception objectiviste des
probabilités). Mathématiquement, on voudrait que
\[\lim_{n \to \infty} \frac{X_1(\omega) + \ldots +X_n(\omega)}{n} = p \text{   pour tout }\omega\in\Omega.\]
C'est en fait complètement faux. En effet, si on considère par exemple
la suite \(\omega_0 = {p,p,p,\ldots}\) qui ne contient que des piles, on
obtient
\[\lim_{n \to \infty} \frac{X_1(\omega_0) + \ldots +X_n(\omega_0)}{n} = 0\]
et plus généralement, on a
\(\lim_{n \to \infty} \frac{1}{n}\sum_{i=1}^n X_i(\omega) = 0\) pour
tout \(\omega\) dans l'ensemble
\(A = \{\omega : \text{ il n'y a qu'un nombre fini de faces}\}\). On
peut même trouver des \(\omega\) pour lesquels la fréquence converge
vers n'importe quel nombre fixé dans \([0,1]\)\,! Bien entendu,
l'événement \(A\) est invraisemblable et on peut montrer qu'il vérifie
\(\mathbb{P}(A) = 0\)\footnote{cf.~\protect\hyperref[IV-pf]{exercice
  plus haut (p.~\pageref*{IV-pf})}}. On verra cependant que la
\textbf{loi des grands nombres} assure~:
\[\mathbb{P}\left(\left\{\omega ; \lim_{n \to \infty} \frac{1}{n}\sum_{i=1}^n X_i(\omega) = p\right\}\right)=1.\]
Ce type de convergence, pour lequel on n'a pas convergence pour tout
\(\omega\), mais seulement pour presque tout \(\omega\) (autrement dit
\(\mathbb{P}\)-presque partout) est typiquement ce qui arrive pour des
variables aléatoires.

On considère une suite \((X_n)_{n\in \mathbb{N}^\ast}\) de variables
aléatoires, définies sur un même espace de probabilité
\((\Omega, \mathcal{A}, \mathbb{P})\), et à valeurs dans
\(\mathbb{R}^d\). On considère également sur le même espace un vecteur
``limite'' \(X\). On notera \(|\cdot|\) la valeur absolue dans
\(\mathbb{R}\) ou la norme euclidienne dans \(\mathbb{R}^d\).
\end{exemple}

\addcontentsline{toc}{subsection}{Modes de convergences}
\begin{definition}[{Modes de
convergences}]\phantomsection\label{IV-modes-de-convergences}

\begin{enumerate}
\def\labelenumi{\arabic{enumi}.}
\tightlist
\item
  La suite \((X_n)_{n\in \mathbb{N}^\ast}\) \emph{converge presque
  sûrement} vers \(X\), ce qui s'écrit \(X_n \to X\) p.s., s'il existe
  un ensemble \(N \in \mathcal{A}\) de probabilité nulle tel que
  \[ X_n (\omega) \xrightarrow[n \to \infty]{}  X(\omega),\,\,\, \forall \omega \notin N.\]
\item
  La suite \((X_n)_{n\in \mathbb{N}^\ast}\) \emph{converge en
  probabilité} vers \(X\), ce qui s'écrit
  \(X_n \xrightarrow{\mathbb{P}} X\), si \(\forall \varepsilon >0\) on a
  \[ \mathbb{P}(|X_n-X| \geq \varepsilon) \xrightarrow[n \to \infty]{}  0. \]
\item
  La suite \((X_n)_{n\in \mathbb{N}^\ast}\) \emph{converge en moyenne}
  (ou dans \(\mathcal{L}^1\)) vers \(X\), ce qui s'écrit
  \(X_n \xrightarrow{\mathcal{L}^1} X\), si \(X_n\) et \(X\) sont dans
  \(\mathcal{L}^1\) et si
  \[ \mathbb{E}(|X_n - X|) \xrightarrow[n \to \infty]{}  0 .\]
\end{enumerate}
\end{definition}

\addcontentsline{toc}{subsection}{Convergence \(\mathcal{L}^p\)}
\begin{remarque}[{Convergence
\(\mathcal{L}^p\)}]\phantomsection\label{IV-convergence-mathcallp}

La définition de la convergence dans \(\mathcal{L}^1\) se généralise aux
ordres supérieurs, pour \(p \in \mathbb{N}^\ast\), on parle de
convergence dans \(\mathcal{L}^p\) (en moyenne quadratique si \(p=2\))
ce qui s'écrit \(X_n \xrightarrow{\mathcal{L}^p} X\), si \(X_n\) et
\(X\) sont dans \(\mathcal{L}^p\) et si
\[ \mathbb{E}(|X_n - X|^p) \xrightarrow[n \to \infty]{}  0 .\]

Ces modes de convergences ne sont pas équivalents comme le montrent les
exemples suivants.
\end{remarque}

\runinhead{Exercice -- En proba et en
moyenne (\(\mathord{\bullet}\))}\phantomsection\addcontentsline{toc}{subsubsection}{En
proba et en moyenne}\label{IV-lien}
Soit \((X_n)_{n\in \mathbb{N}^\ast}\) une suite de variables aléatoires
de Bernoulli à valeurs dans \(\{0,1\}\) telles que
\[ \mathbb{P}(X_n=1) = \frac{1}{n} ;\,\,\,\, \mathbb{P}(X_n = 0) = 1- \frac{1}{n}.\]
Montrer que la suite \((X_n)_{n\in \mathbb{N}^\ast}\) converge en
probabilité et en moyenne vers 0.\solref{IV-answer-lien}

\runinhead{Exercice -- En proba mais pas en
moyenne (\(\mathord{\bullet}\))}\phantomsection\addcontentsline{toc}{subsubsection}{En
proba mais pas en moyenne}\label{IV-lien2}
Soit maintenant une suite \((Y_n)_{n\in \mathbb{N}^\ast}\) de variables
aléatoires de Bernoulli à valeurs dans \(\{0,n^2\}\) telles que
\[ \mathbb{P}(Y_n=n^2) = \frac{1}{n} ;\,\,\,\, \mathbb{P}(Y_n = 0) = 1- \frac{1}{n}.\]
Montrer que la suite \((Y_n)_{n\in \mathbb{N}^\ast}\) converge en
probabilité vers 0 mais pas en moyenne.\solref{IV-answer-lien2}

\runinhead{Exercice -- Presque
sûre (\(\mathord{\bullet}\))}\phantomsection\addcontentsline{toc}{subsubsection}{Presque
sûre}\label{IV-lien3}
Soit \(U\) une variable aléatoire uniforme sur \([0 , 1]\). Posons
\(Z_n = 1_{\{ U \leq \frac{1}{n}\} }\). Alors
\[ \mathbb{P}(Z_n = 1)= \frac{1}{n} ;\,\,\,\, \mathbb{P}(Z_n = 0) = 1 - \frac{1}{n}\]
Montrer que la suite \((Z_n)_{n\in\mathbb{N}^\ast}\) converge presque
sûrement vers 0.\solref{IV-answer-lien3}

Il existe cependant des liens entre ces différents modes de convergences
que l'on explore ci-dessous.

\addcontentsline{toc}{subsection}{p.s. / \(\mathcal{L}^1\)
\(\Rightarrow\) en proba}
\begin{proposition}[{p.s. / \(\mathcal{L}^1\) \(\Rightarrow\) en
proba}]\phantomsection\label{IV-propconv1}

La convergence presque sûre et la convergence en moyenne entraînent la
convergence en probabilité.
\end{proposition}

\begin{preuve}
Soient \(n\in\mathbb{N}^\ast\), \(\varepsilon > 0\) et
\(A_{n,\varepsilon} = \{|X_n - X| > \varepsilon \}\).

\begin{itemize}
\item
  Supposons que \(X_n \to X\) p.s. et soit \(N\) l'ensemble de
  probabilité nulle en dehors duquel on a \(X_n(\omega) \to X(\omega)\).
  Si \(\omega \notin N\), on a \(\omega \notin A_{n,\varepsilon}\) pour
  tout \(n \geq n_0\), où \(n_0\) dépend de \(\omega\) et de
  \(\varepsilon\), ce qui implique que les variables aléatoires
  \(Y_{n,\varepsilon} = 1_{N^c\cap A_{n,\varepsilon}}\) tendent
  simplement vers 0 lorsque \(n \to \infty\). Comme on a aussi
  \(0 \leq Y_{n,\varepsilon} \leq 1\) le théorème de convergence dominée
  implique que
  \(\mathbb{E}(Y_{n,\varepsilon}) \xrightarrow[n \to \infty]{} 0.\) Par
  suite,
  \[\mathbb{P}(A_{n,\varepsilon}) \leq \mathbb{P}(N^c \cap A_{n,\varepsilon}) + \mathbb{P}(N) = \mathbb{P}(N^c \cap A_{n,\varepsilon}) = \mathbb{E}(Y_{n,\varepsilon}) \xrightarrow[n \to \infty]{} 0.\]
\item
  Supposons que \(X_n \xrightarrow{\mathcal{L}^1} X\). Pour
  \(\varepsilon >0\), on a
  \(1_{A_{n,\varepsilon}} \leq \frac{1}{\varepsilon}|X_n - X|\), donc
  \[ \mathbb{P}(A_{n,\varepsilon}) \leq \frac{1}{\varepsilon}\mathbb{E}(|X_n - X|) \xrightarrow[n \to \infty]{} 0.\]
\end{itemize}

\hfill$\blacksquare$
\end{preuve}

La convergence en probabilité n'entraîne pas la convergence en moyenne,
comme nous l'avons vu dans \protect\hyperref[IV-lien2]{l'exemple
précédent (p.~\pageref*{IV-lien2})}, ne serait-ce que parce qu'elle
n'implique pas l'appartenance de \(X_n\) et \(X\) à \(\mathcal{L}^1\).
Si les \(X_n\) ne sont ``pas trop grands'', il y a cependant équivalence
entre les deux modes de convergence. En voici un exemple.

\addcontentsline{toc}{subsection}{Cas borné}
\begin{proposition}[{Cas borné}]\phantomsection\label{IV-propconv2}

S'il existe une constante \(a\) telle que \(|X_n| \leq a\) presque
sûrement, il y a équivalence entre \(X_n \xrightarrow{\mathbb{P}} X\) et
\(X_n \xrightarrow{\mathcal{L}^1} X\).
\end{proposition}

\begin{preuve}
Etant donnée la \protect\hyperref[IV-propconv1]{proposition
précédente (p.~\pageref*{IV-propconv1})}, dont on reprend les notations,
il suffit de montrer que la convergence en probabilité implique la
convergence en moyenne lorsque \(|X_n| \leq a\).

Comme \(|X_n| \leq a\) p.s., on a
\(\{|X| > a +\varepsilon\} \subset A_{n,\varepsilon}\), et donc
\(\mathbb{P}(|X| > a +\varepsilon ) \leq \mathbb{P}(A_{n,\varepsilon})\).
En faisant \(n \to \infty\), on en déduit que
\(\mathbb{P}(|X| > a + \varepsilon) = 0\). Ceci est vrai pour tout
\(\varepsilon\) et donc \[\mathbb{P}(|X|> a) = 0.\] Comme
\(|X_n| \leq a\), on a aussi
\[ |X_n -X| \leq \varepsilon + (X_n + X) 1_{A_{n,\varepsilon}} \leq \varepsilon + 2a 1_{A_{n,\varepsilon}}\]
sur l'ensemble \(\{|X| \leq a \}\) qui est de probabilité 1. On a ainsi
\[ \mathbb{E}(|X_n - X|) \leq \varepsilon +2a \mathbb{P}(A_{n,\varepsilon})\]
On en déduit que
\(\lim \sup_{n \to \infty} \mathbb{E}(|X_n - X|) \leq \varepsilon\), et
comme \(\varepsilon\) est arbitrairement proche de 0, on a le résultat
souhaité.\hfill$\blacksquare$
\end{preuve}

Les rapports entre convergence presque-sûre et convergence en
probabilité sont plus subtils. La première de ces deux convergences est
plus forte que la seconde d'après la
\protect\hyperref[IV-propconv1]{proposition plus
haut (p.~\pageref*{IV-propconv1})}, mais ``à peine plus'', comme le
montre le résultat suivant.

\addcontentsline{toc}{subsection}{Existence d'une sous-suite
convergente}
\begin{proposition}[{Existence d'une sous-suite
convergente}]\phantomsection\label{IV-propconv3}

Si \(X_n \xrightarrow{\mathbb{P}} X\), il existe une sous-suite
\((n_k)\) telle que \(X_{n_k} \to X\) p.s. quand \(k \to \infty\).
\end{proposition}

\begin{preuve}
On remarque d'abord que l'on peut réécrire la définition de la
convergence en probabilité de la manière suivante~:
\(\forall \varepsilon >0\) et \(\delta > 0\), il existe
\(N = N(\delta,\varepsilon)\) tel que
\[n\geq N \Rightarrow \mathbb{P}(|X_n-X| > \varepsilon ) < \delta.\]

Comme la suite \((X_n)_{n\in \mathbb{N}^\ast}\) converge en probabilité
vers \(X\), on peut définir une sous-suite de la manière suivante~:
posons \(n_1 = 1\), et soit
\[n_j = \inf \left\{ n > n_{j-1} ; \mathbb{P}\left(|X_m - X| > \frac{1}{2^j} \right) < \frac{1}{3^j}, \text{ pour } m \geq n \right\}.\]
On a alors
\[\sum_j \mathbb{P}\left(|X_{n_j} - X| > \frac{1}{2^j} \right) < \sum_j \frac{1}{3^j} < \infty\]
et en appliquant le \protect\hyperref[IV-BC]{lemme de
Borel-Cantelli (p.~\pageref*{IV-BC})} aux ensembles
\[A_j =\left\{|X_{n_{j}} - X| > \frac{1}{2^j} \right\},\] on obtient que
la suite \((X_{n_j})_{j\in \mathbb{N}^\ast}\) converge
presque-sûrement.\hfill$\blacksquare$
\end{preuve}

\runinhead{Exercice -- Intervalles de longueur aléatoire
(1) (\(\mathord{\bullet}\))}\phantomsection\addcontentsline{toc}{subsubsection}{Intervalles
de longueur aléatoire (1)}\label{IV-ia}
Soient \(\Omega = \mathbb{R}\) muni de sa tribu borélienne et
\(\mathbb{P}\) la probabilité uniforme sur {[}0, 1{]}. Soit
\(X_n = 1_{A_n}\) , où \(A_n\) est un intervalle de {[}0, 1{]} de
longueur \(\frac{1}{n}\).

Montrer que la suite des \(X_n\) converge en moyenne et en
proba.\solref{IV-answer-ia}

\runinhead{Exercice -- Intervalles de longueur aléatoire
(2) (\(\mathord{\bullet}\mathord{\bullet}\mathord{\bullet}\))}\phantomsection\addcontentsline{toc}{subsubsection}{Intervalles
de longueur aléatoire (2)}\label{IV-ia2}
Montrer que la suite des \(X_n\) ne converge pas presque sûrement.
Extraire une sous-suite qui converge presque
sûrement.\solref{IV-answer-ia2}

\addcontentsline{toc}{subsection}{Continuité}
\begin{proposition}[{Continuité}]\phantomsection\label{IV-propconv4}

Soit \(f\) une fonction continue de \(\mathbb{R}^d\) dans
\(\mathbb{R}\).

\begin{enumerate}
\def\labelenumi{\arabic{enumi}.}
\tightlist
\item
  Si \(X_n \to X\) presque-sûrement, alors \(f(X_n) \to f(X)\)
  presque-sûrement.
\item
  Si \(X_n \xrightarrow{\mathbb{P}} X\), alors
  \(f(X_n) \xrightarrow{\mathbb{P}} f(X)\).
\end{enumerate}
\end{proposition}

\begin{preuve}
\begin{enumerate}
\def\labelenumi{\arabic{enumi}.}
\item
  Soit \(N\) l'ensemble de probabilité nulle en dehors duquel on a
  \(X_n(\omega) \xrightarrow[n \to \infty]{} X(\omega)\). Si
  \(\omega \notin N\), il vient
  \[ \lim_{n \to \infty}f(X_n(\omega)) = f(\lim_{n \to \infty}X_n(\omega)) = f(X(\omega))\]
  par continuité de \(f\), d'où le résultat.
\item
  On remarque d'abord que si \(K >0\) et \(\varepsilon >0\),
  \[\{|f(X_n)-f(X)| \geq \varepsilon\} \subset \{|X| > K\}\cup\{|X| \leq K, |f(X_n)-f(X)| \geq \varepsilon\}.\]
  La fonction \(f\) est uniformément continue sur
  \(\{x : |x| \leq K\}\), donc il existe \(\eta > 0\) tel que
  \(|x-y| <\eta\) et \(|x| \leq 2K\) et \(|y|\leq 2K\) impliquent
  \(|f(x) - f(y)| <\varepsilon\). On a donc
  \[\{|f(X_n)-f(X)| \geq \varepsilon\} \subset \{|X| > K\}\cup\{|X_n-X| \geq \eta\}\]
  d'où
  \[\mathbb{P}(|f(X_n)-f(X)| \geq \varepsilon) \leq \mathbb{P}(|X| > K) + \mathbb{P}(|X_n-X| \geq \eta).\]
  Par hypothèse, il vient
  \[ \lim_{n \to \infty} \mathbb{P}(|f(X_n)-f(X)| \geq \varepsilon) \leq \mathbb{P}(|X| > K) .\]
  Enfin, \(\lim_{K\to +\infty} \mathbb{P}(|X| > K) = 0\) (par
  convergence dominée) et donc la limite ci-dessus est nulle. On a ainsi
  le résultat.
\end{enumerate}

\hfill$\blacksquare$
\end{preuve}

\section{La loi des grands nombres}\label{IV-la-loi-des-grands-nombres}

On présente maintenant l'un des résultats essentiels de la théorie des
probabilités. Ce résultat montre rigoureusement que, quand le nombre de
répétitions de l'expérience tend vers l'infini, la fréquence de
réalisation d'un événement converge p.s. vers la probabilité de
réalisation de cet événement. Ce résultat, appelé \textbf{Loi des grands
nombres}, a d'autres portées fondamentales. Il est en particulier à
l'origine de méthodes de calcul numérique appelées Méthodes de
Monte-Carlo, qui sont extrêmement puissantes et robustes. Elles sont par
exemple très utilisées en Physique, en Mathématiques Financières ou
encore dans les problèmes de quantification d'incertitudes. Une
introduction à ces méthodes sera donnée au chapitre V.

Dans ce paragraphe, on considère une suite
\((X_n)_{n\in\mathbb{N}^\ast}\) de variables aléatoires
\textbf{indépendantes et de même loi} (ou indépendantes et identiquement
distribuées, i.i.d. en abrégé). On considère la ``moyenne'' des \(n\)
premières variables aléatoires~:
\[ M_n = \frac{X_1 + \ldots + X_n}{n},\] et notre objectif est de
montrer que \(M_n\) converge vers l'espérance des \(X_n\) lorsque
celle-ci existe (comme les \(X_n\) ont même loi, cette espérance est la
même pour tout \(n\)).

Nous allons démontrer dans un premier temps la loi des grands nombres
pour des variables aléatoires de carré intégrable.

\addcontentsline{toc}{subsection}{Loi des grands nombres cas
\(\mathcal{L}^2\)}
\begin{theoreme}[{Loi des grands nombres cas
\(\mathcal{L}^2\)}]\phantomsection\label{IV-loi-des-grands-nombres-cas-mathcall2}

Soit \((X_n)_{n\in\mathbb{N}^\ast}\) une suite de variables aléatoires
indépendantes, de même loi et de \textbf{carré intégrable}, et
\(m = \mathbb{E}(X_n)\) leur espérance. Alors la suite
\((M_n)_{n\in\mathbb{N}^\ast}\) définie par
\[M_n = \frac{X_1 + \ldots + X_n}{n}\] converge vers \(m\),
\textbf{presque sûrement et en moyenne}, quand \(n\) tend vers l'infini.
Elle converge donc aussi en probabilité. On a même convergence en
\emph{moyenne quadratique}, à savoir que~:
\[ \mathbb{E}((M_n - m)^2) \xrightarrow[n \to \infty]{} 0.\]
\end{theoreme}

\runinhead{Exercice -- Loi faible des grands
nombres (\(\mathord{\bullet}\))}\phantomsection\addcontentsline{toc}{subsubsection}{Loi
faible des grands nombres}\label{IV-wlln}
Le résultat sur la convergence en probabilité est appelé \emph{loi
faible des grands nombres}. Démontrer ce
résultat.\solref{IV-answer-wlln}

Le résultat est peu informatif et permet d'obtenir certains contrôles
d'erreurs. Le résultat prouvant la convergence presque-sûre est appelé
\emph{loi forte des grands nombres}. Sa preuve est plus délicate et
utilise le \protect\hyperref[IV-BC]{lemme de
Borel-Cantelli (p.~\pageref*{IV-BC})}.

\subsection{Démonstration}\label{IV-duxe9monstration-7}

Notons \(\sigma^2\) la variance des variables \(X_n\), bien définie
puisqu'on les a supposées de carré intégrable. En vertu de la linéarité
de l'espérance, on a
\[ \mathbb{E}(M_n) = m,\,\,\,\, \mathbb{E}((M_n-m)^2) = \mathbb{V}(M_n) = \frac{\sigma^2}{n},\]
d'où la convergence en moyenne quadratique.

Comme \(\mathbb{E}(Y)^2 \leq \mathbb{E}(Y^2)\), on en déduit que
\(\mathbb{E}(|M_n -m|)\to 0\), donc on a aussi la convergence en
moyenne.

La preuve de la convergence presque-sûre est plus délicate.

Quitte à remplacer \(X_n\) par \(X_n - m\) (et donc \(M_n\) par
\(M_n - m\)), nous pouvons supposer que \(m = 0\).

Montrons tout d'abord que la sous-suite
\((M_{n^2})_{n \in \mathbb{N}^\ast}\) converge presque-sûrement vers 0.

La convergence dans \(\mathcal{L}^1\) impliquant la convergence en
probabilité, on sait qu'on peut extraire de \((M_n)_n\) une sous-suite
convergeant p.s. vers 0. Cependant, cela ne suffit pas puisqu'on veut
que la suite \((M_n)_n\) elle-même converge p.s. Pour le montrer, on
construit d'abord une sous suite-particulière qui converge p.s. puis on
traite les termes qui se trouvent entre deux éléments successifs de la
sous-suite.

D'après l'\protect\hyperref[IV-inegbc]{inégalité de
Bienaymé-Chebyshev (p.~\pageref*{IV-inegbc})} et ce qui précède, on a
pour \(q\in \mathbb{N}^\ast\)
\[\mathbb{P}\left(|M_{n^2}|\geq \frac{1}{q}\right) \leq \frac{\sigma^2 q^2}{n^2}\]

Donc si \(A_{n,q} = \{|M_{n^2}|\geq \frac{1}{q}\}\), nous obtenons que
\(\sum_{n\geq 1} \mathbb{P}(A_{n,q}) < \infty\). Posons ensuite
\(B_{n,q} = \cup_{m\geq n} A_{m,q}\) et
\(C_q = \cap_{n \geq 1} B_{n,q} = \lim\sup_n A_{n,q}\). En appliquant le
\protect\hyperref[IV-BC]{lemme de Borel-Cantelli (p.~\pageref*{IV-BC})},
on obtient que \(\mathbb{P}(C_q) = 0\). En conséquence, si on pose
\(N = \cup_{q \in \mathbb{N}^\ast} C_q\), on obtient
\(\mathbb{P}(N) \leq \sum_{q\in\mathbb{N}^\ast} \mathbb{P}(C_q)= 0\).

Si \(\omega \notin N\), alors
\(\omega \in \cap_{q \in \mathbb{N}^\ast} (C_q)^c\). Ainsi,
\(\omega \notin C_q\) pour tout \(q \geq 1\), et donc
\(\omega \notin B_{n,q}\) pour \(n\) assez grand (car \(B_{n,q}\) est
décroissant en \(n\)). Cela signifie que pour tout \(\omega \notin N\),
pour tout \(q \geq 1\), il existe un \(n\) assez grand tel que
\(M_{k^2} \leq \frac{1}{q}\) dès que \(k \geq n\). Autrement dit,
\(M_{n^2} \to 0\) si \(\omega \notin N\), avec \(\mathbb{P}(N) = 0\),
d'où \[ M_{n^2} \xrightarrow[n \to \infty]{} 0 \text{ p.s.}\]

Montrons maintenant que la suite \((M_n)_{n\in\mathbb{N}^\ast}\) tend
presque-sûrement vers 0.

Pour tout entier \(n\), notons \(p(n)\), l'entier tel que
\(p(n)^2 \leq n \leq (p(n)+1)^2\). Alors,
\[ M_n - \frac{p(n)^2}{n}M_{p(n)^2} = \frac{1}{n} \sum_{p = p(n)^2+1}^{n} X_p,\]
et puisque les variables aléatoires de la somme sont indépendantes, il
vient \begin{align*}
\mathbb{E}\left(\left(M_n - \frac{p(n)^2}{n}M_{p(n)^2}\right)^2\right) & = \frac{n-p(n)^2}{n^2}\sigma^2 \\
                                                                 & \leq \frac{2p(n)+1}{n^2} \sigma^2 \\
                                                                 & \leq \frac{2\sqrt{n}+1}{n^2} \sigma^2 \leq \frac{3}{n^{3/2}}\sigma^2
\end{align*}

où on a utilisé le fait que \(p(n) \leq \sqrt{n}\).

En appliquant de nouveau l'\protect\hyperref[IV-inegbc]{inégalité de
Bienaymé-Chebyshev (p.~\pageref*{IV-inegbc})}, on obtient
\[ \mathbb{P}\left(\left|M_n - \frac{p(n)^2}{n}M_{p(n)^2}\right|>a\right) \leq \frac{2\sqrt{n}+1}{n^2} \frac{\sigma^2}{a^2}\]
Comme la série \(\sum_n \frac{2\sqrt{n}+1}{n^2}\) converge, le même
raisonnement que pour \(M_{n^2}\) décrit ci-dessus, montre que
\[ M_n - \frac{p(n)^2}{n}M_{p(n)^2} \to 0 \text{ p.s.}\] Par ailleurs,
on a déjà montré que \(M_{p(n)^2} \to 0\) p.s. et
\(\frac{p(n)^2}{n} \to 1\). On en déduit que \(M_n \to 0\) p.s.

Plus généralement, le théorème suivant donne les hypothèses minimales
assurant la validité de la loi des grands nombres, à savoir que les
\(X_n\) sont dans \(\mathcal{L}_1\) (on se référera par exemple à
\href{https://perso.univ-rennes1.fr/ismael.bailleul/AGREG/COURS/LFGN.pdf}{ce
document en ligne} ou à Jacod and Protter
(\protect\hyperlink{refIV-Jacod}{2003}) pour la démonstration).

\addcontentsline{toc}{subsection}{Loi des grands nombres cas
\(\mathcal{L}^1\)}
\begin{theoreme}[{Loi des grands nombres cas
\(\mathcal{L}^1\)}]\phantomsection\label{IV-lfgn}

Soit \((X_n)_{n\in\mathbb{N}^\ast}\) une suite de variables aléatoires
indépendantes, de même loi et \textbf{intégrables}, et
\(m = \mathbb{E}(X_n)\) leur espérance. Alors la suite
\((M_n)_{n\in\mathbb{N}^\ast}\) définie par
\[M_n = \frac{X_1 + \ldots + X_n}{n}\] converge vers \(m\),
\textbf{presque sûrement et en moyenne}, quand \(n\) tend vers l'infini.
\end{theoreme}

\runinhead{Exercice -- Exponentielle (\(\mathord{\bullet}\))}\phantomsection\addcontentsline{toc}{subsubsection}{Exponentielle}\label{IV-expo}
Soit \((X_i)_{i \in \mathbb{N}^\ast}\) une suite de variables aléatoires
i.i.d. d'espérance \(m\) et \(Y_i = e^{X_i}\). Montrer que~:
\[\left(\prod_{i=1}^n Y_i\right)^{1/n}\] converge presque sûrement vers
une constante à déterminer.\solref{IV-answer-expo}

\runinhead{Exercice -- Variance (\(\mathord{\bullet}\))}\phantomsection\addcontentsline{toc}{subsubsection}{Variance}\label{IV-variance}
Soit \((X_i)_{i \in \mathbb{N}^\ast}\) une suite de variables aléatoires
i.i.d. d'espérance \(m\) et de variance \(\sigma^2\). Montrer que~:
\[\lim_{n \to \infty} \frac{1}{n}\sum_{i=1}^n (X_i - m)^2 \to \sigma^2 \text{   p.s.}\]

\solref{IV-answer-variance}

\chapter{Convergence en loi et théorème central
limite}\label{IV-convergence-en-loi-et-thuxe9oruxe8me-central-limite-1}

Nous allons introduire maintenant une nouvelle notion de convergence de
suites de variables aléatoires. La convergence en loi définie dans ce
paragraphe va concerner les lois des variables aléatoires. Elle
signifiera que les lois sont asymptotiquement ``proches'', sans que les
variables aléatoires elles-mêmes le soient nécessairement.

\section{Convergence en loi}\label{IV-convergence-en-loi}

On considère des vecteurs aléatoires \(X_n\) et \(X\), tous à valeurs
dans le même espace \(\mathbb{R}^d\), mais pouvant éventuellement être
définis sur des espaces de probabilité différents.

\addcontentsline{toc}{subsection}{Convergence en loi}
\begin{definition}[{Convergence en
loi}]\phantomsection\label{IV-defconvloi}

On dit que la suite \((X_n)_{n\in \mathbb{N}^\ast}\) \emph{converge en
loi} vers \(X\) et on écrit \(X_n \xrightarrow{\mathcal{L}} X\), si pour
toute fonction réelle \(f\) \textbf{continue bornée} sur
\(\mathbb{R}^d\),
\[\mathbb{E}(f(X_n)) \xrightarrow[n \to \infty]{} \mathbb{E}(f(X)).\]
\end{definition}

\addcontentsline{toc}{subsection}{Cas fini}
\begin{exemple}[{Cas fini}]\phantomsection\label{IV-cas-fini}

Un cas très simple est celui où toutes les variables aléatoires \(X_n\)
prennent un nombre fini de valeurs \(\{ x_i , 1 \leq i \leq N \}\).
Alors, la suite \((X_n)_{n \in \mathbb{N}^\ast}\) converge en loi vers
\(X\) si et seulement si
\[\lim_{n \to +\infty} \mathbb{P}(X_n = x_i ) = \mathbb{P}(X = x_i),\,\,\, \forall 1 \leq i \leq N\]
Il suffit d'écrire pour \(f\) continue bornée
\[ \mathbb{E}(f (X_n)) = \sum_{i=1}^N f (x_i ) P(X_n = x_i) \]

Dans l'exemple ci-dessus, \(N\) est fini et fixé. Mais nous avons un
résultat analogue (en faisant tendre \(N\) vers l'infini) si les
variables aléatoires ont un nombre dénombrable de valeurs. En
particulier, le cas de la convergence de la loi binomiale vers la loi de
Poisson a été traité en CPGE.
\end{exemple}

\addcontentsline{toc}{subsection}{Univers de définition}
\begin{remarque}[{Univers de
définition}]\phantomsection\label{IV-univers-de-duxe9finition}

Dans la
\protect\hyperref[IV-defconvloi]{définition (p.~\pageref*{IV-defconvloi})},
les v.a. \(X_n\) et \(X\) peuvent être définies sur des univers
distincts puisque seules leurs lois sont en cause. Il arrive même qu'une
suite \(X_n\) converge vers une limite \(X\) qui ne peut pas exister sur
les espaces sur lesquels sont définies les \(X_n\), parce que ceux-ci
sont trop ``petits''~: par exemple, si \(X_n\) est une variable
binomiale à \(n\) modalités, convenablement normalisée, et la limite
\(X\) est gaussienne (on pourra justifier ceci par le
\protect\hyperref[IV-TCL]{théorème central
limite (p.~\pageref*{IV-TCL})}); l'espace naturel sur lequel est définie
\(X_n\) contient \(n+1\) points, et sur un tel espace toutes les v.a.
sont discrètes. La convergence en loi permet donc une sorte de
convergence pour des v.a. pour lesquelles toute autre forme de
convergence serait impossible.
\end{remarque}

\runinhead{Exercice -- Suite
gaussienne (\(\mathord{\bullet}\))}\phantomsection\addcontentsline{toc}{subsubsection}{Suite
gaussienne}\label{IV-suitegauss}
Soit \((X_n)_{n \in \mathbb{N}^\ast}\) et \(X\) des variables aléatoires
de lois respectives \(\mathcal{N}(0,\sigma^2_n)\) et
\(\mathcal{N}(0,\sigma^2)\). On suppose que la suite de réels positifs
\((\sigma_n)_{n\in \mathbb{N}^\ast}\) converge vers \(\sigma > 0\) quand
\(n\) tend vers l'infini. Montrer que la suite
\((X_n)_{n \in \mathbb{N}^\ast}\) converge en loi vers
\(X\).\solref{IV-answer-suitegauss}

La convergence en loi est plus faible que la convergence en probabilité
et donc aussi que les convergences presque-sûre et en moyenne.

\addcontentsline{toc}{subsection}{En proba \(\Rightarrow\) en loi}
\begin{proposition}[{En proba \(\Rightarrow\) en
loi}]\phantomsection\label{IV-en-proba-rightarrow-en-loi}

Soient \((X_n)_{n\in\mathbb{N}\ast}\) et \(X\) des v.a., toutes définies
sur le même espace de probabilité \((\Omega, \mathcal{A}, \mathbb{P})\).
Si \(X_n \xrightarrow{\mathbb{P}} X\), alors
\(X_n\xrightarrow{\mathcal{L}} X\).
\end{proposition}

\begin{preuve}
Soit \(f\) une fonction réelle continue bornée. D'après la
\protect\hyperref[IV-propconv4]{proposition de
continuité (p.~\pageref*{IV-propconv4})}, on a
\(f(X_n) \xrightarrow{\mathbb{P}} f(X)\) et donc \(f(X_n)\) converge
aussi en moyenne vers \(f(X)\) par la
\protect\hyperref[IV-propconv2]{proposition --- cas
borné (p.~\pageref*{IV-propconv2})}. Comme
\(|\mathbb{E}(Y)|\leq \mathbb{E}(|Y|)\) pour toute variable aléatoire
réelle \(Y\), on en déduit
\(\mathbb{E}(f(X_n)) \to \mathbb{E}(f(X))\).\hfill$\blacksquare$
\end{preuve}

On peut finalement résumer les liens d'implications entre les différents
modes de convergence à l'aide de la figure ci-desous.

\begin{figure}
\centering
\includegraphics{cvgces.tex.pdf}
\caption{Relations entre modes de convergence}
\end{figure}

Un moyen efficace de caractériser la convergence en loi des variables
aléatoires réelles passe par l'étude de la suite des fonctions de
répartition.

\addcontentsline{toc}{subsection}{Convergence en loi et fonction de
répartition}
\begin{proposition}[{Convergence en loi et fonction de
répartition}]\phantomsection\label{IV-cvceloifdr}

Soient \((X_n)_{n\in\mathbb{N}^\ast}\) et \(X\) des variables aléatoires
réelles de fonctions de répartition respectives \(F_n\) et \(F\). Pour
que \(X_n \xrightarrow{\mathcal{L}} X\), il faut et il suffit que
\(F_n(x) \xrightarrow[n\to \infty]{} F(x)\) pour tout \(x\) en lequel
\(F\) est continue.

Notons que puisque la fonction \(F\) est continue à droite et
croissante, l'ensemble des points où \(F\) est continue est l'ensemble
\(D = \{x : F (x-) = F (x)\}\), et son complémentaire est au plus
dénombrable. Ainsi, \(D\) est dense\footnote{Soient \(\Omega\) un espace
  topologique et \(A\) une partie de \(\Omega\). On dit que \(A\) est
  \emph{dense} dans \(\Omega\) si l'une des propriétés équivalentes est
  satisfaite~: tout ouvert non vide de \(\Omega\) contient des éléments
  de \(A\)\,; l'adhérence de \(A\) est égale à \(\Omega\)\,; tout point
  de \(\Omega\) est adhérent à \(A\)\,; le complémentaire de \(A\) est
  d'intérieur vide.} dans \(\mathbb{R}\).
\end{proposition}

\begin{preuve}
\begin{enumerate}
\def\labelenumi{\arabic{enumi}.}
\tightlist
\item
  Supposons d'abord que \(X_n \xrightarrow{\mathcal{L}} X\). Soit
  \(a \in \mathbb{R}\) tel que \(F(a-)=F(a)\). Pour tout
  \(p \in \mathbb{N}^\ast\) et tout \(b \in \mathbb{R}\), il existe une
  fonction \(f_{p,b}\) continue bornée sur \(\mathbb{R}\) telle que
  \[ 1_{]-\infty,b]} \leq f_{p,b} \leq 1_{]-\infty,b + \frac{1}{p}]}.\]
  Alors, par définition,
  \(\mathbb{E}(f_{p,b}(X_n)) \to \mathbb{E}(f_{p,b}(X))\) quand \(n\)
  tend vers l'infini. L'inégalité ci-dessus implique que
  \(F_n(a) = \mathbb{P}(X_n \leq a) \leq \mathbb{E}(f_{p,a}(X_n))\) et
  \(\mathbb{E}(f_{p,a}(X)) \leq F(a+1/p)\)\,; donc
  \(\lim\sup_n F_n(a) \leq F(a+1/p)\) pour tout \(p\) et donc on a aussi
  \(\lim\sup_n F_n(a) \leq F(a)\). On a également que
  \(F_n(a) \geq \mathbb{E}(f_{p,a-1/p}(X_n))\) et
  \(\mathbb{E}(f_{p,a-1/p}(X)) \geq F(a-1/p)\)\,; donc
  \(\lim\inf_n F_n(a) \geq F(a-1/p)\) pour tout \(p\) et donc aussi
  \(\lim\inf_n F_n(a) \geq F(a)\), puique \(F(a-)=F(a)\). Ces deux
  résultats impliquent que \(F_n(a) \xrightarrow[n\to \infty]{} F(a)\).
\item
  Inversement, supposons \(F_n(x) \xrightarrow[n\to \infty]{} F(x)\)
  pour tout \(x\in T\), où \(T\) est une partie dense de \(\mathbb{R}\).
  Soit \(f\) une fonction continue bornée sur \(\mathbb{R}\) et
  \(\varepsilon > 0\). Soient \(a,b \in T\) avec
  \(F(a) \leq \varepsilon\) et \(F(b) \geq 1-\varepsilon\). Il existe
  \(n_0\) tel que
  \[ n\geq n_0 \Rightarrow \mathbb{P}(X_n\notin ]a,b]) = 1-F_n(b)+F_n(a) \leq 3\varepsilon.\]
  La fonction \(f\) est uniformément continue sur \([a,b]\), donc il
  existe un nombre fini de points \(a_0 = a < a_1 < \ldots < a_k = b\)
  appartenant tous à \(T\), et tels que
  \(|f(x) - f(a_i)|\leq \varepsilon\) si \(a_{i-1} \leq x \leq a_i\).
  Donc \[ g(x) = \sum_{i=1}^k f(a_i)1_{]a_{i-1},a_i]}(x)\] vérifie
  \(|f-g| \leq \varepsilon\) sur \(]a,b]\). Si \(M = \sup_x |f(x)|\), il
  vient alors
  \[|\mathbb{E}(f(X_n))-\mathbb{E}(g(X_n))| \leq M \mathbb{P}(X_n \notin [a,b]) + \varepsilon,\]
  de même pour \(X\). Enfin,
  \(\mathbb{E}(g(X_n)) = \sum_{i=1}^k f(a_i)(F_n(a_i) - F_n(a_{i-1})\),
  et de même pour \(X\) par définition de \(g\). Comme
  \((F_n(a_i))_{n \in \mathbb{N}^\ast}\) converge vers \(F(a_i)\) pour
  tout \(i\), on en déduit l'existence de \(n_1\) tel que
  \[ n\geq n_1 \Rightarrow |\mathbb{E}(g(X_n)) - \mathbb{E}(g(X))|\leq \varepsilon.\]
  Finalement, on a
  \[ n \geq \sup(n_0,n_1) \Rightarrow |\mathbb{E}(f(X_n))-\mathbb{E}(f(X))|\leq 3\varepsilon + 5 M \varepsilon .\]
  Vu l'arbitraire sur \(\varepsilon\), on en déduit que
  \(\mathbb{E}(f(X_n))\) converge vers \(\mathbb{E}(f(X))\), d'où le
  résultat.
\end{enumerate}

\hfill$\blacksquare$
\end{preuve}

\addcontentsline{toc}{subsection}{Cas où \(X\) admet une densité}
\begin{corollaire}[{Cas où \(X\) admet une
densité}]\phantomsection\label{IV-cas-ouxf9-x-admet-une-densituxe9}

Si la suite \((X_n)_{n\in\mathbb{N}^\ast}\) de variables aléatoires
réelles converge en loi vers \(X\), et si la loi de \(X\) admet une
densité, alors pour tous \(a< b\),
\[ \mathbb{P}(X_n \in ]a,b]) \xrightarrow[n\to \infty]{} \mathbb{P}(X\in]a,b]).\]
\end{corollaire}

\begin{preuve}
La fonction de répartition de \(X\) est alors continue en tout point.
(Mais pas nécessairement celle des variables aléatoires
\(X_n\).)\hfill$\blacksquare$
\end{preuve}

\section{Théorème central
limite}\label{IV-thuxe9oruxe8me-central-limite}

Ce théorème est aussi connu sous le nom de théorème de la limite
centrale. Plus simplement, il apparaît souvent sous l'abréviation TCL.

On considère une suite de variables aléatoire
\((X_n)_{n \in \mathbb{N}^\ast}\) indépendantes, de même loi et de carré
intégrable. On note \(m\) et \(\sigma^2\) l'espérance et la variance
commune aux variables \(X_n\), et \[S_n = X_1 + \ldots + X_n\] ainsi
(\(S_n = n M_n\)). On a vu que la loi des grands nombres assure que
\(M_n\) converge vers \(m\) presque-sûrement et en moyenne. On va
s'intéresser à la vitesse à laquelle cette convergence a lieu.

Pour évaluer cette vitesse, c'est-à-dire trouver un équivalent de
\(\frac{S_n}{n} - m\), on est amené à étudier la limite éventuelle de la
suite \(n^\alpha (\frac{S_n}{n} - m)\) pour différentes valeurs de
\(\alpha\)~: si \(\alpha\) est ``petit'' cette suite va encore tendre
vers 0, et elle va ``exploser'' si \(\alpha\) est ``grand''. On peut
espérer que pour une (et alors nécessairement une seule) valeur de
\(\alpha\), cette suite converge vers une limite qui n'est ni infinie ni
nulle.

Il se trouve que la réponse à cette question a un aspect ``négatif''~:
la suite \(n^\alpha (\frac{S_n}{n} - m)\) ne converge au sens
presque-sûr, ou même en probabilité, pour aucune valeur de \(\alpha\).
Elle a aussi un aspect ``positif''~: cette suite converge, au sens de la
convergence en loi, pour la même valeur \(\alpha = 1/2\) quelle que soit
la loi des \(X_n\), et toujours vers une loi normale.

Ce résultat, qui peut sembler miraculeux, a été énoncé par Laplace
(1749-1827) et démontré beaucoup plus tard par Lyapounov (1901). Il
montre le caractère universel de la loi normale en probabilités (d'où
son nom).

\addcontentsline{toc}{subsection}{central limite (TCL)}
\begin{theoreme}[{central limite (TCL)}]\phantomsection\label{IV-TCL}

Si les \(X_n\) sont des variables aléatoires réelles, indépendantes et
de même loi, de carré intégrable, d'espérance \(m\) et de variance
\(\sigma^2 >0\), alors les variables
\[ \frac{S_n -nm}{\sigma \sqrt{n}}\] convergent en loi vers une variable
aléatoire de loi \(\mathcal{N}(0,1)\).

En d'autres termes, \(\sqrt{n}(M_n - m)\) converge vers une variable
normale de loi \(\mathcal{N}(0,\sigma^2)\).
\end{theoreme}

\begin{preuve}
Soit \(\phi\) la fonction caractéristique de \(X_n - m\), et
\(Y_n = \frac{S_n -nm}{\sigma \sqrt{n}}\). Comme les \(X_n\) sont
indépendantes, la \protect\hyperref[IV-fct_carac_sum]{proposition ---
somme (p.~\pageref*{IV-fct_carac_sum})} entraîne que la fonction
caractéristique de \(Y_n\) est \begin{align*}
    \phi_{Y_n}(u) & = \phi_{\frac{1}{\sigma\sqrt{n}}\sum_{j=1}^n(X_j-m)}(u) \\
                  & = \phi_{\sum_{j=1}^n(X_j-m)}\left(\frac{u}{\sigma\sqrt{n}}\right) \\
                  & = \prod_{j=1}^n \phi \left(\frac{u}{\sigma\sqrt{n}}\right)\\
                  & = \phi\left(\frac{u}{\sigma\sqrt{n}}\right)^n
\end{align*} Comme \(\mathbb{E}(X_n -m) = 0\) et
\(\mathbb{E}((X_n-m)^2) = \sigma^2\), la
\protect\hyperref[IV-fct_carac_deriv]{proposition --- fonction
caractéristique et moments (p.~\pageref*{IV-fct_carac_deriv})} entraîne
que \[\phi'(0) = 0 \text{     et } \phi''(0) = -\sigma^2\] Si on fait le
développement de Taylor à l'ordre 2 au voisinage de zéro, on obtient
\[\phi(u) = 1 - \frac{u^2\sigma^2}{2} + u^2 h(|u|),\] où \(h(u) \to 0\)
quand \(u \to 0\). On a ainsi \begin{align*}
    \phi_{Y_n}(u) & = \phi\left(\frac{u}{\sigma\sqrt{n}}\right)^n \\
                  & = e^{n\log \phi\left(\frac{u}{\sigma\sqrt{n}}\right)} \\
                  & = e^{n\log (1 - \frac{u^2}{2n} + \frac{u^2}{n\sigma^2} h(\frac{u}{\sigma\sqrt{n}}))}
\end{align*} où \(\log\) désigne la valeur principale du logarithme
complexe\footnote{voir par exemple
  https://fr.wikipedia.org/wiki/Logarithme\_complexe} (elle vaut 0 au
point 1 et est continue dans le cercle complexe de centre 1 et de rayon
1/2 et admet le même développement limité au voisinage de \(z=1\) que le
logarithme réel). Comme \(\phi(0) = 1\) et que \(\phi\) est continue en
0, on a que pour \(u\) fixé et \(n\) assez grand,
\[\left| \phi\left(\frac{u}{\sigma\sqrt{n}}\right)-1\right| \leq 1/2.\]
En faisant \(n \to \infty\), on obtient
\[\lim_{n \to \infty} \phi_n(u) = \exp -\frac{u^2}{2}\] Le
\protect\hyperref[IV-levytheorem]{théorème de
Lévy (p.~\pageref*{IV-levytheorem})} implique alors que \(Y_n\) converge
en loi vers \(Z\) de fonction caractéristique \(\phi_Z(u) = e^{-u^2/2}\)
où l'on reconnaît la fonction caractéristique de la loi
\(\mathcal{N}(0,1)\).\hfill$\blacksquare$
\end{preuve}

\addcontentsline{toc}{subsection}{Conséquence}
\begin{remarque}[{Conséquence}]\phantomsection\label{IV-consuxe9quence}

On peut déduire de ce résultat que \(n^\alpha (\frac{S_n}{n} - m)\)
converge vers 0 (resp. \(+\infty\)) en probabilité lorsque
\(\alpha < 1/2\) (resp. \(\alpha > 1/2\)).
\end{remarque}

\subsection{Exemple~: convergence des lois
binomiales}\label{IV-exemple-convergence-des-lois-binomiales}

Supposons que \(S_n\) suive une loi binomiale \(\mathcal{B}(p,n)\). Cela
revient à dire que \(S_n\) a la même loi qu'une somme \(X_1+\ldots+X_n\)
de \(n\) variables aléatoires \(X_i\) indépendantes de loi
\(\mathcal{B}(p,1)\), i.e.~\(\mathbb{P}(X_i= 1) = p\) et
\(\mathbb{P}(X_i= 0) = 1- p\). On a alors \(m=p\) et
\(\sigma^2= p(1-p)\).

On veut calculer \(\mathbb{P}(S_n \leq x)\) pour \(x\) fixé et \(n\)
grand.

Si \(p\) est petit de sorte que \(\theta = np\) ne soit pas trop grand
(en pratique, \(\theta \leq 5\) convient), on peut utiliser
l'approximation par une loi de Poisson, vue en CPGE. Si \(p\) est très
proche de 1, de sorte que \(\theta = n(1 - p)\) soit comme ci-dessus,
alors \(n - S_n\) suit à son tour une loi proche de la loi de Poisson de
paramètre \(\theta\).

Dans les autres cas, on utilise la loi des grands nombres et le théorème
central limite~: \begin{align*}
\frac{S_n}{n} & \xrightarrow{\mathbb{P}} p,\\
\frac{S_n - np}{\sqrt{n p(1-p)}} & \xrightarrow{\mathcal{L}} Y \sim \mathcal{N}(0,1)
\end{align*}

Si on désigne par \(\Phi\) la fonction de répartition de la loi
\(\mathcal{N}(0,1)\), il vient
\[\mathbb{P}(S_n \leq x) \approx \Phi \left(\frac{x - np}{\sqrt{n p(1-p)}}\right)\]

Imaginons que l'on lance 1000 fois une pièce (non truquée). On cherche
la probabilité d'obtenir plus de 545 fois le côté Face. Le calcul exact
utilisant les lois binomiales est extrêmement lourd. Le résultat
ci-dessus nous donne une très bonne approximation. On a
\[ \mathbb{P}(S_{1000} > 545) = \mathbb{P}\left(\frac{S_{1000} - 500}{\sqrt{250}} > \frac{45}{\sqrt{250}}\right) \approx \int_{\frac{45}{\sqrt{250}}}^{+\infty} \frac{1}{\sqrt{2\pi}} e^{-x^2/2} dx.\]

Cette dernière intégrale se calcule numériquement (on trouve encore des
abaques où les valeurs de \(\Phi\) sont tabulées) et on obtient
\[ \mathbb{P}(S_{1000} > 545) \approx 1 - \Phi(2,84) \approx 0,0023.\]

\runinhead{Exercice -- Surbooking (\(\mathord{\bullet}\))}\phantomsection\addcontentsline{toc}{subsubsection}{Surbooking}\label{IV-surbooking}
400 places dans un avion, 420 billets vendus, 8\% de taux de
non-présentation~: quelle probabilité de surbooking\,? On utilisera une
approximation judicieuse.\solref{IV-answer-surbooking}

Le \protect\hyperref[IV-TCL]{théorème (p.~\pageref*{IV-TCL})} admet une
version multidimensionnelle, de preuve similaire. On considère des
vecteurs aléatoires \(X_n\) à valeurs dans \(\mathbb{R}^d\),
indépendants et de même loi, dont les composantes sont de carré
intégrable. On a un vecteur espérance \(m = E(X_n)\), et une matrice de
covariance \(C = (c_{ij} )_{i,j=1,\ldots,d}\) avec
\(c_ij = \text{Cov}(X_i,X_j)\). On peut alors énoncer le TCL
multi-dimensionnel.

\addcontentsline{toc}{subsection}{central limite multi-dimensionnel}
\begin{theoreme}[{central limite
multi-dimensionnel}]\phantomsection\label{IV-central-limite-multi-dimensionnel}

Les vecteurs aléatoires \(\frac{S_n-nm}{\sqrt{n}}\) convergent en loi
vers un vecteur aléatoire gaussien centré (i.e.~d'espérance nulle), de
matrice \(C\).
\end{theoreme}

\addcontentsline{toc}{subsection}{Vitesse de convergence}
\begin{remarque}[{Vitesse de
convergence}]\phantomsection\label{IV-vitesse-de-convergence}

Il est important de noter ici que la vitesse de convergence ne dépend
pas de la dimension des vecteurs \(X_n\).
\end{remarque}

\chapter{Annexe}\label{IV-annexe}

\section{Fonctions
caractéristiques}\label{IV-fonctions-caractuxe9ristiques}

Dans ce paragraphe, nous introduisons un outil important en calcul des
probabilités~: il s'agit de ce que l'on appelle \emph{la fonction
caractéristique} d'une variable aléatoire, et qui dans d'autres branches
des mathématiques s'appelle aussi \emph{la transformée de Fourier}. Elle
nous sera notamment nécessaire, via le théorème de Lévy, pour démontrer
le théorème central limite. L'essentiel de cette section peut être
considéré hors programme, dans le sens où, bien que très utile en
pratique, sa connaissance ne sera pas évaluée à l'examen.

On notera \(< x, y >\) le produit scalaire de deux vecteurs de
\(\mathbb{R}^n\) . Si \(u \in \mathbb{R}^n\) , la fonction (complexe)
\(x \mapsto e^{i < u,x>}\) est continue, de module 1. Donc si \(X\) est
un vecteur aléatoire à valeurs dans \(\mathbb{R}^n\) , nous pouvons
considérer \(e^{i < u,X>}\) comme une variable aléatoire à valeurs
complexes. Ses parties réelle \(Y = \cos(< u, X>)\) et imaginaire
\(Z = \sin(< u, X>)\) sont des variables aléatoires réelles. Ces
variables aléatoires réelles sont de plus bornées par 1, donc elles
admettent une espérance. Il est alors naturel d'écrire que l'espérance
de \(e^{i < u,x>}\) est
\[\mathbb{E}(e^{i < u,X>}) = \mathbb{E}(Y) + i \mathbb{E}(Z) = \mathbb{E}(\cos< u, X>) + i\mathbb{E}(\sin< u, X>) \]

\addcontentsline{toc}{subsection}{Fonction caractéristique}
\begin{definition}[{Fonction
caractéristique}]\phantomsection\label{IV-fonction-caractuxe9ristique}

Si \(X\) est un vecteur aléatoire à valeurs dans \(\mathbb{R}^n\), sa
\emph{fonction caractéristique} est la fonction \(\phi_X\) de
\(\mathbb{R}^n\) dans \(\mathbb{C}\) définie par
\[ \phi_X(u) = \mathbb{E}(e^{i < u,X>}) = \int_{\mathbb{R}^n} e^{i < u,x>} \mathbb{P}_X(dx).\]
\end{definition}

\addcontentsline{toc}{subsection}{Interprétation}
\begin{remarque}[{Interprétation}]\phantomsection\label{IV-interpruxe9tation-1}

La fonction caractéristique ne dépend en fait \textbf{que de la loi
\(\mathbb{P}_X\) de \(X\)}~: c'est la ``transformée de Fourier'' de la
loi \(\mathbb{P}_X\).

Nous verrons que cette fonction porte bien son nom, au sens où elle
caractérise la loi \(\mathbb{P}_X\). C'est une notion qui, de ce point
de vue, généralise la fonction génératrice
\(G_X(s) = \mathbb{E}(s^X), \, s \in [0,1]\), vue en CPGE dans le cas
discret. Elle vérifie \[\phi_X(u) = G_X(e^{iu}) = \mathbb{E}(e^{iuX}) \]
pour une variable \(X\) à valeurs dans \(\mathbb{N}\).
\end{remarque}

\addcontentsline{toc}{subsection}{Propriétés de la fonction
caractéristique}
\begin{proposition}[{Propriétés de la fonction
caractéristique}]\phantomsection\label{IV-propriuxe9tuxe9s-de-la-fonction-caractuxe9ristique}

\(\phi_X\) est de module inférieur à 1, continue, avec
\[ \phi(0) = 1 ;\,\,\, \phi_X(-u) = \overline{\phi_X(u)}.\]
\end{proposition}

\begin{preuve}
\(|z|\) désigne le module d'un nombre complexe \(z\).

Comme \(\mathbb{E}(Y)^2 \leq \mathbb{E}(Y^2)\) pour toute variable
réelle \(Y\), on a~:
\[ |\phi_X(u)|^2 = \mathbb{E}(\cos< u, X>)^2 + \mathbb{E}(\sin< u, X>)^2 \leq \mathbb{E}(\cos^2< u, X> + \sin^2< u, X>) = 1.\]

Pour montrer la continuité, considérons une suite
\(u_p \xrightarrow[p \to \infty]{} u\). Il y a convergence simple de
\(e^{i < u_p,X>}\) vers \(e^{i < u,X>}\). Comme ces variables aléatoires
sont de module inférieur à 1, le théorème de convergence dominée assure
que \(\phi_X(u_p) \xrightarrow[p \to \infty]{} \phi_X(u)\). \(\phi_X\)
est donc continue.\hfill$\blacksquare$
\end{preuve}

\addcontentsline{toc}{subsection}{Transformation linéaire}
\begin{proposition}[{Transformation
linéaire}]\phantomsection\label{IV-fct_carac_vec}

Si \(X\) est un vecteur aléatoire à valeurs dans \(\mathbb{R}^n\), si
\(a \in \mathbb{R}^m\) et \(A\) est une matrice réelle de taille
\(m \times n\), alors
\[ \phi_{a+AX}(u) = e^{i < u,a>} \phi_X (A^t u), \forall u \in \mathbb{R}^m\]
\end{proposition}

\begin{preuve}
Nous avons \(e^{i< u, a + AX>} = e^{i < u,a>}e^{i <A^tu, X>}\). En
effet, \(< u, A X> = < A^t u, X>\). On prend ensuite les espérances pour
obtenir le résultat.\hfill$\blacksquare$
\end{preuve}

\begin{exemple}\phantomsection\label{IV-ex}
\begin{enumerate}
\def\labelenumi{\arabic{enumi}.}
\tightlist
\item
  \(X\) suit une loi binomiale \(\mathcal{B}(n,p)\)~:
  \(\phi_X(u) = (p e^{iu}+ 1- p)^n\).
\item
  \(X\) suit une loi de Poisson \(\mathcal{P}(\theta)\)~:
  \(\phi_X(u) = e^{\theta (e^{iu}-1)}\).
\item
  \(X\) suit une loi uniforme
  \(\mathcal{U}_{[a,b]}\)~:\(\phi_X(u) = \frac{e^{iua} - e^{iub}}{iu(b-a)}\).
\item
  \(X\) suit une loi exponentielle \(\mathcal{E}(\lambda)\),
  \(\lambda > 0\)~: \(\phi_X(u) = \frac{\lambda}{\lambda - iu}\).
\item
  \(X\) suit une loi normale \(\mathcal{N}(0,1)\)~:
  \(\phi_X(u) = e^{-u^2/2}\). (calcul ci-dessous)
\item
  \(X\) suit une loi normale \(\mathcal{N}(\mu,\sigma^2)\)~:
  \(\phi_X(u) = e^{iu\mu -u^2\sigma^2/2}\). (application directe de
  \protect\hyperref[IV-fct_carac_vec]{la proposition
  ci-dessus (p.~\pageref*{IV-fct_carac_vec})})
\end{enumerate}
\end{exemple}

\subsection{Fonction caractéristique d'une variable
gaussienne}\label{IV-fctcaracgauss}

Soit \(X \sim \mathcal{N}(0,1)\). On a \begin{align*}
\phi_X(u) = \mathbb{E}(e^{iuX}) &= \int_\mathbb{R}e^{iux}\frac{1}{\sqrt{2\pi}}e^{-x^2/2}dx \\
                          &= \int_\mathbb{R}\frac{\cos(ux)}{\sqrt{2\pi}}e^{-x^2/2}dx + \int_\mathbb{R}\frac{i\sin(ux)}{\sqrt{2\pi}}e^{-x^2/2}dx
\end{align*} Comme \(x \mapsto \frac{\sin(ux)}{\sqrt{2\pi}}e^{-x^2/2}\)
est impaire et intégrable, son intégrale est nulle, d'où
\[\phi_X(u) = \int_\mathbb{R}\frac{\cos(ux)}{\sqrt{2\pi}}e^{-x^2/2}dx.\]
D'après la \protect\hyperref[IV-fct_carac_deriv]{proposition ---
fonction caractéristique et moments (p.~\pageref*{IV-fct_carac_deriv})},
on peut dériver les deux membres par rapport à \(u\), et on obtient
\[\phi'_X(u) = \frac{1}{\sqrt{2\pi}}\int_\mathbb{R}-x\sin(ux)e^{-x^2/2}dx\]
puis par intégration par parties
\[\phi'_X(u) = -\frac{1}{\sqrt{2\pi}}\int_\mathbb{R}u\cos(ux)e^{-x^2/2}dx = -u\phi_X(u).\]
Ainsi, \(\phi_X\) satisfait à l'équation différentielle
\(\phi'_X(u) = -u\phi_X(u)\), dont la solution générale est
\[\phi_X(u) = C e^{-u^2/2}.\] Comme \(\phi_X(0) = 1,\) on en déduit
finalement que \[\phi_X(u) = e^{-u^2/2}.\]

L'intérêt majeur de la fonction caractéristique réside dans le fait
qu'elle caractérise la loi de la variable aléatoire (d'où son nom).

\addcontentsline{toc}{subsection}{Caractérisation}
\begin{theoreme}[{Caractérisation}]\phantomsection\label{IV-caracfc}

La fonction caractéristique \(\phi_X\) caractérise la loi du vecteur
aléatoire \(X\). Ainsi, si deux vecteurs aléatoires \(X\) et \(Y\) ont
même fonction caractéristique, ils ont même loi.
\end{theoreme}

\begin{preuve}
Soient les fonctions suivantes avec \(\sigma > 0\)~:
\[ f_\sigma (x) = \frac{1}{(2\pi \sigma^2)^{n/2}}\exp\left(-\frac{|x|^2}{2\sigma^2}\right) \text{ et } \widehat{f}_\sigma (u)= \exp\left(-\frac{|u|^2\sigma^2}{2}\right).\]
On note que \(f_\sigma\) est la densité d'un vecteur gaussien \(Z\) de
dimension \(n\), centré et de matrice de covariance \(\sigma^2 I_n\),
autrement dit dont les composantes sont indépendantes, centrées, de
variances \(\sigma^2\).

On a \begin{align*}
\mathbb{E}(e^{i < u,Z>}) = \int f_\sigma (x) e^{i < u,x>} dx &= \int_{\mathbb{R}^n} \prod_{j=1}^n \frac{1}{\sqrt{2\pi \sigma^2}}\exp\left(-\frac{x_j^2}{2\sigma^2}+iu_j x_j\right) dx_1\ldots dx_n \\
                                  &= \prod_{j=1}^n \int_{\mathbb{R}} \frac{1}{\sqrt{2\pi \sigma^2}}\exp\left(-\frac{t^2}{2\sigma^2}+iu_j t\right) dt \\
                                  &= \prod_{j=1}^n e^{-u_j^2\sigma^2/2}\\
                                  & = \widehat{f}_\sigma(u)
\end{align*} d'après l'exemple 5 \href{ex}{ci-dessus} et en utilisant le
théorème de Fubini. On remarque ainsi que
\[ f_\sigma (u-v) = \frac{1}{(2\pi \sigma^2)^{n/2}} \widehat{f}_{\sigma}\left(\frac{u-v}{\sigma^2}\right) = \frac{1}{(2\pi \sigma^2)^{n/2}} \int f_\sigma (x) e^{i < u -v ,x>/\sigma^2} dx.\]

Supposons que \(X\) et \(Y\) admettent la même fonction caractéristique
\(\phi_X = \phi_{Y}\). En utilisant le théorème de Fubini, on a
\begin{align*}
\mathbb{E}_X(f_\sigma (X-v)) &= \int f_\sigma (u-v) \mathbb{P}_X (du) \\
                       &= \int_{\mathbb{R}^n} \frac{1}{(2\pi \sigma^2)^{n/2}} \left( \int_{\mathbb{R}^n} f_\sigma (x) e^{i < u -v ,x>/\sigma^2} dx \right) \mathbb{P}_X (du)\\
                       &= \int_{\mathbb{R}^n} f_\sigma (x) \frac{1}{(2\pi \sigma^2)^{n/2}} \left( \int_{\mathbb{R}^n}  e^{i < u ,x/\sigma^2>} \mathbb{P}_X(du) \right) e^{-i < v ,x>/\sigma^2} dx\\
                       &= \int_{\mathbb{R}^n} f_\sigma (x) \frac{1}{(2\pi \sigma^2)^{n/2}} \phi_X\left(\frac{x}{\sigma^2}\right) e^{-i < v ,x>/\sigma^2} dx,
\end{align*}

et la même égalité reste vraie pour \(\mathbb{P}_{Y}\). On en déduit que
\[\mathbb{E}(g(X)) = \int g(u) \mathbb{P}_X (du) = \int g(u) \mathbb{P}_Y (du) = \mathbb{E}(g(X'))\]
pour toute fonction \(g\) de l'espace vectoriel de fonction engendrées
par \(u \mapsto f_\sigma (u-v)\).

D'après le théorème de Stone-Weiertrass\footnote{voir par exemple
  Simmons (\protect\hyperlink{refIV-Simmons}{2003}) p.160.}, cet espace
est dense dans l'ensemble \(C_0\) des fonctions continues sur
\(\mathbb{R}^n\) et ayant une limite nulle à l'infini, pour la topologie
de la convergence uniforme (la norme est le sup sur \(\mathbb{R}^n\)).

Par suite, \(\mathbb{E}(g(X)) = \mathbb{E}(g(X'))\) pour toute fonction
\(g \in C_0\). Comme l'indicatrice de tout ouvert est limite croissante
de fonctions de \(C_0\), on en déduit que
\(\mathbb{P}_X(A) = \mathbb{E}(1_A (X))\) est égal à
\(\mathbb{P}_{X'}(A) = \mathbb{E}(1_A (X'))\) pour tout ouvert \(A\), ce
qui implique \(P_X = P_{X'}\).\hfill$\blacksquare$
\end{preuve}

La fonction caractéristique offre également un moyen commode de
caractériser l'indépendance.

\begin{corollaire}\phantomsection\label{IV-corollaire}
Soit \(X = (X_1,\ldots,X_n)\) un vecteur aléatoire à valeurs dans
\(\mathbb{R}^n\). Ses composantes \(X_i\) sont indépendantes si et
seulement si pour tous \(u_1, \ldots, u_n \in \mathbb{R}\),
\[ \phi_X(u_1,\ldots,u_n) = \prod_{j=1}^n \phi_{X_j}(u_j) \] où
\(\phi_X\) désigne la fonction caractéristique du vecteur aléatoire
\(X\), et \(\phi_{X_j}\) celle de la composante \(X_j\) pour chaque
\(j\).
\end{corollaire}

\begin{preuve}
On a \(< u,X> = \sum_{j=1}^n u_j X_j\). Si les \(X_i\) sont
indépendantes, et comme \(e^{i < u,x>} = \prod_j e^{i u_j x_j}\), nous
obtenons directement le résultat en utilisant la
\href{Probabilité\%20II.pdf\%20\#indep_fct}{proposition} (valable dans
le cas général) qui montre que l'espérance du produit de fonctions de
variables aléatoires indépendantes est égal au produit des espérances.

Supposons inversement qu'on ait
\(\phi_X (u_1,\ldots,u_n) = \prod_{j=1}^n \phi_{X_j}(u_j)\). On peut
alors construire des variables aléatoires \(X'_j\) indépendantes, telles
que \(X'_j\) et \(X_j\) aient même loi pour tout \(j\) et donc telles
que \(\phi_{X'_j} = \phi_{X_j}\). Si \(X' = (X'_1,\ldots,X'_n)\), on a
donc \(\phi_{X'} = \phi_X\). Donc \(X\) et \(X'\) ont même loi, ce qui
entraîne que pour tous boréliens \(A_j\), on ait
\[ \mathbb{P}(\bigcap_j \{X_j \in A_j\}) = \mathbb{P}(\bigcap_j \{X'_j \in A_j\}) = \prod_j \mathbb{P}(\{X'_j \in A_j\}) = \prod_j \mathbb{P}(\{X_j \in A_j\})\]
d'où l'indépendance.\hfill$\blacksquare$
\end{preuve}

\addcontentsline{toc}{subsection}{Somme}
\begin{proposition}[{Somme}]\phantomsection\label{IV-fct_carac_sum}

Si \(X\) et \(Y\) sont deux vecteurs aléatoires indépendants à valeurs
dans \(\mathbb{R}^n\), la fonction caractéristique de la somme \(X+Y\)
est donnée par \[ \phi_{X+Y} = \phi_X\phi_Y\]
\end{proposition}

\begin{preuve}
Comme \(e^{i< u, X+Y>}=e^{i< u, X>}e^{i< u, Y>}\), il suffit d'appliquer
la \href{Probabilité\%20II.pdf\%20\#indep_fct}{proposition} déjà
utilisée dans la preuve du corollaire ci-dessus.\hfill$\blacksquare$
\end{preuve}

\subsection{Exemples~:}\label{IV-exemples}

Soient \(X\) et \(Y\) deux variables aléatoires réelles indépendantes et
\(Z = X+Y\)~:

\begin{enumerate}
\def\labelenumi{\arabic{enumi}.}
\tightlist
\item
  Si \(X\) suit la loi normale \(\mathcal{N}(m,\sigma^2)\) et \(Y\) suit
  \(\mathcal{N}(m',\sigma'^2)\), alors \(Z\) suit une loi
  \(\mathcal{N}(m+m',\sigma^2+\sigma'^2)\), d'après
  l'\protect\hyperref[IV-ex]{exemple ci-dessus (p.~\pageref*{IV-ex})}
  point 6. et la \protect\hyperref[IV-fct_carac_sum]{proposition
  ci-dessus (p.~\pageref*{IV-fct_carac_sum})}.
\item
  Si \(X\) et \(Y\) suivent des lois de Poisson de paramètres \(\theta\)
  et \(\theta'\), alors \(Z\) suit une loi de Poisson de paramètre
  \(\theta + \theta'\), d'après l'\protect\hyperref[IV-ex]{exemple
  ci-dessus (p.~\pageref*{IV-ex})} point 2. et la
  \protect\hyperref[IV-fct_carac_sum]{proposition
  ci-dessus (p.~\pageref*{IV-fct_carac_sum})}.
\item
  Si \(X\) suit une loi binomiale \(\mathcal{B}(n,p)\) et \(Y\) la loi
  biomiale \(\mathcal{B}(m,p)\), alors \(Z\) suit une loi binomiale
  \(\mathcal{B}(n+m,p)\), d'après l'\protect\hyperref[IV-ex]{exemple
  ci-dessus (p.~\pageref*{IV-ex})} point 1. et la
  \protect\hyperref[IV-fct_carac_sum]{proposition
  ci-dessus (p.~\pageref*{IV-fct_carac_sum})}.
\end{enumerate}

\addcontentsline{toc}{subsection}{Fonction caractéristique et moments}
\begin{proposition}[{Fonction caractéristique et
moments}]\phantomsection\label{IV-fct_carac_deriv}

Soit \(X\) un vecteur aléatoire de \(\mathbb{R}^n\). Si la variable
\(|X|^m\) (où \(|\cdot|\) désigne la norme euclidienne) est intégrable
pour un entier \(m\), la fonction \(\phi_X\) est \(m\) fois continûment
différentiable sur \(\mathbb{R}^n\) et pour tout choix des indices
\(i_1,\ldots, i_m\),
\[ \frac{\partial^m}{\partial u_{i_1}\ldots \partial u_{i_m}} \phi_X(u) = i^m \mathbb{E}(e^{i < u,X >}X_{i_1}\ldots X_{i_m}),\]
où les \(X_j\) sont les composantes de \(X\).
\end{proposition}

\begin{preuve}
Le résultat se démontre par application itérée du théorème de dérivation
sous le signe somme.\hfill$\blacksquare$
\end{preuve}

\begin{remarque}\phantomsection\label{IV-remarque}
En prenant \(u=0\) dans la
\protect\hyperref[IV-fct_carac_sum]{proposition ---
somme (p.~\pageref*{IV-fct_carac_sum})}, la formule permet de calculer
\(\mathbb{E}(X_{i_1}\ldots X_{i_m})\) en fonction des dérivées à
l'origine de \(\phi_X\), autrement dit de calculer tous les moments du
vecteur \(X\), s'ils existent. Par exemple, si \(X\) est à valeurs
réelles et est intégrable (respectivement de carré intégrable), on a
\[\mathbb{E}(X) = i \phi'_X(0), \,\,\,(\text{resp. } \mathbb{E}(X^2) = \phi"_X(0))\]

On donne ici la définition générale d'un vecteur gaussien, qui nous sera
utile pour la démonstration du théorème central limite.
\end{remarque}

\addcontentsline{toc}{subsection}{Vecteur Gaussien (cas général)}
\begin{definition}[{Vecteur Gaussien (cas
général)}]\phantomsection\label{IV-gauss_vec}

Un vecteur aléatoire \(X=(X_1,\ldots,X_n)\) à valeurs dans
\(\mathbb{R}^n\) est dit gaussien si toute combinaison linéaire de ses
composantes, soit \(\sum_{j=1}^n a_j X_j\) où \(a_j \in \mathbb{R}\),
suit une loi normale uni-dimensionnelle (éventuellement dégénérée, par
exemple si on prend \(a_j = 0\) pour tout j).
\end{definition}

\addcontentsline{toc}{subsection}{Cas gaussien}
\begin{theoreme}[{Cas
gaussien}]\phantomsection\label{IV-fct_carac_gauss}

\(X\) est un vecteur gaussien si et seulement si sa fonction
caractéristique s'écrit \[\phi_X(u) = e^{i< u,m> - \frac{1}{2}< u,Cu>}\]
où \(m = \mathbb{E}(X) \in \mathbb{R}^n\) et \(C\) est la matrice de
covariance de \(X\) qui est donc semi-définie positive.
\end{theoreme}

\begin{preuve}
\begin{enumerate}
\def\labelenumi{\arabic{enumi}.}
\tightlist
\item
  Supposons \(\phi_X(u) = e^{i< u,m> - \frac{1}{2}< u,Cu>}\). Pour toute
  combinaison linéaire \(Y = \sum_j a_j X_j = < a, X >\), et pour tout
  \(v \in \mathbb{R}\), on a
  \[ \phi_Y(v) = \phi_X(va) = e^{iv< a ,m> - \frac{v^2}{2}< a,Ca>}\]
  donc \(Y\) suit la loi \(\mathcal{N}(< a, m>, < a, Ca>)\).
\item
  Inversement, soit \(C\) la matrice de covariance de \(X\) et \(m\) son
  vecteur moyenne. Si \(Y = < a, X > = \sum_{j=1}^n a_j X_j\) avec
  \(a\in\mathbb{R}^n\), un calcul simple montre
  \[ \mathbb{E}(Y) = < a, m >, \mathbb{V}(Y) = < a, Ca> \] Par
  hypothèse, \(Y\) suit une loi normale donc vu le point 6. de
  l'\protect\hyperref[IV-ex]{exemple plus haut (p.~\pageref*{IV-ex})},
  sa fonction caractéristique est
  \[ \phi_Y(v) = e^{iv< a ,m> - \frac{v^2}{2}< a,Ca>} \] Mais
  \(\phi_Y(1) = \phi_{< a,X>} (1) = \mathbb{E}(e^{i < a, X>}) = \phi_X(a)\),
  d'où le résultat.
\end{enumerate}

\hfill$\blacksquare$
\end{preuve}

Le théorème suivant caractérise la convergence en loi à l'aide des
fonctions caractéristiques. C'est un critère extrêmement utile dans la
pratique.

\addcontentsline{toc}{subsection}{de Lévy}
\begin{theoreme}[{de Lévy}]\phantomsection\label{IV-levytheorem}

Soit \((X_n)_{n \in \mathbb{N}^\ast}\) une suite de vecteurs aléatoires
à valeurs dans \(\mathbb{R}^d\).

\begin{enumerate}
\def\labelenumi{\arabic{enumi}.}
\item
  Si la suite \((X_n)_{n \in \mathbb{N}^\ast}\) converge en loi vers
  \(X\), alors \(\phi_{X_n}\) converge simplement vers \(\phi_X\).
\item
  Si les \(\phi_{X_n}\) convergent simplement vers une fonction
  (complexe) \(\phi\) sur \(\mathbb{R}^d\) et si cette fonction est
  \textbf{continue} en 0, alors c'est la fonction caractéristique d'une
  variable aléatoire \(X\) et \(X_n \xrightarrow{\mathcal{L}} X\).
\end{enumerate}
\end{theoreme}

\begin{preuve}
\begin{enumerate}
\def\labelenumi{\arabic{enumi}.}
\tightlist
\item
  On remarque que \(\phi_{X_n}(u) = \mathbb{E}(g_u(X_n))\) et
  \(\phi_X(u) = \mathbb{E}(g_u(X))\) où \(g_u\) est la fonction continue
  bornée \(g_u(x) = e^{i < u,x>}\). On applique alors la
  \protect\hyperref[IV-defconvloi]{définition (p.~\pageref*{IV-defconvloi})}.
\item
  Se reporter à Jacod and Protter (\protect\hyperlink{refIV-Jacod}{2003}).
\end{enumerate}

\hfill$\blacksquare$
\end{preuve}

\chapter{Exercices}\label{IV-exercices}

\section{Inégalités de
concentration}\label{IV-inuxe9galituxe9s-de-concentration}

\runinhead{Question
1}\phantomsection\addcontentsline{toc}{subsubsection}{Question
1}\label{IV-idc1}
Soit \((X_n)_n\) une suite de variables aléatoires indépendantes et
identiquement distribuées d'espérance \(\mathbb{E}(X_i) = m\) et de
variance \(\mathbb{V}(X_i) = \sigma^2 \leq 1\). Montrer que pour tout
\(\delta \in \left]0,1\right[\), avec probabilité au moins \(1-\delta\)
on a
\[\left|\frac{1}{n}\sum_{i=1}^n X_i - m \right| \leq \sqrt{\frac{1}{\delta n}}\]

On peut trouver des bornes à décroissance beaucoup plus rapide dans des
cas particuliers.\solref{IV-answer-idc1}

\runinhead{Question
2}\phantomsection\addcontentsline{toc}{subsubsection}{Question
2}\label{IV-idc2}
On suppose désormais que les \(X_i\) sont de loi
\(\mathcal{N}(m, \sigma^2 )\).

\begin{enumerate}
\def\labelenumi{\roman{enumi})}
\item
  On pose pour \(u \in \mathbb{R}: M(u) = \mathbb{E}(e^{u(X_1 -m)})\).
  Calculer \(M (u)\).
\item
  On pose \(S_n = X_1 + \ldots + X_n\) . Montrer que
  \(\forall a \in \mathbb{R}, \forall u \in \mathbb{R}_+\),
  \[\mathbb{P}(S_n - nm > a) \leq e^{-ua} (M(u))^n,\] et que
  \(\forall a \in \mathbb{R}, \forall u \in \mathbb{R}_-\),
  \[\mathbb{P}(S_n - nm < -a) \leq e^{ua} (M(u))^n.\]
\item
  Soit \(Y_n =\frac{S_n}{n}\). Montrer que \(\forall \varepsilon >0\),
  \[\mathbb{P}(|Y_n - m| > \varepsilon) \leq 2 \exp\left(\frac{-n\varepsilon^2}{2\sigma^2}\right).\]
  Cette inégalité est appelée inégalité de Chernov.
\end{enumerate}

On peut dériver ce type d'inégalités pour différentes lois de
probabilité. On voit qu'ici la concentration auprès de la moyenne se
fait à vitesse exponentielle. On a le même type de résultats pour la loi
de Bernoulli par exemple ce qui est très utile en apprentissage
statistique dans les problèmes de classification.\solref{IV-answer-idc2}

\runinhead{Question
3}\phantomsection\addcontentsline{toc}{subsubsection}{Question
3}\label{IV-idc3}

\begin{enumerate}
\def\labelenumi{\roman{enumi})}
\setcounter{enumi}{3}
\item
  On suppose que \(m = 1\) et que \(\sigma^2 = 10\). Quelle taille
  d'échantillon doit-on choisir pour obtenir
  \[\mathbb{P}(|Y_n - m| < \varepsilon) \geq \alpha,\] avec
  \(\alpha = 0,95\) et \(\varepsilon = 0,05\),

  \begin{itemize}
  \tightlist
  \item
    en utilisant l'inégalité de Bienaymé-Chebyshev\,?
  \item
    en utilisant l'inégalité de Chernov\,?
  \end{itemize}
\end{enumerate}

\solref{IV-answer-idc3}

\section{Lemme de Borel-Cantelli}\label{IV-lemme-de-borel-cantelli}

Soit \(A_n\) une suite d'événements sur l'espace probabilisé
\((\Omega,\mathcal{A}, \mathbb{P})\).

\runinhead{Question
1}\phantomsection\addcontentsline{toc}{subsubsection}{Question
1}\label{IV-bc1}
On suppose que \(\sum_{i=1}^n \mathbb{P}(A_n) < \infty\). Montrer que
\(\mathbb{P}(\lim \sup_{n \to \infty} A_n) = \mathbb{P}\left(\bigcap_{n\geq 1} \bigcup_{k \geq n} A_k \right) =0\).\solref{IV-answer-bc1}

\runinhead{Question
2}\phantomsection\addcontentsline{toc}{subsubsection}{Question
2}\label{IV-bc2}
On suppose maintenant que les événements \(A_n\) sont mutuellement
indépendants. Montrer que si \(\sum_{i=1}^n \mathbb{P}(A_n) = \infty\),
alors on a
\(\mathbb{P}(\lim \sup_{n \to \infty} A_n) = 1\).\solref{IV-answer-bc2}

\runinhead{Question
3}\phantomsection\addcontentsline{toc}{subsubsection}{Question
3}\label{IV-bc3}
Donner un exemple où \(\sum_{i=1}^n \mathbb{P}(A_n) = \infty\) et
\(\mathbb{P}(\lim \sup_{n \to \infty} A_n) < 1\) quand les \(A_n\) ne
sont pas indépendants.\solref{IV-answer-bc3}

\section{Convergence vers une
constante}\label{IV-convergence-vers-une-constante}

Soient \(\left(X_n\right)_{n\in\mathbb{N}^\ast}\) une suite de variables
aléatoires réelles, \(X\) une autre variable aléatoire réelle et
\(a \in \mathbb{R}\).

\runinhead{Convergence
\(\mathcal{L}^2\)}\phantomsection\addcontentsline{toc}{subsubsection}{Convergence
\(\mathcal{L}^2\)}\label{IV-cvtocst-l2}
On suppose que \(X\) et chaque \(X_n\), \(n\in\mathbb{N}^\ast\), sont de
carré intégrable.

\begin{enumerate}
\def\labelenumi{\arabic{enumi}.}
\tightlist
\item
  Montrer l'implication
\end{enumerate}

\[X_n \xrightarrow[n\to+\infty]{\mathcal{L}^2} X \Rightarrow \left|\begin{array}{ll} \mathbb{E}\left(X_n\right) \xrightarrow[n\to+\infty]{} \mathbb{E}(X),\\ \mathbb{V}\left(X_n\right) \xrightarrow[n\to+\infty]{} \mathbb{V}(X). \end{array}\right.\]

\begin{enumerate}
\def\labelenumi{\arabic{enumi}.}
\setcounter{enumi}{1}
\tightlist
\item
  Vérifier l'équivalence
  \[X_n \xrightarrow[n\to+\infty]{\mathcal{L}^2} a \Leftrightarrow \left|\begin{array}{ll} \mathbb{E}\left(X_n\right) \xrightarrow[n\to+\infty]{} a,\\ \mathbb{V}\left(X_n\right) \xrightarrow[n\to+\infty]{} 0. \end{array}\right.\]
\end{enumerate}

\solref{IV-answer-cvtocst-l2}

\runinhead{Convergences en loi et
probabilité}\phantomsection\addcontentsline{toc}{subsubsection}{Convergences
en loi et probabilité}\label{IV-cvtocst-loiprob}
Montrer que si \(X_n\) converge en loi vers \(a\) quand \(n\to+\infty\),
alors elle converge aussi en probabilité vers
\(a\).\solref{IV-answer-cvtocst-loiprob}

\section{Fonction de répartition
empirique}\label{IV-fonction-de-ruxe9partition-empirique}

Soient \(X\) une variable aléatoire réelle de fonction de répartition
\(F\) et \(\left(X_n\right)_{n\in\mathbb{N}^\ast}\) une suite de
variables aléatoires indépendantes, de même loi que \(X\). Pour tout
\(n\in\mathbb{N}^\ast\) on définit la fonction de répartition empirique
comme suit~:
\[F_n : x \in \mathbb{R}\mapsto \dfrac{1}{n}\,\sum_{i = 1}^n 1_{]-\infty,x]}(X_i).\]

Son nom est issu de la modélisation probabiliste en statistique. Lorsque
l'on souhaite étudier une variable physique dont on n'est pas à même de
prédire parfaitement les valeurs (e.g.~la hauteur d'eau de la Seine à la
station d'Austerlitz), le statisticien va la considérer comme une
variable aléatoire \(X\). Pour en retrouver les propriétés, il va
observer un certain nombre \(n\in\mathbb{N}^\ast\) de réalisations de
cette variable (en réalisant des mesures sur le terrain), qu'il va à
nouveau considérer comme des variables aléatoires \(X_1,\dots,X_n\) de
même loi que \(X\). Lorsque cette hypothèse est raisonnable, il va
supposer que ces dernières sont même indépendantes. Nous allons voir
dans cet exercice que pour tout \(x\in\mathbb{R}\), \(F_n(x)\), calculée
à partir des données de terrain (d'où le nom attribué à \(F_n\)),
approche bien la quantité théorique \(F(x) = \mathbb{P}(X\leq x)\).

On considère \(x\in\mathbb{R}\) fixé.

\runinhead{Préliminaire}\phantomsection\addcontentsline{toc}{subsubsection}{Préliminaire}\label{IV-fdremp-loi}
Quelle est la loi de la variable aléatoire \(1_{]-\infty,x]}(X)\)\,?
Expliciter son espérance et sa variance.\solref{IV-answer-fdremp-loi}

\runinhead{Biais}\phantomsection\addcontentsline{toc}{subsubsection}{Biais}\label{IV-fdremp-biais}
Soit \(n\in\mathbb{N}^\ast\). Calculer l'espérance de \(F_n(x)\). On dit
que \(F_n(x)\) est un \emph{estimateur sans biais} de
\(F(x)\).\solref{IV-answer-fdremp-biais}

\runinhead{Erreur quadratique
moyenne}\phantomsection\addcontentsline{toc}{subsubsection}{Erreur
quadratique moyenne}\label{IV-fdremp-mse}
Montrer que \(F_n(x) \overset{\mathcal{L}^2}{\longrightarrow} F(x)\)
quand \(n\to+\infty\).\solref{IV-answer-fdremp-mse}

\runinhead{Consistance}\phantomsection\addcontentsline{toc}{subsubsection}{Consistance}\label{IV-fdremp-ps}
Montrer que \(F_n(x) \to F(x)\) p.s. quand \(n\to+\infty\). On dit que
\(F_n(x)\) est un \emph{estimateur fortement consistant} de
\(F(x)\).\solref{IV-answer-fdremp-ps}

\section{Théorème de Weierstrass sur
\([0,1]\)}\label{IV-thuxe9oruxe8me-de-weierstrass-sur-01}

Le théorème de Weierstrass est un résultat classique d'analyse dont
l'énoncé est le suivant~: soient \(I\) un segment de \(\mathbb{R}\) et
\(f : I \to \mathbb{R}\) une fonction continue. Alors il existe une
suite \((P_n)_{n\in\mathbb{N}^\ast}\) de polynômes
\(\mathbb{R}\to\mathbb{R}\) convergeant uniformément vers \(f\) sur
\(I\).

Dans cet exercice, nous allons voir une démonstration constructive de ce
théorème dans le cas particulier où \(I = [0,1]\).

\runinhead{Preuve du
théorème.}\phantomsection\addcontentsline{toc}{subsubsection}{Preuve du
théorème.}\label{IV-weier-thm}
Soit \(\left(X_n\right)_{n\in\mathbb{N}^\ast}\) une suite de variable
aléatoires indépendantes, suivant toutes la même loi de Bernoulli de
paramètre \(x \in ]0,1[\). Pour tout \(n\in\mathbb{N}^\ast\) on pose
\(M_n := \frac{1}{n}\,\sum_{i = 1}^n X_i\). On considère une fonction
\(f : [0,1] \to \mathbb{R}\) continue.

\begin{enumerate}
\def\labelenumi{\arabic{enumi}.}
\item
  Soit \(n\in\mathbb{N}^\ast\). Vérifier que
  \(\mathbb{E}\left( f(M_n) \right)\) peut s'écrire sous forme de
  polynôme.
\item
  Montrer que
  \(\mathbb{E}\left( f(M_n) \right) \xrightarrow[n\to+\infty]{} f(x).\)
\item
  En déduire qu'il existe une suite de polynômes
  \([0,1] \to \mathbb{R}\) qui convergent simplement vers \(f\).
\item
  Montrer que la suite de polynômes exhibée à la question précédente
  converge même uniformément vers \(f\).
\end{enumerate}

\solref{IV-answer-weier-thm}

\section{Théorème de Slutsky}\label{IV-thuxe9oruxe8me-de-slutsky}

\runinhead{Question
1}\phantomsection\addcontentsline{toc}{subsubsection}{Question
1}\label{IV-slut1}
Soient \((X_n)_n\) et \(X\) des variables aléatoires. Montrer que
\(X_n \xrightarrow{\mathcal{L}} X\) si et seulement si
\(\mathbb{E}(f(X_n)) \xrightarrow[n\to\infty]{} \mathbb{E}(f(X))\) pour
toute fonction \(f\) lipschitzienne bornée.\solref{IV-answer-slut1}

\runinhead{Question
2}\phantomsection\addcontentsline{toc}{subsubsection}{Question
2}\label{IV-slut2}
Soient \((X_n)_n\) et \((Y_n)_n\) deux suites de variables aléatoires
définies sur le même espace de probabilité. On suppose que \((X_n)_n\)
converge en loi vers \(X\) et que \((X_n - Y_n)_n\) converge vers 0 en
probabilité. Montrer que \((Y_n)_n\) converge en loi vers
\(X\).\solref{IV-answer-slut2}

\runinhead{Question
3}\phantomsection\addcontentsline{toc}{subsubsection}{Question
3}\label{IV-slut3}
Montrer ce résultat à l'aide du
\protect\hyperref[IV-levytheorem]{théorème de
Lévy (p.~\pageref*{IV-levytheorem})}.

\begin{center}\rule{0.5\linewidth}{0.5pt}\end{center}

\solref{IV-answer-slut3}


\chapter{Références}\label{IV-ruxe9fuxe9rences}

\hypertarget{refs}{}
\begin{cslreferences}
\leavevmode\hypertarget{refIV-Jacod}{}%
Jacod, J., and P. Protter. 2003. \emph{L'essentiel En Théorie Des
Probabilités}. Cassini.
\url{https://hal.archives-ouvertes.fr/hal-00104956}.

\leavevmode\hypertarget{refIV-Simmons}{}%
Simmons, George F. 2003. \emph{Introduction to Topology and Modern
Analysis}. Reprint edition. Malabar, Fla: Krieger Publishing Company.
\end{cslreferences}


\part{Simulation}
\chapter{Objectifs d'apprentissage}\label{V-objectifs-dapprentissage}

Cette section s'efforce d'expliciter et de hiérarchiser les acquis
d'apprentissages associés au chapitre. Ces objectifs sont organisés en
paliers~:

($\mathord{\boldsymbol{\circ}}$) Prérequis ($\mathord{\bullet}$)
Fondamental ($\mathord{\bullet}\mathord{\bullet}$) Standard
($\mathord{\bullet}\mathord{\bullet}\mathord{\bullet}$) Avancé
($\mathord{\bullet}\mathord{\bullet}\mathord{\bullet}\mathord{\bullet}$)
Expert

Sauf mention particulière, la connaissance des démonstrations du
document n'est pas exigible\footnote{l'étude des démonstrations du cours
  peut toutefois contribuer à votre apprentissage, au même titre que la
  résolution d'exercices.}

\runinhead{Intégrale de Monte
Carlo}\phantomsection\label{V-intuxe9grale-de-monte-carlo}

\begin{itemize}
\tightlist
\item
  $\mathord{\bullet}$ connaître le principe de l'intégration par la
  méthode Monte Carlo
\item
  $\mathord{\bullet}$ savoir que cette approche se justifie par la loi
  des grands nombres
\item
  $\mathord{\bullet}$ savoir que l'approximation fournie par le TCL
  offre un contrôle de l'erreur
\end{itemize}

\runinhead{Génération de nombres
pseudo-aléatoires}\phantomsection\label{V-guxe9nuxe9ration-de-nombres-pseudo-aluxe9atoires}

\begin{itemize}
\tightlist
\item
  $\mathord{\bullet}$ connaître le principe de la génération de nombres
  pseudo-aléatoires par la méthode des congruences
\end{itemize}

\runinhead{Méthodes de simulation de
v.a.}\phantomsection\label{V-muxe9thodes-de-simulation-de-v.a.}

\begin{itemize}
\tightlist
\item
  $\mathord{\bullet}$ connaître et savoir implémenter en python la
  méthode d'inversion
\item
  $\mathord{\bullet}$ connaître et savoir implémenter en python la
  méthode de rejet
\item
  $\mathord{\bullet}$ connaître la méthode de Box-Muller
\item
  $\mathord{\bullet}$ connaître et savoir implémenter la simulation de
  vecteurs gaussien par la méthode de Cholesky
\end{itemize}

\runinhead{Echantillonnage
d'importance}\phantomsection\label{V-echantillonnage-dimportance}

\begin{itemize}
\tightlist
\item
  $\mathord{\bullet}$ connaître la définition de la méthode
  d'échantillonnage d'importance
\item
  $\mathord{\bullet}$ savoir qu'un bon choix de densité instrumentale
  permet de minimiser la variance d'estimation
\end{itemize}

\chapter{Intégrale de
Monte-Carlo}\label{V-intuxe9grale-de-monte-carlo-1}

Les méthodes de Monte-Carlo ont été développées initialement par des
physiciens dans les années 1950 (notamment par les travaux de
Metropolis, Ulam, Hastings, Rosenbluth) pour calculer des intégrales
(déterministes) à partir de méthodes probabilistes numériquement assez
économiques. Le nom a été donné en référence au célèbre casino du fait
du caractère aléatoire de ces méthodes.

Les méthodes de simulation sont basées sur la production de nombres
aléatoires, distribués selon une certaine loi de probabilité. Dans de
nombreuses applications, pour une certaine fonction \(h\), on souhaite
calculer, pour une variable aléatoire \(X\) de loi \(\mathbb{P}_X\)
\[\mathcal{I}=\mathbb{E}\left(h(X)\right)=\int h(x)\, \mathbb{P}_X(dx).\]

En général, même si on sait évaluer \(h\) en tout point, on ne peut pas
calculer formellement l'intégrale \(\mathcal{I}\). Le calcul d'intégrale
par la méthode Monte-Carlo consiste dans sa version la plus simple à
générer un \emph{échantillon}
\((X_1,\ldots,X_n) \sim_{i.i.d.}\mathbb{P}_X\), et à estimer
\(\mathcal{I}\) par la moyenne empirique
\[M_n(h)=\frac{1}{n}\sum_{i=1}^{n}h(X_i),\] où i.i.d signifie
indépendant et identiquement distribué. En effet, d'après la
\href{Probabilité\%20III.pdf\%20\#lfgn}{loi forte des grands nombres},
si \(h(x)\) est \(\mathbb{P}_X\)-intégrable, on a l'assurance que quand
\(n \to +\infty\),
\[M_n(h) \rightarrow \int h(x)\mathbb{P}_X(dx) \text{ p.s.}\] Si de plus
\(h(X)^2\) est intégrable, la vitesse de convergence de \(M_n(h)\) peut
être évaluée, puisque la variance
\[\mathbb{V}(M_n(h)) = \frac{1}{n} \int \left(h(x)-\mathcal{I}\right)^2 \mathbb{P}_X(dx)\]
peut également être estimée à partir de l'échantillon
\((X_1,\ldots,X_n)\) par la quantité
\[\sigma_n^2=\frac{1}{n}\sum_{i=1}^{n}\left( h(X_i)-M_n(h)\right)^2.\]
Le \href{Probabilité\%20III.pdf\%20\#TCL}{théorème central limite}
assure alors que pour \(n\) grand,
\[\frac{M_n(h)-\mathcal{I}}{\sigma_n}\] suit approximativement une loi
\(\mathcal{N}(0,1)\)\footnote{ce résultat sera démontré dans le cours de
  science des données au second semestre.}. Cette propriété conduit à la
construction de tests de convergence et de bornes de confiance
asymptotiques pour \(M_n(h)\). Par exemple, on aura
\[\mathbb{P}\left(\mathcal{I} \in \left[M_n(h) - 1.96 \sigma_n,M_n(h) + 1.96 \sigma_n\right]\right) \approx 0.95,\]
où \(1.96 \approx \Phi^{-1}(0,975)\) avec \(\Phi\) la fonction de
répartition de la loi normale centrée réduite.

En outre, cette propriété indique que la vitesse de convergence de
\(M_n(h)\) est de l'ordre de \(\sqrt{n}\), et ce indépendamment de la
dimension du problème. Cela explique la supériorité de cette méthode par
rapport aux méthodes d'intégration numérique déterministes dont les
vitesses de convergence décroissent rapidement (exponentiellement) avec
la dimension du problème.

\runinhead{Exercice -- Calcul numérique d'une
probabilité (\(\mathord{\bullet}\mathord{\bullet}\))}\phantomsection\addcontentsline{toc}{subsubsection}{Calcul
numérique d'une probabilité}\label{V-mc1}
Calculer une estimation de la probabilité de l'événement
\(\sin(X) > 1/2\) pour \(X\) de loi \(\mathcal{N}(0, 1)\).

Indication~: si x est un vecteur,
\texttt{y=(np.sin(x)\textgreater{}1/2)} désigne le vecteur de booléens
dont les coordonnées valent \texttt{True} ou \texttt{False} selon que
les coordonnées correspondantes de \texttt{np.sin(x)} soient ou non
\textgreater{} 1/2. Par ailleurs, les vecteurs de booléens se
convertissent naturellement en vecteurs de 0 et de 1 si on leur applique
une fonction dont les arguments sont des nombres (ex~:
\texttt{np.mean}).\solref{V-answer-mc1}

\chapter{Génération de nombres
pseudo-aléatoires}\label{V-guxe9nuxe9ration-de-nombres-pseudo-aluxe9atoires-1}

Les ordinateurs sont des machines déterministes. Il peut sembler
paradoxal de leur demander de générer des nombres aléatoires\footnote{Von
  Neumann (1951) résume ce problème très clairement~: ``Any one who
  considers arithmetical methods of reproducing random digits is, of
  course, in a state of sin. As has been pointed out several times,
  there is no such thing as a random number --- there are only methods
  of producing random numbers, and a strict arithmetic procedure of
  course is not such a method.''}. En réalité, les algorithmes de
génération de nombres aléatoires vont produire des séquences de nombres
déterministes qui vont avoir l'aspect de l'aléatoire. On s'intéresse ici
à la génération de nombres uniformes sur \(]0,1[\) (on exclut les bornes
par commodité), puisque, comme on le verra dans la suite, il s'agit de
l'ingrédient de base de toutes les méthodes de simulation stochastique.

\addcontentsline{toc}{subsection}{Générateur de nombres uniformes
pseudo-aléatoires}
\begin{definition}[{Générateur de nombres uniformes
pseudo-aléatoires}]\phantomsection\label{V-guxe9nuxe9rateur-de-nombres-uniformes-pseudo-aluxe9atoires}

Un \emph{générateur de nombres uniformes pseudo-aléatoires} est un
algorithme qui étant donné une valeur initiale \(u_0\) et une
transformation \(T\) produit une séquence \(u_i=T^i(u_0)\),
\(i \in\mathbb{N}^\ast\), de valeurs dans \(]0,1[\).

Pour tout \(n\in\mathbb{N}^\ast\), les valeurs \((u_1,\ldots,u_n)\)
reproduisent le comportement d'une suite de variables aléatoires
\((V_1,\ldots,V_n)\) i.i.d de loi uniforme sur {]}0,1{[}, lorsqu'on les
compare au travers d'un ensemble de tests statistiques\footnote{Par
  exemple, la suite de tests
  \href{https://en.wikipedia.org/wiki/Diehard_tests}{Die Hard}, due à
  Marsaglia.}, vérifiant par exemple que la corrélation entre deux
nombres successifs est suffisamment faible.
\end{definition}

\addcontentsline{toc}{subsection}{la méthode des congruences}
\begin{exemple}[{la méthode des
congruences}]\phantomsection\label{V-la-muxe9thode-des-congruences}

Cet algorithme, dû à Lehmer (\protect\hyperlink{refV-Lehmer}{1951}), est
l'un des premiers à avoir été proposé et implémenté. Il repose sur 2
paramètres~:

\begin{itemize}
\tightlist
\item
  le multiplicateur \(a\in\mathbb{N}^\ast\),
\item
  le modulo \(m\in\mathbb{N}^\ast\).
\end{itemize}

Etant donné un entier \(k_0\) tel que \(1 \leq k_0 \leq m-1\), la
séquence de nombres est générée par la transformation suivante~: pour
tout \(n\in\mathbb{N}^\ast\) on définit l'entier
\(k_{n+1} = a\,k_n\,\mathrm{\,mod\,} m\) puis on pose
\[u_{n+1} = \frac{k_{n+1}}{m},\hspace{1em} u_0 = \frac{k_0}{m}.\]

On peut remarquer que l'algorithme va produire une séquence de valeurs
qui sera périodique, c'est-à-dire qu'après un certain nombre
d'itérations, la suite se répétera, et qu'il pourra fournir au plus
\(m-1\) valeurs différentes. Ce trait est commun à tous les générateurs
de nombre aléatoires et est lié aux limitations matérielles des
ordinateurs (on ne peut représenter qu'un nombre fini de nombres). Le
choix des valeurs de \(a\) et \(m\) est par conséquent crucial. Il
existe des critères qui permettent de s'assurer du bon comportement de
cette suite~:

\begin{itemize}
\tightlist
\item
  \(m\) est un nombre premier (le plus grand possible),
\item
  \(p = (m-1) /2\) est un nombre premier,
\item
  \(a^p = -1 \mathrm{\,mod\,} m\).
\end{itemize}

Ce critère assure une période pleine et donc que tous les nombres
\((1/m,\ldots,(m-1)/m)\) seront générés. Cependant, les nombres générés
ne sont pas indépendants, pas même non corrélés. On peut montrer que la
corrélation entre 2 nombres successifs vaut approximativement \(1/a\).
Il convient donc de choisir \(a\) suffisamment grand pour que celle-ci
devienne négligeable. Par exemple, en prenant \(a=1000\) et
\(m=2001179\), on obtient une période de \(2001178\) et une corrélation
de l'ordre de \(10^{-3}\).

Ce générateur de nombres uniformes pseudo-aléatoires, bien que
rudimentaire, a ouvert la voie à des générateurs plus sophistiqués,
toujours basés sur des opérations arithmétiques. Le plus couramment
utilisé à ce jour est l'algorithme de
\href{https://en.wikipedia.org/wiki/Mersenne_Twister}{Mersenne-Twister}.
Il s'agit du générateur par défaut de la plupart des logiciels tels que
Python, R, MATLAB, Julia, MS Excel,\ldots{} Il passe notamment avec
succès toute la batterie de tests Die Hard. En particulier, sa période
vaut \(2^{19937} - 1\).

Pour certains usages, cet algorithme n'est cependant pas recommandé du
fait de sa prédictibilité. C'est notamment un défaut rédhibitoire pour
les applications en cryptographie. Il existe des variantes mieux
adaptées à ce cas de figure. On notera enfin qu'il existe des
générateurs de nombres aléatoires basés sur des phénomènes physiques
comme un bruit électrique ou des phénomènes quantiques et donc
parfaitement imprévisibles.
\end{exemple}

\chapter{Méthodes de simulation de variables aléatoires
réelles}\label{V-muxe9thodes-de-simulation-de-variables-aluxe9atoires-ruxe9elles}

On a vu au chapitre II du cours de Probabilités que l'on pouvait
transformer des variables aléatoires réelles suivant certaines lois pour
en obtenir de nouvelles. Par exemple, si \(X_1,\dots,X_n\) sont
\(n\in\mathbb{N}^\ast\) variables gaussiennes centrées réduites
indépendantes, alors \(X_1^2+\dots,X_n^2\) suit une loi du \(\chi^2\) à
\(n\) degrés de liberté. Dans le même esprit, on va voir ici comment
simuler des v.a.r. de lois diverses à partir de la simulation de
variables uniformes sur \(]0,1[\). On introduit une notation qui sera
utile dans la suite~: pour spécifier que deux v.a.r. \(X\) et \(Y\) ont
même loi, on écrira \(X \overset{\mathcal{L}}{=} Y\).

\section{Méthode d'inversion}\label{V-muxe9thode-dinversion}

L'objectif de ce paragraphe est de définir quand et comment il est
possible de simuler une variable aléatoire réelle \(X\) de fonction de
répartition (f.d.r.) \(F_X\) en transformant la simulation d'une
variable aléatoire \(U\) de loi uniforme sur \(]0,1[\). En d'autres
termes, on cherche à déterminer les conditions sous lesquelles il est
possible d'identifier une fonction borélienne
\(\psi :\,]0,1[ \to \mathbb{R}\) telle que
\(X \overset{\mathcal{L}}{=} \psi(U)\).

Commençons par un cadre simple, où \(F_X\) est \textbf{bijective} d'un
intervalle non vide de \(\mathbb{R}\) sur \(]0,1[\).

\addcontentsline{toc}{subsection}{Cas où \(F_X\) est bijective}
\begin{proposition}[{Cas où \(F_X\) est
bijective}]\phantomsection\label{V-invbij}

Soient \(X\) une variable aléatoire réelle de fonction de répartition
\(F_X\) et \(U\) une variable uniforme sur \(]0,1[\). S'il existe un
intervalle non vide \(]a,b[ \subset \mathbb{R}\) tel que
\(F_X :\, ]a,b[ \to \,]0,1[\) est bijective, de bijection réciproque
\(F_X^{-1} :\, ]0,1[ \to ]a,b[\), alors
\(F_X^{-1}(U) \overset{\mathcal{L}}{=} X\) et
\(F_X(X) \overset{\mathcal{L}}{=} U\).
\end{proposition}

\begin{preuve}
Le premier résultat est immédiat~: pour tout \(x\in\mathbb{R}\), par
croissance de \(F_X\) et donc de \(F_X^{-1}\), on a
\[\mathbb{P}\left(F_X^{-1}(U)\leq x\right) = \mathbb{P}\left(U \leq F_X(x) \right) = F_X(x).\]
Concernant le second, notons \(G\) la fonction de répartition de la
variable aléatoire \(F_X(X)\). La croissance de \(F_X\) garantit que
\(\mathbb{P}\left(F_X(X) \in\, ]0,1[\right) = 1\). Ainsi, pour tout
\(x\in\mathbb{R}\) on a bien
\[G(x) = \left|\begin{array}{ll} 1 & \text{si } x \geq 1,\\ \mathbb{P}\left(F_X(X)\leq x\right) = \mathbb{P}\left( X \leq F_X^{-1}(x) \right) = x & \text{si } 0 < x < 1,\\ 0 & \text{sinon.}\end{array}\right.\]

\hfill$\blacksquare$
\end{preuve}

\runinhead{Exercice -- Exemples d'application
1 (\(\mathord{\bullet}\))}\phantomsection\addcontentsline{toc}{subsubsection}{Exemples
d'application 1}\label{V-exemples1}
Donner un algorithme de simulation d'une v.a.r. \(X\) suivant une loi

\begin{itemize}
\tightlist
\item
  uniforme sur un intervalle \(I \subset \mathbb{R}\),
\item
  Exponentielle de paramètre \(\lambda \in \mathbb{R}_+^\ast\),
\item
  de Cauchy, de densité
  \(x\in\mathbb{R}\mapsto \left(\pi\left(1+x^2\right)\right)^{-1}\),
\item
  de Laplace de paramètres \(\mu \in \mathbb{R}\) et
  \(s\in\mathbb{R}_+^\ast\), de densité
  \(x\in\mathbb{R}\mapsto \frac{1}{2s}\,\exp\left\{-\frac{|x-\mu|}{s}\right\}\),
\item
  Logistique de paramètres \(\mu \in \mathbb{R}\) et
  \(s\in\mathbb{R}_+^\ast\), de fonction de répartition
  \(x\in\mathbb{R}\mapsto \left(1 + \exp\left\{-\frac{x-\mu}{s} \right\}\right)^{-1}\)\,?
\end{itemize}

\solref{V-answer-exemples1}

Dans cette situation idéale, \(\psi = F_X^{-1}\) est une solution à
notre problème. Que se passe-t-il en revanche si \(F_X\) n'est pas
bijective\,?

\runinhead{Exercice -- Pile ou
face (\(\mathord{\bullet}\))}\phantomsection\addcontentsline{toc}{subsubsection}{Pile
ou face}\label{V-pf}
Proposer une méthode pour simuler un tir à pile ou face à partir de la
simulation d'une variable uniforme sur \(]0,1[\).\solref{V-answer-pf}

On a déjà vu au chapitre II que les fonctions de répartition de v.a.r.
possèdent un nombre au plus dénombrable de points de discontinuité. Sur
chaque intervalle où elles sont continues, on peut alors considérer
qu'elles sont bijectives, quitte à réduire les zones de palier à un
point. Cela permet de généraliser la notion de bijection réciproque pour
ces fonctions.

\addcontentsline{toc}{subsection}{Réciproque généralisée}
\begin{definition}[{Réciproque
généralisée}]\phantomsection\label{V-defrecgen}

Soit \(F\) une fonction de répartition. On définit sa \emph{réciproque
généralisée} (aussi appelée \emph{inverse généralisée} ou
\emph{pseudo-inverse}) comme la fonction
\[F^{-} : u \in\, ]0,1[\, \mapsto \inf\left\{ x \in \mathbb{R}: F(x) \geq u \right\} \in \mathbb{R}.\]
\end{definition}

\addcontentsline{toc}{subsection}{Conséquences}
\begin{remarque}[{Conséquences}]\phantomsection\label{V-consuxe9quences}

\begin{itemize}
\tightlist
\item
  Cette fonction est bien définie sur tout \(]0,1[\), car quel que soit
  \(u\) dans cet intervalle, l'ensemble
  \(\left\{ x \in \mathbb{R}: F(x) \geq u \right\}\) est non vide et
  minoré. S'il était vide ou non minoré pour un certain
  \(u_0\in\,]0,1[\), pour tout \(x \in \mathbb{R}\) on aurait dans le
  premier cas \(F(x) < u_0 < 0\) et dans le second
  \(F(x) \geq u_0 > 1\). L'une comme l'autre de ces inégalités est
  impossible pour une fonction de répartition.
\item
  La réciproque généralisée de la f.d.r. \(F_X\) d'une v.a.r. \(X\) est
  aussi appelée \emph{fonction quantile}. On pourra notamment remarquer
  que \(F_X^{-}\left(\frac{1}{2}\right)\) est la médiane de \(X\).
\item
  Lorsque \(F\) réalise une bijection d'un intervalle non vide
  \(I\subset \mathbb{R}\) sur \(]0,1[\), sa réciproque généralisée
  coïncide avec sa bijection réciproque.
\end{itemize}

On a alors le résultat suivant, qui indique que \(\psi = F_X^-\) est une
solution universelle à notre problème. La preuve détaillée est donnée en
Annexe.
\end{remarque}

\addcontentsline{toc}{subsection}{Méthode d'inversion}
\begin{theoreme}[{Méthode d'inversion}]\phantomsection\label{V-invgen}

Soient \(U\) une variable uniforme sur \(]0,1[\) ainsi que \(X\) une
variable aléatoire réelle de fonction de répartition \(F_X\) et de
réciproque généralisée \(F_X^-\). Alors
\(F_X^-(U) \overset{\mathcal{L}}{=} X\).
\end{theoreme}

\runinhead{Exercice -- Exemples d'application
2 (\(\mathord{\bullet}\))}\phantomsection\addcontentsline{toc}{subsubsection}{Exemples
d'application 2}\label{V-exemples2}
Donner un algorithme de simulation d'une v.a.r. \(X\) suivant une loi

\begin{itemize}
\tightlist
\item
  Binomiale de paramètres \(n\in\mathbb{N}^\ast\) et \(p \in\, ]0,1[\),
\item
  Uniforme sur l'union de deux segments non vides et disjoints
  \([a,b], [c,d]\subset\mathbb{R}\), de densité
  \(x\in\mathbb{R}\mapsto (b-a + d-c)^{-1}\,1_{[a,b]\cup[c,d]}(x)\)\,?
\end{itemize}

\solref{V-answer-exemples2}

\subsection{Limitations}\label{V-limitations}

La méthode d'inversion peut sembler universelle pour simuler une v.a.r.
\(X\) à partir de \(U \sim \mathcal{U}_{]0,1[}\). Cependant, elle
nécessite en pratique de disposer d'une expression analytique de \(F_X\)
pour en déduire sa réciproque généralisée. Ce n'est typiquement pas le
cas de nombreuses lois usuelles comme la loi Normale. On va donc
déterminer d'autres procédures pour simuler des variables suivant de
telles lois.

\section{Méthode de rejet}\label{V-muxe9thode-de-rejet}

La méthode de rejet est une alternative populaire à la méthode
d'inversion, lorsque cette dernière ne peut être utilisée directement et
que \textbf{la loi cible possède une densité}. On la doit à Von Neumann
(\protect\hyperlink{refV-vonNeumann}{1951}). Pour en comprendre le
fondement, il faut d'abord introduire une généralisation naturelle de la
loi uniforme dans \(\mathbb{R}\) à \(\mathbb{R}^d\)
(\(d\in\mathbb{N}^\ast\)). On notera \(\lambda\) la mesure de Lebesgue
sur \(\mathbb{R}^d\).

\addcontentsline{toc}{subsection}{Loi uniforme sur un borélien}
\begin{definition}[{Loi uniforme sur un
borélien}]\phantomsection\label{V-loi-uniforme-sur-un-boruxe9lien}

La loi uniforme sur un borélien \(A\subset\mathbb{R}^d\) de volume
\(\lambda(A) > 0\) est une loi de probabilité admettant pour densité
\[f : x\in\mathbb{R}^d \mapsto \dfrac{1_A(x)}{\lambda(A)}.\]
\end{definition}

\runinhead{Exercice -- Loi uniforme sur un
pavé (\(\mathord{\bullet}\))}\phantomsection\addcontentsline{toc}{subsubsection}{Loi
uniforme sur un pavé}\label{V-unipave}
Comment simuler un vecteur aléatoire \((U_1,\dots,U_d)\) de loi uniforme
sur un pavé non vide
\([a_1,b_1]\times\dots\times[a_d,b_d] \subset \mathbb{R}^d\)\,?\solref{V-answer-unipave}

Une propriété fondamentale de cette loi est qu'elle reste stable par
conditionnement, dans le sens suivant.

\addcontentsline{toc}{subsection}{Stabilité par conditionnement}
\begin{proposition}[{Stabilité par
conditionnement}]\phantomsection\label{V-stabunif}

Soit \(U\) un vecteur aléatoire de loi uniforme sur un borélien
\(A\subset\mathbb{R}^d\) de volume \(\lambda(A) > 0\). Alors pour tout
borélien \(B \subset A\) de volume \(\lambda(B) > 0\), la loi
conditionnelle \(\mathbb{P}_{U\mid U\in B}\) de \(U\) sachant que
\(U \in B\) est uniforme sur \(B\).
\end{proposition}

\runinhead{Exercice -- Démonstration (\(\mathord{\bullet}\))}\phantomsection\addcontentsline{toc}{subsubsection}{Démonstration}\label{V-unicond}
\solref{V-answer-unicond}

\runinhead{Exercice -- Simulation (\(\mathord{\bullet}\))}\phantomsection\addcontentsline{toc}{subsubsection}{Simulation}\label{V-simunibor}
Déduire de la \protect\hyperref[V-stabunif]{proposition
précédente (p.~\pageref*{V-stabunif})} une méthode pour simuler un
vecteur aléatoire de loi uniforme sur un borélien
\(B \subset \mathbb{R}^d\) \textbf{borné} et de volume
\(\lambda(B) > 0\).\solref{V-answer-simunibor}

Il est aussi possible de simuler un vecteur aléatoire de loi uniforme
sur certains boréliens non vides et non bornés, comme l'établit la
proposition ci-dessous.

\addcontentsline{toc}{subsection}{Loi uniforme sur le sous-graphe d'une
densité}
\begin{proposition}[{Loi uniforme sur le sous-graphe d'une
densité}]\phantomsection\label{V-simunifdens}

Soient une densité \(f : \mathbb{R}\to \mathbb{R}\), une v.a.r. \(X\) et
une variable \(U\) uniforme sur \(]0,1[\), indépendante de \(X\). On
note
\(A_f := \left\{ (x,y) \in \mathbb{R}\times \mathbb{R}_+ : f(x) \geq y \right\}\)
le domaine limité par le graphe de \(f\) et l'axe des abscisses (souvent
appelé sous-graphe de \(f\)). Si \(X\) est de densité \(f\), alors le
couple \((X,Uf(X))\) suit une loi uniforme sur \(A_f\). Réciproquement, si \((X,Y)\) suit une loi uniforme sur \(A_f\), alors \(X\) est de densité \(f\) et \(Y\) conditionnellement à \(X\) est uniforme sur \(]0,f(X)[\).
\end{proposition}

\begin{preuve}
Commençons par remarquer que \(\lambda(A_f) = 1\). Quel que soit
\((z,v)\in\mathbb{R}^2\), par indépendance de \(X\) et \(U\) et par
Fubini on a \begin{align*}
\mathbb{P}\left( X \leq z, Uf(X) \leq v \right) &= \int_\mathbb{R}\int_\mathbb{R}1_{]-\infty,z]}(x)\,1_{]-\infty,v]}(uf(x))\,1_{]0,1[}(u)\,f(x)\,du\,dx\\
&= \int_{-\infty}^z \left(\int_\mathbb{R}1_{]-\infty,v]\cap ]0,f(x)[}(uf(x))\,f(x)\,du\right)\,dx\\
&= \int_{-\infty}^z \int_{-\infty}^v 1_{]0,f(x)[}(u)\,du\,dx\\
& = \int_{-\infty}^z \int_{-\infty}^v 1_{A_f}(x,u)\,du\,dx.
\end{align*} Ainsi, \((X,Uf(X))\) admet pour densité \(1_{A_f}\), qui
correspond bien à celle d'une loi uniforme sur
\(A_f\).

Réciproquement, supposons que \((X,Y)\) suit une loi uniforme sur
\(A_f\). Alors pour tout \(x\in\mathbb{R}\)
\[\mathbb{P}\left(X \leq x\right) = \int_{-\infty}^x \int_\mathbb{R} 1_{A_f}(u,v)\,du\,dv = \int_{-\infty}^x f(u)\,du,\] ce qui montre que \(X\) est de densité \(f\). De plus, pour tout \(x\in\mathbb{R}\) et \(y\in\mathbb{R}_+\),
\begin{align*}
\mathbb{P}\left(Y \leq y, X \leq x\right) 
&= \int_{-\infty}^x \int_\mathbb{R} 1_{A_f}(u,v)\,1_{]-\infty,y]}(v)\,dv\,du\\
&= \int_{-\infty}^x \int_0^{\min(f(u),y)} dv\,du\\
&= \int_{-\infty}^x \min(f(u),y)\,du.
\end{align*}
ce qui montre que \(Y\) conditionnellement à \(X\) est uniforme sur \(]0,f(X)[\).

\hfill$\blacksquare$
\end{preuve}

Pour simuler un vecteur uniforme \((X,Y)\) sur un ensemble \(A_f\) tel
que défini à la proposition précédente, il suffit donc de simuler une
v.a.r. \(X\) de densité \(f\), puis une variable \(U\) uniforme sur
\(]0,1[\), et de poser \(Y = f(X)U\).

En combinant les résultats précédents, on obtient la méthode de rejet,
illustrée sur la figure ci-dessous, où
\(A_Y := \left\{ (x,y) \in \mathbb{R}\times \mathbb{R}_+ : y \leq f_Y(x) \right\}\).

\subsection{Méthode de rejet}\label{V-muxe9thode-de-rejet-1}

On souhaite simuler une variable aléatoire réelle \(X\) de densité
\(f_X\). Supposons que l'on sait simuler une variable \(U\) uniforme sur
\(]0,1[\), ainsi qu'une v.a.r. \(Y\) (par exemple avec la méthode
d'inversion) de densité \(f_Y\) telle qu'il existe un réel \(a > 0\)
pour lequel on a \(\forall x \in \mathbb{R}\)~:
\(f_X(x) \leq a\,f_Y(x)\). Il suffit alors de suivre l'algorithme
suivant~:

\begin{enumerate}
\def\labelenumi{\arabic{enumi}.}
\item
  simuler \(Y\) et \(U\),
\item
  si \(aUf_Y(Y) > f_X(Y)\), recommencer à l'étape 1.
\item
  poser \(X = Y\).
\end{enumerate}

\begin{figure}
\centering
\includegraphics{MetRej.tex.pdf}
\caption{Méthode de rejet}\label{V-fig-MetRej.tex}
\end{figure}

\subsection{Limitations}\label{V-limitations-1}

La méthode de rejet a l'avantage de permettre de simuler des variables
aléatoires à densité lorsque la méthode d'inversion est inapplicable. Néanmoins, pour pouvoir l'appliquer il faut indiquer une densité instrumentale qui,
multipliée par un réel positif, majore la densité cible, et que l'on
sait simuler. Le taux de rejet, c'est-à-dire la probabilité de
l'événement \(\{aUf_Y(Y) > f_X(Y)\}\), peut parfois être élevé,
notamment lorsque la dimension de \(X\) est grande, ce qui limite
l'efficacité de la méthode.

\runinhead{Exercice -- Taux de
rejet (\(\mathord{\bullet}\mathord{\bullet}\))}\phantomsection\addcontentsline{toc}{subsubsection}{Taux
de rejet}\label{V-tauxrej}
Calculer le taux de rejet de la méthode proposée
ci-dessus.\solref{V-answer-tauxrej}

\section{Simulation de variables aléatoires gaussiennes~:
Box-Muller}\label{V-simulation-de-variables-aluxe9atoires-gaussiennes-box-muller}

On a vu que la méthode d'inversion est inappropriée pour simuler une
variable gaussienne, puisqu'elle requiert l'expression analytique de la
fonction de répartition cible. La méthode de rejet est quant à elle sous-optimale,dans le sens où toutes les variables uniformes générées ne sont pas directement utilisées.
La loi normale étant fondamentale en probabilité, il est nécessaire de disposer d'une méthode de simulation exacte et
efficace.

George E. P. Box et Mervin E. Muller ont proposé en 1958 une telle
méthode (G.E.P. Box and M.E. Muller
(\protect\hyperlink{refV-BoxMuller}{1958})). Elle exploite la propriété
d'invariance par rotation de la densité d'un couple de variables
gaussiennes indépendantes centrées réduites.

\begin{proposition}\phantomsection\label{V-boxmuller}
Soient \(U\) et \(V\) deux variables indépendantes, de loi uniforme sur
\(]0,1[\). Alors les variables aléatoires
\(X = \sqrt{-2\ln(U)}\cos\left(2\pi V\right)\) et
\(Y = \sqrt{-2\ln(U)}\sin\left(2\pi V\right)\) sont indépendantes et
suivent toutes deux une loi normale centrée réduite.
\end{proposition}

\begin{preuve}
On considère \((\widetilde{X},\widetilde{Y})\) un vecteur aléatoire dont
les deux composantes sont indépendantes, de loi normale centrée réduite.
On note ses coordonnées polaires aléatoires \(\widetilde{R}\) et
\(\widetilde{\Theta}\). On a vu dans le cours de Probabilités II que
dans ce cas, \(\widetilde{R}\) et \(\widetilde{\Theta}\) sont
indépendantes, la première de densité
\(f_{\widetilde{R}} : r \in \mathbb{R}\mapsto r\,e^{-\frac{r^2}{2}} 1_{\mathbb{R}_+^\ast}(r)\)
et la seconde de loi uniforme sur \(]0,2\pi]\).

\(X\) et \(Y\) sont les coordonnées cartésiennes obtenues à partir d'un
rayon et d'un angle~: en posant \(R = \sqrt{-2\ln(U)}\) et
\(\Theta = 2\pi V\) on obtient \(X =R\cos(\Theta)\) et
\(Y = R\sin(\Theta)\). Par indépendance de \(U\) et \(V\), \(R\) et
\(\Theta\) sont indépendantes. Pour que le vecteur \((X,Y)\) ait la même
distribution que
\((\widetilde{X},\widetilde{Y}) = \left(\widetilde{R}\cos(\widetilde{\Theta}),\widetilde{R}\sin(\widetilde{\Theta})\right)\),
il suffit de montrer que \(R \overset{\mathcal{L}}{=} \widetilde{R}\) et
\(\Theta \overset{\mathcal{L}}{=} \widetilde{\Theta}\).

\begin{itemize}
\item
  Commençons par étudier la loi de \(R\), de fonction de répartition
  \(F_R\). On remarque que la fonction
  \(u\in\, ]0,1[ \mapsto \sqrt{-2\ln(u)} \in \mathbb{R}_+^\ast\) est
  bijective, strictement décroissante. Ainsi, pour tout
  \(r\in\mathbb{R}_-\) on a \(\mathbb{P}\left(R \leq r \right) = 0\) et
  pour tout \(r \in \mathbb{R}_+^\ast\) on a
  \[\mathbb{P}(R \leq r) = \mathbb{P}\left(\sqrt{-2\ln(U)} \leq r \right) = \mathbb{P}\left(U \geq e^{-\frac{r^2}{2}} \right) = 1 - e^{-\frac{r^2}{2}}.\]
  qui correspond à la fonction de répartition associée à la densité
  \(f_{\widetilde{R}}\).
  \[\int_{-\infty}^r f_{\widetilde{R}}(x)\,dx = \left|\begin{array}{ll}\displaystyle \int_0^r x\,e^{-\frac{x^2}{2}}\,dx = \left[-e^{-\frac{x^2}{2}} \right]_0^r = 1 - e^{-\frac{r^2}{2}}  & \text{si } r>0,\\[1em] 0 & \text{sinon.} \end{array}\right.\]
\item
  En ce qui concerne la loi de \(\Theta\), de fonction de répartition
  \(F_\Theta\), puisque la fonction
  \(v \in\, ]0,1[\, \mapsto 2\pi v \in\, ]0,2\pi[\) est bijective
  strictement croissante, on a que pour tout \(\theta \in \mathbb{R}\)
  \[F_\Theta(\theta) = \left|\begin{array}{ll} 1 & \text{si } \theta \geq 2\pi,\\ \mathbb{P}\left(V\leq \frac{\theta}{2\pi} \right) = \dfrac{\theta}{2\pi} & \text{si } \theta \in\, ]0,2\pi[,\\ 0 & \text{si } \theta \leq 0,\end{array}\right.\]
  qui est bien la fonction de répartition d'une loi uniforme sur
  \(]0,2\pi[\).
\end{itemize}

\hfill$\blacksquare$
\end{preuve}

Cette méthode permet de simuler directement deux variables gaussiennes
centrées réduites indépendantes à partir de deux variables uniformes
indépendantes. Pour simuler une variable gaussienne d'espérance
\(m \in \mathbb{R}\) et de variance \(\sigma^2 \in \mathbb{R}_+^\ast\)
quelconques, on se rappelera que si \(X\) suit une loi normale centrée
réduite, alors \(\sigma X + m\) suit une loi normale d'espérance \(m\)
et de variance \(\sigma^2\). Dans les logiciels de calcul scientifique (python, R, Julia, etc.), la fonction de simulation d'une variable gaussienne repose cependant sur d'autres algorithmes, plus rapides, comme l'algorithme de la \emph{ziggurat} (G. Marsaglia and W. W. Tsang
(\protect\hyperlink{refV-Ziggurat}{2000})) : on recouvre la densité par
des rectangles empilés de même aire, on tire un point uniformément dans
un rectangle, et on ne recourt à un calcul coûteux (exponentielle, rejet) que dans de rares cas.

\chapter{Simulation d'un vecteur gaussien à
densité}\label{V-simulation-dun-vecteur-gaussien-uxe0-densituxe9}

On souhaite simuler un vecteur gaussien \(X = (X_1,\ldots,X_d)\) à
valeurs dans \(\mathbb{R}^d\) d'espérance \(m\) et de matrice de
covariance \(C\) définie positive (et donc inversible) données.

Puisque la matrice \(C\) est définie positive, elle admet une racine
carrée, c'est-à-dire qu'il existe une matrice \(N\) telle que
\(C = N\,N^t\). En effet, on peut par exemple décomposer \(C\) de la
manière suivante~: \[C = V\,D\,V^t,\] où \(V\) est une matrice
orthogonale et \(D\) est la matrice diagonale dont les termes diagonaux
sont les valeurs propres (toutes strictement positives) de \(C\). Il
suffit alors de prendre \(N = V\,D^{1/2}\), où \(D^{1/2}\) est la
matrice diagonale dont les termes diagonaux sont les racines carrées des
valeurs propres.

En pratique, il est coûteux numériquement d'effectuer le calcul des
valeurs propres et des vecteurs propres de \(C\). On va plutôt calculer
sa
\href{https://fr.wikipedia.org/wiki/Factorisation_de_Cholesky}{\emph{décomposition}
ou \emph{factorisation de Cholesky}} qui permet d'écrire \[C = L\,L^t\]
avec \(L\) une matrice triangulaire inférieure.\footnote{Cette
  décomposition est très utile dans la résolution de systèmes linéaires
  de la forme \(A\,x = b\), où \(b\) est connu, \(x\) inconnu et \(A\)
  est définie positive. Cela revient à résoudre \(L\,L^t\,x = b\). On
  pose alors \(y = L^t\,x\) et on résout d'abord \(Ly=b\), ce qui est
  très rapide puisque \(L\) est triangulaire inférieure (on commence par
  \(y_1 = b_1/L_{11}\), puis \(y_2 = (b_2 - L_{21}y_1)/L_{22}\), etc. en
  descendant). On résout ensuite \(L^t\,x = y\), ce qui est aussi très
  rapide pour la même raison (on commence par \(x_n = y_n/L_{nn}\) puis
  on remonte).}

Soit maintenant un autre vecteur gaussien \(Y = (Y_1,\ldots,Y_d)\) à
valeurs dans \(\mathbb{R}^d\) et de matrice de covariance identité,
notée \(I_d\). Autrement dit, les \(Y_i\) sont des variables aléatoires
gaussiennes centrées, réduites et indépendantes.

Alors, le vecteur \(Z = m + L\,Y\) est gaussien, d'espérance \(m\) et de
matrice de covariance \(C\). En effet, \(Z\) est gaussien comme
combinaison linéaire de variables aléatoires gaussiennes,
\(\mathbb{E}(Z) = \mathbb{E}(m + L\,Y) = m\) et
\(\mathbb{V}(Z) = \mathbb{E}((L\,Y)^2) = L I_d L^t = C\).

\chapter{Echantillonnage
d'importance}\label{V-echantillonnage-dimportance-1}

On introduit dans cette section la méthode d'échantillonnage
d'importance (importance sampling en anglais), que l'on appelle aussi,
de manière plus intuitive, échantillonnage préférentiel, pour les lois à
densité. Pour ce faire, nous allons commencer par un exemple qui montre
qu'il peut être plus efficace de simuler des valeurs selon une loi
différente de celle d'intérêt, autrement dit de modifier la
représentation de l'intégrale \(\mathcal{I}\) sous la forme d'une
espérance calculée selon une autre densité.

\addcontentsline{toc}{subsection}{Loi de Cauchy conditionnelle}
\begin{exemple}[{Loi de Cauchy
conditionnelle}]\phantomsection\label{V-loi-de-cauchy-conditionnelle}

Supposons que l'on s'intéresse à calculer la probabilité \(p\) qu'une
variable \(X\) de loi de Cauchy standard soit plus grande que 2 (on peut
le calculer directement et \(p=0.15\))
\[p = \int_2^{+\infty} \frac{1}{\pi(1+x^2)}dx.\] Si on estime \(p\)
directement à partir d'un échantillon \((X_1,\ldots,X_n)\) simulé selon
la loi de Cauchy standard, soit
\[\widehat{p}_1 = \frac{1}{n}\sum_{i=1}^n 1_{X_i > 2},\] la
\emph{variance de l'estimateur}
\(\mathbb{V}(\widehat{p}_1) = p(1-p)/n = 0.127/n\), puisque
\(\widehat{p}_1\) suit une loi binomiale de paramètre \((n,p)\). On peut
réduire cette variance (et donc améliorer la qualité de l'estimateur) en
tirant parti de la symétrie de la densité de la loi de Cauchy, en
formant un second estimateur
\[\widehat{p}_2 = \frac{1}{n}\sum_{i=1}^n 1_{|X_i| > 2},\] dont la
variance vaut \(\mathbb{V}(\widehat{p}_2) = p(1-p/2)/2n = 0.052/n\). La
relative inefficacité de ces méthodes est due au fait que la majeure
partie des valeurs simulées seront en dehors de la zone d'intérêt
\(]2,+\infty[\). En passant par le complémentaire, on peut réécrire
\(p\) comme \[p = \frac{1}{2} - \int_0^2 \frac{1}{\pi(1+x^2)}dx,\] dont
le second terme peut être vu comme l'espérance de
\(h(U) = \frac{2}{\pi(1+U^2)}\) avec \(U\sim \mathcal{U}_{[0,2]}\).
Tirant un échantillon \((U_1,\ldots,U_n)\) i.i.d. de loi uniforme sur
\([0,2]\), on obtient un troisième estimateur~:
\[\widehat{p}_3 = \frac{1}{2} - \frac{1}{n}\sum_{i=1}^n \frac{2}{\pi(1+U_i^2)},\]
dont la variance vaut
\(\mathbb{V}(\widehat{p}_3) = (\mathbb{E}(h(X)^2) - \mathbb{E}(h(U))^2)/n = 0.0285/n\)
(par intégration par parties). Enfin, on peut encore réécrire (voir
Ripley (\protect\hyperlink{refV-ripley}{1987}))
\[p = \int_0^{1/2}\frac{y^{-2}}{\pi(1+y^{-2})}dy,\] qui peut être vue
comme \(\mathbb{E}\left(\frac{V^{-2}}{2\pi(1+V^{-2})}\right)\) avec
\(V\sim \mathcal{U}_{]0,1/2[}\). L'estimateur formé à partir de cette
représentation et d'un échantillon \((V_1,\ldots,V_n)\) i.i.d. de loi
uniforme sur \([0,1/2]\) a une variance de \(0.95\, 10^{-4}/n\). Il est
donc bien plus efficace que \(\widehat{p}_1\) puisqu'il nécessite
environ \(\sqrt{10^3}\approx 32\) fois moins de simulations pour
atteindre la même précision.

On a ainsi vu sur ce cas particulier que l'estimation d'une intégrale de
la forme
\[\mathcal{I}=\mathbb{E}\left(h(X)\right)=\int_{\mathbb{R}^d} h(x) f(x) dx,\]
peut s'écrire de différentes manières, en faisant varier \(h\) et \(f\).
Par conséquent, un estimateur ``optimal'' devrait tenir compte de
l'ensemble de ces possibilités. C'est justement l'idée développée dans
la méthode d'échantillonnage d'importance dont le principe est décrit
dans la définition suivante.
\end{exemple}

\addcontentsline{toc}{subsection}{Echantillonnage d'importance}
\begin{definition}[{Echantillonnage
d'importance}]\phantomsection\label{V-is}

La méthode d'\emph{échantillonnage d'importance} est une évaluation de
\(\mathcal{I}\) basée sur la simulation d'un échantillon
\(X_1, \ldots,X_n\) de loi de densité \(g\) et approximant~:
\[\mathbb{E}_{f}\left(h(X)\right)\approx \frac{1}{n}\sum_{i=1}^{n}\frac{f(X_i)}{g(X_i)}h(X_i),\]
où la notation \(\mathbb{E}_f\) signifie que l'espérance est calculée
avec \(X\sim f\). On appelle souvent les ratios
\(\frac{f(X_i)}{g(X_i)}\) les \emph{poids d'importance} que l'on note
\(w_i\).
\end{definition}

Cette méthode est basée sur la représentation suivante de
\(\mathcal{I}\)~:

\[\mathbb{E}_{f}\left[h(X)\right]=\int h(x)\frac{f(x)}{g(x)}g(x)\,dx\]

que l'on appelle l'\emph{identité fondamentale de l'échantillonnage
d'importance} et l'estimateur converge du fait de la
\protect\hyperref[II-lfgn]{loi forte des grands nombres}.

Cette identité indique qu'une intégrale du type \(\mathcal{I}\) n'est
pas intrinsèquement associée à une loi donnée. L'intérêt de
l'échantillonnage d'importance repose sur le fait qu'il n'y a aucune
restriction sur le choix de la densité \(g\), dite \emph{instrumentale},
que l'on peut donc choisir parmi les densités des lois que l'on sait
simuler aisément. Il y a bien évidemment des choix qui sont meilleurs
que d'autres. Remarquons tout d'abord que bien que l'estimateur proposé
dans la \protect\hyperref[V-is]{définition
ci-dessus (p.~\pageref*{V-is})} converge presque sûrement, sa variance
est finie si
\[\mathbb{E}_g\left(h^2(X)\frac{f^2(X)}{g^2(X)}\right) = \mathbb{E}_f\left(h^2(X)\frac{f(X)}{g(X)}\right) = \int h^2(x)\frac{f^2(x)}{g(x)}\,dx < \infty.\]
On préconise alors l'usage de densités instrumentales \(g\) dont la
queue de distribution est plus épaisse que celle de \(f\) pour éviter
que cette variance puisse être infinie (on notera que cela dépend aussi
de la fonction \(h\) à intégrer). En pratique, on utilise généralement
l'estimateur suivant, de variance finie et qui donne des résultats plus
stables numériquement que celui de la définition~:
\[\frac{\sum_{i=1}^{n}w_ih(X_i)}{\sum_{i=1}^n w_i}\] où on a remplacé
\(n\) par la somme des poids d'importance. Puisque
\(\frac{1}{n}\sum_{i=1^n}w_i = \frac{1}{n}\sum_{i=1}^n\frac{f(x_i)}{g(x_i)}\)
tend vers 1 quand \(n\to\infty\), cet estimateur converge presque
sûrement vers \(\mathbb{E}_{f}\left(h(X)\right)\) par la loi forte des
grands nombres (voir C. P. Robert and G. Casella
(\protect\hyperlink{refV-roca}{2004})).

Parmi les densités \(g\) qui fournissent des estimateurs de variance
finie, il est possible d'exhiber la densité optimale (au sens de la
variance de l'estimateur associé) pour une fonction \(h\) et une densité
\(f\) données.

\addcontentsline{toc}{subsection}{Densité optimale}
\begin{theoreme}[{Densité
optimale}]\phantomsection\label{V-densituxe9-optimale}

Le choix de \(g\) qui minimise la variance de l'estimateur donné dans la
\protect\hyperref[V-is]{définition ci-dessus (p.~\pageref*{V-is})} est
\[g^\ast(x) = \frac{|h(x)|f(x)}{\int |h(x)|f(x) dx}.\]
\end{theoreme}

\begin{preuve}
Notons d'abord que
\[\mathbb{V}_g\left(\frac{h(X)f(X)}{g(X)}\right) = \mathbb{E}_g\left(h^2(X)\frac{f^2(X)}{g^2(X)}\right) - \mathbb{E}_g\left(h(X)\frac{f(X)}{g(X)}\right)^2\]
et le second terme ne dépend pas de \(g\). On minimise donc le premier
terme. D'après \protect\hyperref[II-jensen]{l'inégalité de
Jensen (p.~\pageref*{II-jensen})}, on a
\[\mathbb{E}_g\left(h^2(X)\frac{f^2(X)}{g^2(X)}\right) \geq \mathbb{E}_g\left(|h(X)|\frac{f(X)}{g(X)}\right)^2 = \left(\int|h(x)|f(x)dx\right)^2,\]
qui nous donne une borne inférieure indépendante du choix de \(g\). Elle
est atteinte en prenant \(g=g^\ast\).\hfill$\blacksquare$
\end{preuve}

Ce résultat formel n'a que peu d'intérêt pratique~: le choix optimal de
\(g\) fait intervenir \(\int|h(x)|f(x)dx\), qui est à une valeur absolue
près la quantité que l'on souhaite estimer\,! Il suggère néanmoins de
considérer des densités \(g\) telles que \(|h|f/g\) est quasi constante
et de variance finie. On se reportera au chapitre 3 de C. P. Robert and
G. Casella (\protect\hyperlink{refV-roca}{2004}) pour des exemples où un
bon choix de \(g\) permet des améliorations considérables par rapport à
des estimateurs de Monte-Carlo plus naïfs.

\chapter{Exercices}\label{V-exercices}

\section{Loi uniforme dans un
domaine}\label{V-loi-uniforme-dans-un-domaine}

Ecrire un algorithme pour simuler un point uniforme dans les trois
domaines représentés ci-dessous.

\begin{figure}
\centering
\fbox{[figure manquante : ex1-eps-converted-to.pdf]}
\caption{Domaine A}\label{V-fig-ex1-eps-converted-to.pdf}
\end{figure}

\begin{figure}
\centering
\fbox{[figure manquante : ex2-eps-converted-to.pdf]}
\caption{Domaine B}\label{V-fig-ex2-eps-converted-to.pdf}
\end{figure}

\begin{figure}
\centering
\fbox{[figure manquante : ex3-eps-converted-to.pdf]}
\caption{Domaine C}\label{V-fig-ex3-eps-converted-to.pdf}
\end{figure}

\section{Simulation selon la loi
géométrique}\label{V-simulation-selon-la-loi-guxe9omuxe9trique}

Cette loi correspond à la loi du nombre d'essais avant d'obtenir un
évènement de probabilité \(p\)
\[X\sim{\cal G} (p)~~\mathbb{P}(X=n)=p(1-p)^{n-1}~~n=1,2,\ldots\]

\runinhead{Question
1}\phantomsection\addcontentsline{toc}{subsubsection}{Question
1}\label{V-loigeom1}
Calculer l'espérance et la variance de \(X\)\solref{V-answer-loigeom1}

\runinhead{Question
2}\phantomsection\addcontentsline{toc}{subsubsection}{Question
2}\label{V-loigeom2}
Générer un échantillon de taille 1000 pour \(p=0.6\), en utilisant une
approche génétique. Calculer la moyenne et la variance. Comparer avec
les valeurs théoriques.\solref{V-answer-loigeom2}

\runinhead{Question
3}\phantomsection\addcontentsline{toc}{subsubsection}{Question
3}\label{V-loigeom3}
Montrer que si \(X\sim{\cal E}(\lambda)\), alors
\(\lceil X \rceil \sim{\cal G}(p)\). On précisera la valeur de
\(\lambda\).\solref{V-answer-loigeom3}

\runinhead{Question
4}\phantomsection\addcontentsline{toc}{subsubsection}{Question
4}\label{V-loigeom4}
Donner, implémenter et tester un algorithme pour simuler \({\cal G}(p)\)
directement\solref{V-answer-loigeom4}

\section{Simulation de la loi gaussienne par la méthode de
rejet}\label{V-simulation-de-la-loi-gaussienne-par-la-muxe9thode-de-rejet}

\(X\) est une variable aléatoire de loi gaussienne d'espérance \(m=0\)
et de variance \(\sigma^2=1\), \(\mathcal{N}(0,1)\). Sa densité est
\[f(x) = \frac{1}{\sqrt{2\pi}}\exp\left(-\frac{1}{2}x^2\right),~x\in \mathbb{R}\]

\runinhead{Question
1}\phantomsection\addcontentsline{toc}{subsubsection}{Question
1}\label{V-lgrej1}
Montrer que \(f(x) \leq C \exp \left(-|x|\right)\) où
\(C=\sqrt{\frac{e}{2\pi}}\)\solref{V-answer-lgrej1}

\runinhead{Question
2}\phantomsection\addcontentsline{toc}{subsubsection}{Question
2}\label{V-lgrej2}
On considère
\(g(x) = \frac{1}{2}\exp \left(-|x|\right),~x\in \mathbb{R}\). Montrer
que \(g\) est une densité de probabilité.\solref{V-answer-lgrej2}

\runinhead{Question
3}\phantomsection\addcontentsline{toc}{subsubsection}{Question
3}\label{V-lgrej3}
Proposer un algorithme pour simuler selon \(g\).\solref{V-answer-lgrej3}

\runinhead{Question
4}\phantomsection\addcontentsline{toc}{subsubsection}{Question
4}\label{V-lgrej4}
Implémenter l'algorithme de rejet pour simuler selon
\(f\).\solref{V-answer-lgrej4}

\section{Simulation selon la loi de Wigner}\label{V-wigner}

La loi de Wigner (ou du demi-cercle) est la loi de support \([-2, 2]\)
et de densité \(f(x) = \frac{1}{2\pi}\sqrt{4 - x^2}\). Simuler un grand
nombre de variables aléatoires de loi de Wigner, représenter
l'histogramme associé et le comparer à la
densité.\solref{V-answer-wigner}

\section{Loi des grands nombres et théorème central
limite}\label{V-loi-des-grands-nombres-et-thuxe9oruxe8me-central-limite}

\runinhead{Question
1}\phantomsection\addcontentsline{toc}{subsubsection}{Question
1}\label{V-lgntcl1}
Afin d'illustrer la \href{Probabilité\%20IV.pdf\%20\#lfgn}{Loi des
Grands Nombres}, visualiser la suite
\(S_n = \frac{X_1 + \ldots +X_n}{n}\) pour
\((X_i)_{i \in \mathbb{N}^\ast}\) une suite de variables aléatoires
indépendantes de loi uniforme sur {[}-1, 1{]}.

Indication~: pour \(x= [x_1 , \ldots, x_n]\) vecteur,
\texttt{np.cumsum(x)} est le vecteur
\([x_1 , x_1 + x_2 , x_1 + x_2 + x_3 , \ldots , x_1 + \ldots + x_n]\)
des sommes cumulées des coordonnées de \(x\).\solref{V-answer-lgntcl1}

\runinhead{Question
2}\phantomsection\addcontentsline{toc}{subsubsection}{Question
2}\label{V-lgntcl2}
Faire de même avec des variables aléatoires \(Y_i\) de loi de Cauchy
(c'est-à-dire de densité \(\frac{1}{\pi(1+x^2)}\)). La suite
\(\frac{Y_1 + \ldots +Y_n}{n}\) semble-t-elle converger\,?
Pourquoi\,?\solref{V-answer-lgntcl2}

\runinhead{Question
3}\phantomsection\addcontentsline{toc}{subsubsection}{Question
3}\label{V-lgntcl3}
Calculer la moyenne \(\mu\) et l'écart-type \(\sigma\) des \(X_i\) et
vérifier, afin d'illustrer le
\href{Probabilité\%20IV.pdf\%20\#TCL}{Théorème Central Limite}, que
\(\sqrt{n} (S_n - \mu)/\sigma\) converge en loi, lorsque
\(n \to \infty\), vers une variable de loi \(\mathcal{N}(0, 1)\)~: on se
fixera une grande valeur de n, on simulera un grand nombre de fois
\(\sqrt{n} (S_n - \mu)/\sigma\), on en tracera l'histogramme et on le
comparera à la densité \((2\pi)^-1/2 e^{-x^2}\) la loi
\(\mathcal{N}(0, 1)\).\solref{V-answer-lgntcl3}

\section{Simulation d'un mélange de
gaussiennes}\label{V-simulation-dun-muxe9lange-de-gaussiennes}

On souhaite simuler selon la loi de mélange introduite dans l'exercice
Mélange de loi du chapitre 3.

\runinhead{Question
1}\phantomsection\addcontentsline{toc}{subsubsection}{Question
1}\label{V-melloi1}
Soit \(K\) une v.a.r. de loi~:
\[\mathbb{P}(K=x) = \left\{\begin{array}{ll} 1/2 & \text{ si }x = 1 \\ 1/4 & \text{ si }x = 2 \\ 1/4 & \text{ si }x = 3 \\ 0 & \text{ sinon} \end{array} \right.\]
Implémenter la méthode d'inversion pour simuler
\(K\)\solref{V-answer-melloi1}

\runinhead{Question
2}\phantomsection\addcontentsline{toc}{subsubsection}{Question
2}\label{V-melloi2}
Soit \(X_1\), \(X_2\) et \(X_3\) des variables aléatoires de lois
respectives \(\mathcal{N}(0,1)\), \(\mathcal{N}(5,1/2)\) et
\(\mathcal{N}(8,4)\). Implémenter un algorithme de simulation de
\(X_K\). Comparer l'histogramme obtenu pour un échantillon de taille
1000 avec la densité de \(X_K\).\solref{V-answer-melloi2}

\section{Echantillonnage
d'importance}\label{V-echantillonnage-dimportance-2}

On cherche à évaluer l'espérance d'une variable aléatoire \(X\)
gaussienne centrée réduite (de densité \(f_X\) et de f.d.r. \(F_X\))
sachant qu'elle dépasse la valeur 3.

\runinhead{Question
1}\phantomsection\addcontentsline{toc}{subsubsection}{Question
1}\label{V-echimp1}
Exprimer la densité conditionnelle de
\(X|X>3\).\solref{V-answer-echimp1}

\runinhead{Question
2}\phantomsection\addcontentsline{toc}{subsubsection}{Question
2}\label{V-echimp2}
Calculer la quantité souhaitée à partir d'un échantillon de taille 1000
généré selon \(F_X\). Quelle est la valeur du taux de
rejet\,?\solref{V-answer-echimp2}

\runinhead{Question
3}\phantomsection\addcontentsline{toc}{subsubsection}{Question
3}\label{V-echimp3}
Implémenter l'algorithme d'échantillonnage d'importance en prenant comme
loi instrumentale la loi gaussienne d'espérance 3 et de variance 1.
Quelle est la proportion de poids nuls\,?

\begin{center}\rule{0.5\linewidth}{0.5pt}\end{center}

\solref{V-answer-echimp3}


\chapter{Annexe}\label{V-annexe}

\section{Preuve de la méthode
d'inversion}\label{V-preuve-de-la-muxe9thode-dinversion}

Pour pouvoir démontrer le \protect\hyperref[V-invgen]{théorème de la
méthode d'inversion (p.~\pageref*{V-invgen})}, il faut d'abord établir
un certain nombre de propriétés de la réciproque généralisée d'une
fonction de répartition. Elle peuvent être visualisées sur la figure
ci-dessous.

\begin{proposition}\phantomsection\label{V-proprecgen}
Soit \(F\) une fonction de répartition. Alors sa réciproque généralisée
\(F^-\) satisfait les propriétés suivantes.

\begin{enumerate}
\def\labelenumi{\arabic{enumi}.}
\item
  \(F^-\) est croissante.
\item
  \(\forall\, x \in \mathbb{R}\)~: \(F^- \circ F (x) \leq x\).
\item
  \(\forall\, u \in\, ]0,1[\)~: \(F \circ F^-(u) \geq u\) avec égalité
  si \(u \in F(\mathbb{R})\).
\item
  \(\forall\, (u,x) \in\, ]0,1[\, \times \mathbb{R}\)~:
  \(\left\{F(x) \geq u\right\} \Leftrightarrow \left\{x \geq F^-(u)\right\}\)
  et
  \(\left\{F(x) < u\right\} \Rightarrow \left\{x \leq F^-(u)\right\}\).
\end{enumerate}
\end{proposition}

\begin{preuve}
Pour tout \(u \in\, ]0,1[\) on note
\(F^{-1}\left([u,1]\right) := \left\{x \in \mathbb{R}: F(x) \geq u\right\}\)
l'image réciproque de \([u,1]\) par \(F\).

\begin{enumerate}
\def\labelenumi{\arabic{enumi}.}
\item
  Soit \((u,v) \in\, ]0,1[^2\). Si \(u < v\) alors
  \(F^{-1}\left([v,1]\right) \subset F^{-1}\left([u,1]\right)\) d'où
  \(F^-(u) = \inf F^{-1}\left([u,1]\right) \leq \inf F^{-1}\left([v,1]\right) = F^-(v)\).
  La fonction \(F^-\) est donc bien croissante.
\item
  Soit \(x \in \mathbb{R}\), alors
  \(F^- \circ F(x) = \inf \bigl\{ z \in \mathbb{R}: F(z) \geq F(x) \bigr\} = \inf F^{-1}\left(\left[F(x),1\right]\right)\).
  Comme \(F\) est croissante sur \(\mathbb{R}\) on a
  \([x,+\infty[\, \subseteq F^{-1}\left(\left[F(x),1\right]\right)\)
  donc
  \(F^- \circ F(x) = \inf F^{-1}\left(\left[F(x),1\right]\right) \leq \inf [x,+\infty[\, = x\).
\item
  Soit \(u \in\, ]0,1[\).
\end{enumerate}

\begin{itemize}
\tightlist
\item
  Puisque \(F^{-1}\left([u,1]\right)\) est nécessairement non vide, il
  existe une suite décroissante
  \((x_n)_{n\in\mathbb{N}} \subseteq F^{-1}\left([u,1]\right)\)
  convergeant vers \(\inf F^{-1}\left([u,1]\right) = F^-(u)\). La
  croissance de \(F\) implique que la suite
  \(\bigl(F(x_n)\bigr)_{n\in\mathbb{N}}\) est elle aussi décroissante,
  minorée par \(u\) car
  \((x_n)_{n\in\mathbb{N}} \subseteq F^{-1}\left([u,1]\right)\) donc
  convergente. Sa limite est de même supérieure ou égale à \(u\). Comme
  \(F\) est continue à droite, cette dernière n'est autre que
  \[\lim_{n \rightarrow +\infty} F(x_n) = F\left(\lim_{n\rightarrow+\infty} x_n\right) = F \circ F^-(u).\]
  Nous avons donc bien \(F\circ F^-(u) \geq u\).
\item
  Supposons maintenant que \(u \in F(\mathbb{R})\). Alors
  \(F^{-1}\left(\{u\}\right) := \left\{x \in \mathbb{R}: F(x) = u\right\} \neq \varnothing\).
  Il existe donc une suite décroissante
  \((x_n)_{n\in\mathbb{N}} \subseteq F^{-1}\left(\{u\}\right)\)
  convergeant vers \(\inf F^{-1}\left(\{u\}\right) = F^-(u)\) par
  croissance de \(F\). Comme \(F\) est continue à droite en tout point
  de \(\mathbb{R}\) et \(F(x_n) = u\) pour tout \(n \in \mathbb{N}\), on
  a bien
  \[u = \lim_{n \rightarrow +\infty} F(x_n) = F\left(\lim_{n\rightarrow+\infty} x_n\right) = F \circ F^-(u).\]
\end{itemize}

\begin{enumerate}
\def\labelenumi{\arabic{enumi}.}
\setcounter{enumi}{3}
\tightlist
\item
  Soit \((u,x) \in\, ]0,1[\, \times \mathbb{R}\).
\end{enumerate}

\begin{itemize}
\tightlist
\item
  \textbf{Equivalence.} Supposons \(F(x) \geq u\). On a
  \(F^- \circ F(x) \geq F^-(u)\) par croissance de \(F^-\) (propriété
  1.) et \(F^- \circ F(x) \leq x\) d'après la propriété 2.
  Réciproquement, supposons que \(x \geq F^-(u)\). Alors par croissance
  de \(F\) on a \(F(x) \geq F \circ F^-(u)\) puis
  \(F \circ F^-(u) \geq u\) d'après la propriété 3.
\item
  \textbf{Implication.} Supposons \(F(x) < u\). Tout
  \(z \in F^{-1}\left([u,1]\right)\) vérifie \(F(z) \geq u > F(x)\).
  Comme \(F\) est croissante sur \(\mathbb{R}\) on a donc \(z \geq x\),
  ce qui implique \(x \leq \inf F^{-1}\left([u,1]\right) = F^-(y)\).
\end{itemize}

\hfill$\blacksquare$
\end{preuve}

\begin{figure}
\centering
\includegraphics{RecGen.tex.pdf}
\caption{Réciproque généralisée d'une fonction de
répartition}\label{V-fig-RecGen.tex}
\end{figure}

Nous pouvons maintenant établir la preuve du
\protect\hyperref[V-invgen]{théorème de la méthode
d'inversion (p.~\pageref*{V-invgen})}.

\begin{preuve}
Soient \(x \in \mathbb{R}\) et \((\Omega,\mathcal{A},\mathbb{P})\)
l'espace probabilisé sur lequel sont définies \(U\) et \(X\). D'après la
propriété 4 ci-dessus, on a
\(\left\{\omega \in \Omega : F_X^-(U(\omega)) \leq x \right\} = \left\{\omega \in \Omega : U(\omega) \leq F(x) \right\}\),
d'où
\[\mathbb{P}\left( F^-_X(U) \leq x \right) = \mathbb{P}\left(U \leq F_X(x)\right) = F_X(x).\]

\hfill$\blacksquare$
\end{preuve}

\chapter{Références}\label{V-ruxe9fuxe9rences}

\hypertarget{refs}{}
\begin{cslreferences}
\leavevmode\hypertarget{refV-roca}{}%
C. P. Robert, and G. Casella. 2004. \emph{Monte-Carlo Statistical
Methods}. 2nd edition, Berlin: Springer.

\leavevmode\hypertarget{refV-BoxMuller}{}%
G.E.P. Box, and M.E. Muller. 1958. ``A Note on the Generation of Random
Normal Deviates.'' \emph{Ann. Math. Stat.} 29: 610--11.

\leavevmode\hypertarget{refV-Lehmer}{}%
Lehmer, Derrick H. 1951. ``Mathematical Methods in Large-Scale Computing
Units.'' \emph{Annu. Comput. Lab. Harvard Univ.} 26: 141--46.

\leavevmode\hypertarget{refV-Ziggurat}{}%
G. Marsaglia, and W. W. Tsang. 2000. ``The Ziggurat Method for
Generating Random Variables.'' \emph{Journal of Statistical Software}
5 (8): 1--7. \url{https://doi.org/10.18637/jss.v005.i08}.

\leavevmode\hypertarget{refV-ripley}{}%
Ripley, Brian D. 1987. \emph{Stochastic Simulation}. New York: Wiley.

\leavevmode\hypertarget{refV-vonNeumann}{}%
Von Neumann, John. 1951. ``Various Techniques Used in Connection with
Random Digits.'' \emph{U.S. Nat. Bur. Stand. Math. Ser.} 12: 36--38.
\end{cslreferences}


%% Corrections regroupées en fin de poly. Pour imprimer sans les corrections : \correctionsfalse
%% dans le préambule ; les renvois "(Solution p. x.)" des énoncés disparaissent alors aussi.
\ifcorrections
\part{Corrections des exercices}
\chapter{Solutions}

\section{Exercices essentiels}\label{I-exercices-essentiels}

\runinhead{Propriétés
élémentaires}\phantomsection\addcontentsline{toc}{subsubsection}{Propriétés
élémentaires}\label{I-answer-propelem}

\begin{enumerate}
\def\labelenumi{\arabic{enumi}.}
\item
  \(\mathbb{P}(A) \in [0,1]\) est donné par la définition. Par ailleurs,
  on a \(\Omega = A \cup A^c\) et \(A\cap A^c = \varnothing\). En
  considérant la suite \(A_1 = A\), \(A_2 = A^c\) et
  \(A_n = \varnothing\) pour \(n > 2\), le point 2. de la définition
  nous donne
  \(1 = \mathbb{P}(\Omega) = \mathbb{P}(A) + \mathbb{P}(A^c)\).
\item
  Puisque \(A \subset B\), on a l'union disjointe
  \(B = (B\cap A) \cup (B\cap A^c) = A \cup (B\cap A^c)\). Comme
  \(\mathbb{P}(B\cap A^c) \geq 0\) par définition, on en déduit le
  résultat.
\item
  On a \((A\cup B) = B \cup (A\cap B^c)\), avec
  \(B \cap (A\cap B^c) =\varnothing\) ainsi que
  \(A = (A\cap B^c) \cup (A\cap B)\), où
  \((A\cap B^c) \cap (A\cap B) = \varnothing\). On en déduit
  \(\mathbb{P}(A\cup B) = \mathbb{P}(B) + \mathbb{P}(A\cap B^c)\) et
  \(\mathbb{P}(A) = \mathbb{P}(A\cap B^c) + \mathbb{P}(A\cap B)\), d'où
  le résultat.
\item
  Par récurrence. On a déjà vu que
  \(\mathbb{P}(A\cup B) \leq \mathbb{P}(A) + \mathbb{P}(B)\). Notant
  \(B = \bigcup_{i=1}^n A_i\), alors
  \(\mathbb{P}(\bigcup_{i=1}^{n+1} A_i) \leq \mathbb{P}(B) + \mathbb{P}(A_{n+1})\),
  et l'hypothèse de récurrence nous indique
  \(\mathbb{P}\left(\bigcup_{i=1}^n A_i\right) \leq \sum_{i=1}^n \mathbb{P}(A_i)\),
  d'où le résultat.
\item
  Par récurrence. On a déjà vu que
  \(\mathbb{P}(A \cup B ) = \mathbb{P}(A) + \mathbb{P}(B) - \mathbb{P}(A \cap B)\).
  Notant \(B = \bigcup_{i=1}^n A_i\), alors \begin{align*}
      \mathbb{P}(\bigcup_{i=1}^{n+1} A_i) =& \mathbb{P}(B \cup A_{n+1}) \\
                                  =& \mathbb{P}(B) + \mathbb{P}(A_{n+1}) - \mathbb{P}(B \cap A_{n+1}) \\
                                  =& \sum_{i=1}^n \mathbb{P}(A_i) - \sum_{1 \leq i < j \leq n} \mathbb{P}(A_i \cap A_j) + \ldots + (-1)^n \mathbb{P}\left(\bigcap_{i=1}^n A_i\right) \\
                                   &+ \mathbb{P}(A_{n+1})\\
                                   &-\sum_{i=1}^n \mathbb{P}(A_i\cap A_{n+1}) + \sum_{1 \leq i < j \leq n} \mathbb{P}(A_i \cap A_j) + \ldots + (-1)^{n+1} \mathbb{P}\left(\bigcap_{i=1}^n A_i\cap A_{n+1}\right)\\
                                   =& \sum_{i=1}^{n+1} \mathbb{P}(A_i) - \sum_{1 \leq i < j \leq n +1 } \mathbb{P}(A_i \cap A_j) + \ldots + (-1)^{n+1} \mathbb{P}\left(\bigcap_{i=1}^{n+1} A_i\right).
  \end{align*}
\end{enumerate}

\runinhead{Continuité
monotone}\phantomsection\addcontentsline{toc}{subsubsection}{Continuité
monotone}\label{I-answer-contmon}
On définit une suite \((B_n)_{n\in \mathbb{N}^\ast}\) telle que
\(B_1 = A_1\) et \(B_n = A_n \setminus A_{n-1}\), pour \(n\geq 2\). Les
\(B_n\) ainsi définis sont deux à deux disjoints et
\(\cup B_n = \cup A_n\), donc~:
\[ \mathbb{P}\left(\bigcup_{n\in\mathbb{N}^\ast} A_n\right) = \sum_n \mathbb{P}(B_n) = \lim_{n \to \infty} \sum_{p=1}^n \mathbb{P}(B_p) = \lim_{n \to \infty} \mathbb{P}(A_n).\]

\runinhead{Une définition alternative de la
probabilité}\phantomsection\addcontentsline{toc}{subsubsection}{Une
définition alternative de la probabilité}\label{I-answer-altdef}
Soit \(A_n\) une suite d'éléments de \(\mathcal{A}\) deux-à-deux
disjoints. On définit \(B_n = \cup_{p \leq n} A_n\) et
\(B = \cup_n A_n\). Comme \(\mathbb{P}\) est additive, on a
\(\mathbb{P}(B_n) = \sum_{p \leq n} \mathbb{P}(A_n)\) qui croît vers
\(\sum_n \mathbb{P}(A_n)\) et aussi vers \(\mathbb{P}(B)\) d'après le
point 3.

En considérant les résultats de ces deux exercices, on obtient une
définition alternative de la probabilité en substituant la
\protect\hyperref[I-contmon]{continuité
monotone (p.~\pageref*{I-contmon})} et l'additivité à la propriété de
\protect\hyperref[I-defproba]{\(\sigma\)-additivité (p.~\pageref*{I-defproba})}.

\runinhead{Inégalité}\phantomsection\addcontentsline{toc}{subsubsection}{Inégalité}\label{I-answer-boolinf}
On a déjà vu le cas fini avec \protect\hyperref[I-propelem]{l'inégalité
de Boole (p.~\pageref*{I-propelem})}. On pose
\(B_n = \cup_{i=1}^n A_i\), qui croît vers l'ensemble
\(C = \cup_n A_n\). D'après \protect\hyperref[I-propelem]{l'inégalité de
Boole (p.~\pageref*{I-propelem})}, on a
\[\mathbb{P}(B_n) \leq \sum_{i=1}^n \mathbb{P}(A_i)\] Mais
\(\mathbb{P}(B_n) \to_{n \to \infty} \mathbb{P}(C)\) d'après le
\protect\hyperref[I-contmon]{théorème de continuité
monotone (p.~\pageref*{I-contmon})}, tandis que
\(\sum_{i=1}^n \mathbb{P}(A_i) \to_{n \to \infty} \sum_{n\in\mathbb{N}^\ast} \mathbb{P}(A_n)\).
En passant à la limite, on obtient donc le résultat.

\runinhead{Test}\phantomsection\addcontentsline{toc}{subsubsection}{Test}\label{I-answer-test}
On considère les événements \(A\) ``l'individu est séropositif'', et
\(B\) ``le test de détection donne un résultat positif''. Les données
ci-dessus nous indiquent \(P(A) = 10^{-4}\) d'où
\(\mathbb{P}(A^c) = 0,9999\) , \(\mathbb{P}(B|A) = 0,99\) et
\(\mathbb{P}(B|A^c) = 0,001\). Nous trouvons alors \begin{align*}
\mathbb{P}(A|B) &= \frac{\mathbb{P}(A \cap B)}{\mathbb{P}(B)}\\
        &= \frac{\mathbb{P}(B |A ) \mathbb{P}(A)}{\mathbb{P}(B |A)\mathbb{P}(A) + \mathbb{P}(B|A^c)\mathbb{P}(A^c)}\\
        &= \frac{0,99 \times 10^{-4}}{0,99 \times 10^{-4} + 0,001 \times 0,9999}\\
        &\approx 0,09.
\end{align*} On remarque que contrairement à l'intuition, cette
probabilité est faible.

\runinhead{Démonstration}\phantomsection\addcontentsline{toc}{subsubsection}{Démonstration}\label{I-answer-indep}
Pour \((A^c,B)\), si \(\mathbb{P}(B) = 0\), alors
\(\mathbb{P}(A^c \cap B) = 0 = \mathbb{P}(B)\mathbb{P}(A^c)\). Sinon,
\(\mathbb{P}(A^c \cap B) = \mathbb{P}(B)\mathbb{P}(A^c|B) = \mathbb{P}(B)(1-\mathbb{P}(A|B)) = \mathbb{P}(B)(1-\mathbb{P}(A))=\mathbb{P}(B)\mathbb{P}(A^c)\).

Pour \((A,B^c)\), on inverse les rôles de \(A\) et \(B\) ci-dessus.

Pour \((A^c,B^c)\), on remarque que \(A^c \cap B^c = (A \cup B)^c\),
d'où
\(\mathbb{P}(A^c \cap B^c) = 1-\mathbb{P}(A \cup B) = 1-\mathbb{P}(A) - \mathbb{P}(B) +\mathbb{P}(A)\mathbb{P}(B) = (1-\mathbb{P}(A))(1-\mathbb{P}(B)) = \mathbb{P}(A^c)\mathbb{P}(B^c)\)

\runinhead{Auto-indépendant\,?}\phantomsection\addcontentsline{toc}{subsubsection}{Auto-indépendant\,?}\label{I-answer-autoindep}
\(A\) est indépendant de lui-même ssi
\(\mathbb{P}(A) = \mathbb{P}(A \cap A) = \mathbb{P}(A)\mathbb{P}(A)\),
autrement dit si \(\mathbb{P}(A) = 0\) ou si \(\mathbb{P}(A) = 1\).

\runinhead{Durée de
fonctionnement}\phantomsection\addcontentsline{toc}{subsubsection}{Durée
de fonctionnement}\label{I-answer-ex.expo}

\begin{enumerate}
\def\labelenumi{\arabic{enumi}.}
\tightlist
\item
  \(\mathbb{P}(X \in [50,150]) = \int_{50}^{150} \frac{1}{100}\exp\left(-\frac{x}{100}\right) dx = \exp(-1/2)-\exp(-3/2) \approx 0,38.\)
\item
  \(\mathbb{P}(X \leq 100) = \int_{0}^{100} \frac{1}{100}\exp\left(-\frac{x}{100}\right) dx = 1-e^{-1} \approx 0,63.\)
\item
  \(\mathbb{P}(X = 100) = 0\)
\end{enumerate}

\runinhead{Précipitations}\phantomsection\addcontentsline{toc}{subsubsection}{Précipitations}\label{I-answer-ex.pluie}

\begin{enumerate}
\def\labelenumi{\arabic{enumi}.}
\item
  \[F(x) = \left\{ \begin{array}{ll}
  0 & \text{si } x < 0 \\
  1/2 & \text{si } x = 0 \\ 
  1/2 + 1/2 \int_0^x \frac{1}{10}\exp\left(-\frac{x}{10}\right) & \text{si } x > 0 
  \end{array} \right.\]
\item
  \(F\) n'est pas continue, \(X\) n'admet donc pas de densité
  (\protect\hyperref[I-fdr-va-dens]{voir la
  proposition (p.~\pageref*{I-fdr-va-dens})}).
\end{enumerate}

\section{Indépendance et
conditionnement}\label{I-induxe9pendance-et-conditionnement-1}

\runinhead{Question
1}\phantomsection\addcontentsline{toc}{subsubsection}{Question
1}\label{I-answer-ic-1}
On a
\(1 = \mathbb{P}(A\cup B) = \mathbb{P}(A) + \mathbb{P}(B) - \mathbb{P}(A\cap B)\),
donc \(1 = \mathbb{P}(A) + \mathbb{P}(B) - \mathbb{P}(A)\mathbb{P}(B)\)
si et seulement si \(A\) et \(B\) sont indépendants. Cette égalité se
réécrit \((1-\mathbb{P}(A))(1-\mathbb{P}(B))=0\), qui est vérifiée ssi
\(A\) ou \(B\) est presque sûr.

\runinhead{Question
2}\phantomsection\addcontentsline{toc}{subsubsection}{Question
2}\label{I-answer-ic-2}
Par définition, on a \(A\) et \(B\) indépendants relativement à
\(\mathbb{P}(\,\,|C)\) ssi
\[\mathbb{P}(A\cap B |C) = \mathbb{P}(A|C)\mathbb{P}(B|C) = \mathbb{P}(A)\mathbb{P}(B),\]
car \(A\) et \(B\) sont chacun indépendants de \(C\), d'où
\[\frac{\mathbb{P}(A\cap B \cap C)}{\mathbb{P}(C)} = \mathbb{P}(A)\mathbb{P}(B)\]
par
\protect\hyperref[I-defprobacond]{définition (p.~\pageref*{I-defprobacond})}\,;
et comme \(A\), \(B\) et \(C\) sont indépendants 2 à 2,
\[\mathbb{P}(A\cap B \cap C) = \mathbb{P}(A)\mathbb{P}(B)\mathbb{P}(C)\]
équivaut à \(A\), \(B\) et \(C\) sont mutuellement indépendants.

\runinhead{Question
3}\phantomsection\addcontentsline{toc}{subsubsection}{Question
3}\label{I-answer-ic-3}
On a
\(\mathbb{P}(A^c\cap B^c) = \mathbb{P}((A \cup B)^c) = 1 - \mathbb{P}(A) - \mathbb{P}(B) + \mathbb{P}(A\cap B)\),
donc
\[\mathbb{P}(A\cap B)\mathbb{P}(A^c \cap B^c) = \mathbb{P}(A\cap B) - \mathbb{P}(A\cap B)(\mathbb{P}(A) + \mathbb{P}(B) - \mathbb{P}(A\cap B))\]
Par ailleurs, on a \[\left\{ \begin{array}{ll}
        \mathbb{P}(A \cap B^c)  & = \mathbb{P}(A \setminus (A\cap B)) = \mathbb{P}(A) - \mathbb{P}(A\cap B) \\
        \mathbb{P}(A^c \cap B)  &=  \mathbb{P}(B \setminus (A\cap B)) = \mathbb{P}(B) - \mathbb{P}(A\cap B)
        \end{array}
        \right.\] Donc
\[ \mathbb{P}(A \cap B^c)\mathbb{P}(A^c\cap B) = \mathbb{P}(A)\mathbb{P}(B) -\mathbb{P}(A\cap B)(\mathbb{P}(A) + \mathbb{P}(B)- \mathbb{P}(A\cap B))\]
On a donc bien l'égalité ssi
\(\mathbb{P}(A)\mathbb{P}(B) = \mathbb{P}(A\cap B)\)

\section{Problème de Monty Hall}\label{I-answer-mh}

Considérons les événements \(A\) = `la porte choisie en premier est la
bonne' et \(B\) = `la dernière porte est la bonne'.

On a évidemment \(\mathbb{P}(A) = \frac{1}{3}\), puisque le choix
initial se fait uniformément parmi les 3 portes.

Pour \(B\), on doit tenir compte de l'information donnée par le
présentateur. On a par \protect\hyperref[I-formprobatot]{la formule des
probabilités totales (p.~\pageref*{I-formprobatot})}
\[\mathbb{P}(B) = \mathbb{P}(B|A)\mathbb{P}(A) + \mathbb{P}(B|A^c)\mathbb{P}(A^c) = 0.\frac{1}{3} + 1.\frac{2}{3} = \frac{2}{3}\]
En effet \(\mathbb{P}(B|A^c) = 1\) car si la porte choisie initialement
n'est pas la bonne, c'est nécessairement la dernière.

Il convient donc de changer son choix compte tenu de la nouvelle
information. Pour se convaincre du résultat, on peut refaire le calcul
avec disons 100 portes et le présentateur qui ouvre 98 autres portes
après le choix du candidat.

\section{Mesure de Probabilité}\label{I-mesure-de-probabilituxe9-1}

\runinhead{Question
1}\phantomsection\addcontentsline{toc}{subsubsection}{Question
1}\label{I-answer-mes-proba1}
On vérifie facilement les axiomes~:

\begin{itemize}
\tightlist
\item
  dans \([0,1]\)
\item
  \(Q(\Omega)= 1\)
\item
  la \(\sigma\)-additivité s'obtient en décomposant l'intégrale sur des
  \(A_n\) disjoints
\end{itemize}

\runinhead{Question
2}\phantomsection\addcontentsline{toc}{subsubsection}{Question
2}\label{I-answer-mes-proba2}
Si \(\mathbb{P}(A)=0\), alors
\(Q(A) = \int_{\Omega} X(\omega) 1_A(\omega)\mathbb{P}(d\omega) = \int_{A} X(\omega)\mathbb{P}(d\omega) = 0\)
car on intègre \(X\) sur \(A\) \(\mathbb{P}\)-négligeable.

contre-exemple~:
\((\Omega,\mathcal{A},\mathbb{P}) = (\mathbb{R},\mathcal{B}(\mathbb{R}),\mathbb{P})\)
avec \(\mathbb{P}\) symétrique et qui charge tout \(\mathbb{R}\) (la loi
normale centrée réduite par exemple) et \(X \sim \mathcal{B}_{1/2}\)
(loi de Bernoulli), définie telle que \(X(\omega) = 0\) si
\(\omega \in ]-\infty,0] = A\). On a alors \(\mathbb{P}(A) = 1/2\),
\(Q(A) = 0\) et on vérifie bien les hypothèses.

\runinhead{Question
3}\phantomsection\addcontentsline{toc}{subsubsection}{Question
3}\label{I-answer-mes-proba3}
\(Y\) positive ou bornée donc \(\mathbb{E}_Q(Y)\) est bien définie.

\(\mathbb{E}_Q(Y) = \int_{\Omega} Y(\omega) Q(d\omega) = \int_{\Omega} Y(\omega) X(\omega) \mathbb{P}(d\omega) = \mathbb{E}_\mathbb{P}(XY)\)

\runinhead{Question
4}\phantomsection\addcontentsline{toc}{subsubsection}{Question
4}\label{I-answer-mes-proba4}

\begin{enumerate}
\def\labelenumi{\roman{enumi})}
\item
  \(\mathbb{P}(X > 0) = 1\), donc \(\exists N \in \mathcal{A}\) tq
  \(\mathbb{P}(N) = 0\) et \(\forall \, \omega \in N^c,~ X(\omega) > 0\)

  Par conséquent, d'après la q.3

  \[\mathbb{E}_Q\left(\frac{1}{X}\right) = \int_{N^c} \frac{1}{X(\omega)}X(\omega) \mathbb{P}(d\omega) = 1\]
\item
  Soit \(A \in \mathcal{A}\), on a
  \(R(A) = \mathbb{E}_Q(\frac{1}{X}1_A) = \int_{A \cap N^c} \frac{1}{X(\omega)}X(\omega) \mathbb{P}(d\omega) = \mathbb{P}(A)\)
\item
  On a \(\forall \omega \in N^c, ~ X(\omega) > 0\) et donc
  \(1_{A \cap N^c}(\omega)\,X(\omega) \geq 0\). Par conséquent, si
  \(Q(A\cap N^c) = 0\), alors
  \(\mathbb{E}_Q(A\cap N^c) = \int_{A \cap N^c} X(\omega) \mathbb{P}(d\omega) = 0\)
  ce qui n'a lieu que si \(\mathbb{P}(A\cap N^c) = 0\).
\end{enumerate}

Enfin \(A \cap N \subset N\) et \(\mathbb{P}(N) = 0\) d'où
\(\mathbb{P}(A \cap N) = 0\)

\section{Queues de distributions}\label{I-queues-de-distributions-1}

\runinhead{Etude en
\(+\infty\)}\phantomsection\addcontentsline{toc}{subsubsection}{Etude en
\(+\infty\)}\label{I-answer-dlfdr-pi}

\begin{enumerate}
\def\labelenumi{\arabic{enumi}.}
\tightlist
\item
  Soit \(u>0\). Comme \(h\) est intégrable, on peut écrire
  \begin{align*}
  \int_{-\infty}^{+\infty} h(x)\,dx = \int_{-\infty}^{+\infty} x\,f(x)\,dx & = \int_{-\infty}^{u} x\,f(x)\,dx + \int_{u}^{+\infty} x\,f(x)\,dx\\
  &\geq \int_{-\infty}^{u} x\,f(x)\,dx + u\,\int_{u}^{+\infty} f(x)\,dx\\
  &= \int_{-\infty}^{u} h(x)\,dx + u\,\left(1- F(u)\right).
  \end{align*}
\end{enumerate}

\begin{enumerate}
\def\labelenumi{\arabic{enumi}.}
\setcounter{enumi}{1}
\tightlist
\item
  Soit toujours \(u>0\). Toute fonction de répartition étant à valeurs
  dans \([0,1]\), on a \(u\,\left(1-F(u)\right) \geq 0\). En utilisant
  la question précédente on obtient donc l'encadrement \begin{align*}
  0 \leq u\,\left(1-F(u)\right) & \leq \int_{-\infty}^{+\infty} h(x)\,dx - \int_{-\infty}^u h(x)\,dx\\
  &= \int_{u}^{+\infty} h(x)\,dx = \int_{-\infty}^{+\infty} h(x)\,1_{[u,+\infty]}(x)\,dx.
  \end{align*} Or pour tout \(x\in \mathbb{R}\) on a
  \(\left|h(x)\,1_{[u,+\infty]}(x)\right| \leq \left|h(x)\right|\), qui
  est intégrable. On peut donc appliquer le
  \href{Calcul\%20Intégral\%20II.pdf\%20\#TCD}{théorème de convergence
  dominée}, qui nous donne~: \begin{align*}
  \lim_{u\to+\infty} u\,(1-F(u)) & = \int_{-\infty}^{+\infty} h(x)\,\lim_{u\to+\infty} 1_{[u,+\infty]}(x)\,dx\\
  & = \int_{-\infty}^{+\infty} h(x)\,1_{\varnothing}(x)\,dx = 0.
  \end{align*} On en conclut que \(1-F(x) = o\left(\dfrac{1}{x}\right)\)
  lorsque \(x\to+\infty\).
\end{enumerate}

\runinhead{Etude en
\(-\infty\)}\phantomsection\addcontentsline{toc}{subsubsection}{Etude en
\(-\infty\)}\label{I-answer-dlfdr-mi}

\begin{enumerate}
\def\labelenumi{\arabic{enumi}.}
\setcounter{enumi}{2}
\tightlist
\item
  Prenons maintenant \(u<0\). Comme \(h\) est intégrable, on peut écrire
  \begin{align*}
  \int_{-\infty}^{+\infty} h(x)\,dx = \int_{-\infty}^{+\infty} x\,f(x)\,dx & = \int_{-\infty}^{u} x\,f(x)\,dx + \int_{u}^{+\infty} x\,f(x)\,dx\\
  &\leq u\, \int_{-\infty}^{u} f(x)\,dx + \int_{u}^{+\infty} x\,f(x)\,dx\\
  &= u\,F(u) + \int_{u}^{+\infty} h(x)\,dx.
  \end{align*} Or, comme \(F\) prend ses valeurs dans \([0,1]\),
  \(u\,F(u) \leq 0\). On obtient donc \begin{align*}
  0 \geq u\,F(u) & \geq \int_{-\infty}^{+\infty} h(x)\,dx - \int_{u}^{+\infty} h(x)\,dx\\
  &= \int_{-\infty}^{u} h(x)\,dx = \int_{-\infty}^{+\infty} h(x)\,1_{[-\infty,u]}(x)\,dx.
  \end{align*} Comme précédemment, cette dernière fonction sous
  l'intégrale est dominée par \(-\left|h\right|\) (à gauche) et
  \(\left|h\right|\) (à droite), toutes deux intégrables, donc d'après
  le \href{Calcul\%20Intégral\%20II.pdf\%20\#TCD}{théorème de
  convergence dominée} \begin{align*}
  \lim_{u\to-\infty} u\,F(u) & = \int_{-\infty}^{+\infty} h(x)\,\lim_{u\to-\infty} 1_{[-\infty,u]}(x)\,dx\\
  & = \int_{-\infty}^{+\infty} h(x)\,1_{\varnothing}(x)\,dx = 0.
  \end{align*} On a donc bien \(F(x) = o\left(\dfrac{1}{x}\right)\)
  lorsque \(x\to-\infty\).
\end{enumerate}

\runinhead{Interprétation}\phantomsection\addcontentsline{toc}{subsubsection}{Interprétation}\label{I-answer-dlfdr-interpret}

\begin{enumerate}
\def\labelenumi{\arabic{enumi}.}
\setcounter{enumi}{3}
\tightlist
\item
  L'hypothèse que \(h\) est intégrable nous donne une indication sur la
  vitesse de convergence de \(1-F(x)\) et \(F(x)\) lorsque \(x\) tend
  respectivement vers \(+\infty\) et \(-\infty\)~: elles convergent vers
  \(0\) au moins aussi vite que \(\dfrac{1}{x}\). On dit que les queues
  de la distribution sont ``relativement'' fines.
\end{enumerate}

Pour mieux comprendre l'impact de ce résultat, prenons un exemple
pratique. Considérons un phénomène aléatoire comme la concentration en
polluant dans l'air, modélisé par une loi de probabilité sur
\(\mathbb{R}\) de densité \(f\) nulle sur \(\mathbb{R}_-^\ast\). Etre
capable de déterminer la probabilité que la concentration dépasse un
seuil critique \(x\) est alors très important pour les organismes de
contrôle de la qualité de l'air. Or cet événement correspond
mathématiquement à l'événement \([x,+\infty]\). Dans ce cas,
\(1 - F(x)\) nous donne cette probabilité selon le modèle considéré. Le
résultat nous indique que si l'on choisit un modèle tel que \(h\) est
intégrable, alors cette probabilité décroît plus rapidement que
\(\dfrac{1}{x}\) lorsque le seuil \(x\) augmente. En d'autres termes,
les événements extrêmes (quand \(x\) est grand) restent ``assez'' rares.
L'impact du choix de modèle n'est donc pas négligeable~: si on prend une
distribution à queues trop fines alors que les pics de pollution ne sont
en réalité pas si rares que ça, \(1-F(x)\) risque de sous-estimer le
risque réel de dépassement du seuil\,!

\runinhead{Etude de la
réciproque}\phantomsection\addcontentsline{toc}{subsubsection}{Etude de
la réciproque}\label{I-answer-dlfdr-reciproque}

\begin{enumerate}
\def\labelenumi{\arabic{enumi}.}
\setcounter{enumi}{4}
\item
  La fonction \(F\) considérée est continue et dérivable par morceaux
  sur \(\mathbb{R}\)\,; le seul point où \(F\) n'est pas dérivable est
  \(e\) (le taux d'accroissement n'a pas les mêmes limites à gauche et à
  droite). Elle possède donc \protect\hyperref[I-fdr-va-dens]{une
  densité (p.~\pageref*{I-fdr-va-dens})} qui s'obtient en dérivant
  chaque morceau~: pour tout \(x\in\mathbb{R}\),
  \[f(x) = \left|\begin{array}{ll} e\,\dfrac{\ln(x) + 1}{x^2\,\ln(x)^2} & \text{si } x \geq e,\\ 0 & \text{sinon.}\end{array}\right.\]
\item
  On remarque que pour tout \(x\geq e\),
  \(x\,\left(1-F(x)\right) = \dfrac{e}{\ln(x)} \xrightarrow[x\to+\infty]{} 0\),
  donc le développement asymptotique \((1)\) est bien vérifié. En outre,
  comme \(x\,F(x) = 0\) pour tout \(x < e\), le développement
  asymptotique \((2)\) est aussi bien respecté.
\item
  Pour tout \(x\in\mathbb{R}\),
  \[ h(x) = \left|h(x)\right| = \left|\begin{array}{ll} e\,\dfrac{\ln(x) + 1}{x\,\ln(x)^2} & \text{si } x \geq e,\\ 0 & \text{sinon,} \end{array}\right.\]
  et on pose \(h(\pm\infty) = \lim\limits_{x\to\pm\infty}h(x) = 0\). Par
  conséquent, \(h\) est clairement intégrable sur \([-\infty,e]\), où
  elle est égale presque partout (partout sauf en \(e\)) à la fonction
  nulle. Pour savoir si \(h\) est intégrable sur toute la droite réelle
  il nous faut donc regarder si elle l'est sur \([e,+\infty]\). On
  remarque en outre que pour tout \(x\geq e\),
  \[h(x) \geq \dfrac{1}{x\,\ln(x)}, \] et le terme de droite n'est pas
  intégrable (en effet, une primitive de \(\frac1{x \ln(x)}\) est
  \(\ln(ln(x))\) qui tend vers \(+\infty\) en \(+\infty\)). Par
  conséquent, \(h\) n'est pas intégrable.
\item
  La fonction de répartition \(F\) considérée respecte bien les
  développements asymptotiques \((1)\) et \((2)\), mais la fonction
  \(h\) associée n'est pas intégrable. Nous avons donc montré que la
  réciproque de la propriété est fausse.
\end{enumerate}

\runinhead{Liens entre les développements
asymptotiques}\phantomsection\addcontentsline{toc}{subsubsection}{Liens
entre les développements asymptotiques}\label{I-answer-dlfdr-liensdl}

\begin{enumerate}
\def\labelenumi{\arabic{enumi}.}
\setcounter{enumi}{8}
\tightlist
\item
  On remarque que \(F\) satisfait \((2)\) mais pas \((1)\), tandis que
  \(G\) satisfait \((1)\) mais pas \((2)\). On en conclut que \((1)\)
  n'implique pas \((2)\) et inversement.
\end{enumerate}

\section{Densité et fonction de répartition d'une loi
Normale}\label{I-densituxe9-et-fonction-de-ruxe9partition-dune-loi-normale-1}

\runinhead{Propriétés
générales}\phantomsection\addcontentsline{toc}{subsubsection}{Propriétés
générales}\label{I-answer-fdrgaussprop}

\begin{enumerate}
\def\labelenumi{\arabic{enumi}.}
\tightlist
\item
  Faisons l'étude de \(f\).
\end{enumerate}

\begin{itemize}
\tightlist
\item
  \textbf{Domaine.} \(f : \mathbb{R}\to \mathbb{R}^+\).
\item
  \textbf{Parité.} On remarque que \(f\) est paire~:
  \(\forall\,x\in\mathbb{R}\), \(f(x) = f(-x)\).
\item
  \textbf{Limites aux bornes.}
  \(\lim\limits_{x\to-\infty} f(x) = \lim\limits_{x\to+\infty} f(x) = 0\).
  Ainsi, la parité est aussi vraie aux bornes.
\item
  \textbf{Dérivabilité.} \(f\) est continue sur \(\mathbb{R}\) et
  infiniment dérivable sur \(\mathbb{R}\) en tant que composée de
  fonctions infiniment dérivables. Pour tout \(x\in\mathbb{R}\) sa
  dérivée première s'écrit \(f^\prime(x) = -x\,f(x)\) et a dérivée
  seconde \(f^{\prime\prime}(x) =(x^2 - 1)\,f(x)\).
\item
  \textbf{Variations.} \(f^\prime\) est strictement positive (resp.
  négative) ssi \(-\infty < x < 0\) (resp. \(+\infty > x > 0\)). Elle
  est nulle en \(0\), \(+\infty\) et \(-\infty\). Ainsi, \(f\) est
  strictement croissante sur \(]-\infty,0[\), vaut
  \(\left(2\,\pi\right)^{-1/2}\) en \(0\), puis est strictement
  décroissante sur \(]0,+\infty[\).
\item
  \textbf{Convexité/Concavité.} \(f^{\prime\prime}\) est strictement
  positive (resp. négative) sur \(\mathbb{R}\backslash\,]-1,1[\) (resp.
  \(]-1,1[\)). \(f\) est donc convexe sur
  \(\mathbb{R}\backslash\,]-1,1[\) et concave sur \(]-1,1[\).
\end{itemize}

\begin{enumerate}
\def\labelenumi{\arabic{enumi}.}
\setcounter{enumi}{1}
\tightlist
\item
  Par définition, pour tout \(x\in\mathbb{R}\)
  \[ F(x) = \int_{-\infty}^x \dfrac{1}{\sqrt{2\,\pi}}\,\exp\left\{-\dfrac{u^2}{2}\right\}\,du.\]
  On pose \(F(+\infty) = 1\) et \(F(-\infty) = 0\). Lorsque \(x = 0\),
  comme \(f\) est paire et que son intégrale sur \([-\infty,+\infty]\)
  vaut \(1\) (c'est \protect\hyperref[I-defdens]{une
  densité (p.~\pageref*{I-defdens})}), son intégrale sur \([-\infty,0]\)
  (i.e.~\(F(0)\)) vaut \(\dfrac{1}{2}\). Lorsque \(x > 0\), en
  décomposant l'intégrale puis en utilisant un
  \href{Calcul\%20Intégral\%20I.pdf\%20\#changement-de-variable}{changement
  de variable} on obtient
  \[F(x) = \dfrac{1}{2} + \int_{0}^x \dfrac{1}{\sqrt{2\,\pi}}\,\exp\left\{-\dfrac{u^2}{2}\right\}\,du = \dfrac{1}{2} + \int_{0}^{\frac{x}{\sqrt{2}}} \dfrac{1}{\sqrt{\pi}}\,e^{-v^2}\,dv = \dfrac{1}{2} + \dfrac{1}{2}\,\text{erf}\left(\dfrac{x}{\sqrt{2}}\right).\]
  Lorsque \(x<0\), on peut procéder de la même manière avec \(1 - F(x)\)
  et obtenir la même égalité que ci-dessus.
\end{enumerate}

\runinhead{Encadrement de \(1 - F(x)\) pour tout
\(x > 0\)}\phantomsection\addcontentsline{toc}{subsubsection}{Encadrement
de \(1 - F(x)\) pour tout \(x > 0\)}\label{I-answer-fdrgaussenca}

\begin{enumerate}
\def\labelenumi{\arabic{enumi}.}
\setcounter{enumi}{2}
\item
  Soit \(x>0\). Par définition,
  \[1 - F(x) = \int_{-\infty}^{+\infty} f(u)\,du - \int_{-\infty}^x f(u)\,du = \int_{x}^{+\infty} \dfrac{1}{\sqrt{2\,\pi}}\,\exp\left\{-\dfrac{u^2}{2}\right\}\,du. \]
  Or pour tout \(u\geq x\) on a \(u/x \geq 1\), donc
  \[1 - F(x) \leq \dfrac{1}{\sqrt{2\,\pi}}\,\int_{x}^{+\infty} \dfrac{u}{x}\, \exp\left\{-\dfrac{u^2}{2}\right\}\,du = \dfrac{-1}{x}\,\int_{x}^{+\infty} f^\prime(u)\,du. \]
  Or \(f^\prime\) a pour primitive \(f\) donc on a
  \[1-F(x) \leq -\dfrac{f(+\infty) - f(x)}{x} = \dfrac{f(x)}{x}. \]
\item
  Soit \(x>0\). Par définition,
  \[\left(1 + \dfrac{1}{x^2}\right)\,\left(1-F(x)\right) = \int_x^{+\infty} \left(1 + \dfrac{1}{x^2} \right)\,f(u)\,du.\]
  Or pout tout \(u \geq x\) on a \(\dfrac{1}{u^2} \leq \dfrac{1}{x^2}\),
  d'où
  \[\left(1 + \dfrac{1}{x^2}\right)\,\left(1-F(x)\right) \geq \int_{x}^{+\infty} \left( 1 + \dfrac{1}{u^2} \right)\,f(u)\,du.\]
  Remarquons maintenant que d'après la question 1,
  \(\left(\dfrac{f(x)}{x} \right)^\prime = \dfrac{-x^2\,f(x) - f(x)}{x^2} = -\left(1+\dfrac{1}{x^2}\right)\,f(x)\).
  En d'autres termes, \(A : x\in\mathbb{R}\mapsto -\dfrac{f(x)}{x}\) est
  une primitive de
  \(x\in\mathbb{R}\mapsto \left(1+\dfrac{1}{x^2}\right)\,f(x)\). On
  obtient alors par
  \[\left(1 + \dfrac{1}{x^2}\right)\,\left(1-F(x)\right) \geq - \left(A(+\infty) - A(x)\right) = A(x) = \dfrac{f(x)}{x}.\]
  On en déduit directement l'inégalité de gauche~:
  \[ 1-F(x) \geq f(x)\,\dfrac{x^2}{(x^2+1)\,x} = f(x)\,\dfrac{x}{x^2+1}. \]
\end{enumerate}

\runinhead{Equivalent de \(1-F(x)\) quand
\(x\to+\infty\).}\phantomsection\addcontentsline{toc}{subsubsection}{Equivalent
de \(1-F(x)\) quand \(x\to+\infty\).}\label{I-answer-fdrgaussequi}

\begin{enumerate}
\def\labelenumi{\arabic{enumi}.}
\setcounter{enumi}{4}
\item
  Soit \(x>0\). L'encadrement que nous avons démontré peut être réécrit
  \[ \dfrac{x^2}{1+x^2} \leq \dfrac{1-F(x)}{\dfrac{f(x)}{x}} \leq 1, \]
  et comme \(\dfrac{x^2}{1+x^2} \to 1\) lorsque \(x\to+\infty\), on
  obtient que \(1-F(x) \sim \dfrac{f(x)}{x}\) quand \(x\to+\infty\).
\item
  Pour tout \(x \in [-\infty,+\infty]\) on a
  \(\left|h(x)\right| = \left|x\,f(x)\right| = |x|\,f(x)\) puisque
  \(f\geq 0\). Lorsque \(x \in [0,+\infty]\),
  \(\left|h(x)\right| = h(x)\), or d'après la question 1, \(h\) admet
  pour primitive \(-f\), qui vaut \(0\) en \(+\infty\) et \(-\infty\)
  puis \((2\,\pi)^{-1/2}\) en \(0\). Par conséquent, (d'après
  l'extension du \href{Calcul\%20Intégral\%20I.pdf\%20\#TFC}{théorème
  fondamental du calcul}) l'intégrale de \(\left|h\right|\) sur
  \([0,+\infty]\) existe, et vaut \((2\,\pi)^{-1/2}\). Or
  \(\left|h\right|\) est paire, son intégrale sur \([-\infty,0]\) existe
  donc aussi et vaut \((2\,\pi)^{-1/2}\). On en conclut que \(h\) est
  bien absolument intégrable sur \([-\infty,+\infty]\). En remarquant
  qu'elle est impaire, on obtient immédiatement que son intégrale vaut
  \(0\).
\end{enumerate}

On est donc bien dans le cadre de l'exercice sur les queues de
distributions. L'équivalent que nous avons trouvé implique bien que
\(1-F(x) = o\left(\dfrac{1}{x} \right)\) lorsque \(x\to+\infty\). Il
nous donne juste plus de précisions quant au comportement de \(1-F(x)\)
lorsque \(x\to+\infty\) dans le cas spécifique de la loi Normale centrée
réduite~: on connaît sa vitesse de convergence.

\runinhead{Loi Normale
générale.}\phantomsection\addcontentsline{toc}{subsubsection}{Loi
Normale générale.}\label{I-answer-fdrgaussgen}

\begin{enumerate}
\def\labelenumi{\arabic{enumi}.}
\setcounter{enumi}{6}
\tightlist
\item
  Soit \(x\in\mathbb{R}\). On remarque que
  \(g(x) = \dfrac{1}{\sigma}\,f\left(\dfrac{x-\mu}{\sigma} \right)\).
  Cela implique en particulier que
  \[G(x) =  \int_{-\infty}^x \dfrac{1}{\sigma}\,f\left(\dfrac{u-\mu}{\sigma} \right)\,du = \int_{-\infty}^{\frac{x-\mu}{\sigma}} f(u)\,du = F\left(\dfrac{x-\mu}{\sigma}\right).\]
  On obtient alors l'encadrement suivant, pour tout \(x>\mu\)~:
  \[ \dfrac{x-\mu}{1 + \dfrac{(x-\mu)^2}{\sigma^2}}\,g(x) \leq 1 - G(x) \leq \dfrac{\sigma^2\,g(x)}{x-\mu}.  \]
\end{enumerate}

On en déduit que
\(1-G(x) \sim \dfrac{\sigma^2\,g(x)}{x-\mu} \sim \sigma^2\,\dfrac{g(x)}{x}\)
lorsque \(x\to+\infty\).

\section{Loi de vie et de mort}\label{I-loi-de-vie-et-de-mort-1}

\runinhead{Question
1}\phantomsection\addcontentsline{toc}{subsubsection}{Question
1}\label{I-answer-viemort1}
Par définition, on a
\[ G_{t_0}(t) = \mathbb{P}(T>t+t_0|T>t_0) = \frac{\mathbb{P}(T>t+t_0,T>t)}{\mathbb{P}(T>t)} = \frac{G(t+t_0)}{G(t)}\]

\runinhead{Question
2}\phantomsection\addcontentsline{toc}{subsubsection}{Question
2}\label{I-answer-viemort2}
D'après la question précédente, \(G\) vérifie alors
\(G(t_0+t) = G(t)G(t_0)\) pour tous \(t, t_0 > 0\). Comme \(G\) est
décroissante, continue à droite et tend vers 0 à l'infini, on en déduit
que \(G(t) = e ^{-\theta t}\), pour un \(\theta > 0\). On reconnaît la
fonction de répartition complémentaire d'une loi exponentielle de
paramètre \(\theta\).

\runinhead{Question
3}\phantomsection\addcontentsline{toc}{subsubsection}{Question
3}\label{I-answer-viemort3}

\begin{align*}
D(t) &= \lim_{\Delta t \to 0} \frac{1}{\Delta t}\mathbb{P}(t < T \leq t + \Delta t | T >t)\\
     &= \lim_{\Delta t \to 0} \frac{1}{\Delta t}\frac{\mathbb{P}(t < T \leq t + \Delta t)}{\mathbb{P}(T >t)}\\
     &= \frac{1}{G(t)} \lim_{\Delta t \to 0} \frac{G(t) - G(t + \Delta t)}{\Delta t}\\
     &= -\frac{g(t)}{G(t)}
\end{align*} Ainsi, \(D(t) = \frac{d}{dt}(-\ln G(t))\) et comme
\(G(0) = 1\), alors on a pour \(t >0\)
\[G(t) = \exp\left(-\int_0^tD(s) ds\right)\] Si \(D\) est constant, on
retrouve une loi exponentielle.

\chapter{Solutions}

\section{Exercices essentiels}\label{II-exercices-essentiels}

\runinhead{Approximation par une
constante}\phantomsection\addcontentsline{toc}{subsubsection}{Approximation
par une constante}\label{II-answer-exo-approx}
On a
\[\mathbb{E}((X-a)^2) = \mathbb{E}((X-\mathbb{E}(X) + \mathbb{E}(X) - a)^2 = \mathbb{E}((X - \mathbb{E}(X))^2 ) + (\mathbb{E}(X) - a)^2\]
le terme de gauche ne dépend pas de \(a\) tandis que celui de droite
s'annule en \(a = \mathbb{E}(X)\).

\runinhead{Calcul d'une
covariance}\phantomsection\addcontentsline{toc}{subsubsection}{Calcul
d'une covariance}\label{II-answer-cov-unif}
On a \begin{align*}
\text{Cov}(U,V) &= \mathbb{E}((U-\mathbb{E}(U))(V-\mathbb{E}(V)))\\
          &= \mathbb{E}(UV - \frac{U}{2} - \frac{U}{2} - \mathbb{E}(U)\mathbb{E}(V))\\
          &= \mathbb{E}(U(1-U)) - 1/4 - 1/4 +1/4 \text{ par linéarité et d'après la solution ci-dessous}\\
          &= 1/2 - 1/3 -1/4 = -1/12
\end{align*} et
\[\rho(U,V) = \frac{\text{Cov}(U,V)}{\sqrt{\mathbb{V}(U)}\sqrt{\mathbb{V}(V)}} = -1\]

\runinhead{Coefficient de
corrélation}\phantomsection\addcontentsline{toc}{subsubsection}{Coefficient
de corrélation}\label{II-answer-CS-corr}
On déduit de l'inégalité de
\protect\hyperref[II-CS]{Cauchy-Schwarz (p.~\pageref*{II-CS})} que
\[|\rho(X,Y)| = \left|\frac{\mathbb{E}((X-\mathbb{E}(X))(X-\mathbb{E}(Y)))}{\sqrt{\mathbb{V}(X)\mathbb{V}(Y)}} \right| \leq \frac{\sqrt{\mathbb{E}((X-\mathbb{E}(X))^2)\mathbb{E}((X-\mathbb{E}(Y))^2)}}{\sqrt{\mathbb{V}(X)\mathbb{V}(Y)}} =1\]

\runinhead{Moments d'une variable aléatoire de loi uniforme sur
\([a,b]\)}\phantomsection\addcontentsline{toc}{subsubsection}{Moments
d'une variable aléatoire de loi uniforme sur
\([a,b]\)}\label{II-answer-moments-unif}
Soit \(X \sim \mathcal{U}_{[a,b]}\). Son espérance vaut
\[\mathbb{E}(X) = \int_a^b \frac{x}{b-a} dx = \frac{a+b}{2}\] et puisque
\[\mathbb{E}(X^2) = \int_a^b \frac{x^2}{b-a} dx = \frac{a^2 + ab +b^2}{3},\]
alors sa variance vaut
\[\mathbb{V}(X) = \mathbb{E}(X^2) - \mathbb{E}(X)^2 = \frac{(b-a)^2}{12}.\]

\runinhead{Moments d'une v.a. de loi
exponentielle}\phantomsection\addcontentsline{toc}{subsubsection}{Moments
d'une v.a. de loi exponentielle}\label{II-answer-moments-expo}
Soit \(X \sim \mathcal{E}(\theta)\), alors
\[\mathbb{E}(X) = \int_0^{+\infty} x \theta \exp(-\theta x) dx = [\frac{x}{\theta}\exp(-\theta x)]_0^{+\infty} + \int_0^{+\infty} \exp(-\theta x) dx = \frac{1}{\theta}\]

\runinhead{Moments d'une v.a. de loi
gamma}\phantomsection\addcontentsline{toc}{subsubsection}{Moments d'une
v.a. de loi gamma}\label{II-answer-moments-gamma}
Soit \(X\sim \Gamma(\alpha,\beta)\). On a \begin{align*}
\mathbb{E}(X) &= \frac{1}{\Gamma(\alpha)}\int_0^{+\infty} \theta^\alpha x^\alpha e^{-\theta x} dx\\
        &= \frac{1}{\Gamma(\alpha)}\int_0^{+\infty} \frac{1}{\theta} u^\alpha e^{-u} du \text{ par le changement de variable } u = \theta x\\
        &= \frac{\Gamma(\alpha+1)}{\Gamma(\alpha)} \frac{1}{\theta} = \frac{\alpha}{\theta}
\end{align*} puis \begin{align*}
\mathbb{E}(X^2) &= \frac{1}{\Gamma(\alpha)}\int_0^{+\infty} \theta^\alpha x^{\alpha+1} e^{-\theta x} dx\\
        &= \frac{1}{\Gamma(\alpha)}\int_0^{+\infty} \frac{1}{\theta^2} u^{\alpha+1} e^{-u} du \text{ par le changement de variable } u = \theta x\\
        &= \frac{\Gamma(\alpha+2)}{\Gamma(\alpha)} \frac{1}{\theta^2} = \frac{\alpha(\alpha+1)}{\theta^2}
\end{align*} d'où
\[\mathbb{V}(X) = \mathbb{E}(X^2) - \mathbb{E}(X)^2 = \frac{\alpha}{\theta^2}\]

\runinhead{Moments d'une v.a.
gaussienne}\phantomsection\addcontentsline{toc}{subsubsection}{Moments
d'une v.a. gaussienne}\label{II-answer-moments-gauss}
Pour s'assurer qu'il s'agit bien d'une densité, on remarque que
\(I = \int_{-\infty}^{+\infty} f(x) dx\) vérifie
\[ I^2 = \int_{-\infty}^{+\infty} \int_{-\infty}^{+\infty} f(x) f(y) dxdy = \int_{0}^{2\pi} d\theta\int_{0}^{+\infty}\frac{1}{2\pi} r e^{-r^2/2}dr \]
(en passant en coordonnées polaires dans l'intégrale double). Le calcul
est ensuite aisé et on obtient \(I = 1\).

Pour les moments, on fait d'abord le calcul dans le cas centré réduit
(\(\mu = 0\) et \(\sigma^2 =1\)) puis on s'y ramène par le changement de
variable \(x \mapsto \frac{x-\mu}{\sigma}\). Soit donc
\(X \sim \mathcal{N}(0,1)\), on a \begin{align*}
\mathbb{E}(X) &= \int_{-\infty}^{+\infty} \frac{x}{\sqrt{2\pi}} e^{-x^2/2} dx\\
         &= \left[-\frac{1}{\sqrt{2\pi}} e^{-x^2/2} \right]_{-\infty}^{+\infty} = 0.
\end{align*} et \begin{align*}
\mathbb{V}(X) = \mathbb{E}(X^2) &= \int_{-\infty}^{+\infty} \frac{x^2}{\sqrt{2\pi}} e^{-x^2/2} dx\\
         &= \left[-\frac{x}{\sqrt{2\pi}} e^{-x^2/2} \right]_{-\infty}^{+\infty} +\int_{-\infty}^{+\infty} \frac{1}{\sqrt{2\pi}} e^{-x^2/2} dx \\
         &= 0 + 1.
\end{align*}

\runinhead{\emph{Loi de
Cauchy}}\phantomsection\addcontentsline{toc}{subsubsection}{\emph{Loi de
Cauchy}}\label{II-answer-exo-cauchy}
L'angle \(\theta\) est une variable aléatoire uniforme sur
\([-\pi/2,\pi/2]\), de densité
\(g(\theta) = \frac{1}{\pi}1_{[ -\pi/2,\pi/2 ]}(\theta)\). L'abscisse
\(X\) est donnée par \(X = \tan \theta\), c'est donc une variable
aléatoire, de fonction de répartition \begin{align*}        
F(x) & = \mathbb{P}(X \leq x) \\
     & = \mathbb{P}(\theta \leq \arctan x) \\
     & = \int_{-\infty}^{\arctan x} \frac{1}{\pi}1_{[ -\pi/2,\pi/2 ]}(\theta) d\theta \\
     & = \frac{1}{\pi} \arctan x + \frac{1}{2}.
\end{align*}

\(F\) est de classe \(C^1\) de dérivée
\[f(x) = \frac{1}{\pi(1+x^2)}, x \in \mathbb{R}\]

C'est la densité de la loi de Cauchy. Une variable aléatoire \(X\) de
loi de Cauchy n'admet pas d'espérance. En effet,
\(\frac{x}{\pi(1+x^2)} \sim_{x\to \infty} \frac{1}{x}\) n'est pas
intégrable.

\runinhead{Calcul de l'aire d'un
triangle}\phantomsection\addcontentsline{toc}{subsubsection}{Calcul de
l'aire d'un triangle}\label{II-answer-triangle}
Si la fonction \(1_D\) est intégrable, alors
\protect\hyperref[II-fubiniproba]{le théorème de Fubini est
applicable (p.~\pageref*{II-fubiniproba})}. On a donc \[
a(T) = \int_{\mathbb{R}}\left[ \int_{\mathbb{R}} 1_T(x, y) \, dx \right] \, dy.
\] Or, si \(0 \leq y \leq 1\), \[
1_T(x, y) = \left| 
\begin{array}{rl}
1 & \mbox{si $0 \leq x \leq 1-y$,} \\
0 & \mbox{sinon.}
\end{array}
\right.
\] et dans le cas contraire, \(1_T(x, y) = 0\). On a donc \[
a(T) = \int_0^1\left[ \int_0^{1-y} dx \right] \, dy
=\int_0^1 (1 - y) \, dy= \left[y - \frac{y^2}{2} \right]_0^1 = \frac{1}{2}.
\]

\runinhead{Contre-exemple}\phantomsection\addcontentsline{toc}{subsubsection}{Contre-exemple}\label{II-answer-fubini-counter-example}
Pour tout \(y \in \left]0,1\right]\) (et donc presque tout
\(y \in [0, 1]\)), on a \[
\int_0^1 \frac{x^2 - y^2}{(x^2+y^2)^2} \, dx 
= 
\left[x \mapsto - \frac{x}{x^2+y^2} \right]_0^1
= 
- \frac{1}{1+y^2},
\] par conséquent \[
\int_0^1 \left[ \int_0^1 \frac{x^2 - y^2}{(x^2+y^2)^2} \, dx \right] \, dy
= - \int_0^1 \frac{dy}{1+y^2} \\
= - [y \mapsto \arctan y]_0^1
= - \frac{\pi}{4}.
\] En exploitant la relation \[
\frac{x^2 - y^2}{(x^2+y^2)^2} = \frac{\partial}{\partial y} \left(\frac{y}{x^2+y^2}\right) 
\] on établit de façon similaire que \[
\int_0^1 \left[ \int_0^1 \frac{x^2 - y^2}{(x^2+y^2)^2} \, dy \right] \, dx
= \frac{\pi}{4}.
\]

Or, si \protect\hyperref[II-fubiniproba]{le théorème de
Fubini (p.~\pageref*{II-fubiniproba})} était applicable, on pourrait
intervertir l'ordre d'intégration des variables sans changer le
résultat. Comme cela n'est pas le cas, on en déduit que l'hypothèse
exigée par le théorème de Fubini ne tient pas~: la fonction \[
(x, y) \in [0,1] \times [0,1] \mapsto \frac{x^2 - y^2}{(x^2+y^2)^2}
\] n'est pas intégrable.

\runinhead{Homothétie}\phantomsection\addcontentsline{toc}{subsubsection}{Homothétie}\label{II-answer-h}
Si l'on pose \(D_1 = \mathbb{R}^n\), \(D_2 = \mathbb{R}^n\),
l'application \(h:D_1 \to D_2\) définie par \(h(x) = \alpha x\) est un
\(C^1\)-difféomorphisme. De plus, on a \[
|\det J_{h}(x)| = \left|\det \left[
  \begin{array}{cccc}
  \alpha & 0  & \cdots & 0 \\
  0 & \alpha & \dots & 0 \\
  \vdots & \vdots & \vdots & \vdots \\
  0 & 0 & \dots & \alpha
  \end{array} \right] \right|  = |\alpha^n| = \alpha^n.
\] Par conséquent, le
\protect\hyperref[II-theorem-changement-de-variables]{théorème de
changement de
variables (p.~\pageref*{II-theorem-changement-de-variables})} fournit \[
\int_{\mathbb{R}^n} f(y) \, dy = \int_{\mathbb{R}^n} f(\alpha x) \alpha^n \, dx,
\] soit \[
\int_{\mathbb{R}^n} f(\alpha x) \, dx = \frac{1}{\alpha^n} \int_{\mathbb{R}^n} f(y) \, dy.
\]

\runinhead{Volume et
translation}\phantomsection\addcontentsline{toc}{subsubsection}{Volume
et translation}\label{II-answer-vr}
L'ensemble \(A\) mesurable et de volume fini a une fonction
caractéristique \(1_A\) intégrable et \[
\lambda(A) = \int 1_A(x) \, dx.
\] Soit \(h(x) = x - u\) ou \(u \in \mathbb{R}^3\). La fonction \(h\)
est une bijection de \(\mathbb{R}^3\) sur lui-même, continûment
différentiable ainsi que son inverse, \(h^{-1}(x) = x + u\) et
\(J_h(x) = I\), donc \(\det J_h(x) = 1\). Par
\protect\hyperref[II-theorem-changement-de-variables]{le théorème de
changement de
variables (p.~\pageref*{II-theorem-changement-de-variables})}, la
fonction \(1_{h^{-1}(A)} = 1_{A} \circ h\) est donc intégrable sur
\(\mathbb{R}^3\) et \[
\int 1_{h^{-1}(A)}\, dx = \int_{\mathbb{R}^3} (1_A \circ h) |\det J_h(x)| \, dx = \int_{\mathbb{R}^3} 1_A(x) \, dx = \lambda(A).
\] L'ensemble translaté \(A + u = h^{-1}(A)\) est donc mesurable, de
volume fini égal au volume de \(A\).

\runinhead{Coordonnées
polaires}\phantomsection\addcontentsline{toc}{subsubsection}{Coordonnées
polaires}\label{II-answer-cp}
Les ensembles \(P\) et \(C\) sont ouverts et la fonction \(h\) -- qui
permet de passer des coordonnées polaires aux coordonnées cartésiennes
-- est une bijection de \(P\) dans \(C\). Elle est continûment
différentiable\,; sa matrice jacobienne en \((r, \theta)\) vaut \[
J_{h}(r, \theta) = 
\left[ 
\begin{array}{cr}
\cos \theta & -r \sin \theta \\
\sin \theta & r \cos \theta
\end{array}
\right],
\] et son déterminant satisfait \[
\det J_{h}(r, \theta) = (\cos \theta)(r \cos \theta) - (\sin \theta)(-r\sin \theta)
= r > 0.
\] Le jacobien est donc inversible et \(h^{-1}\) est continûment
différentiable par le théorème d'inversion locale. On peut donc
appliquer \protect\hyperref[II-theorem-changement-de-variables]{le
théorème de changement de
variables (p.~\pageref*{II-theorem-changement-de-variables})} à la
fonction \(f\), ce qui fournit \begin{align*}
\int_C f(x, y) \, d(x, y) &= \int_P f(h(r, \theta)) |\det J_h(r,\theta)| d(r,\theta) \\
&= \int_P g(r, \theta)  r \, d(r,\theta).
\end{align*}

\runinhead{Absence du déterminant
jacobien}\phantomsection\addcontentsline{toc}{subsubsection}{Absence du
déterminant jacobien}\label{II-answer-adj}
La façon la plus rapide de procéder consiste à considérer que
\(f \circ h\) joue le rôle de \(f\) dans
\protect\hyperref[II-theorem-changement-de-variables]{le théorème de
changement de
variables (p.~\pageref*{II-theorem-changement-de-variables})}, que \(h\)
dans notre énoncé désigne \(h^{-1}\) dans ce théorème et que les rôles
de \(D_1\) et \(D_2\) sont intervertis. Une fois que l'on a permuté ces
notations, on réalise que l'intégrale que l'on souhaite calculer est le
membre de gauche de l'équation du théorème de changement de variables,
dont toutes les hypothèses sont par ailleurs satisfaites. Par conséquent
on a \[
\int_{D_1} f (h(x)) \, dx = \int_{D_2} (f \circ h)(h^{-1}(y)) |\det J_{h^{-1}}(y)| \, dy.
\] On peut simplifier \((f \circ h)(h^{-1}(y))\) en \(f(y)\) et
éventuellement exprimer le jacobien de \(h^{-1}\) en fonction du
jacobien de \(h\)~: \(J_{h^{-1}}(y) = [J_{h}(h^{-1}(y))]^{-1}\)\,; on
obtient donc \[
\int_{D_1} f (h(x)) \, dx = 
\int_{D_2} f(y) |\det J_{h^{-1}}(y)| \, dy = 
\int_{D_2} f(y) \frac{1}{|\det J_{h}(h^{-1}(y))|} \, dy.
\]

\runinhead{Lancer de
fléchette}\phantomsection\addcontentsline{toc}{subsubsection}{Lancer de
fléchette}\label{II-answer-exo-flechette}
Les coordonnées cartésiennes de
\(M \in D = \{(x,y) \in \mathbb{R}^2 , x^2+y^2 \leq 1\}\) constituent un
couple de variables aléatoires de densité
\[f_{(X,Y)}(x,y) = \frac{1}{\pi}1_{D} (x,y)\] uniforme sur le disque,
par hypothèse. L'abscisse de \(X\) est distribuée selon la densité
marginale
\[f_X(x) = \int f_{(X,Y)}(x,y) dy = \frac{2}{\pi} (1-x^2)^{1/2} 1_{[-1,1]}(x).\]
La loi de \(Y\) a la même densité.

\runinhead{Lancer de fléchette
(suite)}\phantomsection\addcontentsline{toc}{subsubsection}{Lancer de
fléchette (suite)}\label{II-answer-exo-flechettebis}
La densité des coordonnées cartésiennes ne s'exprime pas comme le
produit des densités marginales\,; elles ne sont pas indépendantes.

\runinhead{Démonstration}\phantomsection\addcontentsline{toc}{subsubsection}{Démonstration}\label{II-answer-demo-cov-indep}
Soit \(X\) et \(Y\) deux v.a. aléatoires indépendantes et de carré
intégrable. On a d'après \protect\hyperref[II-indep_fct]{la
proposition (p.~\pageref*{II-indep_fct})}
\[\mathbb{E}((X-\mathbb{E}(X))(Y-\mathbb{E}(Y))) = \mathbb{E}(X-\mathbb{E}(X))\mathbb{E}(Y-\mathbb{E}(Y)) = 0 = \rho(X,Y)\]

\runinhead{Lancer de fléchette
(fin)}\phantomsection\addcontentsline{toc}{subsubsection}{Lancer de
fléchette (fin)}\label{II-answer-exo-flechetteter}
En coordonnées polaires, le domaine \(D\) est décrit par
\(D = \{(r,\theta) \in [0,1] \times [0,2\pi]\}\). La densité de
\((R,\theta)\) s'écrit ainsi~:
\[f_{(R,\theta)}(r,\theta) = \frac{r}{\pi}1_{D} (r\cos(\theta),r\sin(\theta)) = 2r 1_{[0,1]}(r) \frac{1}{2\pi}1_{[0,2\pi]}(r).\]
Ainsi \(R\) et \(\Theta\) sont indépendantes de densités respectives
\(2r 1_{[0,1]}(r)\) et \(\frac{1}{2\pi}1_{[0,2\pi]}(\theta)\).

\runinhead{Loi
bêta}\phantomsection\addcontentsline{toc}{subsubsection}{Loi
bêta}\label{II-answer-exo-loibeta}
On note d'abord que la dimension de \(Y\) est plus petite que celle de
\(X\). On va donc compléter \(Y\) en prenant par exemple \(Y' = (Y,Z)\),
avec \(Z=U+V\), ce qui correspond à
\(g(u,v) = \left(\frac{u}{u+v},u+v\right)\). Cette application est
bijective de \(A = \left]0,+\infty\right[^2\) dans
\(B = \left]0,1\right[ \times \left]0,+\infty\right[\), d'inverse
\(g^{-1}(y,z) = (yz,z(1-y))\), qui a pour jacobien \(z\).

Comme
\(f_X(u,v) = \frac{\theta^{\alpha+\beta}}{\Gamma(\alpha)\Gamma(\beta)}u^{\alpha-1}v^{\beta-1}e^{-\theta(u+v)}1_A(u,v)\),
on a
\[f_{Y'}(y,z) = \frac{\theta^{\alpha+\beta}}{\Gamma(\alpha)\Gamma(\beta)}z^{\alpha+\beta -1}y^{\alpha-1}(1-y)^{\beta-1}e^{-\theta z}1_B(y,z).\]
On obtient alors la densité de Y en intégrant \(f_{Y'}(y,z)\) par
rapport à \(z\in \left]0,+\infty\right[\)~: \begin{align*}
f_Y(y) &= \int_0^{+\infty} f_{Y'}(y,z) dz\\
       &= \frac{\theta^{\alpha+\beta}}{\Gamma(\alpha)\Gamma(\beta)}y^{\alpha-1}(1-y)^{\beta-1}1_{\left]0,1\right[}(y)\int_0^{+\infty}z^{\alpha+\beta -1}e^{-\theta(z)}dz\\
       &= \frac{\Gamma(\alpha+\beta)}{\Gamma(\alpha)\Gamma(\beta)}y^{\alpha-1}(1-y)^{\beta-1}1_{\left]0,1\right[}(y),
\end{align*} où on a utilisé la définition de la fonction gamma et le
changement de variable linéaire \(z \mapsto \theta z\). On appelle loi
bêta de paramètres \(\alpha\) et \(\beta\) la loi admettant cette
densité. Admettant une grande variété de formes, elle permet de
modéliser de nombreuses distributions à support fini.

La loi bêta apparaît naturellement dans une expérience d'urnes, donnée
par George Pólya dans un article de 1930,
\href{http://www.numdam.org/article/AIHP_1930__1_2_117_0.pdf}{\emph{Sur
quelques points de la théorie des probabilités}}. Il décrit l'expérience
suivante~: on se donne une urne contenant initialement \(r\) boules
rouges et \(b\) boules bleues, on tire une boule dans l'urne, puis on la
remet dans l'urne avec une deuxième boule de même couleur. Alors la
proportion de boules rouges tend vers une variable aléatoire de loi
Beta\((r,b)\), et, inversement, la proportion de boules bleues tend vers
une variable aléatoire de loi Beta\((b,r)\).

Nous obtenons aussi facilement la densité de \(Z\). En effet, on a
\(f_{Y'}(y,z) = f_Y(y)f_Z(z)\) (\(Y\) et \(Z\) sont donc indépendantes),
où
\[f_Z(z) = \frac{\theta^{\alpha+\beta}}{\Gamma(\alpha+\beta)}z^{\alpha+\beta -1}e^{-\theta z}1_{\left]0,+\infty\right[}\]
On a ainsi démontré que si \(U\) et \(V\) sont deux variables aléatoires
indépendantes de lois respectives \(\Gamma(\alpha,\theta)\) et
\(\Gamma(\beta,\theta)\), alors \(U+V\) suit la loi
\(\Gamma(\alpha+\beta,\theta)\) et est indépendante de \(\frac{U}{U+V}\)
qui suit une loi bêta de paramètres \((\alpha,\beta)\).

\runinhead{Divergence de Kullback-Leibler entre deux
gaussiennes}\phantomsection\addcontentsline{toc}{subsubsection}{Divergence
de Kullback-Leibler entre deux gaussiennes}\label{II-answer-exoklgauss}
On a
\[f(x) = \frac{1}{\sqrt{2\pi}\sigma_1} e^{-\frac{(x-\mu_1)^2}{2\sigma_1^2}}\]
et
\[g(x) = \frac{1}{\sqrt{2\pi}\sigma_2} e^{-\frac{(x-\mu_2)^2}{2\sigma_2^2}}\]

On en déduit que \begin{align*}
D_{KL} (f || g) &= \int_{-\infty}^{+\infty} f(x) \log \left(\frac{f(x)}{g(x)}\right) dx \\
                &= \int_{-\infty}^{+\infty} f(x) \left[ \log \left(\frac{\sigma_2}{\sigma_1}\right) + \frac{(x-\mu_2)^2}{2\sigma_2^2} - \frac{(x-\mu_1)^2}{2\sigma_1^2} \right] dx \\
                &= \log \left(\frac{\sigma_2}{\sigma_1}\right) + \frac{1}{2\sigma_2^2} \int_{-\infty}^{+\infty} f(x) (x-\mu_2)^2 dx - \frac{1}{2\sigma_1^2} \int_{-\infty}^{+\infty} f(x) (x-\mu_1)^2 dx
\end{align*} On a~: \begin{align*}
\int_{-\infty}^{+\infty} f(x) (x-\mu_2)^2 dx &= \int_{-\infty}^{+\infty} f(x) ((x-\mu_1)+\mu_1-\mu_2)^2 dx \\
&= \int_{-\infty}^{+\infty} f(x) (x-\mu_1)^2 dx + (\mu_1 - \mu_2)^2 \\
&= \sigma_1^2 + (\mu_1 - \mu_2)^2 
\end{align*} On en déduit la formule annoncée~:
\[D_{KL} (f || g) = \log \left(\frac{\sigma_2}{\sigma_1}\right) + \frac{\sigma_1^2 + (\mu_1 - \mu_2)^2}{2\sigma_2^2} - \frac{1}{2}\]

\section{Durée de vie}\label{II-duruxe9e-de-vie-1}

\runinhead{Question
1}\phantomsection\addcontentsline{toc}{subsubsection}{Question
1}\label{II-answer-dureevie1}
En dérivant, on obtient que la densité de \(T\) est
\(f(t) = t e^{-t^2/2}1_{t>0}\).

L'espérance vaut~:

\begin{align*}
\mathbb{E}(T) &= \int_0^{+\infty} t^2 e^{-t^2/2} dt\\
        &= \left[-t e^{-t^2/2} \right]_0^{+\infty} + \int_0^{+\infty} e^{-t^2/2} dt \text{ par intégration par parties}\\
        &= \frac{1}{2} \int_{-\infty}^{+\infty} e^{-t^2/2} dt \\
        &= \frac{\sqrt{2\pi}}{2}
\end{align*}

\runinhead{Question
2}\phantomsection\addcontentsline{toc}{subsubsection}{Question
2}\label{II-answer-dureevie2}
La probabilité s'écrit~:

\[\mathbb{P}(T\geq 3 |T \geq 1) = \frac{\mathbb{P}(T \geq 3,T \geq 1)}{\mathbb{P}(T \geq 1)} = \frac{\mathbb{P}(T \geq 3)}{\mathbb{P}(T \geq 1)} = \frac{e^{-9/2}}{e^{-1/2}} = e^{-4}\]

\runinhead{Question
3}\phantomsection\addcontentsline{toc}{subsubsection}{Question
3}\label{II-answer-dureevie3}
Les variables aléatoires \(X_1 , \ldots , X_{10}\) sont indépendantes et
suivent une loi de Bernoulli de paramètre~:
\[ \mathbb{P}(T \leq 1) = F (1) = 1 - e^{- 1/2}\]

On en déduit que la loi de N est la loi binomiale de paramètre
\((10, 1 - e^{- 1/2})\).

\runinhead{Question
4}\phantomsection\addcontentsline{toc}{subsubsection}{Question
4}\label{II-answer-dureevie4}
La probabilité que l'équipement en série soit défaillant avant 1 an
vaut~:
\[ P(N \geq 1) = 1 - \mathbb{P}(N = 0) = 1 - e^{-10/2} \approx 9.9 ~10^{-1} = 99\%\]

\runinhead{Question
5}\phantomsection\addcontentsline{toc}{subsubsection}{Question
5}\label{II-answer-dureevie5}
La probabilité que l'équipement en parallèle soit défaillant avant 1 an
vaut~:
\[ \mathbb{P}(N = 10) = (1 - e^{-1/2})^10 \approx 8.9 ~10 ^{-5} \]

\section{Crues centennales de la
Seine}\label{II-crues-centennales-de-la-seine-1}

\runinhead{Question
1}\phantomsection\addcontentsline{toc}{subsubsection}{Question
1}\label{II-answer-crue1}
Soit \(x\in \mathbb{R}\) \begin{align*}
F_Y(x) = F(Y\leq x) &= F(X_1\leq x,\ldots,X_n\leq x)\\
           &= \prod_{i=1}^n F(X_i\leq x) \\
           &= F(x)^n
\end{align*} puisque les \(X_i\) sont indépendantes et de même loi

\runinhead{Question
2}\phantomsection\addcontentsline{toc}{subsubsection}{Question
2}\label{II-answer-crue2}
Soit \(x\in \mathbb{R}\) \begin{align*}
f_Y(x) &= \frac{d}{d x} F_Y(x)\\
       &= \frac{d}{d x} F(x)^n\\
       &= nf(x)F(x)^{n-1}
\end{align*}

\runinhead{Question
3}\phantomsection\addcontentsline{toc}{subsubsection}{Question
3}\label{II-answer-crue3}
On a alors \(f(x) = 1_{[0,1]}(x)\) et
\(F(x) = x 1_{[0,1]}(x) + 1_{[1,+\infty]}\). On en déduit
\[ f_Y(x) = n x^{n-1}1_{[0,1]} \]

\runinhead{Question
4}\phantomsection\addcontentsline{toc}{subsubsection}{Question
4}\label{II-answer-crue4}
On cherche à déterminer \(x\) tel que
\(\mathbb{P}(Y > x) \leq \frac{1}{1000}\). Or,
\[ \mathbb{P}(Y > x) = 1-\mathbb{P}(Y \leq x) = 1 - F_Y(x) = 1 - F(x)^n = 1-x^n \]
Ainsi
\[ 1-x^n \leq \frac{1}{1000} \Leftrightarrow (1-\frac{1}{1000})^{1/n} \leq x \]

\runinhead{Question
5}\phantomsection\addcontentsline{toc}{subsubsection}{Question
5}\label{II-answer-crue5}
On rappelle que la médiane est le réel \(m\) tel que
\(\frac{1}{2} = \mathbb{P}(Y \leq m) = F_Y(m)\).

On cherche donc \(m\) tel que \(m^n = \frac{1}{2}\), soit
\(m = \frac{1}{2^{1/n}}\).

On voit que cette valeur est plus élevée que la médiane des \(X_i\) qui
vaut \(\frac{1}{2}\). Elle tend même vers 1 lorsque \(n\) tend vers
l'infini.

\section{Indépendance et vecteurs
gaussiens}\label{II-induxe9pendance-et-vecteurs-gaussiens-1}

\runinhead{Question
1}\phantomsection\addcontentsline{toc}{subsubsection}{Question
1}\label{II-answer-covindepgauss-joint}
D'après le cours, en interprétant les éléments de \(\mathbb{R}^2\) comme
des vecteurs colonnes, pour tout \(z =(x,y)\in\mathbb{R}^2\) on a \[
f_Z(z) = \dfrac{1}{2\pi\,\sqrt{\det(C)}}\,\exp\left\{ -\dfrac{1}{2}\,\left(z - m\right)^t\,C^{-1}\,\left(z - m\right) \right\}.
\] Or par construction \(C\) peut s'écrire
\[C = \left( \begin{array}{cc} \sigma_X^2 & c\\ c & \sigma_Y^2 \end{array}\right), \]
où \(\sigma_X^2 := \mathbb{V}\left(X\right)\),
\(\sigma_Y^2 := \mathbb{V}\left(Y\right)\) et
\(c = \text{Cov}\left(X,Y\right)\). Comme \(C\) est définie positive
elle est inversible, avec \(\det(C) = \sigma_X^2\,\sigma_Y^2 - c^2 > 0\)
et
\[C^{-1} = \dfrac{1}{\sigma_X^2\,\sigma_Y^2 - c^2}\,\left( \begin{array}{cc} \sigma_Y^2 & -c\\ -c & \sigma_X^2 \end{array}\right).\]
En outre, \(m = \left(m_X, m_Y \right)\) avec \(m_X := \mathbb{E}(X)\)
et \(m_Y := \mathbb{E}(Y)\). On en déduit que \begin{align*}
f_Z(x,y) &= \dfrac{1}{2\pi\,\sqrt{\sigma_X^2\,\sigma_Y^2 - c^2}}\,\exp\biggl\{ \dfrac{-1}{2\,\left(\sigma_X^2\,\sigma_Y^2 - c^2 \right)}\times\\
&\ \ \ \left( \sigma_Y^2\,(x-m_X)^2 - 2c\,(x-m_X)\,(y-m_Y) + \sigma_X^2\,(y-m_Y)^2 \right) \biggr\}\\
&= \dfrac{1}{2\pi\,\sigma_X\,\sigma_Y\,\sqrt{1-\rho^2}}\,\,\exp\biggl\{ \dfrac{-1}{2\,\left( 1 - \rho^2 \right)}\times\\
&\ \ \ \left( \dfrac{(x-m_X)^2}{\sigma_X^2} - 2\rho\,\dfrac{(x-m_X)}{\sigma_X}\,\dfrac{(y-m_Y)}{\sigma_Y} + \dfrac{(y-m_Y)^2}{\sigma_Y^2} \right) \biggr\}.
\end{align*}

\runinhead{Question
2}\phantomsection\addcontentsline{toc}{subsubsection}{Question
2}\label{II-answer-covindepgauss-marg}
Ici, pour \(x,y \in \mathbb{R}^2\), on remarque que~: \begin{align*}
f_Z(x,y) &= \dfrac{1}{2\pi\,\sigma_X\,\sigma_Y\,\sqrt{1-\rho^2}}\,\,\exp\biggl\{ \dfrac{-1}{2\,\left( 1 - \rho^2 \right)}\times\\
&\ \ \ \left( \left(\dfrac{y-m_Y}{\sigma_Y} - \rho\,\dfrac{x-m_X}{\sigma_X}\right)^2 + (1 - \rho^2)\,\dfrac{(x-m_X)^2}{\sigma_X^2} \right) \biggr\}\\
& = \dfrac{1}{\sqrt{2\pi}\,\sigma_X}\,\exp\left\{-\dfrac{(x-m_X)^2}{2\,\sigma_X^2} \right\}\times\\
&\ \ \ \dfrac{1}{\sqrt{2\pi}\,\sigma_Y\,\sqrt{1-\rho^2}}\,\exp\left\{-\dfrac{1}{2\,(1-\rho^2)}\,\left(\dfrac{y-m_Y}{\sigma_Y} - \rho\,\dfrac{x-m_X}{\sigma_X}\right)^2 \right\}.
\end{align*} Ainsi, avec le changement de variable
\(u = \dfrac{y-m_Y}{\sigma_Y}\) on obtient \begin{align*}
f_X(x) &= \int_{-\infty}^{+\infty} f_Z(x,y)\,dy\\
&= \dfrac{1}{\sqrt{2\pi}\,\sigma_X}\,\exp\left\{-\dfrac{(x-m_X)^2}{2\,\sigma_X^2} \right\}\times\\
&\ \ \ \int_{-\infty}^{+\infty} \dfrac{1}{\sqrt{2\pi}\,\sqrt{1-\rho^2}}\,\exp\left\{-\dfrac{1}{2\,(1-\rho^2)}\,\left(u - \rho\,\dfrac{x-m_X}{\sigma_X}\right)^2 \right\}\,du,
\end{align*} et on reconnaît dans cette dernière intégrale la densité
d'une loi Normale d'espérance \(\rho\,\dfrac{x-m_X}{\sigma_X}\) et de
variance \(1-\rho^2\). Cette intégrale vaut donc \(1\) et on conclut que
\[f_X(x) = \dfrac{1}{\sqrt{2\pi}\,\sigma_X}\,\exp\left\{-\dfrac{(x-m_X)^2}{2\,\sigma_X^2} \right\}, \]
qui correspond à la densité d'une loi Normale d'espérance \(m_X\) et de
variance \(\sigma_X^2\).

En procédant de manière symétrique, on obtient de même que \(Y\) suit
une loi Normale d'espérance \(m_Y\) et de variance \(\sigma_Y^2\).

\runinhead{Question
3}\phantomsection\addcontentsline{toc}{subsubsection}{Question
3}\label{II-answer-covindepgauss-indep}
Le premier sens est évident~: si \(X\) et \(Y\) sont indépendantes,
alors nous avons vu dans le cours que
\(\text{Cov}\left(X,Y\right) = 0\), et ce que l'on soit gaussien ou non.

Supposons maintenant que \(\text{Cov}\left(X,Y\right) = 0\), et montrons
l'indépendance entre \(X\) et \(Y\). En reprenant la formule de la
question 1 et en remplaçant \(\rho\) par \(0\), pour tout
\((x,y) \in \mathbb{R}^2\) on a: \begin{align*}
f_Z(x,y) &= \dfrac{1}{2\,\pi\,\sigma_X\,\sigma_Y}\,\exp\left\{ -\dfrac{1}{2}\,\left( \dfrac{(x-\mu_X)^2}{\sigma_X^2} + \dfrac{(y-\mu_Y)^2}{\sigma_Y^2} \right) \right\}\\
&= \dfrac{1}{\sqrt{2\pi}\,\sigma_X}\,\exp\left\{ -\dfrac{(x-m_X)^2}{2\,\sigma_X^2} \right\} \times \dfrac{1}{\sqrt{2\pi}\,\sigma_Y}\,\exp\left\{ -\dfrac{(y-m_Y)^2}{2\,\sigma_Y^2} \right\}\\
&= f_X(x)\times f_Y(y)
\end{align*} d'après la question 2. On en conclut que \(X\) et \(Y\)
sont bien indépendantes.

\section{Symétrie de la
gaussienne}\label{II-symuxe9trie-de-la-gaussienne-1}

\runinhead{Question
1}\phantomsection\addcontentsline{toc}{subsubsection}{Question
1}\label{II-answer-gausssym-loi}
Nous avons admis dans le cours que la fonction de répartition \(F_c\) de
\(X_c\) caractérisait sa loi. Nous allons donc l'expliciter. Soit
\(x \in \mathbb{R}\). Par définition de la fonction de répartition d'une
variable aléatoire et d'après la formule des probabilités totales, on a
\begin{align*}
F_c(x) &= \mathbb{P}\left(X_c \leq x\right) = \mathbb{P}\left( X_c \leq x, |X| > c \right) + \mathbb{P}\left( X_c \leq x, |X| \leq c\right)\\
&= \mathbb{P}\left( X\leq x, |X| > c \right) + \mathbb{P}\left(-X \leq x, |X| \leq c\right)\\
&= \mathbb{P}\left( X\leq x, X > c \right) + \mathbb{P}\left(X \leq x, X < -c \right) + \mathbb{P}\left( X \geq -x, -c \leq X \leq c\right).
\end{align*} Or
\[\mathbb{P}\left( X\leq x, X > c \right) = \left|\begin{array}{ll} \mathbb{P}\left(c < X \leq x \right) = F(x) - F(c) & \text{si } x > c,\\ 0 & \text{si } x\leq c,\end{array}\right.\]
\[ \mathbb{P}\left(X \leq x, X < -c \right) = \mathbb{P}\left( X \leq \min\{-c,x\} \right) = F\left(\min\{-c,x\}\right), \]
\[ \mathbb{P}\left( X \geq -x, -c \leq X \leq c\right) = \left|\begin{array}{ll} \mathbb{P}\left(-c \leq X \leq c \right) = F(c) - F(-c) & \text{si } x > c, \\ \mathbb{P}\left(-x \leq X \leq c \right) = F(c) - F(-x) & \text{si } -c \leq x \leq c,\\ 0 & \text{si } x<-c. \end{array}\right.\]
Ainsi,

\begin{itemize}
\item
  si \(x > c\) alors \(\min\{-c,x\} = -c\) et
  \[F_c(x) = F(x) - F(c) + F\left(\min\{-c,x\}\right) + F(c) - F(-c) = F(x),\]
\item
  si \(-c \leq x \leq c\) alors \(\min\{-c,x\} = -c\) et
  \[F_c(x) = F\left(\min\{-c,x\}\right) + F(c) - F(-x) = 1 - F(c) + F(c) - 1 + F(x) = F(x)\]
  car la densité de la loi Normale centrée réduite est paire, ce qui
  implique que \(F(-x) = 1 - F(x)\) pour tout \(x\in\mathbb{R}\),
\item
  si \(x < -c\) alors \(\min\{-c,x\} = x\) et
  \[F_c(x) =  F\left(\min\{-c,x\}\right) = F(x).\]
\end{itemize}

On en conclut que \(F_c = F\)~: \(X_c\) a la même loi que \(X\). Ce
résultat est dû à la symétrie de la densité de la loi Normale centrée
réduite. On remarque d'ailleurs que le résultat resterait inchangé si
l'on prenait n'importe quelle autre loi de densité paire\,!

\runinhead{Question
2}\phantomsection\addcontentsline{toc}{subsubsection}{Question
2}\label{II-answer-gausssym-cov}
Tout d'abord, par définition,
\[\text{Cov}\left(X,X_c\right) = \mathbb{E}\left( \left(X - \mathbb{E}\left(X\right)\right)\,\left(X_c - \mathbb{E}\left(X_c\right) \right) \right) = \mathbb{E}\left(X\,X_c \right) \]
puisque \(X\) et \(X_c\) sont gaussiennes centrées (d'espérance nulle)
réduites. Or on remarque que l'on peut écrire
\(X_c = X\,1_{\mathbb{R}\backslash[-c,c]}(X) - X\,\,1_{[-c,c]}(X)\). On
déduit par linéarité de l'espérance que \begin{align*}
\mathbb{E}\left(X\,X_c\right) &= \mathbb{E}\left(X^2\,1_{\mathbb{R}\backslash[-c,c]}(X) - X^2\,1_{[-c,c]}(X)\right)\\
&= \mathbb{E}\left(X^2\,\left(1_{\mathbb{R}\backslash[-c,c]}(X) + 1_{[-c,c]}(X) - 1_{[-c,c]}(X)\right) - X^2\,1_{[-c,c]}(X)\right)\\
&= \mathbb{E}\left(X^2\right) - 2\,\mathbb{E}\left(X^2\,1_{[-c,c]}(X)\right).
\end{align*} Etant donné que \(X\) est de loi Normale centrée réduite,
on a directement que \(\mathbb{E}\left(X^2\right)=1\); il suffit donc de
calculer la dernière espérance pour obtenir le résultat. Comme
\(x\mapsto x^2\,1_{[-c,c]}(x)\) est bornée et continue sur
\(\mathbb{R}\backslash\{-c,c\}\) (donc continue presque partout), elle
est intégrable et
\[\mathbb{E}\left(X^2\,1_{[-c,c]}(X)\right) = \int_{-\infty}^{+\infty} x^2\,1_{[-c,c]}\,f(x)\,dx = \int_{-c}^c x^2\,f(x)\,dx. \]
Or on pourra remarquer que \(f^\prime(x) = -x\,f(x)\) puis que
\(f^{\prime\prime}(x) = (x^2-1)\,f(x)\). Puisque \(f\), \(f^\prime\) et
\(f^{\prime\prime}\) sont continues et bornées (voir l'exercice
\emph{Densité et fonction de répartition d'une loi Normale} du chapitre
Probabilités I) elles sont intégrables sur des segments. On en déduit
que \begin{align*}
\mathbb{E}\left(X^2\,1_{[-c,c]}(X)\right) &= \int_{-c}^c (x^2-1)\,f(x)\,dx + \int_{-c}^c f(x)\,dx\\
&= \left[-x\,f(x) \right]_{-c}^c + F(c) - F(-c)\\
&= -c\,\left(f(c) + f(-c)\right) + F(c) - F(-c)\\
&= 2\,\left(F(c) - c\,f(c)\right) - 1
\end{align*} par parité de \(f\). On en conclut que
\(\text{Cov}\left(X,X_c\right) = 3 + 4\,\left(c\,f(c) - F(c)\right)\).

La fonction
\(g : c\in\mathbb{R}_+^\ast \mapsto 3+4\,\left(c\,f(c) - F(c)\right)\)
est continue, infiniment dérivable et pour tout \(c > 0\) on a
\(g^\prime(c) = -4\,c^2\,f(c) < 0\). Par conséquent, \(g\) est
strictement décroissante, avec \(\lim\limits_{c\to0^+} g(c) = 1\) et
\(\lim\limits_{c\to+\infty} g(c) = -1\). Le théorème des valeurs
intermédiaires nous permet de conclure qu'il existe bien un unique
\(c_0 > 0\) tel que \(g(c_0) = 0\).

\runinhead{Question
3}\phantomsection\addcontentsline{toc}{subsubsection}{Question
3}\label{II-answer-gausssym-indep}
Les variables aléatoires \(X\) et \(X_{c_0}\) sont indépendantes ssi
pour tous ensembles \(A,B \in \mathcal{E}_\mathbb{R}\) on a
\[\mathbb{P}\left(X \in A, X_{c_0} \in B \right) = \mathbb{P}\left(X \in A\right)\,\mathbb{P}\left(X_{c_0} \in B \right).\]

Or si l'on prend \(A = B = ]-\infty, x]\) où \(x < -c_0\), alors
\begin{align*}
\mathbb{P}\left( X \leq x, X_{c_0} \leq x \right) &= \mathbb{P}\left( X \leq x, X < -c_0 \right) + \mathbb{P}\left( X \leq x, X > c_0 \right)\\
&\ \ \ + \mathbb{P}\left( X\leq x, X \geq -x, -c_0 \leq X \leq c_0 \right)\\
&= \mathbb{P}\left( X \leq x \right) +0 + 0 = F(x)\\
& \neq F(x)\,F(x) = \mathbb{P}\left( X \leq x \right)\,\mathbb{P}\left(X_{c_0} \leq x \right).
\end{align*}

On en conclut que bien que \(\text{Cov}\left(X,X_{c_0}\right) = 0\),
\(X\) et \(X_{c_0}\) ne sont pas indépendantes.

\runinhead{Question
4}\phantomsection\addcontentsline{toc}{subsubsection}{Question
4}\label{II-answer-gausssym-vecgauss}
Supposons que \((X,X_{c_0})\) soit un vecteur gaussien. Alors, puisque
\(\text{Cov}\left(X,X_{c_0}\right) = 0\), \(X\) et \(X_{c_0}\) sont
nécessairement indépendantes. Or nous venons de voir que ce n'était pas
le cas; nous avons donc contradiction. On en conclut que \((X,X_{c_0})\)
n'est pas gaussien.

\section{Loi du \(\chi^2\)}\label{II-loi-du-chi2-1}

\runinhead{A 1 degré de
liberté}\phantomsection\addcontentsline{toc}{subsubsection}{A 1 degré de
liberté}\label{II-answer-khi2-1dl}

\begin{enumerate}
\def\labelenumi{\arabic{enumi}.}
\tightlist
\item
  Soit \(h : \mathbb{R}\to \mathbb{R}\) une fonction mesurable telle que
  \(h\times f\) est intégrable. Alors en opérant le changement de
  variable \(u = x^2\) on a \begin{align*}
  \mathbb{E}\left( h(Y) \right) &= \mathbb{E}\left(h(X^2)\right) = \int_{\mathbb{R}} h(x^2)\,f(x)\,dx\\
  & = \int_{-\infty}^{0} h(x^2)\,f(x)\,dx + \int_{0}^{+\infty} h(x^2)\,f(x)\,dx\\
  &= \int_{0}^{+\infty} h(u)\,f\left(-\sqrt{u}\right)\,\dfrac{1}{2\,\sqrt{u}} \,du + \int_{0}^{+\infty} h(u)\,f\left(\sqrt{u}\right)\,\dfrac{1}{2\,\sqrt{u}} \,du.
  \end{align*} Or \(f\) est paire, donc pour tout \(u \in \mathbb{R}^+\)
  on a \(f\left(-\sqrt{u}\right) = f\left(\sqrt{u}\right)\) et
  \begin{align*}
  \mathbb{E}\left( h(Y) \right) &= \int_{\mathbb{R}_+} h(u)\,f\left(\sqrt{u}\right)\,\dfrac{1}{\sqrt{u}} \,du.
  \end{align*} Par propriété \(Y\) possède donc une densité \(f_Y\)
  valant pour tout \(x\in\mathbb{R}\)
  \[f_Y(x) = \left|\begin{array}{ll} f\left(\sqrt{x}\right)\,\dfrac{1}{\sqrt{x}} = \dfrac{x^{-\frac{1}{2}}}{2^{\frac{1}{2}}\,\Gamma\left(\frac{1}{2}\right)}\,\exp\left\{-\dfrac{x}{2}\right\} & \text{si } x > 0,\\ 0 & \text{si } x \leq 0. \end{array}\right. \]
\end{enumerate}

\textbf{Remarque.} La fonction \(f_Y\) est bien une densité~: elle est
continue par morceaux donc mesurable, puis positive, et son intégrale
sur \(\mathbb{R}\) vaut bien \(1\)~: \begin{align*}
\int_{\mathbb{R}} f_Y(x)\,dx &= \dfrac{1}{2\,\Gamma\left(\frac{1}{2}\right)}\,\int_{0}^{+\infty} \left(\dfrac{x}{2}\right)^{\frac{1}{2}-1}\,\exp\left\{-\dfrac{x}{2}\right\}\,dx \\
&= \dfrac{1}{\Gamma\left(\frac{1}{2}\right)}\,\int_{0}^{+\infty} u^{\frac{1}{2}-1}\,e^{-u}\,du = \dfrac{\Gamma\left(\frac{1}{2}\right)}{\Gamma\left(\frac{1}{2}\right)} = 1.
\end{align*}

\runinhead{A n degrés de
liberté}\phantomsection\addcontentsline{toc}{subsubsection}{A n degrés
de liberté}\label{II-answer-khi2-ndl}

\begin{enumerate}
\def\labelenumi{\arabic{enumi}.}
\setcounter{enumi}{1}
\tightlist
\item
  On note \((\mathcal{P}_n)\) la propriété à démonter au rang \(n\).
\end{enumerate}

\textbf{Initialisation.} Nous avons vu à la question précédente que
\((\mathcal{P}_1)\) est vraie.

\textbf{Hérédité.} Soit maintenant \(n \geq 2\) et supposons
\((\mathcal{P}_{n-1})\) vraie. Alors \(Y = \sum_{i = 1}^{n-1}Y_i + Y_n\)
est la somme de deux variables aléatoires indépendantes, dont la
première suit un \(\chi_{n-1}^2\) de densité notée \(f_{n-1}\) et la
seconde un \(\chi_1^2\) de densité notée \(f_1\). Elle possède donc une
densité, égale au produit de convolution de \(f_{n-1}\) et \(f_1\)~:
pour tout \(x\in\mathbb{R}\)
\[f_Y(x) = \int_{\mathbb{R}} f_1(x-u)\,f_{n-1}(u)\,du.\] Or pour tout
\(x,u\in\mathbb{R}\) on a \(f_1(x-u) = 0\) ssi \(x \leq u\) et
\(f_{n-1}(u) = 0\) ssi \(u \leq 0\). Donc leur produit est strictement
positif ssi \(x > u > 0\), nul sinon. La densité \(f_Y\) est par
conséquent nulle sur \(\mathbb{R}^-\) et nous allons maintenant exhiber
son expression sur \(\mathbb{R}_+^\ast\). Soit \(x > 0\), alors
\begin{align*}
f_Y(x) &= \int_{0}^{x} \dfrac{(x-u)^{-\frac{1}{2}}}{\sqrt{2}\,\Gamma\left(\frac{1}{2}\right)}\,\exp\left\{-\dfrac{x-u}{2}\right\}\,\dfrac{u^{\frac{n-1}{2}-1}}{2^{\frac{n-1}{2}}\,\Gamma\left(\frac{n-1}{2}\right)}\,\exp\left\{-\dfrac{u}{2}\right\} \,du\\
&= \dfrac{\exp\left\{-\dfrac{x}{2} \right\}}{2^{\frac{n}{2}}\,\Gamma\left(\frac{1}{2}\right)\Gamma\left(\frac{n-1}{2}\right)} \,\int_0^x \left(x-u\right)^{-\frac{1}{2}}\,u^{\frac{n-3}{2}}\,du.
\end{align*} Avec le changement de variable \(v=\dfrac{u}{x}\) on
obtient \begin{align*}
f_Y(x) &= \dfrac{x^{\frac{n}{2}-1}\,\exp\left\{-\dfrac{x}{2} \right\}}{2^{\frac{n}{2}}\,\Gamma\left(\frac{1}{2}\right)\Gamma\left(\frac{n-1}{2}\right)} \,\int_0^{1} v^{\frac{n-3}{2}}\,\left(1 - v\right)^{-\frac{1}{2}} \,dv.
\end{align*} Or \(f_Y\) est une densité, son intégrale vaut donc \(1\)~:
\begin{align*}
\int_{\mathbb{R}} f_Y(x)\,dx &= \int_0^{+\infty} \dfrac{x^{\frac{n}{2}-1}\,\exp\left\{-\dfrac{x}{2} \right\}}{2^{\frac{n}{2}}\,\Gamma\left(\frac{1}{2}\right)\Gamma\left(\frac{n-1}{2}\right)}\,\left(\int_0^{1} v^{\frac{n-3}{2}}\,\left(1 - v\right)^{-\frac{1}{2}} \,dv\right)\,dx\\
&= \left(\int_0^{1} v^{\frac{n-3}{2}}\,\left(1 - v\right)^{-\frac{1}{2}} \,dv\right)\,\left(\int_0^{+\infty} \dfrac{x^{\frac{n}{2}-1}\,\exp\left\{-\dfrac{x}{2} \right\}}{2^{\frac{n}{2}}\,\Gamma\left(\frac{1}{2}\right)\Gamma\left(\frac{n-1}{2}\right)}\,dx\right)\\
& \left(\int_0^{1} v^{\frac{n-3}{2}}\,\left(1 - v\right)^{-\frac{1}{2}} \,dv\right)\,\left(\int_0^{+\infty} \dfrac{x^{\frac{n}{2}-1}\,\exp\left\{-x \right\}}{\Gamma\left(\frac{1}{2}\right)\Gamma\left(\frac{n-1}{2}\right)}\,dx\right)\\
&= \left(\int_0^{1} v^{\frac{n-3}{2}}\,\left(1 - v\right)^{-\frac{1}{2}} \,dv\right)\,\left( \dfrac{\Gamma\left(\frac{n}{2}\right)}{\Gamma\left(\frac{1}{2}\right)\Gamma\left(\frac{n-1}{2}\right)}\right) = 1.
\end{align*}

On en déduit que
\[ \int_0^{1} v^{\frac{n-3}{2}}\,\left(1 - v\right)^{-\frac{1}{2}} \,dv = \dfrac{\Gamma\left(\frac{1}{2}\right)\Gamma\left(\frac{n-1}{2}\right)}{\Gamma\left(\frac{n}{2}\right)}, \]
et finalement que pour tout \(x\in\mathbb{R}\)
\[f_Y(x) = \left|\begin{array}{ll} \dfrac{x^{\frac{n}{2}-1}}{2^{\frac{n}{2}}\,\Gamma\left(\frac{n}{2}\right)}\,\exp\left\{-\dfrac{x}{2}\right\} & \text{si } x> 0,\\ 0 & \text{sinon;}\end{array}\right.\]
\((\mathcal{P}_n)\) est donc vraie.

\textbf{Conclusion.} La propriété est vraie pour tout
\(n\in\mathbb{N}^\ast\).

\section{Combinaisons linéaires de variables Gaussiennes
indépendantes}\label{II-combinaisons-linuxe9aires-de-variables-gaussiennes-induxe9pendantes-1}

\runinhead{Préliminaires}\phantomsection\addcontentsline{toc}{subsubsection}{Préliminaires}\label{II-answer-CLIGauss-pre}
Cette question a été traitée de manière générale dans le cours. Nous en
proposons une preuve alternative, basée sur le calcul de la fonction de
répartition de la variable aléatoire \(s\,X + m\), qui caractérise sa
loi. Elle dépend clairement des valeurs de \(s\).

\begin{itemize}
\item
  Si \(s = 0\), alors \(s\,X + m\) est toujours égale à \(m\)~: sa loi
  est une masse de Dirac en \(\{m\}\) et pour tout \(x\in\mathbb{R}\) on
  a \(\mathbb{P}\left(s\,X +m \leq x\right) = 1_{[m,+\infty[}(x)\).
\item
  Si \(s\neq 0\), alors pour tout \(x\in\mathbb{R}\) on a
  \[\mathbb{P}\left( s\,X + m \leq x \right) = \left|\begin{array}{ll}\displaystyle\mathbb{P}\left( X \leq \dfrac{x-m}{s} \right) = \int_{-\infty}^{\frac{x-m}{s}} f(u)\,du & \text{si } s>0,\\[1em] \displaystyle \mathbb{P}\left( X \geq \dfrac{x-m}{s} \right) = \int_{\frac{x-m}{s}}^{+\infty} f(u)\,du & \text{si } s<0,\end{array}\right.\]
  qui en posant le changement de variable \(v = s\,x+m\) donne
  \[\mathbb{P}\left( s\,X + m \leq x \right) = \int_{-\infty}^{x} \dfrac{1}{|s|}\,f\left(\dfrac{u-m}{s} \right)\,dx.\]
  Ainsi, \(s\,X+m\) admet une densité, qui pour tout \(x\in\mathbb{R}\)
  est égale à
  \[ \dfrac{1}{|s|}\,f\left(\dfrac{u-m}{s}\right) = \dfrac{1}{\sqrt{2\,\pi}\,|s|}\,\exp\left\{-\dfrac{(x-m)^2}{2\,s^2} \right\}.\]
  On reconnaît la densité d'une loi Normale d'espérance \(m\) et de
  variance \(s^2\).
\end{itemize}

\runinhead{Combinaisons
linéaires}\phantomsection\addcontentsline{toc}{subsubsection}{Combinaisons
linéaires}\label{II-answer-CLIGauss-cl}

\begin{enumerate}
\def\labelenumi{\arabic{enumi}.}
\setcounter{enumi}{1}
\tightlist
\item
  Commençons par supposer que pour tout \(n\in\mathbb{N}^\ast\),
  \(a\in\mathbb{R}^n\) est tel qu'aucune de ses composantes n'est nulle.
  Nous allons montrer par récurrence sur \(n\) que \(S_n^a\) suit une
  loi Normale d'espérance nulle et de variance \(\sum_{i = 1}^n a_i^2\).
  On note cette propriété \((\mathcal{P}_n)\).
\end{enumerate}

\textbf{Initialisation.}

\begin{itemize}
\item
  Si \(n = 1\) et \(a_1 \neq 0\), alors \(S_1^a = a_1\,X_1\) suit une
  loi Normale centrée de variance \(a_1^2\) d'après la question 1;
  \((\mathcal{P}_1)\) est donc vraie.
\item
  Si \(n=2\) et \(a_1,a_2 \neq 0\), alors d'après le cours
  \(S_2^a = a_1\,X_1 + a_2\,X_2\) admet une densité, notée \(f_n^a\),
  égale au produit de convolution des densités \(f_1\) de \(a_1\,X_1\)
  et \(f_2\) de \(a_2\,X_2\). En outre, d'après la question 1, pour tout
  \(x\in\mathbb{R}\) on a
  \(f_1(x) = \dfrac{1}{|a_1|}\,f\left(\dfrac{x}{a_1} \right)\) et
  \(f_2(x) = \dfrac{1}{|a_2|}\,f\left(\dfrac{x}{a_2} \right)\). Par
  conséquent, pour tout \(x\in\mathbb{R}\), \begin{align*}
  f_2^a(x) &= \int_{\mathbb{R}} f_1(x-u)\,f_2(u)\,du = \int_\mathbb{R}\dfrac{1}{|a_1|\,|a_2|}\,f\left(\dfrac{x - u}{a_1}\right)\,f\left(\dfrac{u}{a_2}\right)\,du\\
  &= \dfrac{1}{\sqrt{a_1^2 +a_2^2}}\,f\left(\dfrac{x}{\sqrt{a_1^2+a_2^2}}\right)\, \int_{\mathbb{R}} \dfrac{\sqrt{a_1^2+a_2^2}}{|a_1|\,|a_2|}\,\dfrac{f\left(\dfrac{x - u}{a_1}\right)}{f\left(\dfrac{x}{\sqrt{a_1^2+a_2^2}}\right)}\,f\left(\dfrac{u}{a_2}\right)\,du.
  \end{align*} Or pour tout \(x,u\in\mathbb{R}\) on a \begin{align*}
  \dfrac{f\left(\dfrac{x - u}{a_1}\right)}{f\left(\dfrac{x}{\sqrt{a_1^2+a_2^2}}\right)}\,f\left(\dfrac{u}{a_2}\right) & = \dfrac{1}{\sqrt{2\pi}}\,\dfrac{\exp\left\{-\dfrac{(x-u)^2}{2\,a_1^2} \right\}}{\exp\left\{ -\dfrac{x^2}{2\,(a_1^2+a_2^2)} \right\}}\,\exp\left\{ -\dfrac{u^2}{2\,a_2^2} \right\}\\
  & = \dfrac{1}{\sqrt{2\pi}}\,\exp\left\{-\dfrac{1}{2}\,\left( \dfrac{(x-u)^2}{a_1^2} + \dfrac{u^2}{a_2^2} - \dfrac{x^2}{a_1^2+a_2^2} \right) \right\}\\
  &= \dfrac{1}{\sqrt{2\pi}}\,\exp\left\{ -\dfrac{\left(a_2^2\,x - (a_1^2 + a_2^2)\,u\right)^2}{2\,a_1^2\,a_2^2\,(a_1^2+a_2^2)} \right\}\\
  &= \dfrac{1}{\sqrt{2\pi}}\,\exp\left\{-\dfrac{\left(u - \dfrac{a_2^2\,x}{a_1^2 + a_2^2}\right)^2}{2\,\dfrac{a_1^2\,a_2^2}{a_1^2+a_2^2}}\right\},
  \end{align*} qui multiplié par
  \(\dfrac{\sqrt{a_1^2+a_2^2}}{|a_1|\,|a_2|}\) correspond à la densité
  d'une loi Normale d'espérance \(\dfrac{a_2^2\,x}{a_1^2+a_2^2}\) et de
  variance \(\dfrac{a_1^2\,a_2^2}{a_1^2+a_2^2}\). La précédente
  intégrale vaut donc \(1\) et \((\mathcal{P}_2)\) est vraie.
\end{itemize}

\textbf{Héritage.} Soit maintenant \(n\geq 2\), et supposons
\((\mathcal{P}_{n-1})\) vraie. Alors
\(S_n^a = S_{n-1}^{a_{-n}} + a_n\,X_n\), où
\(a_{-n} := (a_1,\dots,a_{n-1})\). Or \(S_{n-1}^{a_{-n}}\) et
\(a_n\,X_n\) sont des variables gaussiennes centrées, de variances
respectives \(\sum_{i = 1}^{n-1} a_i^2\) d'après \((\mathcal{P}_{n-1})\)
et \(a_n^2\) d'après la question 1. Par \((\mathcal{P}_2)\), \(S_n^a\)
suit donc une loi Normale centrée de variance \(\sum_{i = 1}^n a_i^2\).

\textbf{Conclusion.} On en conclut que \((\mathcal{P}_n)\) est vraie
pour tout \(n\in\mathbb{N}^\ast\).

\begin{enumerate}
\def\labelenumi{\arabic{enumi}.}
\setcounter{enumi}{2}
\tightlist
\item
  Calculons cette covariance~: \begin{align*}
  \text{Cov}\left(S_n^a, S_n^b\right) &= \text{Cov}\left(\sum_{i =1}^n a_i\,X_i, \sum_{j = 1}^n b_j\,X_j\right)\\
  &= \sum_{i = 1}^n a_i\,b_i\,\mathbb{V}\left(X_i\right) + \sum_{1\leq i\neq j \leq n} a_i\,b_j\,\text{Cov}\left(X_i,X_j\right).
  \end{align*} Or par hypothèse \(\mathbb{V}\left(X_i\right) = 1\) pour
  tout \(i\in\{1,\dots,n\}\) et par indépendance on a
  \(\text{Cov}\left(X_i,X_j\right) = 0\) pour tous
  \(i,j \in \{1,\dots,n\}\) tels que \(i\neq j\). Ainsi,
  \[ \text{Cov}\left(S_n^a, S_n^b\right) = \sum_{i = 1}^n a_i\,b_i \]
  qui est nulle ssi \(a\) et \(b\) sont orthogonaux.
\end{enumerate}

\section{Loi de Maxwell}\label{II-loi-de-maxwell-1}

\runinhead{Question
1}\phantomsection\addcontentsline{toc}{subsubsection}{Question
1}\label{II-answer-Maxwell-1}
La densité du triplet \((X, Y, Z)\) étant invariante par rotation, elle
est de la forme \(f_{(X,Y,Z)} (x, y, z) = \phi(x^2 + y^2 + z^2)\). De
plus, les variables \(X, Y\) et \(Z\) étant indépendantes, on a
\[f_{(X,Y,Z)} (x, y, z) = f_X (x)f_Y (y)f_Z (z) \; \mbox{ pour tout $(x, y, z) \in \mathbb{R}^3$.} \]
Soit \((x_0,y_0,z_0) \in \mathbb{R}^3\)\,; comme
\(\phi(x_0^2 + y_0^2 + z_0^2) > 0\) on a
\(f_{(X, Y, Z)}(x_0, y_0, z_0) > 0\) et donc \(f_X(x_0) \neq 0\),
\(f_Y(y_0) \neq 0\) et \(f_Z(z_0) \neq 0\). Pour tout
\(x \in \mathbb{R}^3\), comme \[
f_X(x) f_Y(y_0) f_Z(z_0) = \phi(x^2 + y_0^2 + z^2) = \phi(y_0^2 + x^2 + z_0^2)
= f_X(y_0) f_Y(x) f_Z(z_0),
\] on a l'égalité \[
f_X(x) = \frac{f_X(y_0)}{f_Y(y_0)} f_Y(x).
\] Comme on a par ailleurs \[
\int_{-\infty}^{+\infty} f_X(x) \, dx = \int_{-\infty}^{+\infty} f_Y(x) \, dx = 1,
\] il est nécessaire que le ratio \({f_X(y_0)}/{f_Y(y_0)}\) soit égal à
\(1\). On a donc \(f_X = f_Y\)\,; la preuve que \(f_Y = f_Z\) s'obtient
de manière similaire.

Pour finir, puisque pour tout \(x \in \mathbb{R}\) on a
\(f(x)^3 = f_{(X, Y, Z)}(x, x, x) = \phi(3 x^2) > 0\), la fonction \(f\)
est bien positive en tout point.

\runinhead{Question
2}\phantomsection\addcontentsline{toc}{subsubsection}{Question
2}\label{II-answer-Maxwell-2}
L'égalité
\[f(x)f(y)f(z) = \phi(x^2 + y^2 + z^2 )  \; \text{ pour tout } (x, y, z) \in \mathbb{R}^3\]
implique \(f (x)f'(y)f (z) = 2y\phi' (x^2 + y^2 + z^2 )\) et
\(f' (x)f (y)f (z) = 2x\phi' (x^2 + y^2 + z^2 ).\) On en déduit que
\[xf (x)f' (y)f (z) = f' (x)yf (y)f (z)  \; \text{ pour tout } (x, y, z) \in \mathbb{R}^3.\]

Soit \((x, y_0, z_0) \in \mathbb{R}^3\) avec \(y_0 \neq 0\)\,; on a
\(f(x) > 0\), \(f(y_0) > 0\) et \(f(z_0) > 0\). En posant \[
c = \frac{f'(y_0) f(z_0)}{y_0 f(y_0) f(z_0)},
\] on voit que
\[f' (x) = cxf(x) \; \mbox{pour tout $x \in \mathbb{R}$.}\] On peut
réécrire cette équation différentielle sous la forme \[
(\ln f(x))' = cx.
\] En intégrant les deux membres de l'équation entre \(0\) et \(x\), on
obtient \[
\ln f(x) - \ln f(0) = c\frac{x^2}{2}
\] et donc \[
f(x) = f(0) e^{c x^2/2},
\] ce qui est bien la forme cherchée \(f(x) = a e^{cx^2/2}\).

\runinhead{Question
3}\phantomsection\addcontentsline{toc}{subsubsection}{Question
3}\label{II-answer-Maxwell-3}
Comme \(\int_\mathbb{R}f (x) dx = 1\), on a nécessairement \(c < 0\) et
on en déduit que
\[f (x) = \frac{1}{\sqrt{2\pi \sigma^2}}\exp\left(-\frac{x^2}{2\sigma^2}\right).\]
Les variables aléatoires \(X\), \(Y\) et \(Z\) suivent donc loi
gaussienne centrée et la covariance du vecteur est \(\sigma^2 I_3\), où
\(I_3\) est la matrice identité.

\runinhead{Question
4}\phantomsection\addcontentsline{toc}{subsubsection}{Question
4}\label{II-answer-Maxwell-4}
On a
\[\frac{1}{2} m\mathbb{E}(\|V \|^2 ) = \frac{1}{2} m\mathbb{E}(V_1^2 +V_2^2 + V_3^2).\]
Comme \(\mathbb{E}(V_i^2) = \sigma^2\), on en déduit que
\(E_c = \frac{3}{2}m \sigma^2\) et par conséquent que
\[\sigma^2 = \frac{kT}{m}.\]

\runinhead{Question
5}\phantomsection\addcontentsline{toc}{subsubsection}{Question
5}\label{II-answer-Maxwell-5}
À l'aide de la méthode de la fonction muette, on montre que \(V_1^2\)
suit une loi gamma de paramètres \((1/2,\frac{1}{2\sigma^2})\). En
effet, en appliquant le changement de variable \(y=x^2\) sur
\(\mathbb{R}_-^\ast\) et \(\mathbb{R}_+^\ast\), on obtient~:
\[f_{V_1^2} (y) = \left(f(-\sqrt{y}) + f(\sqrt{y})\right) \frac{1}{2\sqrt{y}} 1_{]0,+\infty[}(x) = f(\sqrt{y}) \frac{1}{\sqrt{y}} 1_{]0,+\infty[}(x)\]
par parité de \(f\), soit \begin{align*}
f_{V_1^2} (y) &= \frac{1}{\sqrt{\pi}}\left(\frac{1}{2\sigma^2}\right)^{1/2} y^{-1/2} \exp\left(-\frac{y}{2\sigma^2}\right)\\
              &= \frac{1}{\Gamma(1/2)}\left(\frac{1}{2\sigma^2}\right)^{1/2} y^{-1/2} \exp\left(-\frac{y}{2\sigma^2}\right)
\end{align*} On déduit de l'indication que
\(\|V\|^2 = V_1^2 +V_2^2 + V_3^2\) suit une loi gamma de paramètres
\((3/2,\frac{1}{2\sigma^2})\). Sa densité est
\[f_{\|V\|^2}(x) = \frac{1}{\sqrt{2\pi}\sigma^3}\sqrt{x} e^{- x/2\sigma^2}1_{]0,+\infty[}(x).\]

À l'aide de la méthode de la fonction muette, en appliquant le
changement de variable \(v = \sqrt{x}\) sur \(\mathbb{R}_+^\ast\) et
avec \(\sigma^2 = \frac{kT}{m}\), on montre enfin que \(\|V\|\) est une
variable de densité~:

\[f_{\|V\|}(v) = \frac{\sqrt{2}}{\sqrt{\pi}}\left(\frac{m}{kT}\right)^{3/2}v^2 e^{- mv^2/2kT} 1_{]0,+\infty[}(v).\]

Il s'agit de la densité de la loi de Maxwell.

\chapter{Solutions}

\runinhead{Dans un triangle
(1)}\phantomsection\addcontentsline{toc}{subsubsection}{Dans un triangle
(1)}\label{III-answer-etb1}
La densité marginale de \(X\) est donnée par
\(f_X(x) = \int f_{X,Y}(x,y) dy = 1_{]0,1[}(x)\) et pour
\(x \in ]0,1[\), \[f_{Y|X=x} (y) = \frac{1}{x} 1_{]0,x[}(y)\] Ainsi
\(X\) est uniformément distribué sur \(]0,1[\), et la loi de \(Y\)
sachant \(X =x\) est uniforme sur \(]0,x[\) pour \((0 < x < 1)\).

\runinhead{Dans un triangle
(2)}\phantomsection\addcontentsline{toc}{subsubsection}{Dans un triangle
(2)}\label{III-answer-etb2}
Pour un tel \(x\), l'espérance conditionnelle \(\mathbb{E}(Y|X=x)\) vaut
ainsi \(x/2\) et nous obtenons \(\mathbb{E}(Y|X) = \frac{X}{2}\).

\runinhead{Auto-conditionnement}\phantomsection\addcontentsline{toc}{subsubsection}{Auto-conditionnement}\label{III-answer-autocond}
On a \(\psi(y) = \mathbb{E}(Y|Y=y) = y\) et donc
\(\mathbb{E}(Y|Y) = \psi(Y) = Y\) p.s.

\runinhead{Espérance conditionnelle d'un produit de
variables}\phantomsection\addcontentsline{toc}{subsubsection}{Espérance
conditionnelle d'un produit de variables}\label{III-answer-prod}
On a \(\mathbb{E}(XY|X = x) = x\mathbb{E}(Y|X=x)\), d'où
\(\mathbb{E}(XY|X) = X\mathbb{E}(Y|X)\) p.s.

\runinhead{Une deuxième
approche}\phantomsection\addcontentsline{toc}{subsubsection}{Une
deuxième approche}\label{III-answer-detail}
On remarque qu'on peut réécrire l'erreur comme
\[\ell(a,b)= \mathbb{E}(Y^2) -2a\mathbb{E}(XY) -2b \mathbb{E}(Y) +a^2\mathbb{E}(X^2) +2ab\mathbb{E}(X) +b^2.\]
Cette fonction est convexe et par conséquent admet un minimum global
atteint au point \((a^*, b^*)\in\mathbb{R}^2\) qui annule les dérivées
dérivées partielles de \(\ell\) en \(a\) et \(b\). Ainsi, on a
\[\frac{\partial \ell(a^*,b^*)}{\partial b} = -2\mathbb{E}(Y^2) + 2a^* \mathbb{E}(X) + 2b^*=0\]
d'où \(b^* = \mathbb{E}(Y) - a^* \mathbb{E}(X)\) Et, \[\begin{aligned}
\frac{\partial \ell(a^*,b^*)}{\partial a} & = -2\mathbb{E}(XY) + 2a \mathbb{E}(X^2) + 2b \mathbb{E}(X)\\
                                   & = -2\mathbb{E}(XY) + 2a \mathbb{E}(X^2) + 2\mathbb{E}(X)\mathbb{E}(Y) - 2 a^* \mathbb{E}(X^2)\\
                                   & = -2\text{Cov}(X,Y) + a^*\mathbb{V}(X)=0
\end{aligned}\]

d'où
\(a^* = \frac{\text{Cov}(X,Y)}{\mathbb{V}(X)} = \rho(X,Y)\frac{\sigma_Y}{\sigma_X}.\)

\runinhead{Variance
totale}\phantomsection\addcontentsline{toc}{subsubsection}{Variance
totale}\label{III-answer-vartot}

\begin{align*}
\mathbb{V}(Y) =& \mathbb{E}((Y-\mathbb{E}(Y))^2) = \mathbb{E}(\mathbb{E}((Y-\mathbb{E}(Y))^2|X)) \text{   par la formule de l'espérance totale}\\
=& \mathbb{E}(\mathbb{E}((Y-\mathbb{E}(Y|X)+\mathbb{E}(Y|X)-\mathbb{E}(Y))^2|X)) \\
=& \mathbb{E}(\mathbb{E}((Y-\mathbb{E}(Y|X))^2|X)) + \mathbb{E}(\mathbb{E}((\mathbb{E}(Y|X)-\mathbb{E}(Y))^2|X)) \\
 &+ 2\mathbb{E}(\mathbb{E}((Y-\mathbb{E}(Y|X))(\mathbb{E}(Y|X)-\mathbb{E}(Y))|X))\\
=& \mathbb{E}(\mathbb{V}(Y|X)) + \mathbb{E}((\mathbb{E}(Y|X)-\mathbb{E}(Y))^2) + 2\mathbb{E}((\mathbb{E}(Y|X)-\mathbb{E}(Y))\mathbb{E}((Y-\mathbb{E}(Y|X))|X))\\
=& \mathbb{E}(\mathbb{V}(Y|X)) + \mathbb{V}(\mathbb{E}(Y|X)) \text{   car    } \mathbb{E}((Y-\mathbb{E}(Y|X))|X) = 0 
\end{align*}

\section{Couple de variables}\label{III-couple-de-variables-1}

\runinhead{Question
1}\phantomsection\addcontentsline{toc}{subsubsection}{Question
1}\label{III-answer-ex1-1}
On voit d'abord que la densité du couple \((X,Y)\) vaut~:
\[f_{X,Y}(x,y) = f_{X|Y=y}(x)f_Y(y) = xe^{-xy}1_{\mathbb{R}_+}(x) 1_{[1,+\infty[} (y)\]
Soit \(h\) une fonction continue bornée sur \(\mathbb{R}^2_+\). Le
changement de variable \((x, y) \mapsto (t = xy , y)\) de jacobien
\(y\), donne alors que
\[\mathbb{E}(h(T, Y )) = \mathbb{E}(h(XY, Y )) = \int_1^{+\infty}\int_0^{+\infty} h(t,y) e^{-t} \frac{t}{y^2} dt dy\]
et donc la densité du couple \((T, Y)\) vaut
\[f_{T,Y} = e^{-t} \frac{t}{y^2}1_{[1,+\infty[} (y)1_{\mathbb{R}_+}(t)\]
Elle s'écrit comme produit d'une fonction de \(t\) et d'une fonction de
\(y\). On en déduit que \(T\) et \(Y\) sont indépendantes et que \(T\) a
pour densité \(te^{-t} 1_{\mathbb{R}_+}(t)\).

\runinhead{Question
2}\phantomsection\addcontentsline{toc}{subsubsection}{Question
2}\label{III-answer-ex1-2}
La loi marginale de \(X\) a pour densité
\(f_X (x) = \int_1^{+\infty} x e^{-xy} dy = e^{-x}\). Ainsi \(X\) suit
une loi exponentielle de paramètre 1 et la loi conditionnelle de \(Y\)
sachant \(X = x\) admet la densité~:
\[f_{Y|X=x}(y) = \frac{f_{X,Y}(x,y)}{f_X(x)} = x e^{-x(y-1)} 1_{[1,+\infty[} (y)\]
pour x \textgreater{} 0.

\runinhead{Question
3}\phantomsection\addcontentsline{toc}{subsubsection}{Question
3}\label{III-answer-ex1-3}
On en déduit que
\[\mathbb{E}(Y|X=x ) = \int_1^{+\infty} y x e^{-x(y-1)} dy = \frac{x+1}{x}1_{\mathbb{R}_+}(x)\]
par intégration par parties. Ainsi \(\mathbb{E}(Y|X) = \frac{X+1}{X}\).

\section{Mélanges de lois}\label{III-muxe9langes-de-lois-1}

\runinhead{Question
1}\phantomsection\addcontentsline{toc}{subsubsection}{Question
1}\label{III-answer-melloi1}
Soit \(B\) un borélien. Par indépendance de \(K\) avec \(X_i\), on a
\[\mathbb{P}(X \in B \mid K = i) = \mathbb{P}(X_i \in B \mid K = i) = \mathbb{P}(X_i \in B).\]
La loi de \(X\) sachant \(\{K = i\}\) est donc la même que celle de
\(X_i\), d'où
\[f_{X\mid K = i} : x\in\mathbb{R}\mapsto f_i(x) = \dfrac{1}{\sqrt{2\pi}\sigma_i}\,\exp\left\{- \dfrac{(x-m_i)^2}{2\sigma_i^2} \right\}.\]

\runinhead{Question
2}\phantomsection\addcontentsline{toc}{subsubsection}{Question
2}\label{III-answer-melloi2}
Soit \(B\) un borélien. D'après la formule des probabilités totales et
la question précédente, on a
\[\mathbb{P}(X\in B) = \sum_{i = 1}^n p_i\,\mathbb{P}(X \in B \mid K = i) = \sum_{i = 1}^n p_i\,\mathbb{P}(X_i \in B).\]
La variable aléatoire \(X\) admet donc une densité, qui vaut
\[f_X : x\in\mathbb{R}\mapsto \sum_{i = 1}^n p_i\,f_i(x).\]

\runinhead{Question
3}\phantomsection\addcontentsline{toc}{subsubsection}{Question
3}\label{III-answer-melloi3}
D'après la question précédente, \(X\) a pour espérance \begin{align*}
\mathbb{E}(X) &= \int_\mathbb{R}x\,f_X(x)\,dx = \int_\mathbb{R}x \sum_{i = 1}^n p_i\, f_i(x)\,dx =  \sum_{i = 1}^n p_i \int_\mathbb{R}x\,f_i(x)\,dx\\
& = \sum_{i = 1}^n p_i\,m_i.
\end{align*} Quant à la variance de \(X\), en utilisant l'égalité
\(\sum_{i= 1}^n p_i = 1\), elle vaut \begin{align*}
\mathbb{V}(X) &= \mathbb{E}\left(X^2\right) - \mathbb{E}(X)^2 = \int_\mathbb{R}x^2\,f_X(x)\,dx - \left(\sum_{i = 1}^n p_i\,m_i\right)^2\\
&= \sum_{i = 1}^n p_i (\sigma_i^2 + m_i^2) - \sum_{i = 1}^n p_i\left(\sum_{j = 1}^n p_j\,m_j\right)^2\\
&= \sum_{i = 1}^n p_i\,\sigma_i^2 + \sum_{i = 1}^n p_i \left(m_i - \sum_{j=1}^n p_j\,m_j \right)^2.
\end{align*} On retrouve bien la forme désirée, avec la dispersion des
espérances
\[\bar{\sigma}^2 := \sum_{i = 1}^n p_i \left(m_i - \sum_{j=1}^n p_j\,m_j \right)^2.\]

\runinhead{Question
4}\phantomsection\addcontentsline{toc}{subsubsection}{Question
4}\label{III-answer-melloi4}
Si l'on souhaite approcher la loi de \(X\) avec une unique Gaussienne,
et non un mélange, les questions précédentes suggèrent de prendre celle
d'espérance \(\sum_{i =1}^n p_i\,m_i\) et de variance
\(\sum_{i = 1}^n p_i\,\sigma^2_i + \bar{\sigma}^2\). Voir figure
ci-dessous.

\begin{figure}
\centering
\includegraphics{PdfMelGauss.tex.pdf}
\caption{Illustration}
\end{figure}

\section{Lois conjuguées}\label{III-answer-loiconj}

On considère dans tout cet exercice \(B_1\) et \(B_2\) des Boréliens.

\runinhead{Question
1}\phantomsection\addcontentsline{toc}{subsubsection}{Question
1}\label{III-answer-loiconj-expexp}
D'après les hypothèses on a \begin{align*}
\mathbb{P}_{X,Y}(B_1\times B_2) &= \int_{B_2} \left(\int_{B_1} d\mathbb{P}_{X\mid Y = y}(x) \right) d\mathbb{P}_Y(y) \hspace{1em}\text{par Fubini conditionnel,}\\
&= \int_{B_2} \left(\int_{B_1} y\,e^{-y x}\,1_{\mathbb{R}_+^\ast}(x)\,dx \right) \lambda\,e^{-\lambda y}\,1_{\mathbb{R}_+^\ast}(y)\,dy\\
&= \int_{B_1} \int_{B_2} \lambda\, y\,e^{-(x+\lambda)\,y}1_{\mathbb{R}_+^\ast}(x)\,1_{\mathbb{R}_+^\ast}(y) \,dy\, dx \hspace{1em}\text{par Fubini.}
\end{align*} Le vecteur aléatoire \((X,Y)\) possède donc une densité
jointe
\[f_{X,Y} : (x,y) \in\mathbb{R}^2 \mapsto \lambda\, y\,e^{-(x+\lambda)\,y}1_{\mathbb{R}_+^\ast}(x)\,1_{\mathbb{R}_+^\ast}(y).\]
La variable aléatoire \(X\) a donc aussi une densité~: pour tout
\(x\in\mathbb{R}\) \begin{align*}
f_X(x) &= \int_{\mathbb{R}} f_{X,Y}(x,y)\,dy = \int_\mathbb{R}\lambda\, y\,e^{-(x+\lambda)\,y}1_{\mathbb{R}_+^\ast}(x)\,1_{\mathbb{R}_+^\ast}(y)\,dy\\
&= \left|\begin{array}{ll}\displaystyle \dfrac{\lambda}{x+\lambda} \int_{0}^{+\infty} y\,(x+\lambda)\,e^{-(x+\lambda)\,y}\,dy & \text{si } x>0,\\[1em] 0 & \text{sinon.} \end{array}\right.
\end{align*} On reconnaît dans cette dernière intégrale la formule de
l'espérance d'une loi Exponentielle de paramètre \(x+\lambda\), et on en
déduit que pour tout \(x\in\mathbb{R}\)
\[f_X(x) = \dfrac{\lambda}{(x+\lambda)^2}\,1_{\mathbb{R}_+^\ast}(x).\]
Pour tout \(x\in\mathbb{R}_+^\ast\) la variable \(Y\) sachant
\(\{X=x\}\) admet donc aussi une densité, que l'on explicite avec la
formule de Bayes~: pour tout \(y\in\mathbb{R}\) \begin{align*}
f_{Y\mid X=x}(y) &= \dfrac{f_{X,Y}(x,y)}{f_X(x)} = \dfrac{\lambda\, y\,e^{-(x+\lambda)\,y}\,1_{\mathbb{R}_+^\ast}(y)}{\dfrac{\lambda}{(x+\lambda)^2}}\\
&= (x+\lambda)^2\,y\,e^{-(x+\lambda)\,y}\,1_{\mathbb{R}_+\ast}(y).
\end{align*} Comme \(\Gamma(2) = 1\), on reconnaît ici la densité d'une
loi Gamma d'indice \(2\) et de paramètre d'échelle \(x+\lambda\).

\runinhead{Question
2}\phantomsection\addcontentsline{toc}{subsubsection}{Question
2}\label{III-answer-loiconj-gampoi}
D'après les hypothèses, en procédant comme précédemment, on a
\begin{align*}
\mathbb{P}_{X,Y}(B_1\times B_2) &= \int_{B_2} \left(\int_{B_1} d\mathbb{P}_{X\mid Y = y}(x) \right) d\mathbb{P}_Y(y)\\
&= \int_{B_2} \left(\sum_{x\in B_1} \dfrac{y^{x}}{x!}\,e^{-y}\,1_{\mathbb{N}}(x) \right) \dfrac{\theta^{\alpha}}{\Gamma(\alpha)}\,y^{\alpha-1}e^{-\theta y}\,1_{\mathbb{R}_+}(y)\,dy\\
&= \sum_{x\in B_1\cap\mathbb{N}} \left(\dfrac{1}{x!} \int_{B_2\cap\mathbb{R}_+} \dfrac{\theta^{\alpha}}{\Gamma(\alpha)}\,y^{x+\alpha-1}e^{-(\theta+1) y}\,dy\right)\\
&= \sum_{x\in B_1\cap\mathbb{N}} \biggl(\dfrac{\Gamma(x+\alpha)\,\theta^\alpha}{x!\,\Gamma(\alpha)\,(\theta+1)^{x+\alpha}}\\
&\hspace{5em} \times \int_{B_2\cap\mathbb{R}_+} \dfrac{(\theta+1)^{x+\alpha}}{\Gamma(x+\alpha)}\,y^{x+\alpha-1}e^{-(\theta+1) y}\,dy\biggr).
\end{align*} On reconnaît dans cette dernière intégrale la densité d'une
loi Gamma d'indice \(x+\alpha\) et de paramètre d'échelle \(\theta+1\),
qui correspond exactement à la loi conditionnelle de \(Y\) sachant
\(\{X=x\}\) pour \(x\in\mathbb{N}\). En effet, on a d'une part
\begin{align*}
\mathbb{P}_X(B_1) &= \mathbb{P}_{X,Y}(B_1 \times \mathbb{R}) = \sum_{x\in B_1\cap\mathbb{N}} \dfrac{\theta^\alpha}{\Gamma(\alpha)}\, \dfrac{\Gamma(x+\alpha)}{x!\,(\theta+1)^{x+\alpha}},
\end{align*} ce qui donne bien pour tout \(x\in\mathbb{N}\)~:
\begin{align*}
\mathbb{P}_{Y\mid X=x} (B_2) &= \mathbb{P}\left(Y \in B_2 \mid X = x\right) = \dfrac{\mathbb{P}_{X,Y}\left(\{x\}\times B_2\right)}{\mathbb{P}_X(\{x\})}\\
&= \int_{B_2\cap\mathbb{R}_+} \dfrac{(\theta+1)^{x+\alpha}}{\Gamma(x+\alpha)}\,y^{x+\alpha-1}\,e^{-(\theta+1)y}\,dy.
\end{align*}

\section{Randomisation}\label{III-answer-randomize}

En termes probabilistes et selon les notations de l'exercice, il s'agit
de calculer \(\mathbb{E}(T_N - T_0)\), où la variable aléatoire \(T_N\)
peut s'écrire en fonction d'une somme aléatoire de variables aléatoires
indépendantes~: \[T_N = \sum_{i = 1}^N X_i + T_0.\] Comme la boutique
ferme au bout d'un certain temps, toutes les variables aléatoires
figurant dans l'équation précédente sont bornées, donc intégrables. On
peut ainsi calculer \(\mathbb{E}(T_N-T_0)\) à l'aide de la formule de
l'espérance totale~:
\[\mathbb{E}\left(T_N - T_0\right) = \mathbb{E}\left(\mathbb{E}\left(T_N \mid N\right)\right) - T_0.\]
Pour tout \(n\in\mathbb{N}^\ast\) l'énoncé suggère que \(N\) est
indépendante de \(X_1,\dots,X_n\), elles-mêmes indépendantes et de même
loi que \(X\), d'où~:
\[\mathbb{E}\left(T_n \mid N = n\right) = \sum_{i = 1}^n \mathbb{E}(X_i\mid N = n) = \sum_{i = 1}^n \mathbb{E}(X_i) = n\mathbb{E}(X).\]
Ainsi, en posant
\(\psi : n \in\mathbb{N}^\ast \mapsto n \mathbb{E}(X)\), on obtient
\[\mathbb{E}\left(T_N - T_0\right) = \mathbb{E}\left(\psi(N)\right) - T_0 = \mathbb{E}(N)\mathbb{E}(X) - T_0.\]
C'était prévisible~: en posant arbitrairement \(T_0 = 0\), le temps
d'attente moyen est le temps d'attente moyen entre deux arrivées,
multiplié par le rang moyen du gagnant. Si la loterie dépendait des
temps d'arrivées, par exemple en faisant gagner le premier client qui
arrive au moins 10 minutes après le client précédent, \(\psi\), et donc
le résultat, seraient différents.

\section{Etats cachés --- indépendance
conditionnelle}\label{III-etats-cachuxe9s-induxe9pendance-conditionnelle-1}

\runinhead{Question
1}\phantomsection\addcontentsline{toc}{subsubsection}{Question
1}\label{III-answer-ec1}
Soit \(n\in\mathbb{N}^\ast\). Quels que soient \(x\in\mathbb{R}\) et
\(B\) borélien on a \begin{align*}
\mathbb{P}_{Y_n \mid X_n = x}(B) &= \mathbb{E}\left(1_B(X_n+\epsilon_n) \mid X_n = x \right)\\
&= \int_\mathbb{R}1_B(x+y)\,d\mathbb{P}_{\epsilon_n\mid X_n = x}(y)\\
& = \int_\mathbb{R}1_B(x+y)\,f_\epsilon(y)\,dy \hspace{1em}\text{par indépendance de $X_n$ et $\epsilon_n$}\\
&= \int_B f_\epsilon(y-x)\,dy.
\end{align*} Ainsi, \(\mathbb{P}_{Y_n\mid X_n =x}\) admet bien une
densité~:
\[f_{Y_n\mid X_n=x} : y \in\mathbb{R}\mapsto f_\epsilon(y - x).\]

\runinhead{Question
2}\phantomsection\addcontentsline{toc}{subsubsection}{Question
2}\label{III-answer-ec2}
Soient \(n\in\mathbb{N}^\ast\), \((x_1,\dots,x_n)\in\mathbb{R}^n\) et
\(B_1,\dots,B_n\) des boréliens. Pour simplifier les écritures, on note
\(x_{1:n}\) tout vecteur \((x_1,\dots,x_n)\) de \(\mathbb{R}^n\). Alors
\begin{align*}
&\mathbb{P}_{Y_{1:n}\mid X_{1:n}=x_{1:n}}(B_1\times\dots\times B_n)
=\mathbb{E}\left(\prod_{i = 1}^n 1_{B_i}(X_i + \epsilon_i) \Bigm\vert X_{1:n} = x_{1:n} \right)\\
&= \int_{\mathbb{R}^n} \prod_{i = 1}^n 1_{B_i}(x_i + y_i)\,d\mathbb{P}_{\epsilon_{1:n}\mid X_{1:n} = x_{1:n}}(y_{1:n})\\
&= \int_{\mathbb{R}^n} \prod_{i = 1}^n 1_{B_i}(x_i + y_i)\,d\mathbb{P}_{\epsilon_{1:n}}(y_{1:n}) \ \text{par indépendance des $\epsilon_i$ et $X_j$,}\\
&= \prod_{i = 1}^n \int_{\mathbb{R}} 1_{B_i}(x_i + y_i)\,f_\epsilon(x_i)\,dy_{i} \ \text{par Fubini et indépendance et même loi des $\epsilon_i$,}\\
&= \prod_{i = 1}^n \int_{\mathbb{R}} 1_{B_i}(y_i)\,f_\epsilon(y_i - x_i)\,dy_i\\
& = \prod_{i = 1}^n \int_{\mathbb{R}} 1_{B_i}(y_i)\,f_{Y_i\mid X_i = x_i}(y_i)\,dy_i\ \text{par la question 1,}\\
&= \prod_{i = 1}^n \mathbb{P}_{Y_i\mid X_i = x_i}(B_i).
\end{align*} Les \(n\) relevés de température sont donc bien
indépendants conditionnellement aux états cachés.

\section{Covariance totale}\label{III-answer-covtot}

Tout d'abord, par linéarité de l'espérance conditionnelle on a~:
\begin{align*}
\text{Cov}(X,Y \mid Z) &= \mathbb{E}\Bigl( \bigl( X - \mathbb{E}(X\mid Z) \bigr)\bigl( Y - \mathbb{E}(Y\mid Z) \bigr) \Bigm| Z  \Bigr)\\
&= \mathbb{E}\Bigl( XY - X\mathbb{E}(Y\mid Z) - Y\mathbb{E}(X\mid Z) + \mathbb{E}(X\mid Z)\mathbb{E}(Y\mid Z) \Bigm| Z  \Bigr)\\
&= \mathbb{E}(XY \mid Z) - \mathbb{E}(X\mid Z)\mathbb{E}(Y\mid Z).
\end{align*}

En utilisant la formule de l'espérance totale et la linéarité de
l'espérance, on obtient alors \begin{align*}
\text{Cov}(X,Y) &= \mathbb{E}(XY) - \mathbb{E}(X)\mathbb{E}(Y)\\
&= \mathbb{E}\bigl( \mathbb{E}(XY \mid Z) \bigr) - \mathbb{E}\bigl( \mathbb{E}(X \mid Z) \bigr)\mathbb{E}\bigl( \mathbb{E}(Y \mid Z) \bigr)\\
&= \mathbb{E}\bigl( \mathbb{E}(XY \mid Z) - \mathbb{E}(X\mid Z)\mathbb{E}(Y\mid Z) \bigr)\\
&\ \ \  + \mathbb{E}\bigl( \mathbb{E}(X\mid Z)\mathbb{E}(Y\mid Z) \bigr) - \mathbb{E}\bigl( \mathbb{E}(X \mid Z) \bigr)\mathbb{E}\bigl( \mathbb{E}(Y \mid Z) \bigr)\\
&= \mathbb{E}\bigl(\text{Cov}(X,Y\mid Z)\bigr) + \text{Cov}\bigl( \mathbb{E}(X\mid Z), \mathbb{E}(Y\mid Z) \bigr).
\end{align*}

\section{Débruitage (formule de
Tweedie)}\label{III-duxe9bruitage-formule-de-tweedie-1}

\runinhead{Question
1}\phantomsection\addcontentsline{toc}{subsubsection}{Question
1}\label{III-answer-tweedie1}
\(f_Y(y) = \int_{\mathbb{R}^d} f_{X,Y}(x,y) dx = \int_{\mathbb{R}^d} f_{Y|X=x}(y) f_X(x) dx = \int_{\mathbb{R}^d} f_Z(y-x) f_X(x) dx\)
où
\(f_Z(y-x) = f_{Y|X=x}(y) = (2 \pi)^{-d/2} (\det \Sigma)^{1/2}\exp\left(-\frac12 (y-x)^t\Sigma^{-1}(y-x)\right)\).
On notera que l'on retrouve le produit de convolution de \(f_X\) et
\(f_Z\) comme densité de \(X+Z\).

\runinhead{Question
2}\phantomsection\addcontentsline{toc}{subsubsection}{Question
2}\label{III-answer-tweedie2}
On remarque d'abord que \(\nabla \log f_Y = \frac{\nabla f_Y}{f_Y}\).
Ensuite, par application du résultat de l'exercice B.8.5 du poly de
maths 1 (intégrales à paramètres) dont on vérifie aisément les
hypothèses, on a
\[\nabla f_Y(y) = \int_{\mathbb{R}^d} \nabla f_{Y|X=x}(y) f_X(x) dx\] et
\begin{align*}
\nabla f_{Y|X=x}(y) &= (2 \pi)^{-d/2} (\det \Sigma)^{1/2}\exp\left(-\frac12 (y-x)^t\Sigma^{-1}(y-x)\right) \Sigma^{-1}(x-y)\\
&= f_{Y|X=x}(y) \Sigma^{-1}(x-y)
\end{align*} Finalement,
\[\nabla \log f_Y(y) = \frac{\int_{\mathbb{R}^d} f_{Y|X=x}(y) \Sigma^{-1}(x-y) f_X(x) dx}{f_Y(y)}\]

\runinhead{Question
3}\phantomsection\addcontentsline{toc}{subsubsection}{Question
3}\label{III-answer-tweedie3}
En utilisant l'égalité
\(f_{X,Y}(x,y) = f_{Y|X=x}f_X(x) = f_{X|Y=y}(x)f_Y(y)\), on obtient
\begin{align*}
\nabla \log f_Y(y) &= \frac{\int_{\mathbb{R}^d} f_{Y|X=x}(y) \Sigma^{-1}(x-y) f_X(x) dx}{f_Y(y)}\\
&= \int_{\mathbb{R}^d} f_{X|Y=y}(x) \Sigma^{-1}(x-y) dx \\
& = \Sigma^{-1}\left(\int_{\mathbb{R}^d} x f_{X|Y=y}(x) dx -y \right)
\end{align*} d'où \(\mathbb{E}(X|Y=y) = y + \Sigma\nabla\log f_Y(y)\).

\section{Non-réponse}\label{III-non-ruxe9ponse-1}

\runinhead{Question
1}\phantomsection\addcontentsline{toc}{subsubsection}{Question
1}\label{III-answer-nonrep1}
L'espérance conditionnelle de \(Y\) sachant \(X\) peut s'écrire comme la
solution au problème de minimisation
\[\min_{\phi(X)\in L^2_X} \mathbb{E}\left(\left(Y-\phi(X)\right)^2\right).\]
Or pour \(\phi(X)\in L^2_X\) on a ici
\[\mathbb{E}\left(\left(Y-\phi(X)\right)^2\right) = \mathbb{E}\left( \left(Y - \phi(1)\right)^2 1_{\{1\}}(X) \right) + \mathbb{E}\left( \left(Y - \phi(0)\right)^2 1_{\{0\}}(X) \right),\]
il suffit donc de résoudre pour tout \(x\in\{0,1\}\)
\[\min_{\lambda \in \mathbb{R}} \mathbb{E}\left(\left(Y-\lambda\right)^2 1_{\{x\}}(X)\right).\]
Soit \(x\in\{0,1\}\) et posons
\(J_x : \lambda\in\mathbb{R}\mapsto \mathbb{E}\left(\left(Y-\lambda\right)^2 1_{\{x\}}(X)\right)\).
Alors pour tout \(\lambda\in\mathbb{R}\)
\[J_x(\lambda) = \mathbb{E}\left(Y^21_{\{x\}}(X)\right) + \lambda^2\,\mathbb{P}(X=x) -2\lambda\,\mathbb{E}\left(Y1_{\{x\}}(X)\right)\]
et sa dérivée
\[J_x^\prime(\lambda) = 2\lambda\,\mathbb{P}(X=x) -2\,\mathbb{E}\left(Y1_{\{x\}}(X)\right)\]
s'annule en
\[\lambda_x := \dfrac{\mathbb{E}\left(Y1_{\{x\}}(X)\right)}{\mathbb{P}(X=x)} = \mathbb{E}(Y\mid X = x).\]
On en conclut que
\[\mathbb{E}(Y\mid X) = \mathbb{E}(Y\mid X = 1)1_{\{1\}}(X) + \mathbb{E}(Y\mid X = 0)1_{\{0\}}(X).\]
Or on remarque que \(1_{\{1\}}(X) = X\) et \(1_{\{0\}}(X) = 1 - X\), ce
qui fait de \(\mathbb{E}(Y\mid X)\) une fonction affine de \(X\). Elle
est par définition la meilleure approximation de \(Y\) par une fonction
de \(X\), elle coïncide donc avec l'approximation affine de \(Y\) par
\(X\):
\[\mathbb{E}(Y\mid X) = m + \dfrac{\rho\,\sigma}{\sqrt{p(1-p)}}\,(X - p).\]

\runinhead{Question
2}\phantomsection\addcontentsline{toc}{subsubsection}{Question
2}\label{III-answer-nonrep2}
D'après la question précédente, on a
\(\mathbb{E}(Y\mid X) = m_0 + (m_1 - m_0) X\), la meilleure
approximation affine de \(Y\) par \(X\). Ainsi, \(m_0\) et \(m_1\)
satisfont
\[\left|\begin{array}{l} m_1 - m_0 = \dfrac{\rho\sigma}{\sqrt{p(1-p)}},\\[1em] m_0 = m - (m_1-m_0)p,  \end{array}\right. \Leftrightarrow \left|\begin{array}{l} m_1 = m + \rho\sigma\sqrt{\dfrac{1-p}{p}},\\[1em] m_0 = m - \rho\sigma\sqrt{\dfrac{p}{1-p}}.  \end{array}\right.\]

\runinhead{Question
3}\phantomsection\addcontentsline{toc}{subsubsection}{Question
3}\label{III-answer-nonrep3}
Par la formule de la variance totale et d'après la question 1, on a
\begin{align*}
\sigma^2 &= \mathbb{V}\left(Y\right) = \mathbb{E}\bigl(\mathbb{V}\left(Y\mid X\right)\bigr) + \mathbb{V}\bigl(\mathbb{E}(Y\mid X)\bigr)\\
&= p\,\sigma^2_1 + (1-p)\,\sigma^2_0 + \dfrac{\rho^2\sigma^2}{p\,(1-p)}\mathbb{V}(X)\\
&= p\,\sigma^2_1 + (1-p)\,\sigma^2_0 + \rho^2\sigma^2.
\end{align*} Cette égalité se simplifie et donne bien
\[\sigma^2 = \dfrac{(1-p)\,\sigma_0^2 + p\,\sigma_1^2}{1-\rho^2}.\]

\runinhead{Question
4}\phantomsection\addcontentsline{toc}{subsubsection}{Question
4}\label{III-answer-nonrep4}
Lorsque \(X\) et \(Y\) sont non corrélées, i.e.~\(\rho = 0\), on obtient
\(m_0 = m_1 = m\) puis \(\sigma^2 = (1-p)\,\sigma_0^2 + p\,\sigma_1^2\).
En d'autres termes, \(\mathbb{E}(Y\mid X) = m\) est une variable
aléatoire constante, et
\(\mathbb{E}\bigl(\mathbb{V}(Y\mid X)\bigr) = \sigma^2\). Dans ce cas,
la non-réponse n'affecte pas l'espérance, mais potentiellement la
variance (la dispersion du temps de travail peut être différente chez
les répondants et les non-répondants). Ces deux propriétés sont encore
vraies en cas d'indépendance entre \(X\) et \(Y\), puisque
l'indépendance implique la non corrélation, mais nous avons de plus
\(\mathbb{V}(Y\mid X) = \sigma^2 = \sigma^2_1 = \sigma^2_0\); la
variable aléatoire \(\mathbb{V}(Y\mid X)\) est elle aussi constante.
Cette fois-ci, la dispersion est la même chez les répondants et les
non-répondants~: la non-réponse n'affecte pas la variance.

\chapter{Solutions}

\runinhead{Démonstration}\phantomsection\addcontentsline{toc}{subsubsection}{Démonstration}\label{IV-answer-bienaymuxe9cheby}
C'est une application immédiate de
\protect\hyperref[IV-inegmarkov]{l'inégalité de
Markov (p.~\pageref*{IV-inegmarkov})} à \((X-\mathbb{E}(X))\) avec
\(p =2\).

\runinhead{Pile ou face
infini}\phantomsection\addcontentsline{toc}{subsubsection}{Pile ou face
infini}\label{IV-answer-pf}
En suivant l'indication, on a \(\mathbb{P}(A_n) = p\), d'où
\(\sum_{n=1}^\infty \mathbb{P}(A_n) = \infty\). Le lemme de
Borel-Cantelli nous indique alors que
\(\mathbb{P}(\lim \sup_{n \to \infty} A_n) = \mathbb{P}\left(\cap_{n\geq 1} \cup_{k \geq n} A_n \right) = 1\).
Autrement dit, on a presque sûrement un nombre infini de faces. En
passant au complémentaire, on en déduit que l'événement \(A\) est de
probabilité nulle.

\runinhead{En proba et en
moyenne}\phantomsection\addcontentsline{toc}{subsubsection}{En proba et
en moyenne}\label{IV-answer-lien}
Pour tout \(\varepsilon \in ]0,1[\), la probabilité
\(\mathbb{P}(|X_n|\geq \varepsilon) = \frac{1}{n}\) tend vers 0 quand
\(n\) tend vers l'infini. Ainsi, la suite
\((X_n)_{n\in \mathbb{N}^\ast}\) tend vers \(X=0\) en probabilité. Comme
\(\mathbb{E}(X_n) = \frac{1}{n}\), elle tend également en moyenne vers
0.

\runinhead{En proba mais pas en
moyenne}\phantomsection\addcontentsline{toc}{subsubsection}{En proba
mais pas en moyenne}\label{IV-answer-lien2}
Par le même argument que ci-dessus, nous voyons que la suite
\((Y_n)_{n\in \mathbb{N}^\ast}\) converge en probabilité vers 0, mais
comme \(\mathbb{E}(Y_n) = n\), la suite ne converge pas en moyenne vers
0 (ni vers aucune autre limite finie).

\runinhead{Presque
sûre}\phantomsection\addcontentsline{toc}{subsubsection}{Presque
sûre}\label{IV-answer-lien3}
Si \(\omega \in \{U > 0\}\) est fixé, alors
\(\exists n_0 \in \mathbb{N}\) tel que \(U(\omega) > \frac{1}{n_0}\), et
donc tel que \(Z_n(\omega) = 0\) pour tout \(n \geq n_0\). Comme
\(\mathbb{P}(U>0)=1\), ceci montre que la suite
\((Z_n)_{n\in \mathbb{N}^\ast}\) converge presque sûrement vers 0.

\runinhead{Intervalles de longueur aléatoire
(1)}\phantomsection\addcontentsline{toc}{subsubsection}{Intervalles de
longueur aléatoire (1)}\label{IV-answer-ia}
\(\mathbb{E}(X_n) = \frac{1}{n}\), et la suite \(X_n\) tend vers
\(X = 0\) en moyenne, et donc en probabilité.

\runinhead{Intervalles de longueur aléatoire
(2)}\phantomsection\addcontentsline{toc}{subsubsection}{Intervalles de
longueur aléatoire (2)}\label{IV-answer-ia2}
Supposons que les \(A_n\) soient placés bout-à-bout, en recommençant en
0 chaque fois qu'on arrive au point 1. Il est clair que l'on parcourt
indéfiniment l'intervalle {[}0, 1{]} (car la série de terme général
\(1/n\) diverge). Ainsi la suite numérique \(X_n (\omega)\) ne converge
pour aucun \(\omega\), et on n'a pas \(X_n \to X\) presque-sûrement.
Cependant comme la série \(\sum_n 1/n^2\) converge, il s'en suit que
\(X_{n^2} \to X = 0\) presque-sûrement. En effet, posant
\(B_n =\{\omega \in \Omega, X_{n^2}(\omega)=1\}\), on a
\(\mathbb{P}(B_n) = \frac{1}{n^2}\) et
\(\sum_{n=1}^{+\infty} \mathbb{P}(B_n) = \frac{\pi^2}{6}\). Le
\protect\hyperref[IV-BC]{lemme de Borel Cantelli (p.~\pageref*{IV-BC})}
nous indique alors que \(\mathbb{P}(B) = 0\) où
\(B = \cap_{n\geq 1}\cup_{k\geq n} B_n\) et donc
\(X_{n^2}(\omega) \to 0,\, \forall \omega \notin B\).

\runinhead{Loi faible des grands
nombres}\phantomsection\addcontentsline{toc}{subsubsection}{Loi faible
des grands nombres}\label{IV-answer-wlln}
Par l'inégalité de Bienaymé-Chebyshev, on a
\[\mathbb{P}(|M_n-m| > \epsilon) \leq \frac{\mathbb{V}(M_n)}{\epsilon} = \frac{\sigma^2}{n\epsilon},\]
où \(\sigma^2 = \mathbb{V}(X_i), \forall i \in \mathbb{N}^\ast\).

\runinhead{Exponentielle}\phantomsection\addcontentsline{toc}{subsubsection}{Exponentielle}\label{IV-answer-expo}
On a
\(\log\left(\prod_{i=1}^n Y_i\right)^{1/n} = \frac{1}{n}\sum_{i=1}^n X_i \to m\)
p.s. par la \protect\hyperref[IV-lfgn]{loi forte des grands
nombre (p.~\pageref*{IV-lfgn})}. Puisque l'exponentielle est continue,
la \protect\hyperref[IV-propconv4]{proposition de
continuité (p.~\pageref*{IV-propconv4})} nous indique que
\(\left(\prod_{i=1}^n Y_i\right)^{1/n} \to e^m\) p.s.

\runinhead{Variance}\phantomsection\addcontentsline{toc}{subsubsection}{Variance}\label{IV-answer-variance}
Les variables aléatoires \(Y_i = (X_i - m)^2\) sont i.i.d. telles que
\(\mathbb{E}(Y_i)= \sigma^2\). Par la \protect\hyperref[IV-lfgn]{loi
forte des grands nombre (p.~\pageref*{IV-lfgn})}, on a
\[\lim_{n \to \infty} \frac{1}{n} \sum_{i=1}^n (X_i - m)^2 \to \sigma^2  \text{   p.s.}\]

\runinhead{Suite
gaussienne}\phantomsection\addcontentsline{toc}{subsubsection}{Suite
gaussienne}\label{IV-answer-suitegauss}
Soit \(f\) une fonction continue bornée sur \(\mathbb{R}\). On a
\begin{align*}
\mathbb{E}(f(X_n)) &= \int_\mathbb{R}f(y) \frac{1}{\sqrt{2\pi}\sigma_n} \exp\left(-\frac{y^2}{2\sigma^2_n}\right) dy \\
             &\xrightarrow[n \to \infty]{} \int_\mathbb{R}f(y) \frac{1}{\sqrt{2\pi}\sigma} \exp\left(-\frac{y^2}{2\sigma^2}\right) dy
\end{align*} où l'on a utilisé le théorème de convergence dominée
appliqué à la suite de fonctions
\(\frac{1}{\sqrt{2\pi}\sigma_n} \exp\left(-\frac{y^2}{2\sigma^2_n}\right)\).

\runinhead{Surbooking}\phantomsection\addcontentsline{toc}{subsubsection}{Surbooking}\label{IV-answer-surbooking}
Le nombre de passager \(S_{420}\) se présentant effectivement à
l'embarquement suit une loi binomiale \(\mathcal{B}(420,0.92)\). On veut
donc calculer
\(\mathbb{P}(S_{420} > 400) = 1 - \mathbb{P}(S_{420}<400)\). Le TCL nous
fournit l'approximation
\[\frac{S_{420} - np}{\sqrt{n p (1-p)}} = \frac{S_{420} - 386.4}{5.56} \approx Y \sim \mathcal{N}(0,1)\]
Ainsi, \(\mathbb{P}(S_{420}<400) \approx \Phi(2.44) \approx 0.99\) La
probabilité de surbooking est donc de 1\%.

\section{Inégalités de
concentration}\label{IV-inuxe9galituxe9s-de-concentration-1}

\runinhead{Question
1}\phantomsection\addcontentsline{toc}{subsubsection}{Question
1}\label{IV-answer-idc1}
L'inégalité de Chebyshev nous donne pour tout \(a>0\)
\[\mathbb{P}\left(\left|\frac{1}{n}\sum_{i=1}^n X_i - m \right| > a\right) \leq \frac{\mathbb{V}(X_1)}{n a^2}\]
Prenant, \(\frac{\mathbb{V}(X_1)}{n a^2} = \delta\), on obtient
\(a = \frac{1}{\sqrt{n\delta}}\), d'où le résultat. On retrouve au
passage la même vitesse de convergence que celle donnée par le TCL.

\runinhead{Question
2}\phantomsection\addcontentsline{toc}{subsubsection}{Question
2}\label{IV-answer-idc2}

\begin{enumerate}
\def\labelenumi{\roman{enumi})}
\tightlist
\item
  \(M(u) = \mathbb{E}(e^{u(X_1 -m)}) = \phi_{X_1}(-iu) = e^{u^2 \sigma^2/2}\),
  où \(\phi_{X_1}\) est la fonction caractéristique de \(X_1\).
\item
  On a \(e^{u(M_n-nm)} \geq e^{ua} 1_{S_n -nm \geq a}\), d'où
  \[\mathbb{P}(S_n -nm \geq a) \leq \frac{\mathbb{E}(e^{u(M_n-nm)})}{e^{ua}} = e^{-ua}M(u)^n\]
  par indépendance des \(X_i\).
\item
  \begin{align*}
  \mathbb{P}(|Y_n - m| \geq \varepsilon) &= \mathbb{P}(S_n - nm \geq n\varepsilon) + P(S_n - nm \leq -n\varepsilon) \\
                         &\leq e^{-nu\varepsilon} (M (u))^n + e^{-n(-u)(-\varepsilon)} (M (-u))^n = 2e^{-nu\varepsilon}(M (u))^n \\
                         &\leq 2e^{-nu\varepsilon}e^{nu^2\sigma^2/2}, \,\,\, \forall u \geq 0
  \end{align*} La meilleure majoration va être obtenue en minimisant
  l'exposant, c'est-à-dire pour \(u = \varepsilon/\sigma^2\). Nous en
  déduisons l'inégalité de Chernov.
\end{enumerate}

\runinhead{Question
3}\phantomsection\addcontentsline{toc}{subsubsection}{Question
3}\label{IV-answer-idc3}
Par l'inégalité de Bienaymé-Chebyshev,
\(\mathbb{P}(|Y_n - m| > \varepsilon) \leq \frac{\mathbb{V}(Y_n)}{\varepsilon^2} = \frac{\sigma^2}{n\varepsilon^2}\).

Ainsi, \(\mathbb{P}(|Y_n - m| \leq \varepsilon)\leq \alpha\) dès que
\(1-\frac{\sigma^2}{n\varepsilon^2}\). Avec \(\varepsilon = 0,05\), il
vient \(n \geq 80000\).

Par l'inégalité de Chernov, nous obtenons
\(2e^{- n\varepsilon^2/20} \leq 0, 05\). Il vient que \(n \geq 29512\).
Pour avoir une évaluation du même ordre de la probabilité cherchée, nous
pouvons donc prendre un échantillon beaucoup plus petit si nous
utilisons l'inégalité de Chernov. A taille d'échantillon fixée, nous
aurons une meilleure évaluation avec cette inégalité.

\section{Lemme de Borel-Cantelli}\label{IV-lemme-de-borel-cantelli-1}

\runinhead{Question
1}\phantomsection\addcontentsline{toc}{subsubsection}{Question
1}\label{IV-answer-bc1}
On voit dans un premier temps que
\(\bigcap_{n\geq 0} \bigcup_{k \geq n} A_k \in \mathcal{A}\) par unions
et intersections dénombrables. On a
\[\mathbb{P}(\lim \sup_n A_n ) = \lim_{p \to \infty} \mathbb{P}(\cup_{n\geq p} A_n) \leq \lim_{p \to \infty} \sum_{n \geq p} \mathbb{P}(A_n),\]
où on remarque que les deux suites sont décroissantes.

Si la série \(\sum_n \mathbb{P}(A_n)\) est convergente, le reste de
cette série tend vers 0 et l'inégalité implique que
\(\mathbb{P}(\lim \sup_n A n) = 0\).

\runinhead{Question
2}\phantomsection\addcontentsline{toc}{subsubsection}{Question
2}\label{IV-answer-bc2}
Supposons maintenant que les \(A_n\) soient indépendants et que la série
diverge. Soit \(m\) un nombre entier. Nous avons \begin{align*}
\mathbb{P}(\cup_{i=p}^m A_i) &= 1- \mathbb{P}(\cap_{i=p}^m A_i^c) = 1- \prod_{i=p}^m \mathbb{P}(A_i^c) \text{ par indépendance}\\
& = 1- \prod_{i=p}^m (1-\mathbb{P}(A_i)) \geq 1- e^{-\sum_{i=p}^m\mathbb{P}(A_i)}
\end{align*} du fait de l'inégalité \(1 - x \leq e^{-x}\) pour
\(x \geq 0\). Ainsi,
\[ \mathbb{P}(\cup_{i=p}^{\infty} A_i) \geq 1- e^{-\sum_{i=p}^{\infty} \mathbb{P}(A_i)} = 1 \]
et l'on conclut finalement que pour tout \(p\),
\(\mathbb{P}(\cup_{i=p}^{\infty} A_i) = 1\), ce qui implique finalement
que \(\mathbb{P}(\lim \sup_{n \to \infty} A_n) = 1\).

\runinhead{Question
3}\phantomsection\addcontentsline{toc}{subsubsection}{Question
3}\label{IV-answer-bc3}
Prendre tous les \(A_n\) égaux à un même événement \(A\) de probabilité
\(\mathbb{P}(A) \in \left]0,1\right[\).

\section{Convergence vers une
constante}\label{IV-convergence-vers-une-constante-1}

\runinhead{Convergence
\(\mathcal{L}^2\)}\phantomsection\addcontentsline{toc}{subsubsection}{Convergence
\(\mathcal{L}^2\)}\label{IV-answer-cvtocst-l2}

\begin{enumerate}
\def\labelenumi{\arabic{enumi}.}
\tightlist
\item
  Supposons \(X_n \overset{\mathcal{L}^2}{\longrightarrow} X\) quand
  \(n\to+\infty\),
  i.e.~\(\mathbb{E}\left(\left(X_n - X\right)^2\right) \xrightarrow[n\to+\infty]{} 0\).
\end{enumerate}

Par positivité de la variance, on a tout d'abord
\[\mathbb{E}\left(\left(X_n - X\right)^2\right) \geq \left(\mathbb{E}\left(X_n-X\right)\right)^2 \geq 0,\]
qui garantit que
\(\left|\mathbb{E}\left(X_n-X\right)\right| = \left|\mathbb{E}(X_n) - \mathbb{E}(X)\right| \to 0\)
quand \(n\to+\infty\).

Ensuite, par inégalité triangulaire on a \begin{align*}
\left|\mathbb{V}\left(X_n\right) - \mathbb{V}\left(X\right)\right| &= \left|\mathbb{E}\left(X_n^2\right) - \mathbb{E}(X_n)^2 - \mathbb{E}\left(X^2\right) + \mathbb{E}(X)^2 \right|\\
& \leq \left|\mathbb{E}\left(X_n^2 - X^2\right)\right| + \left|\mathbb{E}(X_n)^2 - \mathbb{E}(X)^2\right|.
\end{align*}

Or, comme nous avons vu que
\(\mathbb{E}(X_n) \xrightarrow[n\to+\infty]{} \mathbb{E}(X)\), on a
directement que
\(\left|\mathbb{E}(X_n)^2 - \mathbb{E}(X)^2\right| \xrightarrow[n\to+\infty]{} 0\).
Par ailleurs, \begin{align*}
\left|\mathbb{E}\left(X_n^2 - X^2\right)\right| &= \left|\mathbb{E}\left( \left(X_n-X\right)^2 \right) - 2\mathbb{E}\left(X^2\right) + 2\mathbb{E}\left(X\,X_n\right)\right|\\
&= \left|\mathbb{E}\left( \left(X_n-X\right)^2 \right) + 2\mathbb{E}\left(X\,(X_n-X)\right)\right|\\
&\leq \mathbb{E}\left( \left(X_n-X\right)^2 \right) + 2\left|\mathbb{E}\left(X\,(X_n-X)\right)\right|.
\end{align*} Par l'inégalité de Cauchy-Schwartz on a
\[\left|\mathbb{E}\left(X\,(X_n-X)\right)\right|^2 \leq \mathbb{E}\left(X^2\right)\,\mathbb{E}\left( \left(X_n-X\right)^2 \right),\]
donc
\[\left|\mathbb{E}\left(X_n^2 - X^2\right)\right| \leq \mathbb{E}\left( \left(X_n-X\right)^2 \right) + 2\sqrt{\mathbb{E}\left(X^2\right)}\,\sqrt{\mathbb{E}\left( \left(X_n-X\right)^2 \right)},\]
qui tend bien vers \(0\) quand \(n\to+\infty\) par hypothèse. On en
conclut que
\(\left|\mathbb{V}\left(X_n\right) - \mathbb{V}\left(X\right)\right|\xrightarrow[n\to+\infty]{} 0.\)

\begin{enumerate}
\def\labelenumi{\arabic{enumi}.}
\setcounter{enumi}{1}
\tightlist
\item
  Par définition, \(X_n \overset{\mathcal{L}^2}{\longrightarrow} a\)
  quand \(n\to+\infty\) ssi
  \(\mathbb{E}\left(\left(X_n - a\right)^2\right) \xrightarrow[n\to+\infty]{} 0\).
\end{enumerate}

Le premier sens de l'équivalence est immédiatement obtenu par la
question précédente, en prenant \(X = a\) presque-sûrement.

Réciproquement, pour tout \(n\in\mathbb{N}^\ast\), par linéarité de
l'espérance on a \begin{align*}
\mathbb{E}\left( \left(X_n-a\right)^2 \right) &= \mathbb{E}\left( X_n^2 + a^2 - 2aX_n \right) = \mathbb{E}\left(X_n^2\right) + a^2 - 2a\mathbb{E}\left(X_n\right)\\
&= \mathbb{V}\left(X_n\right) + a^2 - 2a\mathbb{E}\left(X_n\right) + \mathbb{E}\left(X_n\right)^2\\
&= \mathbb{V}\left(X_n\right) + \left(\mathbb{E}(X_n) - a\right)^2.
\end{align*} On en déduit immédiatement le deuxième sens de
l'équivalence.

\runinhead{Convergences en loi et
probabilité}\phantomsection\addcontentsline{toc}{subsubsection}{Convergences
en loi et probabilité}\label{IV-answer-cvtocst-loiprob}
Supposons que \(X_n \overset{\mathcal{L}}{\longrightarrow} a\) quand
\(n\to+\infty\). Cela revient à dire que \(X_n\) tend en loi vers la
variable aléatoire \(X = a\) p.s. quand \(n\to+\infty\). Dans ce cas,
\(X\) a pour fonction de répartition
\(x\in\mathbb{R}\mapsto 1_{[a,+\infty[}(x)\) et pour tout
\(x \in \mathbb{R}\backslash\{a\}\), on a
\[\mathbb{P}\left(X_n \leq x \right) \xrightarrow[n\to+\infty]{} 1_{[a,+\infty[}(x).\]

Soit maintenant \(\varepsilon > 0\). D'après la croissance de la
fonction de répartition et l'hypothèse de convergence en loi, on a
\begin{align*}
\mathbb{P}\left(\left|X_n - a\right| \geq \varepsilon \right) &= \mathbb{P}\left(X_n \geq a + \varepsilon \right) + \mathbb{P}\left(X_n \leq a - \varepsilon \right)\\
&= 1 - \mathbb{P}\left(X_n < a + \varepsilon \right) + \mathbb{P}\left(X_n \leq a - \varepsilon \right)\\
&\leq 1 - \mathbb{P}\left(X_n \leq a + \frac{\varepsilon}{2} \right) + \mathbb{P}\left(X_n \leq a - \varepsilon \right)\\
& \xrightarrow[n\to+\infty]{} 1 - 1_{[a,+\infty[}\left(a+\frac{\varepsilon}{2}\right) + 1_{[a,+\infty[}\left(a-\varepsilon\right) = 0.
\end{align*}

Ainsi, on a bien \(X_n \overset{\mathbb{P}}{\longrightarrow} a\) quand
\(n\to+\infty\).

\section{Fonction de répartition
empirique}\label{IV-fonction-de-ruxe9partition-empirique-1}

\runinhead{Préliminaire}\phantomsection\addcontentsline{toc}{subsubsection}{Préliminaire}\label{IV-answer-fdremp-loi}
La variable aléatoire \(1_{]-\infty,x]}(X)\) est à valeurs dans
\(\{0,1\}\) et
\[\mathbb{P}\left( 1_{]-\infty,x]}(X) = 1\right) = F(x) = 1 - \mathbb{P}\left( 1_{]-\infty,x]}(X) = 0\right).\]
On reconnaît une loi de Bernoulli de paramètre \(F(x)\). On en déduit
directement que \(\mathbb{E}(1_{]-\infty,x]}(X)) = F(x)\) et
\(\mathbb{V}(1_{]-\infty,x]}(X)) = F(x)\,(1-F(x))\).

\runinhead{Biais}\phantomsection\addcontentsline{toc}{subsubsection}{Biais}\label{IV-answer-fdremp-biais}
Sous les hypothèses de l'exercice, on a \begin{align*}
\mathbb{E}\left(F_n(x) \right) &= \dfrac{1}{n}\,\sum_{i = 1}^n \mathbb{E}\left( 1_{]-\infty,x]}(X_i) \right) \hspace{1em}\text{par linéarité de l'espérance,}\\
& = \dfrac{1}{n}\,\sum_{i = 1}^n \mathbb{E}\left( 1_{]-\infty,x]}(X) \right) \hspace{1em}\text{car $X_1,\dots,X_n$ ont même loi que $X$,}\\
& = \dfrac{1}{n}\,\sum_{i = 1}^n F(x) \hspace{1em}\text{d'après la première question,}\\
& = F(x).
\end{align*}

\runinhead{Erreur quadratique
moyenne}\phantomsection\addcontentsline{toc}{subsubsection}{Erreur
quadratique moyenne}\label{IV-answer-fdremp-mse}
Soit \(n\in\mathbb{N}^\ast\). D'après la question précédente, on a

\(\mathbb{E}\left( \left(F_n(x) - F(x)\right)^2 \right) = \mathbb{E}\left( \left(F_n(x) - \mathbb{E}(F_n(x))\right)^2 \right) = \mathbb{V}\left(F_n(x)\right).\)

Or \begin{align*}
\mathbb{V}\left(F_n(x)\right) &= \dfrac{1}{n^2}\,\sum_{i = 1}^n \mathbb{V}\left(1_{]-\infty,x]}(X_i)\right) \hspace{1em}\text{par indépendance des variables,}\\
&= \dfrac{1}{n^2}\,\sum_{i = 1}^n \mathbb{V}\left(1_{]-\infty,x]}(X)\right) \hspace{1em}\text{car $X_1,\dots,X_n$ ont même loi que $X$,}\\
&= \dfrac{1}{n^2}\,\sum_{i = 1}^n F(x)\,(1-F(x)) \hspace{1em}\text{d'après la première question,}\\
&= \dfrac{1}{n}\, F(x)\,(1-F(x)) \xrightarrow[n\to+\infty]{} 0.
\end{align*}

On en conclut que
\(F_n(x) \overset{\mathcal{L}^2}{\longrightarrow} F(x)\) quand
\(n\to+\infty\).

\runinhead{Consistance}\phantomsection\addcontentsline{toc}{subsubsection}{Consistance}\label{IV-answer-fdremp-ps}
On remarque que pour tout \(n\in\mathbb{N}^\ast\), \(F_n(x)\) n'est
autre que la moyenne de \(n\) variables aléatoires indépendantes de même
loi de Bernoulli de paramètre \(F(x)\). La loi forte des grands nombres
nous assure donc qu'elle converge presque-sûrement vers l'espérance de
cette loi, qui n'est autre que \(F(x)\).

\addcontentsline{toc}{subsection}{Pour aller plus loin}
\begin{remarque}[{Pour aller plus
loin}]\phantomsection\label{IV-pour-aller-plus-loin}

On peut en fait aller plus loin et montrer que l'on a la convergence
presque sûre uniformément sur \(\mathbb{R}\) voir par exemple
\href{http://math.univ-lyon1.fr/~gelineau/devagreg/Theoreme_Dini.pdf}{ce
document}, c'est le théorème de Glivenko-Cantelli très utile en
statistiques.
\end{remarque}

\section{Théorème de Weierstrass sur
\([0,1]\)}\label{IV-thuxe9oruxe8me-de-weierstrass-sur-01-1}

\runinhead{Preuve du
théorème.}\phantomsection\addcontentsline{toc}{subsubsection}{Preuve du
théorème.}\label{IV-answer-weier-thm}

\begin{enumerate}
\def\labelenumi{\arabic{enumi}.}
\tightlist
\item
  Comme \(X_1,\dots,X_n\) sont indépendantes de même loi de Bernoulli de
  paramètre \(x\in]0,1[\), \(M_n\) est à valeurs dans
  \(\bigl\{\frac{k}{n} : k \in \{0,\dots,n\} \bigr\}\) et pour tout
  \(k\in\{0,\dots,n\}\) on a
  \[\mathbb{P}\left(M_n = \dfrac{k}{n}\right) = \binom{n}{k} x^k\,(1-x)^{n-k}.\]
  Ainsi,
  \[\mathbb{E}\left(f(M_n)\right) = \sum_{k=0}^{n} \binom{n}{k} f\left(\dfrac{k}{n}\right)\, x^k\,(1-x)^{n-k}.\]
\end{enumerate}

Les polynômes de la forme
\(B_{n,k} : u\in [0,1] \mapsto \binom{n}{k} u^k\,(1-u)^{n-k}\) sont
appelés \emph{polynômes de Bernstein}.

\begin{enumerate}
\def\labelenumi{\arabic{enumi}.}
\setcounter{enumi}{1}
\tightlist
\item
  Comme une variable de Bernoulli admet une espérance égale à son
  paramètre, la loi (forte) des grands nombres nous assure que
  \(M_n \to x\) p.s. quand \(n\to+\infty\). Puisque \(f\) est continue,
  on a de même \(f\left(M_n\right) \to f(x)\) p.s. quand
  \(n\to+\infty\).
\end{enumerate}

Par ailleurs, la continuité de \(f\) sur \([0,1]\) nous assure que \(f\)
est bornée (l'image d'un compact par une fonction continue
\(\mathbb{R}\to \mathbb{R}\) est un compact). Ainsi, pour tout
\(n\in\mathbb{N}^\ast\), la variable aléatoire \(f\left(M_n\right)\)
l'est également. Cela nous permet d'appliquer le théorème de convergence
dominée de la question 1~: en notant \((\Omega,\mathcal{A},\mathbb{P})\)
l'espace probabilisé sur lequel sont définies nos variables aléatoires,

\[\mathbb{E}\left( f\left(M_n\right) \right) = \int_{\Omega} f\left(M_n(\omega)\right)\,\mathbb{P}(d\omega) \xrightarrow[n\to+\infty]{} \int_{\Omega} f(x)\,\mathbb{P}(d\omega) = f(x).\]

\begin{enumerate}
\def\labelenumi{\arabic{enumi}.}
\setcounter{enumi}{2}
\item
  En combinant les résultats des deux questions précédentes, on obtient
  que la suite \(\left(P_n\right)_{n\in\mathbb{N}^\ast}\) des polynômes
  définis pour tout \(n\in\mathbb{N}^\ast\) par
  \[P_n : x \in ]0,1[ \mapsto \sum_{k=0}^{n} \binom{n}{k} f\left(\dfrac{k}{n}\right)\, x^k\,(1-x)^{n-k}\]
  converge simplement vers \(f(x)\). On étend simplement ce résultat à
  \([0,1]\) en remarquant que pour tout \(n\in\mathbb{N}^\ast\) on a
  directement l'égalité
  \[P_n(0) = f\left(\dfrac{0}{n}\right) = f(0) \text{ et } P_n(1) = f\left(\dfrac{n}{n}\right) = f(1).\]
  On a même toujours l'égalité
  \(P_n(x) = \mathbb{E}\left( f\left(M_n\right) \right)\) pour
  \(x\in[0,1]\) si l'on remarque qu'une loi de Bernoulli de paramètre
  \(a\in \{0,1\}\) n'est autre d'une Dirac en \(\{a\}\), ce qui signifie
  que \(M_n = a\) p.s. et donc que
  \(\mathbb{E}\left(f(M_n)\right) = f(a) = P_n(a)\).
\item
  On souhaite montrer que
  \[\forall\,\varepsilon > 0 \ \exists\,N_\varepsilon\in\mathbb{N}^\ast : \left(n\geq N_\varepsilon\right) \Rightarrow \left(\forall\,x\in[0,1] : \left|P_n(x) - f(x)\right|\leq \varepsilon\right).\]
\end{enumerate}

Remarquons tout d'abord que la fonction \(f\) étant bornée, elle admet
un maximum sur \([0,1]\), que l'on note \(K\). Puisqu'elle est continue
sur un segment réel, elle est aussi uniformément continue~:
\[\forall\,\varepsilon > 0\ \ \exists\,\delta_{\varepsilon} > 0 : \forall\,(x,y) \in ]0,1[^2,\ \ |x-y| < \delta_\varepsilon \Rightarrow \left|f(x) - f(y)\right| < \varepsilon.\]
Soient \(\varepsilon > 0\), puis \(\delta = \delta_{\varepsilon/2} > 0\)
pour lequel l'implication précédente est vraie avec
\(\dfrac{\varepsilon}{2}\). Alors pour tous \(n\in\mathbb{N}^\ast\) et
\(x\in]0,1[\) on a \begin{align*}
\left|\mathbb{E}\left( f(M_n) \right) - f(x)\right| & \leq \mathbb{E}\left( \left|f(M_n) - f(x) \right| \right) \hspace{1em}\text{par l'inégalité de Jensen,}\\
& = \mathbb{E}\left( \left|f(M_n) - f(x) \right|\,1_{[0,\delta[}\left(\left|M_n - x\right|\right) \right)\\
&\ \ \ + \mathbb{E}\left( \left|f(M_n) - f(x) \right|\,1_{[\delta,+\infty[}\left(\left|M_n - x\right|\right) \right)\\
&\leq \dfrac{\varepsilon}{2} + 2|K|\,\mathbb{P}\left(\left|M_n-x\right|\geq \delta\right) \hspace{1em}\text{par inégalité triangulaire,}\\
&\leq \dfrac{\varepsilon}{2} + 2|K|\,\dfrac{\mathbb{V}(M_n)}{\delta^2} \hspace{1em}\text{par Bienaymé-Chebyshev.}
\end{align*} Or pour tout \(n\in\mathbb{N}^\ast\) \begin{align*}
\mathbb{V}(M_n) &= \mathbb{V}\left(\dfrac{1}{n}\,\sum_{i = 1}^n X_i \right)\\
&= \dfrac{1}{n^2}\,\sum_{i = 1}^n \mathbb{V}(X_i) \hspace{1em}\text{par indépendance de $X_1,\dots,X_n$,}\\
&= \dfrac{1}{n^2}\,\sum_{i = 1}^n x(1-x) \hspace{1em}\text{car $X_1,\dots,X_n$ suivent la même loi de Bernoulli},\\
&= \dfrac{x(1-x)}{n} < \dfrac{1}{4n} \hspace{1em}\text{car $x \in ]0,1[$.}
\end{align*} Par conséquent, pour tout \(n\in\mathbb{N}^\ast\),
\[\left|\mathbb{E}\left( f(M_n) \right) - f(x)\right| < \dfrac{\varepsilon}{2} + \dfrac{|K|}{2n\delta^2}.\]
En prenant \(N_\varepsilon\) le plus petit entier naturel supérieur ou
égal à \(\dfrac{|K|}{\varepsilon\,\delta^2}\), on obtient alors que pour
tout \(n\geq N_\varepsilon\),
\(\left|\mathbb{E}\left( f(M_n) \right) - f(x)\right| < \varepsilon.\)
Comme \(\varepsilon\) et \(N_\varepsilon\) sont les mêmes quel que soit
\(x \in [0,1]\), on obtient bien la convergence uniforme désirée.

\section{Théorème de Slutsky}\label{IV-thuxe9oruxe8me-de-slutsky-1}

\runinhead{Question
1}\phantomsection\addcontentsline{toc}{subsubsection}{Question
1}\label{IV-answer-slut1}
La démonstration est analogue à celle de la
\protect\hyperref[IV-cvceloifdr]{la proposition sur la convergence des
f.d.r. (p.~\pageref*{IV-cvceloifdr})}, où on peut voir que l'on peut
remplacer les fonctions continues bornées \(f_{p,b}\) approchant
\(1_{\left]-\infty,b\right]}\) par des fonctions lipschitziennes
bornées. Voir la démonstration du théorème 18.7 dans Jacod and Protter
(\protect\hyperlink{refIV-Jacod}{2003}) pour une version complète.

\runinhead{Question
2}\phantomsection\addcontentsline{toc}{subsubsection}{Question
2}\label{IV-answer-slut2}
Du fait du résultat de la question 1, il suffit de montrer que
\(\lim_{n \to \infty} \mathbb{E}(f (Y_n)) = \mathbb{E}(f (X))\), pour
toute fonction lipschitzienne bornée. Soit \(f\) une telle fonction. On
a alors \(|f (x)| \leq k\) et \(|f (x) -f (y)| \leq C|x - y|\), pour des
constantes \(k\) et \(C\).

Soit \(\varepsilon > 0\) donné. On a \begin{align*}
|\mathbb{E}(f (Y_n )) - \mathbb{E}(f (X_n ))| &\leq \mathbb{E}(|f (Y_n) - f (X_n )|) \\
                                  &\leq \mathbb{E}(|f (Y_n) - f (X_n )|(1_{|Y_n-X_n|\leq \varepsilon} + 1_{|Y_n-X_n| > \varepsilon}))\\
                                  &\leq C \varepsilon + 2k P(|X_n - Y_n | > \varepsilon)
\end{align*} Le deuxième terme du membre de droite tend vers 0 quand
\(n\) tend vers l'infini, car \((X_n - Y_n )_n\) converge en probabilité
vers 0. Comme \(\varepsilon\) est arbitrairement petit, nous en
déduisons que
\(\lim_{n\to\infty}|\mathbb{E}(f (Y_n )) - \mathbb{E}(f (X_n ))| = 0\),
d'où le résultat.

\runinhead{Question
3}\phantomsection\addcontentsline{toc}{subsubsection}{Question
3}\label{IV-answer-slut3}
Soit \(u \in R^d\) . On a~:
\[|\phi_{Y_n} (u) - \phi_X (u)| \leq |\phi_{Y_n} (u) - \phi_{X_n} (u)| + |\phi_{X_n} (u) - \phi_{X} (u)|\]
D'une part, le \protect\hyperref[IV-levytheorem]{théorème de
Lévy (p.~\pageref*{IV-levytheorem})} partie 1. montre que
\(|\phi_{X_n}(u) - \phi_X (u)|\) tend vers 0 quand \(n \to \infty\).
D'autre part, \begin{align*}
|\phi_{Y_n} (u) - \phi_{X_n} (u)|  & = |\mathbb{E}(e^{i< u, Y_n>} - e^{i< u, X_n>})| \\
& = |\mathbb{E}(e^{i<u,X_n>} (e^{i<u,Y_n-X_n>} - 1))|\leq \mathbb{E}(|e^{i<u,Y_n-X_n>} - 1|)
\end{align*} tend vers 0 quand \(n \to \infty\) d'après
\protect\hyperref[IV-propconv2]{la proposition --- cas
borné (p.~\pageref*{IV-propconv2})} appliquée aux variables aléatoires
\(e^{i<u,Y_n -X_n>}-1\) dont la convergence en proba vers 0 est assurée
par la \protect\hyperref[IV-propconv4]{propriété de
continuité (p.~\pageref*{IV-propconv4})}.

\chapter{Solutions}

\runinhead{Calcul numérique d'une
probabilité}\phantomsection\addcontentsline{toc}{subsubsection}{Calcul
numérique d'une probabilité}\label{V-answer-mc1}
cf.~notebook

\runinhead{Exemples d'application
1}\phantomsection\addcontentsline{toc}{subsubsection}{Exemples
d'application 1}\label{V-answer-exemples1}
On suppose que \(U \sim \mathcal{U}_{]0,1[}\)

\begin{itemize}
\item
  Uniforme sur un intervalle \(I \subset \mathbb{R}\)~:

  Si \(I = ]a,b[,~a< b\), alors la fonction de répartition de
  \(X \sim \mathcal{U}_I\) est
  \(F_X(x) = \frac{x-a}{b-a} 1_{]a,b[}(x) + 1_{[b,+\infty[}(x)\), d'où
  par inversion \(X = a + (b-a)U\)
\item
  Exponentielle de paramètre \(\lambda \in \mathbb{R}_+^\ast\)~:

  la fonction de répartition de \(X \sim \mathcal{E}(\lambda)\) est
  \(F_X (x) = 1-\exp(-\lambda x)\), d'où par inversion
  \(X = -\frac{\ln(1-U)}{\lambda}\). Pour simplifier, on remarquera que
  si \(U \sim \mathcal{U}_{]0,1 [ }\) , alors
  \(1- U \sim \mathcal{U}_{]0,1[}\), d'où
  \(X = -\frac{\ln(U)}{\lambda}\)
\item
  de Cauchy, de densité
  \(x\in\mathbb{R}\mapsto \left(\pi\left(1+x^2\right)\right)^{-1}\),

  la loi de Cauchy admet la fonction de répartition
  \(F(x) = \frac{1}{\pi} \arctan (x) + \frac{1}{2}\), d'où
  \(X = \tan(\pi U -1/2)\)
\item
  de Laplace de paramètres \(\mu \in \mathbb{R}\) et
  \(s\in\mathbb{R}_+^\ast\), de densité
  \(x\in\mathbb{R}\mapsto \frac{1}{2s}\,\exp\left\{-\frac{|x-\mu|}{s}\right\}\),

  la loi de Laplace admet la fonction de répartition
  \[F(x) = \left\{ \begin{array}{ll}
    \frac{1}{2} \exp(x) & \text{ si } x < 0,\\
    1 - \frac{1}{2} \exp(-x) & \text{ si } x \geq 0,
    \end{array} \right. = \frac{1}{2}\left(1 + \text{sgn}(x)(1-\exp(-|x|))\right)\]

  d'où \(X = \text{sgn}(U - 1/2)\ln(1-2|U-1/2|)\), où sgn est la
  fonction signe.
\item
  Logistique de paramètres \(\mu \in \mathbb{R}\) et
  \(s\in\mathbb{R}_+^\ast\), de fonction de répartition
  \(x\in\mathbb{R}\mapsto \left(1 + \exp\left\{-\frac{x-\mu}{s} \right\}\right)^{-1}\)~:

  par inversion,
  \(X = s \left(-\ln\left(\frac{1}{U}-1\right)+ \mu\right)\)
\end{itemize}

\runinhead{Pile ou
face}\phantomsection\addcontentsline{toc}{subsubsection}{Pile ou
face}\label{V-answer-pf}
Il suffit d'associer à ``pile'' un événement portant sur
\(U \sim \mathcal{U}_{]0,1[}\) de la bonne probabilité, \(p \in ]0,1[\),
par exemple \(\{U \leq p\}\). On obtient ainsi l'algorithme~:

\begin{enumerate}
\def\labelenumi{\arabic{enumi}.}
\tightlist
\item
  Générer \(u\) selon \(\mathcal{U}_{]0,1[}\)
\item
  Renvoyer ``pile'' si \(u \leq p\), ``face'' sinon
\end{enumerate}

\runinhead{Exemples d'application
2}\phantomsection\addcontentsline{toc}{subsubsection}{Exemples
d'application 2}\label{V-answer-exemples2}

\begin{itemize}
\item
  Binomiale de paramètres \(n\in\mathbb{N}^\ast\) et \(p \in\, ]0,1[\),

  La fonction de répartition de la loi binomiale est donnée par
  \[F(x) = \begin{cases}
    0&\text{si } \;x<0\\
    \sum_{{k=0}}^{{\lfloor x\rfloor }}{n \choose k}p^{k}(1-p)^{{n-k}} & \text{si } \;0\leq x<n\\
    1&\text{si} \;x \geq n
    \end{cases}\]

  où \({\lfloor x\rfloor }\) est la partie entière de \(x\).\\
  On peut alors découper \(]0,1[\) en intervalles de la forme
  \(]\sum_{{k=0}}^{{m}}{n \choose k}p^{k}(1-p)^{{n-k}},\sum_{{k=0}}^{{m+1}}{n \choose k}p^{k}(1-p)^{{n-k}}[\)
  pour \(m \in \{0,n-1\}\) et renvoyer \(m+1\) si un \(u\) généré selon
  \(\mathcal{U}_{] 0,1[}\) appartient à l'un de ces intervalles, 0
  sinon.

  Il est cependant beaucoup plus simple de remarquer qu'une variable
  aléatoire de loi \(\mathcal{B}(n,p)\) s'obtient comme la somme de
  \(n\) variables de Bernoulli que l'on peut simuler comme
  \protect\hyperref[V-answer-pf]{ci-dessus (p.~\pageref*{V-answer-pf})}.
\end{itemize}

\begin{itemize}
\item
  Uniforme sur l'union de deux segments non vides et disjoints
  \([a,b], [c,d]\subset\mathbb{R}\), de densité
  \(x\in\mathbb{R}\mapsto (b-a + d-c)^{-1}\,1_{[a,b]\cup[c,d]}(x)\)

  La fonction de répartition d'une telle variable est
  \[F(x) = \frac{x-a}{b-a + d-c} 1_{]a,b[}(x) + \frac{b-a}{b-a + d-c} 1_{[b,c]}(x) + \frac{x-c+b-a}{b-a + d-c} 1_{]c,d[}(x) + 1_{[d,+\infty[}(x)\]
  Pour simuler cette variable, on peut donc utiliser l'algorithme
  suivant~:

  \begin{enumerate}
  \def\labelenumi{\arabic{enumi}.}
  \tightlist
  \item
    Générer \(u\) selon \(\mathcal{U}_{]0,1[}\)
  \item
    Retourner \[\begin{cases} 
      a + u(b-a + d -c) & \text{ si } u \in ]0,\frac{b-a}{b-a + d-c}[ \\
      c + (u(b-a + d -c) - (b-a)) \text{ sinon }
      \end{cases}\]
  \end{enumerate}
\end{itemize}

\runinhead{Loi uniforme sur un
pavé}\phantomsection\addcontentsline{toc}{subsubsection}{Loi uniforme
sur un pavé}\label{V-answer-unipave}
On note \([a_1,b_1]\times\dots\times[a_d,b_d] \subset \mathbb{R}^d\). On
peut voir que la densité d'une v.a. uniforme sur un tel pavé s'écrit
\(f(x) = \frac{1}{b_1 - a_1} \cdots \frac{1}{b_d - a_d} 1_{[a_1,b_1]\times\dots\times[a_d,b_d]}(x)\)
soit comme le produit des densités de ses coordonnées qui sont donc
indépendantes. On peut donc simuler la variable d'intérêt avec cet
algorithme~:

\begin{enumerate}
\def\labelenumi{\arabic{enumi}.}
\tightlist
\item
  Générer \(u_1,\ldots,u_d\) indépendamment selon
  \(\mathcal{U}_{]0,1[}\)
\item
  Retourner \((a_1 + u_1(b_1-a_1), \ldots, a_d + u_d(b_d-a_d))\)
\end{enumerate}

\runinhead{Démonstration}\phantomsection\addcontentsline{toc}{subsubsection}{Démonstration}\label{V-answer-unicond}
Soit \(C \in \mathcal{B}(B)\) (considérer les boréliens de \(B\) suffit
à caractériser la loi conditionnelle à l'événement \(U \in B\) puisque
si \(C \cap B = \varnothing\), alors nécessairement,
\(\mathbb{P}(U \in C |U \in B) = 0\)). On a \begin{align*}
\mathbb{P}(U \in C |U \in B) & = \frac{\mathbb{P}(U\in C, U \in B)}{\mathbb{P}(U \in B)} = \frac{\mathbb{P}(U\in C)}{\mathbb{P}(U \in B)} \\
& = \frac{\int_C \frac{dx}{\lambda(A)}}{\int_B \frac{dx}{\lambda(A)}} \text{ puisque } U \sim \mathcal{U}_A \\
& = \frac{\lambda(C)}{\lambda(B)} = \int_C \frac{1}{\lambda(B)} dx
\end{align*} on reconnaît bien la loi uniforme sur \(B\).

\runinhead{Simulation}\phantomsection\addcontentsline{toc}{subsubsection}{Simulation}\label{V-answer-simunibor}
On prend un pavé \(R\) tel que \(B \subset R\) puis on applique
l'algorithme suivant~:

\begin{enumerate}
\def\labelenumi{\arabic{enumi}.}
\tightlist
\item
  générer \(u\) selon \(\mathcal{U}_R\)
\item
  Si \(u \notin R\), retour en 1.
\item
  retourner \(u\)
\end{enumerate}

\runinhead{Taux de
rejet}\phantomsection\addcontentsline{toc}{subsubsection}{Taux de
rejet}\label{V-answer-tauxrej}
Puisque le couple \((Y, Uf_Y(Y))\) est uniforme sur \(A_f\), alors le
couple \((Y, aUf_Y(Y))\) est uniforme sur \(A_{af}\), d'où
\begin{align*}
\mathbb{P}(\{aUf_Y(Y) > f_X(Y)\}) & = \mathbb{E}(1_{\{aUf_Y(Y) > f_X(Y)\}}) \\
& = \mathbb{E}(\mathbb{E}(1_{\{aUf_Y(y) > f_X(y)\}}|Y=y)) \\
& = \int_R \int_0^{af_Y(y)} 1_{\{auf_Y(Y) > f_X(Y)\}} \frac{1}{af_Y(y)} du f_Y(y) dy \\
& = \int_R \left(1 - \frac{f_X(y)}{af_Y(y)}\right) f_Y(y) dy \\
& = 1 - \frac{1}{a}
\end{align*}

\section{Loi uniforme dans un
domaine}\label{V-loi-uniforme-dans-un-domaine-1}

cf.~notebook

\section{Simulation selon la loi
géométrique}\label{V-simulation-selon-la-loi-guxe9omuxe9trique-1}

\runinhead{Question
1}\phantomsection\addcontentsline{toc}{subsubsection}{Question
1}\label{V-answer-loigeom1}
\(\mathbb{P}(X=n) = p (1-p)^{n-1} = (1-q)q^{n-1} \text{, où } p \text{ est la probabilité de succès, } n=1,2,\ldots\)
\begin{align*}
\mathbb{E}[X]&=\sum_{i=1}^{\infty}i\mathbb{P}(X=i)\\
&=\sum_{i=1}^{\infty} i (1-q)q^{i-1}\\
&=(1-q)\sum_{i=1}^{\infty} iq^{i-1}\\
&=(1-q)\sum_{i=1}^{\infty} \frac{\mathrm{d}}{\mathrm{d} q} q^i\\
&=(1-q)\frac{\mathrm{d}}{\mathrm{d} q}\left(\sum_{i=1}^{\infty}  q^i\right)\\
&=(1-q)\frac{\mathrm{d}}{\mathrm{d} q}\left(\frac{q}{1-q}\right), \text{ car } \sum_{i=1}^{\infty}  q^i = q \sum_{i=0}^{\infty}  q^i = \frac{q}{1-q}\\
&=(1-q)\frac{1}{(1-q)^2} \\
&=\frac{1}{(1-q)}\\
&=\frac{1}{p}
\end{align*} \bigskip Pour la variance
\(\mathbb{V}[X] = \mathbb{E}[X^2]-\mathbb{E}[X]\) \begin{align*}
\mathbb{E}[X^2]&=\sum_{i=1}^{\infty}i^2\mathbb{P}(X=i)\\
&=\sum_{i=1}^{\infty} i^2 (1-q)q^{i-1}\\
&=(1-q)\sum_{i=1}^{\infty} i^2q^{i-1}\\
&=(1-q)\sum_{i=1}^{\infty}\left( i\left((i+1)q^{i-1} - q^{i-1}\right)\right)\\
&=(1-q)\sum_{i=1}^{\infty}\left( i(i+1)q^{i-1} - iq^{i-1}\right)\\
&=(1-q)\sum_{i=1}^{\infty}\left( \frac{\mathrm{d}^2}{\mathrm{d} q^2} q^{i+1} - \frac{\mathrm{d}}{\mathrm{d} q} q^i\right)\\
&=(1-q)\left( \frac{\mathrm{d}^2}{\mathrm{d} q^2}\left(\sum_{i=1}^{\infty}q^{i+1}\right)-\frac{\mathrm{d}}{\mathrm{d} q}\left(\sum_{i=1}^{\infty}  q^i\right)\right)
\end{align*} On a déjà calculé le second terme ci-dessus. \begin{align*}
\frac{\mathrm{d}^2}{\mathrm{d} q^2}\sum_{i=1}^{\infty}q^{i+1} &= \frac{\mathrm{d}^2}{\mathrm{d} q^2}q \sum_{i=1}^{\infty}q^{i}\\
&=\frac{\mathrm{d}^2}{\mathrm{d} q^2}\frac{q^2}{1-q}
\end{align*} Calcul des dérivées~: \begin{align*}
\frac{\mathrm{d}}{\mathrm{d} q}\frac{q^2}{1-q} &= \frac{2q(1-q)+q^2}{(1-q)^2}\\
&= \frac{2q-q^2}{(1-q)^2}
\end{align*}

\begin{align*}
\frac{\mathrm{d}}{\mathrm{d} q}\frac{2q-q^2}{(1-q)^2}&=\frac{2(1-q)(1-q)^2+2(1-q)(2q-q^2)}{(1-q)^4}\\
&= \frac{2(1-q)^2+2(2q-q^2)}{(1-q)^3}\\
&= \frac{2 - 4q +2q^2 + 4q -2q^2}{(1-q)^3}\\
&= \frac{2}{(1-q)^3}
\end{align*} On en déduit \begin{align*}
\mathbb{E}[X^2]&=(1-q)\left(\frac{2}{(1-q)^3}-\frac{1}{(1-q)^2}\right)\\
&=\frac{2}{(1-q)^2}-\frac{1}{(1-q)}\\
&=\frac{2-(1-q)}{(1-q)^2}\\
&=\frac{1+q}{(1-q)^2}
\end{align*} Et finalement \begin{align*}
\mathbb{V}[X] &= \mathbb{E}[X^2]-\mathbb{E}[X]\\
&= \frac{1+q}{(1-q)^2}-\frac{1}{(1-q)^2}\\
&=\frac{q}{(1-q)^2}\\
&=\frac{1-p}{p^2}
\end{align*}

\runinhead{Question
2}\phantomsection\addcontentsline{toc}{subsubsection}{Question
2}\label{V-answer-loigeom2}
cf.~notebook

\runinhead{Question
3}\phantomsection\addcontentsline{toc}{subsubsection}{Question
3}\label{V-answer-loigeom3}

\begin{align*}
\mathbb{P}(\lceil X \rceil = n) & = \mathbb{P}(n-1 < X \leq n) \\
& = F_X(n) - F_X(n-1) \\
& = e^{-\lambda (n-1)}(1-e^{-\lambda})
\end{align*} d'où \(\lambda = -\ln(1-p)\)

\runinhead{Question
4}\phantomsection\addcontentsline{toc}{subsubsection}{Question
4}\label{V-answer-loigeom4}
cf.~notebook

\section{Simulation de la loi gaussienne par la méthode de
rejet}\label{V-simulation-de-la-loi-gaussienne-par-la-muxe9thode-de-rejet-1}

\(X\) est une variable aléatoire de loi gaussienne d'espérance \(m=0\)
et de variance \(\sigma^2=1\), \(\mathcal{N}(0,1)\). Sa densité est
\[f(x) = \frac{1}{\sqrt{2\pi}}\exp\left(-\frac{1}{2}x^2\right),~x\in \mathbb{R}\]

\runinhead{Question
1}\phantomsection\addcontentsline{toc}{subsubsection}{Question
1}\label{V-answer-lgrej1}
Soit \(x \in \mathbb{R}\) \begin{align*}
x^2 & \geq 2|x| - 1 \\
\frac{1}{\sqrt{2\pi}}e^{-x^2/2} & \leq \sqrt{\frac{e}{2\pi}} e^{-|x|}
\end{align*}

\runinhead{Question
2}\phantomsection\addcontentsline{toc}{subsubsection}{Question
2}\label{V-answer-lgrej2}
\(g\) est bien positive et son intégrale sur \(\mathbb{R}\) vaut
\[\int_\mathbb{R}g(x) dx = \int_\mathbb{R}\frac{1}{2} e^{-|x|} dx = \int_{\mathbb{R}^+} e^{-x} dx = 1\]

\runinhead{Question
3}\phantomsection\addcontentsline{toc}{subsubsection}{Question
3}\label{V-answer-lgrej3}
voir dans \protect\hyperref[V-exemples1]{exemples d'application
1 (p.~\pageref*{V-exemples1})}.

\runinhead{Question
4}\phantomsection\addcontentsline{toc}{subsubsection}{Question
4}\label{V-answer-lgrej4}
cf.~notebook

\section{Simulation de la loi de Wigner}\label{V-answer-wigner}

cf.~notebook

\section{Loi des grands nombres et théorème central
limite}\label{V-loi-des-grands-nombres-et-thuxe9oruxe8me-central-limite-1}

\runinhead{Question
1}\phantomsection\addcontentsline{toc}{subsubsection}{Question
1}\label{V-answer-lgntcl1}
cf.~notebook

\runinhead{Question
2}\phantomsection\addcontentsline{toc}{subsubsection}{Question
2}\label{V-answer-lgntcl2}
cf.~notebook

\runinhead{Question
3}\phantomsection\addcontentsline{toc}{subsubsection}{Question
3}\label{V-answer-lgntcl3}
cf.~notebook

\section{Simulation d'un mélange de
gaussiennes}\label{V-simulation-dun-muxe9lange-de-gaussiennes-1}

\runinhead{Question
1}\phantomsection\addcontentsline{toc}{subsubsection}{Question
1}\label{V-answer-melloi1}
cf.~notebook

\runinhead{Question
2}\phantomsection\addcontentsline{toc}{subsubsection}{Question
2}\label{V-answer-melloi2}
cf.~notebook

\section{Echantillonnage
d'importance}\label{V-echantillonnage-dimportance-3}

\runinhead{Question
1}\phantomsection\addcontentsline{toc}{subsubsection}{Question
1}\label{V-answer-echimp1}
On a
\(F_{X|X>3}(x) = \mathbb{P}(X \leq x |X > 3) = \frac{\mathbb{P}(X\leq x, X>3)}{1-F_X(3)} = \frac{F_X(x)-F_X(3)}{1-F_X(3)}1_{]3,+\infty[}(x)\)
d'où \(f_{X|X>3}(x) = \frac{f_X(x)}{1-F_X(3)}1_{]3,+\infty[}(x)\)

\runinhead{Question
2}\phantomsection\addcontentsline{toc}{subsubsection}{Question
2}\label{V-answer-echimp2}
cf.~notebook

\runinhead{Question
3}\phantomsection\addcontentsline{toc}{subsubsection}{Question
3}\label{V-answer-echimp3}
cf.~notebook

\fi

\end{document}
